% ============================================================================
% Fichier     : NRAC_main.tex
% Titre       : Normes, gestion des risques et audit de cybersécurité
% Auteur      : MINKA MI NGUIDJOÏ Thierry Emmanuel
% Institution : Chaire UNESCO d'Afrique Centrale Accès TIC — ENSP / LIMSI
% Version     : 1.0 -- novembre 2026 (traité complet : P00 à C16, annexes A à F)
% ============================================================================

\documentclass[12pt, twoside, openright, a4paper]{book}

% --- Feuille de style unifiée ---
\usepackage{Style_NRAC}
\raggedbottom

% --- Bibliographie ---
\addbibresource{bibliographie/NRAC_references.bib}

% ============================================================================
% MÉTADONNÉES PDF
% ============================================================================
\hypersetup{
    pdftitle={Normes, gestion des risques et audit de cybersécurité},
    pdfauthor={MINKA MI NGUIDJOÏ Thierry Emmanuel},
    pdfsubject={MSN-2S3-SEC08 · Master Professionnel en Sécurité des Systèmes Numériques · Chaire UNESCO d'Afrique Centrale Accès TIC},
    pdfkeywords={ISO/IEC 27001, ISO/IEC 27005, EBIOS RM, FAIR, audit de cybersécurité, GRC, NEK}
}

\begin{document}

% --------------------------------------------------------------------------
% ÉLÉMENTS PRÉLIMINAIRES (pages romaines)
% --------------------------------------------------------------------------
\frontmatter

    % ==============================================================================
% Fichier : front/00_page_titre.tex
% Titre   : Page de titre - Normes, Risques et Audit de Cybersécurité
% Auteur  : MINKA MI NGUIDJOÏ Thierry Emmanuel
% Version : 1.0 -- octobre 2026
% ==============================================================================

\begin{titlepage}
\thispagestyle{empty}

% --- Fond décoratif : bandeaux bleu marine et filets or ---
\begin{tikzpicture}[remember picture, overlay]
    \fill[CouleurPrimaire]
        (current page.north west) rectangle
        ([yshift=-3.6cm] current page.north east);
    \fill[CouleurAccent]
        ([yshift=-3.6cm] current page.north west) rectangle
        ([yshift=-4.0cm] current page.north east);
    \fill[CouleurPrimaire]
        (current page.south west) rectangle
        ([yshift=2.5cm] current page.south east);
    \fill[CouleurAccent]
        ([yshift=2.5cm] current page.south west) rectangle
        ([yshift=2.9cm] current page.south east);
    % En-tête institutionnel, posé dans le bandeau supérieur
    \node[align=center, text=white] at ([yshift=-1.8cm] current page.north) {%
        {\fontsize{11}{13}\selectfont
        \textsc{République du Cameroun}}\\[0.3ex]
        {\color{CouleurAccent!60!white}\rule{0.4\paperwidth}{0.4pt}}\\[0.6ex]
        {\fontsize{10}{12}\selectfont\textbf{\institutionprimaire}}\\[0.2ex]
        {\fontsize{8}{10}\selectfont \tutelle}\\[0.2ex]
        {\fontsize{9}{11}\selectfont \unitepedago}};
    % Logos : enspy_logo.jpg à gauche, unesco_logo.jpeg à droite
    \node[fill=white, rounded corners=4pt, inner sep=2pt]
        at ([xshift=1.75cm, yshift=-1.8cm] current page.north west)
        {\includegraphics[height=2.6cm]{figures/logos/enspy_logo.jpg}};
    \node[fill=white, rounded corners=4pt, inner sep=2pt]
        at ([xshift=-1.75cm, yshift=-1.8cm] current page.north east)
        {\includegraphics[height=2.6cm]{figures/logos/unesco_logo.jpeg}};
    % Année et version, posées dans le bandeau inférieur
    \node[align=center] at ([yshift=1.25cm] current page.south) {%
        {\color{white}\bfseries\large Année académique \anneescolaire}\\[0.3ex]
        {\color{CouleurAccent!70!white}\small
        Version 1.0 | Document pédagogique de l'UE \codeUE}};
\end{tikzpicture}

\vspace*{0.6cm}

% --- Titre principal (sur deux lignes) ---
\begin{center}
    {\color{CouleurPrimaire}\rule{0.85\textwidth}{1.2pt}}\\[1.2ex]
    {\fontsize{26}{32}\selectfont\bfseries\color{CouleurPrimaire}
        Normes, Gestion des Risques\\[0.4ex]
        et Audit de Cybersécurité}\\[0.8ex]
    {\color{CouleurAccent}\rule{0.85\textwidth}{2pt}}\\[0.5ex]
    {\color{CouleurPrimaire}\rule{0.85\textwidth}{0.5pt}}
\end{center}

\vspace{0.8cm}

\begin{center}
    {\fontsize{13}{16}\selectfont\itshape\color{CouleurSecondaire}
        Du référentiel à la preuve~:\\
        piloter, mesurer et auditer la sécurité}
\end{center}

\vspace{0.5cm}

% --- Cartouche sigle et code UE ---
\begin{center}
    \begin{tikzpicture}
        \node[draw=CouleurPrimaire, fill=CouleurPrimaire, rounded corners=6pt,
              inner xsep=14pt, inner ysep=6pt]{%
            \color{white}\bfseries\large \sigleCours{} - UE~\codeUE
        };
    \end{tikzpicture}
\end{center}

\vspace{1.0cm}

% --- Schéma de la discipline : normes, risques, contrôles, audit, amélioration,
%     lus par les trois mesures µ, U, H ---
\begin{center}
\begin{tikzpicture}[
    etape/.style = {draw=CouleurBordDef, fill=CouleurDef, rounded corners=3pt,
        minimum width=2.5cm, minimum height=0.8cm, align=center,
        font=\small\bfseries\color{CouleurPrimaire}, inner sep=2pt},
    arrow/.style = {-{Stealth[length=4pt]}, thick, color=CouleurSecondaire}
]
\node[etape] (n1) at (0, 0)      {Normes};
\node[etape] (n2) at (3.2, 0)    {Risques};
\node[etape] (n3) at (6.4, 0)    {Contrôles};
\node[etape] (n4) at (6.4, -1.4) {Audit};
\node[etape] (n5) at (3.2, -1.4) {Plan d'action};
\node[etape] (n6) at (0, -1.4)   {Amélioration};
\draw[arrow] (n1) -- (n2); \draw[arrow] (n2) -- (n3); \draw[arrow] (n3) -- (n4);
\draw[arrow] (n4) -- (n5); \draw[arrow] (n5) -- (n6); \draw[arrow] (n6) -- (n1);
\node[draw=CouleurPrimaire, fill=CouleurPrimaire, rounded corners=4pt,
      align=center, font=\footnotesize\bfseries\color{white}, inner xsep=8pt]
      at (3.2, -2.8)
      {\OUI{} / \NON{} / \NEK{} \quad$\longrightarrow$\quad
       $\mu$ (conformité) · $U$ (incertitude) · $H$ (entropie)};
\end{tikzpicture}
\end{center}

\vspace{0.6cm}

\begin{center}
    {\large\bfseries\color{CouleurPrimaire}
        \auteurcours}\\[0.5ex]
    {\normalsize\color{CouleurSecondaire}
       Enseignant Chercheur | \laboratoire{} | ENSPY}\\[0.2ex]
    {\small\color{CouleurBordInfo}
        \institutionprimaire}\\[0.2ex]
    {\small\color{CouleurBordInfo}
        \href{mailto:maletyon@proton.me}{maletyon@proton.me}}
\end{center}

\end{titlepage}

    \cleardoublepage

    % ============================================================================
% Fichier  : front/01_avant_propos.tex — Avant-propos
% Version  : 1.0 -- novembre 2026 (proposition à relire par l'auteur)
% ============================================================================

\chapter*{Avant-propos}
\addcontentsline{toc}{chapter}{Avant-propos}
\markboth{Avant-propos}{Avant-propos}

Une banque peut être certifiée et compromise la même année. Ce paradoxe
apparent, qui ouvre ce livre, n'a rien d'exceptionnel~: il naît chaque fois
qu'une organisation confond ce qu'elle a fait vérifier avec ce qu'elle est
réellement, et ce qu'elle a regardé avec ce qui existe. Les normes, les
méthodes de gestion des risques et l'audit sont les outils dont on dispose
pour réduire cet écart~; ils peuvent aussi le masquer, quand on les emploie
pour rassurer plutôt que pour savoir.

Ce cours s'adresse aux étudiants du master professionnel en management des
systèmes numériques, spécialité sécurité des systèmes numériques, de la
Chaire UNESCO d'Afrique centrale, et aux praticiens qui exercent ou
exerceront les métiers de responsable de la sécurité des systèmes
d'information, d'auditeur en cybersécurité, de gestionnaire des risques
cyber, de responsable GRC ou de consultant en mise en œuvre de systèmes de
management. Il suppose les bases de la cybersécurité~; il ne suppose aucune
connaissance préalable des normes, de la gestion des risques ni de l'audit.

Il a trois partis pris. Le premier est de relier, plutôt que de juxtaposer,
les normes, le risque et l'audit~: une même organisation les vit ensemble,
et les échecs naissent souvent entre eux. Le deuxième est d'ancrer chaque
notion dans le contexte camerounais et africain~: le droit de la
cybersécurité, la régulation bancaire de la CEMAC, les organismes publics
et leurs contrôleurs, au côté des normes internationales. Le troisième est
de mesurer ce que l'on ne sait pas~: chaque cote, chaque conclusion, chaque
chiffre du livre dit sur quoi il repose, et l'instrument \OUI{}~/ \NON{}~/
\NEK{}, issu du programme de recherche de l'auditologie, donne à cette
exigence une forme mesurable.

Les trois organisations qui accompagnent le lecteur, la BANQUE SANAGA,
LiZbiZ et l'ANPP, sont fictives, comme leurs dirigeants, leurs prestataires,
les dates et les montants. Les mécanismes, les normes et les textes cités
sont exacts, avec, pour chaque texte, le degré de vérification indiqué à
l'annexe~\ref{ann:textes}. Les normes sous licence sont paraphrasées, jamais
citées.

Le livre est accompagné d'un livret des corrigés, qui explique et annote
chaque exercice, et de NDAP SUKULU, sur lequel les fiches d'auto-évaluation
et les ateliers du cours permettent de pratiquer chaque chapitre.

\bigskip
\hfill\begin{minipage}{0.5\textwidth}
\raggedleft
\IfFileExists{../../Logos/signature_avant_propos.png}{\includegraphics[height=1.6cm]{../../Logos/signature_avant_propos.png}}{\textit{MINKA MI NGUIDJOÏ Thierry Emmanuel}}\\
Yaoundé, novembre 2026
\end{minipage}

    \cleardoublepage

    % ============================================================================
% Fichier  : front/02_mode_emploi.tex — Mode d'emploi
% Version  : 1.0 -- novembre 2026
% ============================================================================
% Règle des diapositives (GCI, ACSI) : 30 h en présentiel = 15 séances de
% 2 h ; fiches et ateliers hors des 30 h. Répartition des séances indicative.
% ============================================================================

\chapter*{Mode d'emploi du cours}
\addcontentsline{toc}{chapter}{Mode d'emploi du cours}
\markboth{Mode d'emploi}{Mode d'emploi}

\section*{Structure générale}

Le cours compte un prologue et six parties, soit 16~chapitres, pour
30~heures en présentiel réparties en quinze séances de deux heures, qui
couvrent l'ensemble des parties. Les fiches de questions-réponses par
chapitre et les dix laboratoires de 60~minutes de NDAP SUKULU sont des
travaux à faire à la maison, hors de ces 30~heures. Les corrigés des
exercices sont réunis dans un livret séparé.

\begin{tableaucours}{cL{6.4cm}cL{3.2cm}}{%
    \tableauHeader{Partie} &
    \tableauHeader{Thème} &
    \tableauHeader{Chap.} &
    \tableauHeader{Niveau}
}
0 & Prologue~: certifiée, et pourtant compromise & P & introductif \\
I & Fondements~: gouvernance, risque et incertitude & 1--2 & fondamental \\
II & Méthodes~: gérer les risques, conduire l'audit & 3--7 & fondamental à
intermédiaire \\
III & Cadre normatif et juridique & 8--10 & fondamental \\
IV & Pratique~: SMSI, audit technique et GRC & 11--14 & intermédiaire à
expert \\
V & Défis avancés & 15 & expert \\
VI & Recherche~: mesurer ce que l'on ne sait pas & 16 & recherche \\
\end{tableaucours}

\section*{Les quinze séances}

\begin{tableaucours}{cL{6.8cm}L{4.2cm}}{%
    \tableauHeader{Séance} &
    \tableauHeader{Contenu} &
    \tableauHeader{Travail à la maison}
}
1 & Prologue~: le certificat, le périmètre, l'inconnu & Fiche P \\
2 & Chap.~1~: gouverner la cybersécurité & Fiche 1 \\
3 & Chap.~2~: le risque cyber et sa cote & Fiche 2, atelier TD01 \\
4 & Chap.~3~: ISO~31000 et ISO/IEC~27005 & Fiche 3 \\
5 & Chap.~4~: EBIOS Risk Manager & Fiche 4, ateliers TD02 et TD03 \\
6 & Chap.~5~: quantifier avec FAIR & Fiche 5, atelier TD04 \\
7 & Chap.~6~: conduire un audit & Fiche 6 \\
8 & Chap.~7~: décomposer, coter, conclure & Fiche 7, atelier TD05 \\
9 & Chap.~8~: la famille ISO/IEC~27000 & Fiche 8, atelier TD06 \\
10 & Chap.~9~: référentiels et correspondances & Fiche 9, atelier TD07 \\
11 & Chap.~10~: droit et déontologie & Fiche 10 \\
12 & Chap.~11~: mettre en œuvre et auditer un SMSI & Fiche 11, atelier TD08 \\
13 & Chap.~12~: l'audit technique & Fiche 12, atelier TD09 \\
14 & Chap.~13 et~14~: la GRC en pratique, les cas intégrateurs & Fiches 13
et 14, atelier TD10 \\
15 & Chap.~15 et~16~: risques émergents, mesurer l'inconnu & Fiches 15 et
16 \\
\end{tableaucours}

Cette répartition est indicative~; l'enseignant l'adapte au calendrier et au
rythme du groupe.

\section*{Les dix laboratoires de NDAP SUKULU}

Les dix laboratoires, NRAC-LAB-01 à NRAC-LAB-10, sont présentés dans la
section « NDAP SUKULU, la salle de classe compagnon », qui suit ce mode
d'emploi.

\section*{Le gabarit des chapitres}

Chaque chapitre suit le même plan~: une épigraphe, le niveau et les
prérequis, les objectifs, une situation d'ouverture tirée de l'un des trois
cas, cinq à huit sections, ce qu'il faut retenir, des points de contrôle,
trois à quatre exercices gradués, des lectures, et une transition.

\section*{Les encadrés}

\begin{tableaucours}{L{3.6cm}L{8.6cm}}{%
    \tableauHeader{Encadré} &
    \tableauHeader{Rôle}
}
Objectifs & Ce que l'on doit savoir faire à la fin du chapitre \\
Définition & Les notions, en gras dans le texte et indexées \\
Violation typique & Une erreur fréquente, et la règle qui l'évite \\
Information & Un complément, un renvoi, une nuance \\
Opération (SANAGA, LIZBIZ, ANPP) & Le cas~: les constatations, puis
l'analyse, avec un niveau de confiance \\
Lecture & Un calcul de $\mu$, $U$ et $H$, et sa lecture \\
Fiche de risque & Un risque décrit, coté et décidé \\
Papier de travail & Ce que le dossier doit contenir pour être refait par un
autre \\
Métier & Ce que la notion change pour chacun des cinq métiers du cours \\
Texte juridique & Un article cité dans son texte officiel, et sa portée \\
Ce qu'il faut retenir & La synthèse du chapitre \\
Exercice & Trois niveaux de difficulté~: \niveaubase, \niveaumoyen,
\niveauexpert \\
\end{tableaucours}

\section*{Les conventions}

\begin{itemize}
    \item \textbf{Les cotes}~: \OUI{} (établi par une preuve), \NON{}
          (absent ou contredit par une preuve), \NEK{} (connaissance
          insuffisante). Aucun barème ne donne de demi-point à \NEK{}.
    \item \textbf{Les mesures}~: le taux de conformité $\mu$, le taux
          d'incertitude $U$ et l'entropie $H$, en bits, se lisent toujours
          ensemble~; règle de décision~: $U$ d'abord, $\mu$ ensuite, $H$ pour
          contrôler (Chapitre~\ref{chap:conclure}).
    \item \textbf{L'échelle de cotation} des risques est celle du
          Chapitre~\ref{chap:risque} (quatre niveaux de vraisemblance et
          d'impact, criticité de 1 à 16).
    \item \textbf{Constatation et analyse} sont toujours séparées~; l'analyse
          porte un niveau de confiance.
    \item \textbf{Les textes}~: seuls les articles lus dans leur version
          officielle sont cités entre guillemets~; les normes sous licence
          sont paraphrasées~; le degré de vérification figure à
          l'annexe~\ref{ann:textes}.
    \item \textbf{Les travaux de recherche} sont cités avec leur statut~:
          prépublication déposée, ou travail en préparation
          (Chapitre~\ref{chap:auditologie}).
    \item \textbf{Les cas} sont fictifs~: noms en majuscules, montants en
          francs CFA, villes réelles.
\end{itemize}

    \cleardoublepage

    % ==============================================================================
% Fichier : front/03_ndap_sukulu.tex
% Titre   : NDAP SUKULU, l'environnement compagnon du cours
% Auteur  : MINKA MI NGUIDJOÏ Thierry Emmanuel
% Version : 1.0 -- novembre 2026
% ==============================================================================
% Texte commun aux cours de l'auteur ; seul le tableau des labs est propre au
% cours (fichier front/03b_labs.tex).
% ==============================================================================

\chapter*{NDAP SUKULU, la salle de classe compagnon}
\addcontentsline{toc}{chapter}{NDAP SUKULU, la salle de classe compagnon}
\markboth{NDAP SUKULU}{NDAP SUKULU}

\begin{center}
    \includegraphics[width=0.62\textwidth]{figures/logos/logo_ndap_sukulu.png}
\end{center}

\section*{Un nom venu du pays bassa}

\begin{quote}\itshape
Le soir tombe sur Éséka. Un enfant revient des champs avec son grand-père et
lui demande pourquoi il doit marcher si loin, chaque matin, pour aller à
l'école. Le vieil homme s'arrête au bord du chemin et répond, en bassa~:
« \textup{Ndap sukulu}, la salle de classe, ce ne sont pas les murs. Ce sont
les murs qui viennent après. La salle de classe, c'est l'endroit où quelqu'un
t'attend pour t'apprendre ce qu'il sait. » L'enfant insiste~: « Et qui
m'attend~? » Le grand-père sourit~: « \textup{Malet yon}. Ton enseignant. »
\end{quote}

En bassa, la langue paternelle de l'auteur, \textbf{\textit{ndap sukulu}}
désigne la salle de classe et \textbf{\textit{malet yon}} signifie « ton
enseignant ». L'environnement compagnon de ce cours porte ce nom parce qu'il
prolonge la salle de classe au-delà des heures de cours~: il attend
l'étudiant le soir, le week-end, à distance, avec des questions, des pièces à
examiner et une correction immédiate. Il est édité par \textbf{MaletYon},
« ton enseignant », et mis à la disposition des étudiants et des praticiens
qui souhaitent se perfectionner (\href{mailto:maletyon@proton.me}{maletyon@proton.me}).

\section*{Ce que l'on y trouve}

\begin{itemize}
    \item \textbf{Une fiche d'auto-évaluation par chapitre.} Chaque fiche
          puise vingt questions dans une banque de quarante (QCM, vrai ou
          faux, textes à trous). Elle dure trente minutes~; la copie validée
          est verrouillée et notée aussitôt. Le seuil de réussite est de
          60~\%.
    \item \textbf{Des laboratoires pratiques.} Chaque lab place l'étudiant
          devant les pièces d'une affaire du fil rouge (tableurs, journaux,
          contrats, procès-verbaux) dans un bureau d'analyste virtuel,
          ouvert dans le navigateur. Il dure soixante minutes, se termine
          par des points de contrôle notés automatiquement et renvoie aux
          sections du livre qu'il mobilise.
    \item \textbf{Un suivi pour l'enseignant}~: rapport de classe, notes par
          fiche et par lab, journal d'activité.
\end{itemize}

Les fiches et les labs sont des travaux à faire hors des heures de cours.
Ils ne remplacent ni le livre ni les exercices de fin de chapitre~: ils en
vérifient la maîtrise.

\section*{Comment y travailler}

\begin{enumerate}
    \item Se connecter au portail NDAP SUKULU avec ses identifiants et le
          code de cours \textbf{\codeUE}.
    \item Pour un lab, commencer par vérifier l'empreinte SHA-256 des pièces
          (\texttt{sha256sum -c MANIFESTE\_SHA256.txt}), puis travailler sur
          la copie du dossier \texttt{travail}~: les pièces originales sont
          en lecture seule et l'activité est journalisée, comme dans une
          vraie mission.
    \item Lire le rappel de cours affiché par le lab et relire les sections
          citées avant de répondre.
\end{enumerate}

\begin{boxinfo}
\textbf{Si le portail n'est pas encore ouvert.} En attendant le déploiement
de NDAP SUKULU, les fiches de questions sont fournies au format texte dans le
dossier \texttt{fiches} du cours, prêtes à être distribuées. Les corrections
des fiches figurent dans le livret de correction, à la suite des corrigés des
exercices de chaque chapitre.
\end{boxinfo}

\section*{Les laboratoires du cours}

% Tableau des labs du cours NRAC dans NDAP SUKULU (propre au cours).
Le cours compte dix laboratoires, disponibles sous le code \codeUE. Le dernier
rassemble le dispositif de GRC de la BANQUE SANAGA et sert de synthèse.

\begin{tableaucours}{lL{6.2cm}L{1.6cm}L{2.0cm}}{%
    \tableauHeader{Lab} &
    \tableauHeader{Sujet} &
    \tableauHeader{Chapitre} &
    \tableauHeader{Opération}
}
NRAC-LAB-01 & Coter le registre des risques et sa base de preuve & 2 & SANAGA \\
NRAC-LAB-02 & EBIOS Risk Manager~: cadrage et sources de risque & 4 & LIZBIZ \\
NRAC-LAB-03 & EBIOS Risk Manager~: écosystème et scénario opérationnel & 4 & LIZBIZ \\
NRAC-LAB-04 & Quantifier un risque cyber avec FAIR & 5 & SANAGA \\
NRAC-LAB-05 & Décomposer, coter, conclure & 7 & SANAGA \\
NRAC-LAB-06 & Confronter une déclaration d'applicabilité aux preuves & 8 & SANAGA \\
NRAC-LAB-07 & Une base commune, plusieurs référentiels & 9 & ANPP \\
NRAC-LAB-08 & Auditer un SMSI & 11 & SANAGA \\
NRAC-LAB-09 & Audit technique sur exports figés & 12 & LIZBIZ \\
NRAC-LAB-10 & Synthèse GRC~: obligations, matrice, indicateurs & 13 & SANAGA \\
\end{tableaucours}
\vspace{-0.8em}
\begin{center}{\small\itshape Les laboratoires du cours dans NDAP SUKULU}\end{center}


    \cleardoublepage

    \tableofcontents
    \cleardoublepage

    \listoffigures
    \listoftables
    \cleardoublepage

% CORPS (pages arabes)
% Numérotation des chapitres CONTINUE sur tout le cours : le prologue est le
% chapitre 0, puis C01 à C16, un numéro par fichier.
% --------------------------------------------------------------------------
\mainmatter

    % ======================================================================
    % PARTIE 0 — PROLOGUE OPÉRATIONNEL
    % Standards de référence : Un cas qui échoue : un SMSI certifié, compromis hors périmètre.
    % Niveau : introductif
    % ======================================================================
    \part*{Prologue Opérationnel}
    \addcontentsline{toc}{part}{Prologue Opérationnel}

        \setcounter{chapter}{-1}
        % ============================================================================
% Fichier  : partie0/P00_certifiee_compromise.tex
% Titre    : Certifiée, et pourtant Compromise
% Auteur   : MINKA MI NGUIDJOÏ Thierry Emmanuel
% Version  : 1.0 -- novembre 2026
% Corrigés : corriges/Corr_P00.tex
% ============================================================================
% Faits du cas (bible des cas, Claude outputs/02_bible_des_cas.md) :
% - corpus : BANQUE SANAGA, Douala, 38 agences, DG M. FOTSO, COBAC ; bascule
%   ratée du système central en mars 2026 (GCI P00) ;
% - propres à SEC08 : certificat ISO/IEC 27001:2022 de juin 2025 (CERTIQUAL
%   AFRIQUE, organisme fictif) sur le seul centre de données principal de
%   Douala ; rançongiciel du 14/09/2026 entré par l'agence de Ngaoundéré,
%   180 postes dans 9 agences, 3 jours, rançon de 90 M FCFA non payée, coût
%   410 M FCFA ; Mme ABENA, RSSI.
% ISO/IEC 27001:2022, article 4.3 : paraphrasé (texte sous licence), jamais
% cité.
% ============================================================================

\chapter[Certifiée, et pourtant Compromise]{Prologue : Certifiée, et pourtant Compromise}
\label{chap:prologue}
\chaptermark{Prologue}

\epigraph{%
    « Un certificat dit ce qui a été audité.\\
    Il ne dit rien de ce qui ne l'a pas été. »%
}{Principe du cours}

\niveauChapitre{Introductif}{%
    Aucun prérequis de gestion des risques ni d'audit. Les notions de base de
    la cybersécurité (MSN-1S1-SEC01) suffisent.%
}

\begin{boxobjectifs}
\begin{itemize}
    \item Reconstituer, à partir de ses traces, la compromission d'une banque
          certifiée ISO/IEC~27001.
    \item Distinguer un certificat, la conformité à une norme et la sécurité
          effective.
    \item Lire le périmètre d'un système de management de la sécurité et
          repérer un périmètre de complaisance.
    \item Reconnaître qu'un risque jamais apprécié n'est pas un risque faible,
          mais une inconnue.
    \item Lire un ensemble de contrôles par trois mesures~: $\mu$, $U$ et $H$.
    \item Énoncer les cinq réflexes du cours et situer ses parties, ses
          ateliers et ses cas.
\end{itemize}
\end{boxobjectifs}

\bigskip

% ============================================================================
\section*{Avant de commencer : un certificat au mur, des guichets fermés}
% ============================================================================

Douala, jeudi 17~septembre 2026. Dans le hall du siège de la BANQUE
SANAGA\footnote{La BANQUE SANAGA, ses dirigeants, ses prestataires,
l'organisme certificateur, les dates et les montants de ce cas sont fictifs.
Les mécanismes d'attaque, les normes et les textes cités sont exacts.}, un
cadre doré présente le certificat ISO/IEC~27001 obtenu quinze mois plus tôt.
Dans les étages, la cellule de crise en est à sa troisième nuit. Neuf
agences sur trente-huit ont gardé leurs guichets fermés pendant trois jours~;
180 postes de travail ont été chiffrés par un rançongiciel~; les attaquants
réclament 90~millions de francs CFA.

Le directeur général, M.~FOTSO, pose la question que tout le monde se pose~:
« Nous sommes certifiés. Comment est-ce possible~? » La responsable de la
sécurité des systèmes d'information, Mme~ABENA, répond par une seule
phrase~: « Le certificat porte sur le centre de données. L'attaque est
passée par une agence. »

\medskip

Ce prologue reconstitue l'incident, montre pourquoi il ne contredit pas le
certificat, et en tire les réflexes qui structurent tout le cours. Il ne
suppose aucune connaissance préalable des normes ni de l'audit~: les notions
sont nommées ici et définies dans les chapitres qui suivent.

\begin{boxalert}
\textbf{Violation typique.} Présenter un certificat ISO/IEC~27001 comme la
preuve que « la banque est sécurisée ».

\smallskip
\textbf{Règle~:} un certificat atteste qu'un système de management, sur un
périmètre donné, satisfait les exigences d'une norme à la date de l'audit.
On lit toujours le périmètre avant de lire le certificat.
\end{boxalert}

\separateur

% ============================================================================
\section{Ce qui s'est passé}
\label{sec:pro-faits}
% ============================================================================

\subsection{La chronologie}

La RSSI a reconstitué l'incident à partir des journaux des postes, de la
messagerie et du partage de fichiers, puis l'a présenté au comité des
risques. Chaque ligne du Tableau~P.1 renvoie à une trace.

\begin{tableaucours}{L{2.8cm}L{6.4cm}L{3.0cm}}{%
    \tableauHeader{Date} &
    \tableauHeader{Événement} &
    \tableauHeader{Trace}
}
Juin 2025 & Certificat ISO/IEC~27001:2022 délivré par CERTIQUAL AFRIQUE pour
trois ans. Périmètre~: le centre de données principal de Douala &
Certificat et rapport d'audit de certification \\
Novembre 2025 & Revue annuelle des risques~: aucun scénario ne porte sur le
réseau des agences, hors périmètre & Registre des risques 2025 \\
Lundi 14 sept. 2026, 10~h~12 & Un agent de l'agence de Ngaoundéré ouvre la
pièce jointe d'un faux avis de virement & Journal de la messagerie \\
14 sept., 10~h~14 & Un programme malveillant s'installe sur le poste~; il
récupère les identifiants enregistrés dans le navigateur & Journal du
poste \\
14 -- 15 sept. & Avec ces identifiants, l'attaquant accède au partage de
fichiers commun aux agences et s'y déplace & Journal du serveur de
fichiers \\
Mardi 15 sept., 2~h~40 & Chiffrement simultané de 180 postes dans 9
agences~; message de rançon & Journaux des postes \\
15 -- 17 sept. & Guichets fermés dans les 9 agences~; reconstruction des
postes à partir d'images propres & Rapport de la cellule de crise \\
18 sept. & Réouverture des agences~; la rançon n'est pas payée & Décision
du directeur général \\
\end{tableaucours}
\vspace{-0.8em}
\begin{center}{\small\itshape Tableau~P.1. Chronologie de l'incident
(reconstitution de la RSSI)}\end{center}

Le centre de données certifié n'a pas été atteint~: il n'était pas relié au
partage de fichiers des agences. Les comptes, les opérations et le système
central sont restés intacts. Le coût de l'incident, chiffré par la
direction financière, atteint pourtant \textbf{410~millions de francs CFA}~:
arrêt des guichets, reconstruction des postes, prestataire de réponse à
incident, communication auprès de la clientèle.

\subsection{Constater, puis analyser}

Comme dans toute investigation, on sépare ce qui est établi de ce que l'on
en déduit.

\begin{boxscenario}{Opération SANAGA. Le rançongiciel}
\textbf{Constatations.}
\begin{itemize}
    \item Le certificat couvre le seul centre de données principal de
          Douala~; les 38 agences, la messagerie et les postes de travail sont
          hors périmètre.
    \item Le registre des risques 2025 ne contient aucun scénario sur les
          agences.
    \item L'attaque est entrée par un courriel ouvert en agence, puis s'est
          propagée par un partage de fichiers commun aux agences.
    \item Les identifiants volés n'étaient protégés par aucun second facteur
          d'authentification.
\end{itemize}

\textbf{Analyse} (confiance élevée). Le certificat n'est pas faux~: le
système de management certifié couvre bien le centre de données, et celui-ci
n'a pas été touché. Ce qui a échoué, c'est le choix du périmètre et
l'appréciation des risques, qui s'arrêtait à ses frontières. La banque a
protégé ce qu'elle avait décidé d'auditer, et laissé dans l'ombre ce qu'elle
n'avait pas regardé.
\end{boxscenario}

% ============================================================================
\section{Certificat, conformité, sécurité}
\label{sec:pro-certificat}
% ============================================================================

\subsection{Trois choses différentes}

\begin{boxdef}
\termecle{Système de management de la sécurité de
l'information}\index{SMSI} (\sigle{SMSI})~: ensemble de politiques, de
processus, de rôles et de contrôles par lequel une organisation pilote la
sécurité de son information sur un périmètre défini, selon un cycle
d'amélioration continue. L'ISO/IEC~27001 en fixe les
exigences~\cite{iso27001_2022}.

\medskip
\termecle{Certificat}\index{certificat}~: attestation délivrée par un
organisme tiers accrédité, après un audit, qu'un système de management
satisfait les exigences d'une norme, sur un périmètre et pour une durée
donnés.

\medskip
\termecle{Conformité}\index{conformité}~: satisfaction d'une exigence
précise (d'une norme, d'une loi, d'un contrat), établie par une preuve.

\medskip
\termecle{Sécurité effective}\index{sécurité effective}~: capacité réelle de
l'organisation à résister aux attaques qu'elle subit, et à s'en relever.
\end{boxdef}

\begin{tableaucours}{L{2.6cm}L{4.4cm}L{5.2cm}}{%
    \tableauHeader{Notion} &
    \tableauHeader{Ce qu'elle établit} &
    \tableauHeader{Ce qu'elle n'établit pas}
}
Certificat & Un SMSI conforme à la norme sur le périmètre, à la date de
l'audit & La sécurité hors périmètre~; la sécurité après l'audit \\
Conformité & Qu'une exigence précise est satisfaite & Que l'exigence suffit à
couvrir le risque \\
Sécurité effective & Que les attaques réelles échouent ou sont contenues &
Rien sans mesure~: elle se démontre par des tests et des incidents \\
\end{tableaucours}
\vspace{-0.8em}
\begin{center}{\small\itshape Tableau~P.2. Trois notions à ne pas
confondre}\end{center}

\subsection{Le périmètre, première décision de sécurité}

L'ISO/IEC~27001 demande à l'organisation de fixer les limites de son SMSI en
tenant compte de son contexte, des attentes des parties intéressées, et des
interfaces et dépendances entre ses propres activités et celles des
autres~\cite{iso27001_2022}. Le périmètre n'est donc pas libre~: il doit être
défendable. Celui de la BANQUE SANAGA avait été choisi pour obtenir
rapidement le certificat, en excluant ce qui était difficile à maîtriser.

\begin{boxdef}
\termecle{Périmètre de complaisance}\index{périmètre de complaisance}~:
périmètre de SMSI choisi pour faciliter la certification plutôt que pour
couvrir les risques de l'organisation, au point d'exclure des activités ou
des actifs dont dépend le périmètre retenu.
\end{boxdef}

Les agences de la banque n'étaient pas un sujet à part~: elles reçoivent les
clients, saisissent les opérations et sont reliées au système central. Un
attaquant installé en agence se trouve à quelques pas de ce que le
certificat protège. Le Chapitre~\ref{chap:smsi} montrera comment construire
un périmètre défendable, et comment l'auditeur le conteste.

\begin{boxinfo}
\textbf{Et le régulateur~?} Le règlement COBAC R-2016/04 impose aux
établissements de crédit un dispositif de contrôle interne qui couvre
l'ensemble de leurs risques, y compris les risques opérationnels liés aux
systèmes d'information~\cite{cobac2016r04}. Un certificat sur une partie du
système ne dispense pas la banque de maîtriser l'autre. Le
Chapitre~\ref{chap:droit} présente ces obligations.
\end{boxinfo}

% ============================================================================
\section{Un Risque jamais Apprécié n'est pas un Risque Faible}
\label{sec:pro-risque}
% ============================================================================

\subsection{Le trou dans le registre}

Le registre des risques de 2025 recensait des scénarios sur le centre de
données~: panne électrique, défaillance du prestataire d'hébergement,
intrusion par Internet. Aucun ne portait sur les agences. Lors de la revue,
le comité des risques a lu le registre et conclu qu'aucun risque
intolérable n'était ouvert.

Le raisonnement est faux, et c'est l'erreur la plus répandue en gestion des
risques~: un risque absent du registre a été lu comme un risque nul. Il
n'était pas nul~; il était \emph{inconnu}. La bonne cote n'était pas
« faible », mais « non évalué », et cette cote aurait dû déclencher une
appréciation.

\begin{boxalert}
\textbf{Violation typique.} Lire l'absence d'un scénario dans le registre
comme la preuve que le risque est maîtrisé.

\smallskip
\textbf{Règle~:} ce qui n'a pas été apprécié est coté \NEK{} (connaissance
insuffisante), pas « faible ». Une cote de vraisemblance sans preuve est une
opinion~; elle se signale comme telle.
\end{boxalert}

\subsection{Une fiche de risque}

Voici ce qu'aurait dû contenir le registre. Le Chapitre~\ref{chap:risque}
en définira chaque rubrique, et le Chapitre~\ref{chap:ebios} montrera comment
construire le scénario avec la méthode EBIOS Risk Manager.

\begin{boxobjet}{R-AG-01 · Rançongiciel en agence}
\textbf{Actif essentiel~:} le service au guichet des agences (versements,
retraits, ouvertures de compte).

\textbf{Actifs supports~:} postes de travail des agences, messagerie,
partage de fichiers commun, comptes des agents.

\textbf{Source de risque~:} groupe cybercriminel motivé par le gain.

\textbf{Événement redouté~:} indisponibilité des guichets pendant plusieurs
jours dans plusieurs agences.

\textbf{Scénario~:} courriel piégé ouvert en agence, vol d'identifiants,
déplacement par le partage commun, chiffrement des postes.

\textbf{Vraisemblance~:} \NEK{} en 2025 (scénario jamais apprécié)~; après
l'incident, élevée.

\textbf{Impact~:} perte d'exploitation, coûts de reconstruction, atteinte à la
réputation, compte rendu au régulateur.

\textbf{Mesures existantes~:} antivirus sur les postes~; aucun second facteur~;
partage de fichiers ouvert à toutes les agences~; sauvegardes des postes non
testées.

\textbf{Risque résiduel~:} intolérable, tant que le partage n'est pas
cloisonné et que l'authentification n'est pas renforcée.
\end{boxobjet}

% ============================================================================
\section{Lire la Sécurité par Trois Mesures}
\label{sec:pro-mesures}
% ============================================================================

\subsection{Deux colonnes, un certificat}

Après l'incident, la direction de l'audit interne, dirigée par
Mme~EKAMBI, évalue dix contrôles clés de sécurité de deux façons~: sur le
centre de données certifié, puis sur le réseau des agences. Chaque contrôle
est coté \OUI{} (en place et prouvé), \NON{} (absent ou défaillant) ou
\NEK{} (impossible à établir faute de preuve).

\begin{tableaucours}{L{6.6cm}cc}{%
    \tableauHeader{Contrôle} &
    \tableauHeader{Centre de données} &
    \tableauHeader{Agences}
}
Second facteur d'authentification pour les accès à distance & \OUI{} & \NON{} \\
Correctifs de sécurité appliqués dans le délai fixé & \OUI{} & \NEK{} \\
Sauvegardes hors ligne testées & \OUI{} & \NON{} \\
Cloisonnement du réseau & \OUI{} & \NON{} \\
Comptes d'administration nominatifs et revus & \OUI{} & \NEK{} \\
Journalisation centralisée et revue & \OUI{} & \NEK{} \\
Filtrage des pièces jointes de la messagerie & \OUI{} & \OUI{} \\
Sensibilisation des agents dans l'année & \NEK{} & \NEK{} \\
Plan de réponse aux incidents testé & \NON{} & \NON{} \\
Protection des postes (antivirus) & \OUI{} & \OUI{} \\
\textbf{Total (\OUI{} / \NON{} / \NEK{})} & \textbf{8 / 1 / 1} & \textbf{2 / 4 / 4} \\
\end{tableaucours}
\vspace{-0.8em}
\begin{center}{\small\itshape Tableau~P.3. Dix contrôles clés, deux
périmètres}\end{center}

\subsection{Trois mesures}

Le cours lit un ensemble de cotes par trois mesures, définies au
Chapitre~\ref{chap:conclure}~:
\begin{itemize}
    \item le \termecle{taux de conformité}\index{taux de conformité}~$\mu$, la
          part des \OUI{}~;
    \item le \termecle{taux d'incertitude}\index{taux d'incertitude}~$U$, la
          part des \NEK{}~;
    \item l'\termecle{entropie}\index{entropie}~$H$, la dispersion des
          réponses, en bits, de 0 à $\log_2 3 \approx 1{,}585$~:
          $H = -\sum_m p(m)\log_2 p(m)$.
\end{itemize}

\begin{boxlecture}{Deux périmètres de la BANQUE SANAGA}
\textbf{Constatation.}
\begin{itemize}
    \item Centre de données (8/1/1)~: $\mu = 0{,}80$, $U = 0{,}10$,
          $H = -(0{,}8\log_2 0{,}8 + 2 \times 0{,}1\log_2 0{,}1) \approx 0{,}92$~bit.
    \item Agences (2/4/4)~: $\mu = 0{,}20$, $U = 0{,}40$,
          $H = -(0{,}2\log_2 0{,}2 + 2 \times 0{,}4\log_2 0{,}4) \approx 1{,}52$~bit.
    \item Ensemble (10/5/5)~: $\mu = 0{,}50$, $U = 0{,}25$, $H = 1{,}50$~bit.
\end{itemize}

\textbf{Analyse} (confiance élevée sur le calcul, moyenne sur la lecture).
Le certificat décrit la première ligne seulement~: un périmètre largement
conforme et bien connu. Les agences combinent une conformité faible et une
forte incertitude~: quatre contrôles sur dix ne peuvent pas même être
évalués, faute de trace. C'est l'endroit où l'attaque est entrée. Lue
d'un bloc, la banque paraît « à moitié conforme »~; lue par périmètre, elle
montre un îlot maîtrisé dans une zone inconnue.

\medskip
Une cote mérite enfin d'être relue~: le filtrage des pièces jointes est
coté \OUI{} en agence, alors que l'attaque est entrée par une pièce jointe.
Il n'y a pas de contradiction~: \OUI{} dit que le contrôle est en place et
prouvé, pas qu'il est infaillible. Un filtre arrête les pièces jointes
connues comme malveillantes~; celle du 14~septembre ne l'était pas encore.
C'est le troisième réflexe du cours~: une mesure se juge à son
fonctionnement face aux attaques réelles.

\medskip
\hfill\lecture{0{,}20}{0{,}40}{1{,}52}
\end{boxlecture}

\begin{boxinfo}
\textbf{Une règle à retenir dès maintenant.} On ne conclut jamais sur une
seule mesure. Un taux de conformité de 0,50 ne dit pas si la moitié du
système est solide et l'autre fragile, ou si l'on ignore l'état de la
moitié du système~: $U$ et $H$ le disent. Le Chapitre~\ref{chap:conclure}
en tirera la règle de décision~: $U$ d'abord, $\mu$ ensuite, $H$ pour
contrôler.
\end{boxinfo}

% ============================================================================
\section{Cinq Réflexes}
\label{sec:pro-reflexes}
% ============================================================================

L'incident se résume en cinq réflexes. Chacun fait l'objet d'au moins un
chapitre.

\begin{tableaucours}{cL{4.0cm}L{6.0cm}c}{%
    \tableauHeader{N°} &
    \tableauHeader{Réflexe} &
    \tableauHeader{Ce qu'il évite} &
    \tableauHeader{Chapitre}
}
1 & Lire le périmètre avant le certificat & Croire protégé ce qui n'a jamais
été audité & \ref{chap:smsi} \\
2 & Ce qui n'est pas apprécié est inconnu, pas faible & Un registre qui
rassure parce qu'il est incomplet & \ref{chap:risque} \\
3 & Une mesure se prouve en fonctionnement & Des contrôles qui existent sur
le papier & \ref{chap:audit} \\
4 & Mesurer l'incertitude & Une conformité moyenne qui cache des zones
inconnues & \ref{chap:conclure} \\
5 & Relier norme, risque et audit & Trois exercices séparés qui ne se
parlent pas & \ref{chap:grc} \\
\end{tableaucours}
\vspace{-0.8em}
\begin{center}{\small\itshape Tableau~P.4. Les cinq réflexes du
cours}\end{center}

\begin{boxprocedure}
\textbf{Le minimum d'une cote documentée.} Pour qu'une cote de contrôle ou
de risque puisse être défendue devant un comité, un auditeur ou le
régulateur, le papier de travail qui la porte indique au moins~:
\begin{itemize}
    \item le contrôle ou le scénario concerné, et son périmètre~;
    \item le critère (la norme, la politique ou la loi qui fixe l'attendu)~;
    \item la cote (\OUI{}, \NON{} ou \NEK{}) et la preuve qui l'appuie~;
    \item pour un \NEK{}, la raison pour laquelle on ne peut pas savoir~;
    \item la date, l'auteur et le réviseur.
\end{itemize}
\end{boxprocedure}

\subsection{Ce que l'incident dit à chaque métier}

\begin{boxmetier}{RSSI}
\textbf{Ce que Mme~ABENA retient}

Le périmètre du SMSI était une décision de sécurité, et non une formalité
de certification. Elle propose au comité d'étendre le SMSI aux agences en
deux ans, de cloisonner le partage de fichiers et d'imposer un second
facteur pour tous les accès. Elle sait que le prochain audit de
surveillance de l'organisme certificateur ne lui sera d'aucun secours sur
ces sujets, puisqu'ils sont hors périmètre.
\end{boxmetier}

\begin{boxmetier}{Auditeur en cybersécurité}
\textbf{Contester le périmètre}

L'auditeur ne se contente pas de vérifier ce qui est inclus~: il examine
ce qui est exclu, et la justification de l'exclusion. Un périmètre qui
exclut une activité dont dépend le périmètre retenu est une constatation
d'audit en soi (Chapitre~\ref{chap:smsi}).
\end{boxmetier}

\begin{boxmetier}{Gestionnaire des risques cyber}
\textbf{Apprécier hors des frontières}

L'appréciation des risques ne s'arrête pas au périmètre du certificat~: elle
porte sur les activités de l'organisation. Le gestionnaire des risques ajoute
le scénario R-AG-01, le chiffre (Chapitre~\ref{chap:fair}) et le présente
au comité avec ses options de traitement.
\end{boxmetier}

\begin{boxmetier}{Responsable GRC}
\textbf{Un seul dispositif}

Le registre des risques, la déclaration d'applicabilité du SMSI et le
programme d'audit doivent couvrir le même univers. Le responsable GRC
vérifie qu'aucun pan de l'activité n'échappe aux trois à la fois
(Chapitre~\ref{chap:grc-pratique}).
\end{boxmetier}

\begin{boxmetier}{Consultant en mise en œuvre de SMSI}
\textbf{Conseiller un périmètre défendable}

Le consultant qui a accompagné la certification a proposé le périmètre le
plus facile à certifier. Son rôle était aussi de dire ce que ce choix
laissait sans protection. Un périmètre étroit peut être une étape~; il doit
alors être présenté comme tel, avec un calendrier d'extension.
\end{boxmetier}

% ============================================================================
\section{Le Cours et son Fil Rouge}
\label{sec:pro-cours}
% ============================================================================

\subsection{Une discipline, une mesure}

Ce cours enseigne à appliquer les normes de sécurité, à gérer les risques
et à auditer la cybersécurité. Il les relie par un même fil rouge~: la
mesure de ce que l'on sait et de ce que l'on ignore. L'instrument de
réponse \OUI{}~/ \NON{}~/ \NEK{} et les mesures $\mu$, $U$ et $H$ viennent du
programme de recherche de l'auditologie~\cite{minka2026auditologie,minka2026nek}.
Le cours les applique à l'audit, comme les autres cours du corpus, et aussi
à l'appréciation des risques, où une vraisemblance cotée sans preuve est un
\NEK{} déguisé.

\begin{boxinfo}
\textbf{Ce que dit la recherche.} L'article \emph{Vers l'auditologie} pose
l'\emph{observabilité partielle} comme premier axiome~: l'auditeur ne voit
jamais tout le système, il en voit des traces. Le périmètre d'un certificat
en est une illustration~: il délimite ce qui a été observé. L'article sur
l'incertitude épistémique en tire l'instrument trichotomique et ses
mesures. Les travaux sur la décomposition, le suivi dans le temps et le jeu
entre auditeur et audité sont en préparation~; le Chapitre~\ref{chap:auditologie}
les présente comme tels.
\end{boxinfo}

\subsection{Trois cas}

\begin{tableaucours}{L{3.2cm}L{5.2cm}L{3.8cm}}{%
    \tableauHeader{Cas} &
    \tableauHeader{Situation} &
    \tableauHeader{Thèmes dans ce cours}
}
BANQUE SANAGA (Douala) & Banque de 38 agences, certifiée ISO/IEC~27001 sur
son centre de données, compromise en agence & SMSI, périmètre,
quantification, régulation bancaire \\
LiZbiZ (Douala) & Logistique, quatre centres de données, douze agences,
écart de stock et fraude & EBIOS Risk Manager, audit technique,
décomposition \\
ANPP (Yaoundé) & Agence publique de paiement & Audit de sécurité
réglementaire, correspondances de référentiels \\
\end{tableaucours}
\vspace{-0.8em}
\begin{center}{\small\itshape Tableau~P.5. Les trois cas fil rouge}\end{center}

\subsection{Le parcours}

\begin{tableaucours}{L{3.8cm}cL{6.8cm}}{%
    \tableauHeader{Partie} &
    \tableauHeader{Chapitres} &
    \tableauHeader{Question traitée}
}
I. Fondements & 1 à 2 & Qu'est-ce que gouverner la sécurité~? Qu'est-ce
qu'un risque cyber, et que sait-on de lui~? \\
II. Méthodes & 3 à 7 & Comment apprécier, traiter et chiffrer un risque~?
Comment conduire un audit et conclure~? \\
III. Cadre normatif & 8 à 10 & Quelles normes, quels référentiels, quel
droit, quelle déontologie~? \\
IV. Pratique & 11 à 14 & Comment mettre en œuvre et auditer un SMSI,
auditer la technique, faire vivre la GRC~? \\
V. Défis avancés & 15 & Tiers, nuage, intelligence artificielle, audit
continu \\
VI. Recherche & 16 & Mesurer ce que l'on ne sait pas \\
\end{tableaucours}
\vspace{-0.8em}
\begin{center}{\small\itshape Tableau~P.6. Les six parties du cours}\end{center}

Le cours compte 30~heures en présentiel, réparties en quinze séances de
deux heures qui couvrent le prologue et les six parties. Les fiches de
questions-réponses de chaque chapitre et les dix ateliers d'environ
90~minutes de NDAP SUKULU, sous le code \codeUE{} (registre des risques,
EBIOS Risk Manager en deux séances, quantification FAIR, décomposition et
cotation, déclaration d'applicabilité, table de correspondance des
référentiels, audit d'un SMSI, audit technique, synthèse), sont des
travaux à faire à la maison, hors de ces 30~heures. Les corrigés des
exercices de fin de chapitre sont réunis dans un livret séparé.

\begin{boxinfo}
\textbf{Ce que le cours ne fait pas.} Il ne refait pas les autres
enseignements du master~: le test d'intrusion relève de l'UE
MSN-2S3-SEC03, la forensique de l'UE MSN-2S3-SEC07, le droit du numérique et
des données personnelles de l'UE MSN-2S3-JUR01, l'audit de protection des
données de l'UE MSN-2S4-SEC02. Il y renvoie chaque fois que nécessaire.
\end{boxinfo}

\separateur

\begin{boxresume}
\textbf{Ce qu'il faut retenir de ce prologue~:}
\begin{itemize}
    \item \textbf{Un certificat dit ce qui a été audité}~: un SMSI conforme,
          sur un périmètre, à une date. Il ne dit rien du reste.
    \item \textbf{Certificat, conformité, sécurité effective} sont trois
          notions distinctes~; seule la dernière se mesure aux attaques.
    \item \textbf{Le périmètre est une décision de sécurité}~: un périmètre
          de complaisance exclut ce dont dépend ce qu'il protège.
    \item \textbf{Un risque jamais apprécié est inconnu, pas faible}~: il se
          cote \NEK{} et appelle une appréciation.
    \item \textbf{Trois mesures}, $\mu$, $U$ et $H$, lisent un ensemble de
          contrôles~; lues par périmètre, elles montrent où se cache
          l'incertitude.
    \item \textbf{Cinq réflexes}~: lire le périmètre, coter l'inconnu, prouver
          le fonctionnement, mesurer l'incertitude, relier norme, risque et
          audit.
\end{itemize}
\end{boxresume}

\begin{boxinfo}
\textbf{Points de contrôle.} Vous avez compris ce prologue si vous savez~:
\begin{itemize}
    \item expliquer pourquoi l'incident ne contredit pas le certificat~;
    \item dire ce qu'est un périmètre de complaisance, sur un exemple~;
    \item justifier la cote \NEK{} d'un scénario jamais apprécié~;
    \item calculer $\mu$, $U$ et $H$ pour 2, 4 et 4 réponses.
\end{itemize}
\end{boxinfo}

% ============================================================================
\section*{Exercices}
% ============================================================================

\begin{boxexo}[title={\ding{45}\quad Exercice~P.1. Certificat, conformité ou sécurité~? \niveaubase}]
Pour chacune des phrases suivantes, dites si elle porte sur le certificat,
la conformité ou la sécurité effective, et si elle est exacte.
\begin{enumerate}[(a)]
    \item « Notre certificat ISO/IEC~27001 prouve que nos agences sont
          protégées. »
    \item « L'audit de certification a constaté que la politique de sauvegarde
          du centre de données était appliquée. »
    \item « Aucun des trois exercices de simulation d'attaque de l'année n'a
          permis d'atteindre le système central. »
    \item « La politique de mots de passe satisfait l'exigence de longueur de
          la politique de sécurité du groupe. »
\end{enumerate}
\end{boxexo}

\begin{boxexo}[title={\ding{45}\quad Exercice~P.2. Lire un périmètre \niveaumoyen}]
Le périmètre d'un SMSI est ainsi rédigé~: « La gestion du système
d'information du centre de données principal de Douala, à l'exclusion des
postes de travail et de la messagerie. »
\begin{enumerate}[(a)]
    \item Citez trois activités de la banque dont dépend ce périmètre et qui
          en sont exclues.
    \item Pour chacune, décrivez un scénario d'attaque qui, partant de
          l'activité exclue, atteindrait le périmètre.
    \item Proposez une rédaction plus défendable, et un calendrier
          d'extension réaliste.
\end{enumerate}
\end{boxexo}

\begin{boxexo}[title={\ding{45}\quad Exercice~P.3. Mesurer deux périmètres \niveaumoyen}]
Six mois après l'incident, la même évaluation des dix contrôles donne, pour
les agences, 6~\OUI{}, 3~\NON{} et 1~\NEK{}.
\begin{enumerate}[(a)]
    \item Calculez $\mu$, $U$ et $H$ pour les agences.
    \item Comparez avec la situation de septembre (2/4/4). Qu'est-ce qui a
          le plus changé~: la conformité ou la connaissance~?
    \item Le nombre de \NON{} est passé de 4 à 3 seulement. Est-ce une
          mauvaise nouvelle~? Justifiez.
\end{enumerate}
\end{boxexo}

\begin{boxexo}[title={\ding{45}\quad Exercice~P.4. La note au conseil d'administration \niveauexpert}]
Le conseil d'administration demande à Mme~ABENA une note d'une page.
\begin{enumerate}[(a)]
    \item Expliquez, sans jargon, pourquoi le certificat n'a pas empêché
          l'incident, et pourquoi il reste valable.
    \item Rédigez une fiche de risque complète pour le scénario « vol
          d'identifiants de messagerie d'un dirigeant », sur le modèle de
          R-AG-01.
    \item Proposez trois décisions au conseil, avec leur coût approximatif
          en ordre de grandeur et le risque qu'elles traitent.
    \item Indiquez ce que le prochain audit devra vérifier pour que le
          conseil sache si ces décisions ont été appliquées.
\end{enumerate}
\end{boxexo}

\bigskip

\begin{boxinfo}
\textbf{Pour aller plus loin.}
\begin{itemize}
    \item L'ISO/IEC~27001:2022~\cite{iso27001_2022}, articles~4.1 à 4.3~:
          le contexte de l'organisation et le périmètre du SMSI.
    \item L'article \emph{Vers l'auditologie}~\cite{minka2026auditologie},
          sur l'observabilité partielle.
    \item L'article sur l'incertitude épistémique~\cite{minka2026nek}, qui
          définit la réponse \NEK{} et les mesures $\mu$, $U$, $H$.
\end{itemize}
\end{boxinfo}

\begin{boxinfo}
\textbf{Transition.} Le prologue a montré une banque qui gérait la sécurité
de son centre de données et ignorait celle de ses agences. Avant d'apprendre
les méthodes, il faut définir ce que signifie gouverner la cybersécurité, et
qui en porte la responsabilité. C'est l'objet du Chapitre~\ref{chap:grc}.
\end{boxinfo}


    % ======================================================================
    % PARTIE 1 — FONDEMENTS : GOUVERNANCE, RISQUE ET INCERTITUDE
    % Standards de référence : ISO/IEC 27000, ISO 31000, modèle des trois lignes de l'IIA ; article A1.
    % Niveau : fondamental
    % ======================================================================
    \part{Fondements : Gouvernance, Risque et Incertitude}
    \label{part:fondements}

        % ============================================================================
% Fichier  : partie1/P01_C01_gouverner_cybersecurite_grc.tex
% Titre    : Gouverner la Cybersécurité : du Contrôle à la GRC
% Auteur   : MINKA MI NGUIDJOÏ Thierry Emmanuel
% Version  : 1.0 -- novembre 2026
% Corrigés : corriges/Corr_C01.tex
% ============================================================================
% Faits du cas (bible des cas) : BANQUE SANAGA, 38 agences, DG M. FOTSO,
% RSSI Mme ABENA rattachée au directeur des risques, audit interne dirigé par
% Mme EKAMBI ; certificat ISO/IEC 27001:2022 sur le seul centre de données ;
% rançongiciel du 14/09/2026. Faits nouveaux introduits ici (reportés dans
% la bible) : le périmètre du SMSI a été proposé par le consultant, accepté
% par le DSI et signé par le DG sans passer par le conseil ; comité de
% sécurité sans procès-verbaux ; évaluation de gouvernance 2025 : 2/4/2.
% Sources : modèle des trois lignes (IIA, 2020) ; ISO/IEC 27001:2022, art. 5
% et 6.1.2 (paraphrasés) ; ISO/IEC 27014:2020 (processus de gouvernance,
% paraphrasés, à vérifier sur le texte) ; NIST CSF 2.0, fonction GOVERN ;
% COBIT 2019 (domaine EDM) ; loi n° 2010/012, art. 26 al. 1 (texte français
% officiel, Ann. D d'ACSI) ; COBAC R-2016/04 (paraphrasé, sans numéro
% d'article).
% ============================================================================

\chapter{Gouverner la Cybersécurité : du Contrôle à la GRC}
\label{chap:grc}
\chaptermark{Gouverner la cybersécurité}

\epigraph{%
    « La sécurité se décide au conseil,\\
    elle s'exécute dans les services,\\
    elle se vérifie ailleurs. »%
}{Principe du cours}

\niveauChapitre{Fondamental}{%
    Maîtrise du Chapitre~\ref{chap:prologue}%
}

\begin{boxobjectifs}
\begin{itemize}
    \item Définir gouvernance, gestion des risques et conformité, et leur
          articulation.
    \item Distinguer la gouvernance de la sécurité de sa gestion
          opérationnelle.
    \item Appliquer le modèle des trois lignes à la sécurité.
    \item Situer les rôles du conseil, de la direction, du comité de
          sécurité, du RSSI, du propriétaire de risque et du délégué à la
          protection des données.
    \item Choisir et justifier le rattachement d'un RSSI.
    \item Construire et vérifier une matrice RACI de la sécurité.
    \item Évaluer un dispositif de gouvernance par $\mu$, $U$ et $H$.
\end{itemize}
\end{boxobjectifs}

\bigskip

% ============================================================================
\section*{Avant de commencer : qui a décidé du périmètre~?}
% ============================================================================

Douala, 6~octobre 2026. Trois semaines après le rançongiciel, le comité
des risques du conseil d'administration de la BANQUE SANAGA\footnote{Comme
dans tout le cours, la BANQUE SANAGA, ses dirigeants, les dates et les
montants sont fictifs~; les normes et les textes cités sont exacts.} se
réunit en séance extraordinaire. Le président du comité ne demande pas
comment l'attaque s'est produite~: la RSSI l'a déjà expliqué. Il pose une
autre question~: « Qui a décidé que les agences resteraient hors du
périmètre du certificat~? »

Personne ne répond tout de suite. On reconstitue~: en 2024, le consultant qui a
accompagné la certification a proposé le centre de données seul~; le
directeur des systèmes d'information l'a accepté, parce que le calendrier
était serré~; le directeur général a signé la déclaration de périmètre, au
milieu d'autres documents. Le conseil d'administration ne l'a jamais vue.
Une décision qui engageait la sécurité de trente-huit agences a été prise
sans que l'organe qui répond des risques de la banque en soit informé.

\medskip

Ce chapitre explique ce qu'est gouverner la cybersécurité, qui en porte la
responsabilité et comment on le vérifie. Il pose le cinquième réflexe du
prologue~: relier norme, risque et audit au lieu de les mener comme trois
exercices séparés.

\begin{boxalert}
\textbf{Violation typique.} Traiter une décision de sécurité structurante
(périmètre, acceptation d'un risque, budget) comme une décision technique,
prise et signée sans que l'instance qui répond des risques en soit saisie.

\smallskip
\textbf{Règle~:} toute décision qui fixe le niveau de risque accepté par
l'organisation remonte à l'instance qui en répond. Le technicien propose,
la direction arbitre, le conseil approuve et surveille.
\end{boxalert}

\separateur

% ============================================================================
\section{Du Contrôle à la GRC}
\label{sec:c01-grc}
% ============================================================================

\subsection{Trois âges de la sécurité}

La sécurité des systèmes d'information n'a pas toujours été une affaire de
gouvernance. On peut distinguer trois âges, qui coexistent encore dans
beaucoup d'organisations.

\begin{tableaucours}{L{2.6cm}L{4.6cm}L{5.0cm}}{%
    \tableauHeader{Âge} &
    \tableauHeader{Question posée} &
    \tableauHeader{Limite}
}
Technique & Quels outils installer (pare-feu, antivirus, sauvegardes)~? &
Personne ne dit quels risques ces outils traitent, ni lesquels ils
laissent \\
Contrôle & Les contrôles prévus sont-ils appliqués~? & On vérifie ce qui
est prévu, pas ce qui manque \\
Gouvernance & Quels risques l'organisation accepte-t-elle, qui en répond,
et comment le sait-on~? & Exige que la direction s'en saisisse
réellement \\
\end{tableaucours}
\vspace{-0.8em}
\begin{center}{\small\itshape Tableau~1.1. Trois âges de la sécurité des
systèmes d'information}\end{center}

La BANQUE SANAGA était à l'âge du contrôle~: ses contrôles du centre de
données étaient documentés et appliqués, et l'audit de certification l'a
confirmé. Elle n'était pas à l'âge de la gouvernance~: personne, au niveau
du conseil, n'avait décidé quels risques la banque acceptait de laisser
hors du périmètre.

\subsection{Gouvernance, risque, conformité}

\begin{boxdef}
\termecle{Gouvernance de la sécurité}\index{gouvernance de la sécurité}~:
ensemble des responsabilités et des pratiques par lesquelles l'organe de
direction fixe les objectifs de sécurité de l'organisation, décide du
niveau de risque acceptable, attribue les rôles et surveille les
résultats~\cite{iso27014_2020}.

\medskip
\termecle{Gestion des risques}\index{gestion des risques}~: activités
coordonnées pour identifier, analyser, évaluer et traiter les risques, puis
en suivre l'évolution. Le Chapitre~\ref{chap:risque} en définit le
vocabulaire.

\medskip
\textbf{Conformité}~: au sens du prologue, satisfaction prouvée d'une
exigence~; appliquée ici à l'ensemble des obligations qui s'imposent à
l'organisation (lois, règlements, contrats, normes auxquelles elle
s'engage) et à ses propres politiques.

\medskip
\termecle{Gouvernance, risque et conformité}\index{GRC} (\sigle{GRC})~:
approche qui pilote la
gouvernance, la gestion des risques et la conformité comme un seul
dispositif, avec un référentiel commun d'objectifs, de risques, de
contrôles et de preuves.
\end{boxdef}

Les trois termes ne désignent pas trois services. Ils désignent trois
questions, qui se posent sur le même objet~:
\begin{itemize}
    \item la \textbf{gouvernance} répond à «~que voulons-nous, et jusqu'où
          acceptons-nous le risque~?~»~;
    \item la \textbf{gestion des risques} répond à «~qu'est-ce qui peut nous
          empêcher d'y parvenir, et que faisons-nous~?~»~;
    \item la \textbf{conformité} répond à «~respectons-nous ce qui nous est
          imposé et ce que nous nous sommes imposé~?~».
\end{itemize}

L'\textbf{audit} n'est pas une quatrième lettre~: il donne une assurance
indépendante sur les trois. C'est pourquoi le cours l'étudie à part
(Chapitre~\ref{chap:audit}).

\begin{figure}[H]
\centering
\begin{tikzpicture}[
    bloc/.style={draw=CouleurPrimaire, thick, rounded corners,
                 fill=CouleurPrimaire!8, minimum width=3.4cm,
                 minimum height=1.3cm, align=center, font=\small},
    audit/.style={draw=CouleurAccent, thick, rounded corners,
                  fill=CouleurAccent!12, minimum width=11.6cm,
                  minimum height=0.9cm, align=center, font=\small},
    fleche/.style={->, thick, CouleurSecondaire}]
    \node[bloc] (G) at (0,0) {\textbf{Gouvernance}\\objectifs, appétence,\\rôles};
    \node[bloc] (R) at (4.1,0) {\textbf{Risque}\\apprécier,\\traiter, suivre};
    \node[bloc] (C) at (8.2,0) {\textbf{Conformité}\\obligations,\\preuves};
    \draw[fleche] (G) -- node[above, font=\scriptsize]{fixe le seuil} (R);
    \draw[fleche] (R) -- node[above, font=\scriptsize]{choisit les mesures} (C);
    \draw[fleche] (C.south) to[bend left=25]
        node[below, font=\scriptsize]{rend compte} (G.south);
    \node[audit] (A) at (4.1,-2.6) {\textbf{Audit}~: assurance indépendante
        sur les trois};
\end{tikzpicture}
\caption{L'articulation de la gouvernance, du risque et de la conformité}
\label{fig:c01-grc}
\end{figure}

\begin{boxinfo}
\textbf{GRC n'est pas un logiciel.} Le marché vend des « outils GRC ». Ils
peuvent aider à tenir les registres, mais un outil ne remplace ni une
décision du conseil sur l'appétence, ni un propriétaire de risque qui
signe. Une organisation qui achète l'outil sans définir les rôles obtient
des registres mieux présentés, et aussi incomplets. Le
Chapitre~\ref{chap:grc-pratique} montre comment faire vivre la GRC au
quotidien.
\end{boxinfo}

% ============================================================================
\section{Gouverner n'est pas Gérer}
\label{sec:c01-gouverner}
% ============================================================================

\subsection{Deux niveaux de décision}

La distinction entre gouvernance et gestion vient de la gouvernance des
entreprises~; les référentiels du numérique l'ont reprise. COBIT sépare
ainsi les objectifs de gouvernance, réunis dans le domaine
\emph{évaluer, diriger, surveiller}, des objectifs de gestion, qui
planifient, construisent, exploitent et contrôlent~\cite{isaca2018cobitintro}.
L'ISO/IEC~27014 applique la même logique à la sécurité de l'information~:
l'organe de direction \emph{évalue} la situation et les risques,
\emph{dirige} en fixant objectifs et politiques, \emph{surveille} les
résultats, et \emph{communique} avec les parties
prenantes~\cite{iso27014_2020}.

\begin{tableaucours}{L{3.0cm}L{4.6cm}L{4.6cm}}{%
    \tableauHeader{Processus} &
    \tableauHeader{Ce que fait l'organe de direction} &
    \tableauHeader{À la BANQUE SANAGA, avant l'incident}
}
Évaluer & Apprécie la situation de sécurité et les risques au regard de la
stratégie & Le conseil ne recevait aucun rapport de sécurité \\
Diriger & Fixe les objectifs, l'appétence, les politiques et les moyens &
La politique existait~; l'appétence au risque cyber n'était pas fixée \\
Surveiller & Vérifie que les objectifs sont atteints, par des indicateurs
et par l'audit & Le seul retour était le certificat \\
Communiquer & Rend compte aux actionnaires, au régulateur, aux clients &
Communication de crise improvisée en septembre 2026 \\
\end{tableaucours}
\vspace{-0.8em}
\begin{center}{\small\itshape Tableau~1.2. Les processus de gouvernance de
la sécurité (d'après l'ISO/IEC~27014)}\end{center}

\begin{boxdef}
\termecle{Organe de direction}\index{organe de direction}~: personne ou
groupe qui répond, en dernier ressort, de la performance et de la
conformité de l'organisation~: le conseil d'administration dans une
société anonyme, le conseil d'administration ou l'autorité de tutelle dans
un établissement public.

\medskip
\termecle{Direction}\index{direction} (ou management)~: personnes qui
dirigent l'organisation au quotidien, sous l'autorité de l'organe de
direction~: le directeur général et ses directeurs.
\end{boxdef}

\subsection{Ce que la norme attend de la direction}

L'ISO/IEC~27001 consacre son article~5 au \emph{leadership}. En substance,
la direction doit démontrer qu'elle s'implique dans le système de
management~: elle établit une politique de sécurité compatible avec la
stratégie, intègre ses exigences dans les processus de l'organisation,
fournit les ressources, et attribue les responsabilités et les autorités
des rôles de sécurité~\cite{iso27001_2022}. L'article~6.1.2 ajoute que
chaque risque identifié doit avoir un \emph{propriétaire}.

La norme n'exige pas que le conseil d'administration approuve le
périmètre. Elle exige que la direction s'implique~; elle ne dit pas
comment l'organe de direction la surveille. C'est là que la gouvernance
d'entreprise prend le relais~: pour une banque, le règlement COBAC
R-2016/04 confie à l'organe délibérant la responsabilité du dispositif de
contrôle interne et de maîtrise des risques, et à l'organe exécutif sa mise
en œuvre~\cite{cobac2016r04}. Le Chapitre~\ref{chap:droit} précisera ces
obligations.

\subsection{Le cadre NIST~: gouverner d'abord}

Dans sa version~2.0, publiée en 2024, le cadre de cybersécurité du NIST a
ajouté une sixième fonction, \emph{Gouverner}, aux cinq fonctions
historiques (identifier, protéger, détecter, répondre,
rétablir)~\cite{nist2024csf}. Elle couvre six catégories~: le contexte de
l'organisation~; la stratégie de gestion des risques~; les rôles, les
responsabilités et les autorités~; la politique~; la supervision~; la
gestion des risques de la chaîne d'approvisionnement. Le cadre la place au
centre des cinq autres~: sans elle, elles n'ont pas d'objectif. Le
Chapitre~\ref{chap:referentiels} étudie ce cadre et ses correspondances
avec l'ISO/IEC~27001.

\begin{boxscenario}{Opération SANAGA. Une décision orpheline}
\textbf{Constatations.}
\begin{itemize}
    \item La déclaration de périmètre du SMSI porte la seule signature du
          directeur général.
    \item Les procès-verbaux du conseil et de son comité des risques de 2024
          et 2025 ne mentionnent ni le périmètre ni le certificat, sauf pour
          annoncer son obtention.
    \item Aucun document ne fixe l'appétence de la banque au risque cyber.
\end{itemize}

\textbf{Analyse} (confiance élevée). La direction a satisfait l'exigence
de la norme~: elle s'est impliquée et a signé. Mais l'organe de direction
n'a ni évalué, ni dirigé, ni surveillé cette décision. On peut discuter,
comme le prologue, de la prise en compte des interfaces et des dépendances
du périmètre~; mais le défaut premier n'est pas de conformité à
l'ISO/IEC~27001, il est de gouvernance. Un audit de certification pouvait
contester le périmètre~; il ne pouvait pas corriger le silence du conseil,
puisqu'il ne porte pas sur son fonctionnement.
\end{boxscenario}

% ============================================================================
\section{Le Modèle des Trois Lignes appliqué à la Sécurité}
\label{sec:c01-trois-lignes}
% ============================================================================

\subsection{Le modèle}

L'Institut des auditeurs internes (\sigle{IIA}) a publié en 2020 le
\termecle{modèle des trois lignes}\index{modèle des trois
lignes}~\cite{iia2020troislignes}. Il remplace l'ancien modèle des
« trois lignes de défense »~: le mot \emph{défense} a disparu, parce que
la gestion des risques sert aussi à saisir des opportunités, et parce que
les lignes doivent coopérer, pas seulement se protéger les unes des autres.
Le modèle distingue des \emph{rôles}, pas des services~: une même
direction peut tenir un rôle de première ligne pour une activité et de
deuxième ligne pour une autre.

\begin{tableaucours}{L{2.7cm}L{4.5cm}L{5.0cm}}{%
    \tableauHeader{Rôle} &
    \tableauHeader{Ce qu'il fait en sécurité} &
    \tableauHeader{À la BANQUE SANAGA}
}
Organe de gouvernance & Fixe l'appétence, approuve la stratégie de
sécurité, surveille & Conseil d'administration, comité des risques du
conseil \\
Première ligne & Exploite les systèmes et applique les mesures de sécurité
au quotidien & DSI (exploitation, réseau, postes), directeurs métier,
chefs d'agence \\
Deuxième ligne & Fixe le cadre, conseille, surveille et remet en cause la
première ligne & RSSI, direction des risques, conformité, DPD \\
Troisième ligne & Donne une assurance indépendante à l'organe de
gouvernance & Direction de l'audit interne (Mme~EKAMBI) \\
Prestataires d'assurance externes & Assurance complémentaire &
CERTIQUAL AFRIQUE, commissaires aux comptes \\
\end{tableaucours}
\vspace{-0.8em}
\begin{center}{\small\itshape Tableau~1.3. Le modèle des trois lignes
appliqué à la sécurité de la BANQUE SANAGA}\end{center}

Le superviseur, la COBAC, n'apparaît pas dans le tableau~: il n'est pas
un prestataire au service de la banque, mais une autorité extérieure qui
contrôle l'ensemble du dispositif, les trois lignes comprises.

\subsection{Où placer le RSSI~?}

Le RSSI est l'acteur le plus difficile à placer, parce qu'on lui confie
souvent des tâches des deux premières lignes à la fois. S'il fixe la
politique, surveille les indicateurs et rend compte à la direction, il
tient un rôle de deuxième ligne. S'il administre lui-même le pare-feu, il
tient un rôle de première ligne. S'il fait les deux, il surveille son
propre travail.

\begin{boxalert}
\textbf{Violation typique.} Le RSSI administre les équipements de
sécurité, puis rédige le rapport qui dit au comité qu'ils fonctionnent.

\smallskip
\textbf{Règle~:} celui qui exploite une mesure n'en atteste pas
l'efficacité. Si l'organisation est trop petite pour séparer les rôles,
elle compense par une revue indépendante (audit interne, auditeur
externe) et le dit dans son rapport.
\end{boxalert}

\subsection{Ce que le modèle ne dit pas}

Le modèle des trois lignes est un modèle de rôles. Il ne dit pas si les
lignes couvrent tout l'univers des risques. À la BANQUE SANAGA, les trois
lignes existaient et fonctionnaient~; mais aucune ne regardait les agences
sous l'angle de la cybersécurité~: la DSI les exploitait sans appréciation
des risques, la RSSI concentrait ses efforts sur le périmètre certifié, et
l'audit interne n'avait pas inscrit la sécurité des agences à son
programme. Trois lignes bien rangées peuvent laisser le même trou. C'est
ce que mesure l'évaluation de la Section~\ref{sec:c01-mesurer}.

% ============================================================================
\section{Les Acteurs et leurs Responsabilités}
\label{sec:c01-acteurs}
% ============================================================================

\subsection{Huit acteurs}

\begin{tableaucours}{L{3.2cm}L{9.0cm}}{%
    \tableauHeader{Acteur} &
    \tableauHeader{Responsabilité principale}
}
Conseil et comité des risques & Fixe l'appétence au risque cyber, approuve
la stratégie et les décisions structurantes (périmètre, risques acceptés
au-delà d'un seuil), reçoit les rapports du RSSI et de l'audit interne \\
Direction générale & Met en œuvre la stratégie, fournit les ressources,
arbitre entre les directions, signe la politique de sécurité \\
Comité de sécurité & Instance de direction qui réunit les directeurs
métier, la DSI, le RSSI et les risques~; suit les projets de sécurité, les
incidents et les indicateurs, et prépare les décisions \\
RSSI & Propose la politique, coordonne l'appréciation des risques,
surveille les mesures, conseille, rend compte \\
DSI & Exploite les systèmes et les mesures techniques, applique la
politique dans les projets et l'exploitation \\
Propriétaire de risque & Directeur qui a l'autorité et les moyens de
traiter un risque~; décide du traitement et accepte le risque résiduel
dans la limite de l'appétence \\
Audit interne & Évalue le dispositif de façon indépendante et rend compte
au conseil \\
DPD & Veille à la conformité des traitements de données personnelles,
conseille, sert de point de contact avec l'autorité de protection \\
\end{tableaucours}
\vspace{-0.8em}
\begin{center}{\small\itshape Tableau~1.4. Les acteurs de la gouvernance
de la sécurité}\end{center}

\begin{boxdef}
\termecle{Propriétaire de risque}\index{propriétaire de risque}~:
personne qui a la responsabilité et l'autorité de gérer un
risque~\cite{iso27000_2018}. Ce n'est pas le RSSI~: c'est en général le
directeur dont l'activité est exposée, parce que c'est lui qui peut
dégager les moyens de traitement et qui supporte les conséquences.

\medskip
\termecle{Appétence au risque}\index{appétence au risque}~: niveau de
risque qu'une organisation est prête à accepter pour atteindre ses
objectifs. Elle est fixée par l'organe de direction~\cite{coso2017erm}~; le
Chapitre~\ref{chap:risque} montre comment l'exprimer.
\end{boxdef}

\subsection{Le propriétaire du risque R-AG-01}

La fiche du prologue décrivait le risque. Il lui manquait les rubriques de
gouvernance~: qui en répond, qui peut l'accepter, qui le surveille.

\begin{boxobjet}{R-AG-01 · Rubriques de gouvernance}
\textbf{Propriétaire du risque~:} le directeur du réseau d'agences. C'est
lui qui dirige les agences, arbitre leur budget et supporte la fermeture
des guichets.

\textbf{Contributeurs~:} DSI (postes, messagerie, partage de fichiers),
RSSI (analyse et mesures), direction de la communication (crise).

\textbf{Instance d'acceptation~:} le comité des risques du conseil. Tant
que le conseil n'a pas fixé l'appétence, aucun directeur ne peut accepter
seul un risque de cette ampleur~; une fois l'appétence fixée, le comité
reste compétent pour tout résiduel qui la dépasse.

\textbf{Surveillance~:} la RSSI suit trois indicateurs (part des comptes
d'agence avec second facteur, délai d'application des correctifs, succès
des tests de restauration) et les présente chaque trimestre au comité de
sécurité.

\textbf{Assurance~:} la direction de l'audit interne inscrit le réseau
d'agences à son programme 2027.
\end{boxobjet}

\subsection{Rattacher le RSSI}

Le rattachement du RSSI n'est pas un détail d'organigramme~: il détermine
ce qu'il peut dire, et à qui. Aucune norme n'impose un rattachement
unique~; chaque option a ses forces et ses risques.

\begin{tableaucours}{L{3.0cm}L{4.6cm}L{4.6cm}}{%
    \tableauHeader{Rattachement} &
    \tableauHeader{Force} &
    \tableauHeader{Risque}
}
DSI & Proximité technique, mise en œuvre rapide & Le RSSI surveille son
propre supérieur~; la sécurité passe après les délais des projets \\
Direction des risques & Intégration au dispositif de maîtrise des risques~;
deuxième ligne naturelle & Distance avec la technique~; risque d'une
sécurité réduite à des registres \\
Direction générale & Autorité, accès direct aux arbitrages & Charge de la
direction générale~; peut isoler le RSSI des autres directions \\
Conformité & Lien avec les obligations légales & Une sécurité réduite à la
conformité, aveugle aux menaces non réglementées \\
\end{tableaucours}
\vspace{-0.8em}
\begin{center}{\small\itshape Tableau~1.5. Les options de rattachement du
RSSI}\end{center}

À la BANQUE SANAGA, la RSSI est rattachée au directeur des risques. C'est
un choix défendable~: il place la sécurité en deuxième ligne, à distance
de la DSI. Mais ce rattachement ne valait rien sans un accès au comité des
risques du conseil~: les rapports de la RSSI s'arrêtaient au directeur des
risques. Le critère décisif n'est pas l'organigramme, c'est le
\emph{droit d'être entendu} par l'instance qui décide.

\begin{boxinfo}
\textbf{Le délégué à la protection des données.} Le DPD et le RSSI
travaillent sur des objets voisins~: les données personnelles sont aussi
des informations à protéger. Leurs missions diffèrent pourtant~: le RSSI
protège l'information de l'organisation contre les menaces~; le DPD veille
aux droits des personnes dont les données sont traitées. Confondre les
deux fonctions crée un conflit d'intérêts, car le RSSI peut avoir à
défendre un traitement que le DPD devrait contester. La place du DPD dans
la loi n°~2024/017 relative à la protection des données à caractère
personnel est étudiée dans l'UE MSN-2S3-JUR01~; le cours n'y revient que
pour situer les rôles.
\end{boxinfo}

% ============================================================================
\section{La Matrice RACI de la Sécurité}
\label{sec:c01-raci}
% ============================================================================

\subsection{Quatre lettres}

\begin{boxdef}
\termecle{Matrice RACI}\index{matrice RACI}~: tableau qui croise des
activités et des acteurs, et indique pour chaque case le rôle de l'acteur~:
\begin{itemize}
    \item \textbf{R} (\emph{Responsible}, réalise)~: exécute l'activité~;
    \item \textbf{A} (\emph{Accountable}, répond)~: rend compte du résultat,
          approuve, et ne peut pas déléguer cette responsabilité~;
    \item \textbf{C} (\emph{Consulted}, consulté)~: donne un avis avant la
          décision~;
    \item \textbf{I} (\emph{Informed}, informé)~: est tenu au courant après.
\end{itemize}
\end{boxdef}

\begin{boxprocedure}
\textbf{Vérifier une matrice RACI.} Le papier de travail qui examine une
matrice contrôle, ligne par ligne~:
\begin{itemize}
    \item \textbf{un seul A} par activité~: deux personnes qui répondent,
          c'est personne~;
    \item \textbf{au moins un R}~: une activité sans exécutant n'est pas
          faite~;
    \item \textbf{un A qui a l'autorité}~: celui qui répond doit pouvoir
          décider et dégager les moyens~;
    \item \textbf{pas de R et de A d'attestation sur la même tête}~: celui
          qui exécute une mesure n'atteste pas son efficacité~;
    \item \textbf{les décisions de risque au bon niveau}~: l'acceptation
          d'un risque au-delà de l'appétence remonte à l'organe de
          gouvernance~;
    \item \textbf{la preuve}~: pour chaque A, un document signé (procès-verbal,
          décision, rapport) montre que le rôle a été tenu.
\end{itemize}
\end{boxprocedure}

\subsection{La matrice de la BANQUE SANAGA, après l'incident}

La RSSI propose au comité de sécurité la matrice du Tableau~1.6. Les
colonnes sont~: CA (conseil d'administration et son comité des risques),
DG, CS (comité de sécurité), RSSI, DSI, PR (propriétaire de risque, ici le
directeur métier concerné), AI (audit interne) et DPD.

\begin{tableaucours}{L{4.6cm}cccccccc}{%
    \tableauHeader{Activité} &
    \tableauHeader{CA} & \tableauHeader{DG} & \tableauHeader{CS} &
    \tableauHeader{RSSI} & \tableauHeader{DSI} & \tableauHeader{PR} &
    \tableauHeader{AI} & \tableauHeader{DPD}
}
Fixer l'appétence au risque cyber & A & R & C & C & -- & C & I & -- \\
Approuver la politique de sécurité & I & A & C & R & C & C & I & C \\
Fixer le périmètre du SMSI & A & R & C & R & C & C & I & I \\
Apprécier les risques d'une activité & -- & I & I & R & C & A & I & C \\
Accepter un risque résiduel sous le seuil & -- & I & I & C & -- & A/R & I & -- \\
Accepter un risque résiduel au-delà du seuil & A & R & C & C & -- & C & I & -- \\
Exploiter les mesures techniques & -- & -- & I & C & A/R & -- & -- & -- \\
Surveiller l'efficacité des mesures & I & I & I & A/R & C & I & -- & -- \\
Notifier un incident majeur au régulateur & I & A & I & R & C & C & I & C \\
Auditer le dispositif de sécurité & I & I & I & C & C & C & A/R & -- \\
\end{tableaucours}
\vspace{-0.8em}
\begin{center}{\small\itshape Tableau~1.6. Matrice RACI de la sécurité de
la BANQUE SANAGA (proposition de la RSSI)}\end{center}

Deux lignes demandent une explication. Pour « fixer le périmètre du
SMSI », le DG et la RSSI réalisent (R) et le conseil répond (A)~: c'est
précisément la décision qui avait échappé au conseil. Pour « auditer le
dispositif », l'audit interne répond seul~; la RSSI n'est que consultée,
parce qu'elle ne peut pas auditer ce qu'elle surveille. La mention
«~A/R~» signale un acteur qui répond de l'activité et la réalise par ses
propres équipes~: elle est normale pour l'exploitation (DSI) ou l'audit
(audit interne), à condition que l'attestation d'efficacité vienne d'un
autre acteur.

\begin{boxinfo}
\textbf{Une matrice se lit aussi par colonne.} Une colonne sans aucun A
désigne un acteur qui ne répond de rien~; une colonne chargée de A désigne
un goulet d'étranglement. Dans le Tableau~1.6, la colonne CA porte trois A,
tous sur des décisions de niveau de risque~: c'est le signe que le conseil
gouverne. Un conseil sans A en sécurité ne gouverne pas la sécurité.
\end{boxinfo}

% ============================================================================
\section{Mesurer la Gouvernance}
\label{sec:c01-mesurer}
% ============================================================================

\subsection{Huit questions}

Une gouvernance se décrit dans des documents~; elle se vérifie dans leur
usage. Après l'incident, Mme~EKAMBI évalue la gouvernance de la sécurité
telle qu'elle fonctionnait en 2025, par huit questions. Chaque réponse est
cotée \OUI{} (prouvé), \NON{} (absent ou contredit par les pièces) ou
\NEK{} (impossible à établir faute de pièce).

\begin{tableaucours}{cL{6.6cm}cL{3.4cm}}{%
    \tableauHeader{N°} &
    \tableauHeader{Question} &
    \tableauHeader{Cote} &
    \tableauHeader{Pièce}
}
G1 & Une politique de sécurité est-elle approuvée par la direction
générale~? & \OUI{} & Politique signée, 2024 \\
G2 & Le périmètre du SMSI a-t-il été approuvé par le conseil ou son comité
des risques~? & \NON{} & Procès-verbaux 2024--2025 \\
G3 & L'appétence au risque cyber est-elle fixée par le conseil~? & \NON{} &
Aucun document \\
G4 & Chaque risque du registre a-t-il un propriétaire désigné~? & \NON{} &
Registre 2025, sans colonne propriétaire \\
G5 & Le RSSI est-il indépendant de la DSI~? & \OUI{} & Organigramme,
fiche de poste \\
G6 & Le comité de sécurité s'est-il réuni au moins deux fois en 2025~? &
\NEK{} & Aucun procès-verbal retrouvé \\
G7 & Un rapport de sécurité a-t-il été présenté au comité des risques du
conseil en 2025~? & \NEK{} & Ordre du jour sans annexe \\
G8 & L'audit interne a-t-il couvert la sécurité dans son programme
2024--2025~? & \NON{} & Programmes d'audit \\
\end{tableaucours}
\vspace{-0.8em}
\begin{center}{\small\itshape Tableau~1.7. Évaluation de la gouvernance de
la sécurité en 2025}\end{center}

\begin{boxlecture}{Gouvernance de la sécurité de la BANQUE SANAGA en 2025}
\textbf{Constatation.} 2~\OUI{}, 4~\NON{}, 2~\NEK{} sur 8 questions~:
\[
\mu = \tfrac{2}{8} = 0{,}25~;\quad U = \tfrac{2}{8} = 0{,}25~;\quad
H = -\left(2 \times 0{,}25\log_2 0{,}25 + 0{,}5\log_2 0{,}5\right)
  = 0{,}5 + 0{,}5 + 0{,}5 = 1{,}50~\text{bit}.
\]

\textbf{Analyse} (confiance élevée sur le calcul et sur les \NON{},
moyenne sur la lecture des \NEK{}). La gouvernance existait sur le papier
(politique, RSSI indépendant) et manquait dans les décisions~: périmètre,
appétence, propriétaires, assurance. Les deux \NEK{} portent sur la
\emph{surveillance}~: rien ne prouve que l'organe de direction ait été
informé. Un \NEK{} sur la surveillance n'est pas neutre~: un conseil qui
ne peut pas prouver qu'il a été informé ne peut pas prouver qu'il a
gouverné.

\medskip
\hfill\lecture{0{,}25}{0{,}25}{1{,}50}
\end{boxlecture}

\subsection{Ce que la mesure ajoute}

Dans l'évaluation des dix contrôles clés, le centre de données certifié
obtenait une conformité de 0,80 (prologue, Tableau~P.3). La gouvernance
qui l'encadrait obtient 0,25.
Les deux chiffres ne se contredisent pas~: ils ne mesurent pas le même
objet. Le premier dit que les contrôles du périmètre sont appliqués~; le
second dit que personne, au niveau qui décide, n'a choisi ce périmètre en
connaissance de cause. Le cinquième réflexe du cours, relier norme, risque
et audit, consiste à lire ces deux niveaux ensemble.

% ============================================================================
\section{Le Cadre Camerounais}
\label{sec:c01-cadre}
% ============================================================================

Deux textes donnent à la gouvernance de la sécurité une assise juridique
au Cameroun. Le Chapitre~\ref{chap:droit} les étudie en détail.

\begin{boxjuris}
\textbf{\loicam{n°~2010/012}, article~26, alinéa~1}~\cite{loi2010012}~:

\medskip
\textit{« Les exploitants des systèmes d'information prennent toutes les
mesures techniques et administratives afin de garantir la sécurité des
services offerts. A cet effet, ils se dotent de systèmes normalisés leur
permettant d'identifier, d'évaluer, de traiter et de gérer continûment les
risques liés à la sécurité des systèmes d'information dans le cadre des
services offerts directement ou indirectement. »}

\medskip
\textbf{Portée pour la gouvernance.} La loi impose aux exploitants de
systèmes d'information un dispositif \emph{normalisé} et \emph{permanent}
de gestion des risques de sécurité, pour les services fournis
« directement ou indirectement ».

\medskip
\textbf{Analyse} (lecture du cours, confiance moyenne en l'absence de
jurisprudence connue). Le texte ne prévoit pas de restreindre ce
dispositif à une partie des services. Un dispositif qui n'apprécie les
risques que sur le centre de données répond donc mal à cette obligation,
quel que soit le certificat.
\end{boxjuris}

Pour les établissements de crédit, le règlement COBAC R-2016/04 organise
le contrôle interne et la maîtrise des risques, y compris les risques liés
aux systèmes d'information, et en répartit la responsabilité entre
l'organe délibérant et l'organe exécutif~\cite{cobac2016r04}. Le défaut de
gouvernance de la BANQUE SANAGA n'est donc pas seulement une faiblesse de
bonne pratique~: il expose la banque à une observation du superviseur.

% ============================================================================
\section{Ce que la Gouvernance demande à chaque Métier}
\label{sec:c01-metiers}
% ============================================================================

\begin{boxmetier}{RSSI}
\textbf{Obtenir le droit d'être entendu}
\begin{enumerate}
    \item Demander un point permanent à l'ordre du jour du comité des
          risques du conseil, au moins deux fois par an.
    \item Présenter peu d'indicateurs, reliés aux risques et à l'appétence,
          et dire ce qu'on ignore (les \NEK{}) autant que ce qu'on sait.
    \item Se dégager de l'exploitation des outils de sécurité, ou faire
          revoir par un tiers ce qu'on exploite soi-même.
\end{enumerate}
\end{boxmetier}

\begin{boxmetier}{Responsable GRC}
\textbf{Un seul référentiel}
\begin{enumerate}
    \item Tenir à jour la matrice RACI et la faire approuver par le comité
          de sécurité.
    \item Vérifier que chaque risque du registre a un propriétaire, et que
          chaque acceptation au-delà du seuil porte la décision de
          l'instance compétente.
    \item Relier registre des risques, déclaration d'applicabilité et
          programme d'audit, pour qu'aucun pan de l'activité n'échappe aux
          trois (Chapitre~\ref{chap:grc-pratique}).
\end{enumerate}
\end{boxmetier}

\begin{boxmetier}{Auditeur en cybersécurité}
\textbf{Auditer la gouvernance, pas seulement les contrôles}

L'auditeur ne s'arrête pas aux mesures techniques. Il demande qui a décidé
du périmètre, qui a accepté chaque risque résiduel, et quelle pièce le
prouve. Une décision structurante sans instance compétente est une
constatation d'audit, même si chaque contrôle est en place.
\end{boxmetier}

\begin{boxmetier}{Gestionnaire des risques cyber}
\textbf{Faire signer les propriétaires}

Le gestionnaire des risques anime l'appréciation, mais ne possède pas les
risques. Il obtient de chaque directeur concerné qu'il reconnaisse son
risque, choisisse un traitement et signe l'acceptation du résiduel. Un
risque que personne n'a signé reste, en pratique, à la charge du RSSI, qui
n'a pas les moyens de le traiter.
\end{boxmetier}

\separateur

\begin{boxresume}
\textbf{Ce qu'il faut retenir de ce chapitre~:}
\begin{itemize}
    \item \textbf{GRC}~: trois questions sur un même objet~: ce que l'on
          veut et jusqu'où l'on accepte le risque, ce qui peut l'empêcher,
          ce que l'on doit respecter. L'audit donne une assurance
          indépendante sur les trois.
    \item \textbf{Gouverner n'est pas gérer}~: l'organe de direction évalue,
          dirige, surveille et communique~; la direction met en œuvre.
    \item \textbf{Trois lignes}~: l'exploitation, la surveillance et
          l'assurance indépendante sont des rôles distincts~; celui qui
          exploite une mesure n'en atteste pas l'efficacité.
    \item \textbf{Le propriétaire de risque} est le directeur qui a
          l'autorité de le traiter, pas le RSSI.
    \item \textbf{Le rattachement du RSSI} compte moins que son droit
          d'être entendu par l'instance qui décide.
    \item \textbf{RACI}~: un seul A par activité, des décisions de risque au
          bon niveau, une preuve pour chaque A.
    \item \textbf{La gouvernance se mesure}~: à la BANQUE SANAGA, 0,25 de
          conformité et 0,25 d'incertitude en 2025, alors que les
          contrôles du périmètre certifié obtenaient 0,80.
\end{itemize}
\end{boxresume}

\begin{boxinfo}
\textbf{Points de contrôle.} Vous avez compris ce chapitre si vous savez~:
\begin{itemize}
    \item dire, sur un exemple, ce qui relève de la gouvernance, de la
          gestion des risques et de la conformité~;
    \item placer les acteurs de la sécurité d'une organisation dans le
          modèle des trois lignes~;
    \item repérer les défauts d'une matrice RACI~;
    \item expliquer pourquoi un défaut de gouvernance échappe à un audit de
          certification.
\end{itemize}
\end{boxinfo}

% ============================================================================
\section*{Exercices}
% ============================================================================

\begin{boxexo}[title={\ding{45}\quad Exercice~1.1. Gouvernance, risque ou conformité~? \niveaubase}]
Classez chacune des activités suivantes dans la gouvernance, la gestion des
risques ou la conformité, et indiquez l'acteur qui en répond à la BANQUE
SANAGA.
\begin{enumerate}[(a)]
    \item Le conseil décide qu'aucune agence ne doit rester fermée plus de
          24~heures pour une cause informatique.
    \item La RSSI anime un atelier qui identifie les scénarios de
          rançongiciel sur le réseau des agences.
    \item La direction juridique vérifie que la banque a bien transmis à la
          COBAC le rapport demandé après la bascule de mars 2026.
    \item Le directeur du réseau d'agences choisit de cloisonner le partage
          de fichiers plutôt que d'accepter le risque.
    \item Le comité des risques du conseil reçoit le rapport trimestriel
          de sécurité.
\end{enumerate}
\end{boxexo}

\begin{boxexo}[title={\ding{45}\quad Exercice~1.2. Les trois lignes de LiZbiZ \niveaumoyen}]
Chez LiZbiZ, M.~KOUNOU, responsable informatique, détient tous les droits
de l'application intégrée~; il administre aussi le pare-feu et rédige
chaque année le rapport de sécurité présenté au directeur général,
M.~ETEME. Mme~NKENG est contrôleuse interne. Il n'y a pas d'audit interne.
\begin{enumerate}[(a)]
    \item Placez chaque acteur dans le modèle des trois lignes, en
          distinguant les rôles qu'il tient effectivement.
    \item Identifiez les deux conflits de rôles les plus graves et
          expliquez le risque que chacun fait courir.
    \item Proposez une organisation réaliste pour une entreprise de cette
          taille, sans créer de direction de l'audit interne.
\end{enumerate}
\end{boxexo}

\begin{boxexo}[title={\ding{45}\quad Exercice~1.3. Corriger une matrice RACI \niveaumoyen}]
Une première version de la matrice de la BANQUE SANAGA contenait les lignes
suivantes (colonnes~: CA, DG, CS, RSSI, DSI, PR, AI)~:
\begin{itemize}
    \item Fixer le périmètre du SMSI~: --, A, --, R, A, --, --~;
    \item Accepter un risque résiduel au-delà du seuil~: --, I, --, A, --,
          C, --~;
    \item Exploiter les mesures techniques~: --, --, --, R, R, --, --~;
    \item Surveiller l'efficacité des mesures~: --, --, I, A, R, --, --~;
    \item Auditer le dispositif de sécurité~: --, I, --, A, --, --, R.
\end{itemize}
\begin{enumerate}[(a)]
    \item Pour chaque ligne, relevez les règles de vérification violées.
    \item Proposez une version corrigée de chaque ligne.
    \item Quelle ligne, si elle était restée en l'état, aurait reproduit le
          défaut à l'origine de l'incident de septembre~2026~?
\end{enumerate}
\end{boxexo}

\begin{boxexo}[title={\ding{45}\quad Exercice~1.4. Un an après \niveauexpert}]
En octobre 2027, Mme~EKAMBI reprend les huit questions du Tableau~1.7. Le
conseil a approuvé un périmètre étendu aux agences et fixé une appétence~;
le registre a une colonne propriétaire, mais trois risques sur
vingt n'ont pas de propriétaire désigné~; le comité de sécurité s'est
réuni trois fois, procès-verbaux à l'appui~; deux rapports ont été
présentés au comité des risques~; l'audit interne a programmé une mission
sur les agences pour 2028. La politique et le rattachement du RSSI n'ont
pas changé.
\begin{enumerate}[(a)]
    \item Cotez les huit questions et justifiez chaque cote.
    \item Calculez $\mu$, $U$ et $H$, et comparez avec 2025.
    \item La question G4 est-elle bien posée~? Proposez une formulation qui
          permette une cote plus informative, et dites ce qu'elle changerait.
    \item Rédigez, en une demi-page, le mandat du comité de sécurité~:
          composition, fréquence, attributions, ce qu'il décide et ce qu'il
          remonte au conseil.
\end{enumerate}
\end{boxexo}

\bigskip

\begin{boxinfo}
\textbf{Pour aller plus loin.}
\begin{itemize}
    \item Le modèle des trois lignes de l'IIA~\cite{iia2020troislignes}.
    \item L'ISO/IEC~27014:2020~\cite{iso27014_2020}, sur la gouvernance de
          la sécurité de l'information.
    \item Le cadre de cybersécurité du NIST, version~2.0~\cite{nist2024csf},
          et sa fonction \emph{Gouverner}.
    \item L'ISO/IEC~27001:2022~\cite{iso27001_2022}, article~5
          (leadership).
\end{itemize}
\end{boxinfo}

\begin{boxinfo}
\textbf{Transition.} Gouverner, c'est décider du niveau de risque que l'on
accepte. Encore faut-il savoir ce qu'est un risque cyber, comment on le
décrit et comment on le cote, et reconnaître quand une cote n'est qu'une
opinion. C'est l'objet du Chapitre~\ref{chap:risque}.
\end{boxinfo}

        % ============================================================================
% Fichier  : partie1/P01_C02_risque_cyber_incertitude.tex
% Titre    : Le Risque Cyber : Notions, Mesure, Incertitude
% Auteur   : MINKA MI NGUIDJOÏ Thierry Emmanuel
% Version  : 1.0 -- novembre 2026
% Corrigés : corriges/Corr_C02.tex
% ============================================================================
% Faits du cas (bible des cas) : BANQUE SANAGA (C01, P00) ; R-AG-01 ; FAIR
% 0,2 à 0,6 événement par an ; coût de l'incident 410 M FCFA, 3 jours ;
% prestataire DIGICORE sans clause d'audit ; aucun contrôle sur les
% changements (GCI). LiZbiZ : registre R1 à R8 et appétence (ACSI C03,
% repris sans modification), constats de l'audit des quatre centres.
% Faits nouveaux (reportés dans la bible) : registre 2025 de la SANAGA, six
% scénarios R-DC-01 à R-DC-06 ; revue des 12 cotes : 6 OUI / 2 NON / 4 NEK ;
% appétence provisoire adoptée en novembre 2026 (convention du corpus).
% Échelles : convention du corpus (ACSI, Tableau 3.2), précisée ici par des
% bornes de fréquence. Convention de réduction par étape du scénario :
% propre à SEC08.
% Sources : ISO 31000:2018 (définition reformulée) ; ISO/IEC 27005:2022 ;
% guide EBIOS RM (ANSSI, 2018) ; Cox (2008) sur les matrices de risque.
% ============================================================================

\chapter{Le Risque Cyber : Notions, Mesure, Incertitude}
\label{chap:risque}
\chaptermark{Le risque cyber}

\epigraph{%
    « Une cote sans preuve est une opinion\\
    qui a pris l'apparence d'un chiffre. »%
}{Principe du cours}

\niveauChapitre{Fondamental}{%
    Maîtrise du Chapitre~\ref{chap:grc}%
}

\begin{boxobjectifs}
\begin{itemize}
    \item Définir risque, actif, menace, vulnérabilité, source de risque,
          objectif visé, événement redouté et scénario.
    \item Expliquer ce qui distingue le risque cyber des autres risques
          opérationnels.
    \item Coter vraisemblance, impact et criticité sur une échelle commune,
          et en connaître les limites.
    \item Distinguer risque inhérent, risque résiduel et appétence, et ne
          réduire un risque que par des mesures testées.
    \item Distinguer l'incertitude aléatoire de l'incertitude épistémique.
    \item Reconnaître une cote sans preuve comme une incertitude (\NEK{}) et
          mesurer la base de preuve d'un registre par $\mu$, $U$ et $H$.
\end{itemize}
\end{boxobjectifs}

\bigskip

% ============================================================================
\section*{Avant de commencer : sur quoi repose ce « 1 »~?}
% ============================================================================

Douala, novembre 2026. Mme~ABENA relit le registre des risques que le
comité des risques avait approuvé un an plus tôt\footnote{Comme dans tout
le cours, la BANQUE SANAGA, ses dirigeants, les dates et les montants sont
fictifs~; les normes et les méthodes citées sont exactes.}. Il compte six
scénarios, tous sur le centre de données. L'un d'eux, « intrusion par
Internet », porte une vraisemblance de 1 sur 4~: rare. La RSSI demande au
responsable de l'exploitation d'où vient ce chiffre.

« Nous n'avons jamais été attaqués », répond-il. Elle lui demande alors qui
lit les journaux du pare-feu d'accès à Internet. Personne~: installé au
siège, hors du périmètre certifié, il conserve ses journaux trente jours,
puis les efface~; la journalisation du centre de données ne voit que ce
qui a franchi ce pare-feu. La banque ne savait pas si elle avait été
attaquée. Le « 1 » ne mesurait pas une fréquence d'attaques~; il mesurait
l'absence de regard.

\medskip

Le prologue a montré qu'un risque jamais apprécié n'est pas faible, mais
inconnu. Ce chapitre va plus loin~: même un risque apprécié et coté peut
reposer sur une inconnue. Il définit le vocabulaire du risque cyber, montre
comment on le cote, puis apprend à coter la cote elle-même.

\begin{boxalert}
\textbf{Violation typique.} Coter la vraisemblance d'une attaque à partir
de l'absence d'incident connu, alors que rien n'aurait permis de détecter
l'incident.

\smallskip
\textbf{Règle~:} l'absence de preuve d'attaque n'est une preuve d'absence
que si l'on regardait. Une cote fondée sur un silence non surveillé est
\NEK{}.
\end{boxalert}

\separateur

% ============================================================================
\section{Définir le Risque Cyber}
\label{sec:c02-definir}
% ============================================================================

\subsection{L'effet de l'incertitude sur les objectifs}

\begin{boxdef}
\termecle{Risque}\index{risque}~: effet de l'incertitude sur les objectifs.
L'effet est un écart, positif ou négatif, par rapport à ce qui est
attendu. On décrit un risque par ses sources, les événements possibles,
leurs conséquences et leur vraisemblance~\cite{iso31000_2018}.
\end{boxdef}
{\small\itshape Définition reformulée d'après l'ISO~31000:2018.}

\medskip

Un risque se rattache toujours à un \emph{objectif}. « Les postes de
Ngaoundéré sont chiffrés » n'est un risque qu'au regard d'un objectif, par
exemple «~servir les clients au guichet tous les jours ouvrés~». C'est
pourquoi le conseil, qui fixe les objectifs, fixe aussi le niveau de risque
acceptable (Chapitre~\ref{chap:grc}).

\begin{boxdef}
\termecle{Risque cyber}\index{risque cyber}~: risque dont la source est
une action, malveillante ou accidentelle, sur les systèmes d'information
ou les données, et dont les conséquences touchent les activités de
l'organisation~: interruption, perte financière, fuite d'information,
sanction, atteinte à la réputation.
\end{boxdef}

\subsection{Les deux portes d'entrée de l'ISO/IEC~27005}

L'ISO/IEC~27005:2022, guide de la gestion des risques de sécurité de
l'information, admet deux manières d'identifier les
risques~\cite{iso27005_2022}~: à partir des \emph{actifs} (on recense ce
qui a de la valeur, puis les menaces et les vulnérabilités qui le
concernent), ou à partir des \emph{événements} (on part des scénarios
stratégiques, des sources de risque et de leurs conséquences sur
l'activité). La première est exhaustive mais lourde~; la seconde va
droit aux scénarios qui comptent, au risque d'en oublier. Le
Chapitre~\ref{chap:iso27005} les met en œuvre~; ce chapitre en pose le
vocabulaire commun.

% ============================================================================
\section{Les Composantes d'un Risque}
\label{sec:c02-composantes}
% ============================================================================

\subsection{Huit notions}

\begin{boxdef}
\termecle{Valeur métier}\index{valeur métier} (ou actif essentiel)~:
activité, processus ou information qui a de la valeur pour l'organisation
(le service au guichet, le fichier des clients).

\medskip
\termecle{Bien support}\index{bien support} (ou actif support)~:
composant sur lequel repose une valeur métier~: matériel, logiciel,
réseau, local, personne, prestataire.

\medskip
\termecle{Menace}\index{menace}~: cause potentielle d'un incident (un
rançongiciel, une erreur d'administration, une coupure électrique).

\medskip
\termecle{Vulnérabilité}\index{vulnérabilité}~: faiblesse d'un bien
support qu'une menace peut exploiter (un partage de fichiers ouvert à
tous, un compte sans second facteur).

\medskip
\termecle{Source de risque}\index{source de risque}~: élément, personne ou
groupe susceptible d'engendrer le risque~; pour une attaque, l'attaquant
et ses capacités.

\medskip
\termecle{Objectif visé}\index{objectif visé}~: ce que la source de
risque cherche à obtenir (un gain financier, de l'information, une
nuisance).

\medskip
\termecle{Événement redouté}\index{événement redouté}~: atteinte à une
valeur métier, exprimée du point de vue de l'organisation (indisponibilité
des guichets pendant plusieurs jours).

\medskip
\termecle{Scénario}\index{scénario de risque}~: enchaînement qui mène
une source de risque, par des biens supports et leurs vulnérabilités,
jusqu'à l'événement redouté.
\end{boxdef}

Les notions de valeur métier, bien support, source de risque, objectif
visé et événement redouté sont celles de la méthode EBIOS Risk Manager,
publiée par l'Agence nationale de la sécurité des systèmes d'information
(\sigle{ANSSI}, France)~\cite{anssi2018ebios}. Le cours les retient parce
qu'elles séparent nettement ce que l'on protège (la valeur métier) de ce
par quoi on est attaqué (les biens supports). D'autres textes emploient
d'autres mots pour les mêmes idées~; le Tableau~2.1 donne les plus
courants.

\begin{tableaucours}{L{3.4cm}L{4.0cm}L{4.8cm}}{%
    \tableauHeader{Ce cours} &
    \tableauHeader{EBIOS Risk Manager} &
    \tableauHeader{Autres formulations courantes}
}
Valeur métier & Valeur métier & Actif essentiel, actif primaire \\
Bien support & Bien support & Actif support \\
Source de risque & Source de risque (SR) & Attaquant, menace au sens
large \\
Objectif visé & Objectif visé (OV) & Motivation de l'attaquant \\
Événement redouté & Événement redouté & Événement, conséquence \\
Scénario & Scénario stratégique, puis opérationnel & Scénario de risque,
scénario d'incident \\
Impact & Gravité & Conséquence \\
\end{tableaucours}
\vspace{-0.8em}
\begin{center}{\small\itshape Tableau~2.1. Correspondance des
vocabulaires (simplifiée)}\end{center}

\subsection{Le scénario R-AG-01, décomposé}

\begin{figure}[H]
\centering
\begin{tikzpicture}[
    etape/.style={draw=CouleurPrimaire, thick, rounded corners,
                  fill=CouleurPrimaire!8, text width=2.6cm,
                  minimum height=1.9cm, align=center, font=\scriptsize},
    fleche/.style={->, thick, CouleurSecondaire}]
    \node[etape] (S) at (0,0) {\textbf{Source de risque}\\groupe
        cybercriminel\\[2pt]\textbf{Objectif visé}\\extorsion};
    \node[etape] (B) at (3.3,0) {\textbf{Biens supports}\\messagerie,
        postes, partage de fichiers, comptes};
    \node[etape] (E) at (6.6,0) {\textbf{Événement redouté}\\guichets
        indisponibles plusieurs jours};
    \node[etape] (V) at (9.9,0) {\textbf{Valeur métier}\\service au
        guichet des agences};
    \draw[fleche] (S) -- (B);
    \draw[fleche] (B) -- (E);
    \draw[fleche] (E) -- (V);
    \node[font=\scriptsize, align=center, text width=6cm] at (3.3,-1.6)
        {vulnérabilités~: pas de second facteur,\\partage ouvert à toutes
        les agences, sauvegardes non testées};
\end{tikzpicture}
\caption{Les composantes du scénario R-AG-01}
\label{fig:c02-scenario}
\end{figure}

La décomposition dit où agir. On peut viser la source (on ne la contrôle
pas), les vulnérabilités des biens supports (second facteur,
cloisonnement), ou les conséquences sur la valeur métier (sauvegardes
restaurables, procédure de guichet en mode dégradé). Les mesures se
classent ainsi par \emph{étape du scénario}~: entrée, propagation,
récupération. Cette classification servira à coter le risque résiduel
(Section~\ref{sec:c02-residuel}).

% ============================================================================
\section{Ce qui rend le Risque Cyber Particulier}
\label{sec:c02-specificites}
% ============================================================================

Le risque cyber se gère avec les outils des risques opérationnels, mais il
en déjoue plusieurs hypothèses. Chacune a une conséquence sur la manière
de le coter.

\begin{tableaucours}{L{3.0cm}L{4.8cm}L{4.4cm}}{%
    \tableauHeader{Particularité} &
    \tableauHeader{Ce qu'elle signifie} &
    \tableauHeader{Conséquence pour la cote}
}
Un adversaire qui s'adapte & L'attaquant choisit sa cible et contourne
les mesures connues~; la fréquence dépend de la défense & Une fréquence
passée ne prédit pas la fréquence future \\
Peu d'historique & Les incidents sont rares, mal déclarés, et changent de
forme chaque année & Les cotes reposent souvent sur l'avis d'expert \\
Des incidents invisibles & Un attaquant peut rester des semaines dans un
réseau sans être détecté & L'absence d'incident connu ne vaut que si l'on
surveillait \\
Des effets simultanés & Un même scénario touche d'un coup de nombreux
sites (9 agences en une nuit) & L'impact se cote sur le scénario entier,
pas par site \\
Des pertes très dispersées & La plupart des incidents coûtent peu~;
quelques-uns coûtent énormément & Une moyenne trompe~; il faut un
intervalle (Chapitre~\ref{chap:fair}) \\
\end{tableaucours}
\vspace{-0.8em}
\begin{center}{\small\itshape Tableau~2.2. Cinq particularités du risque
cyber}\end{center}

\begin{boxinfo}
\textbf{Un risque qui dépend de la défense.} Une coupure électrique ne
devient pas plus fréquente parce que la banque a installé un groupe
électrogène. Une attaque, si~: quand une cible se protège, l'attaquant en
cherche une autre, ou une autre porte. La vraisemblance d'un scénario cyber
n'est donc jamais une donnée de la nature~; c'est une hypothèse sur un
adversaire, qui se révise.
\end{boxinfo}

% ============================================================================
\section{Coter un Risque}
\label{sec:c02-coter}
% ============================================================================

\subsection{Une échelle commune}

Pour comparer des risques de natures différentes, et pour que le conseil
puisse arbitrer, on les cote sur une même échelle. Le cours reprend la
convention à quatre niveaux du corpus, et la précise par des bornes de
fréquence~; chaque organisation fixe la sienne.

\begin{tableaucours}{cL{4.2cm}L{7.4cm}}{%
    \tableauHeader{Niv.} &
    \tableauHeader{Vraisemblance} &
    \tableauHeader{Impact (le critère le plus défavorable l'emporte)}
}
1 & Rare~: moins d'une fois en 10~ans & Mineur~: perte inférieure à
5~millions FCFA, interruption de moins de 4~heures \\
2 & Possible~: une fois en 3 à 10~ans & Modéré~: de 5 à 50~millions FCFA,
interruption de moins de 24~heures \\
3 & Probable~: une fois en 1 à 3~ans & Majeur~: de 50 à 500~millions FCFA,
interruption de 1 à 3~jours, réserve du commissaire aux comptes ou compte
rendu exigé par le régulateur \\
4 & Quasi certain~: au moins une fois par an & Critique~: plus de
500~millions FCFA, interruption de plus de 3~jours, sanction d'un
régulateur \\
\end{tableaucours}
\vspace{-0.8em}
\begin{center}{\small\itshape Tableau~2.3. Échelles de cotation
(convention du cours)}\end{center}

\begin{boxdef}
\termecle{Criticité}\index{criticité} (ou score)~: produit
$\text{vraisemblance} \times \text{impact}$, de 1 à 16. Elle sert à
classer les risques et à les situer par rapport à l'appétence.
\end{boxdef}

L'incident de septembre 2026 a coûté 410~millions FCFA et fermé les
guichets trois jours~: il se situe au niveau~3 d'impact. Le scénario
redouté de R-AG-01 va plus loin~: une propagation à l'ensemble du réseau,
au-delà de trois jours. Son impact se cote~4. On cote le scénario que l'on
redoute, pas seulement celui que l'on a vécu.

\subsection{L'appétence, traduite sur l'échelle}

En novembre 2026, le comité des risques du conseil adopte, à titre
provisoire, la règle d'appétence commune au corpus~:
\begin{itemize}
    \item une criticité résiduelle de \textbf{9 ou plus} est
          \emph{intolérable} et doit être traitée sans délai~;
    \item de \textbf{5 à 8}, le risque est \emph{traité dans l'année}~;
    \item de \textbf{1 à 4}, il est \emph{accepté sous surveillance}, par
          une décision écrite du propriétaire~;
    \item tout risque d'impact~4 est présenté au comité, quelle que soit sa
          criticité.
\end{itemize}

La dernière règle n'est pas un détail. Elle corrige le principal défaut
des matrices, qu'il faut maintenant examiner.

\subsection{Ce que la matrice cache}

La matrice vraisemblance~$\times$ impact est l'outil le plus répandu de la
gestion des risques. Elle est aussi l'un des plus critiqués. Cox a montré
qu'elle peut classer deux risques dans le mauvais ordre, et donner la même
couleur à des risques dont les pertes attendues diffèrent de plusieurs
ordres de grandeur~\cite{cox2008matrices}. Trois raisons~:
\begin{itemize}
    \item \textbf{des échelles ordinales}~: le niveau~4 n'est pas « deux
          fois » le niveau~2~; leurs bornes (une fois par an, une fois en
          trois ans) sont dans un rapport de 1 à 3 au moins, et les bornes
          d'impact dans un rapport de 1 à 10~;
    \item \textbf{un produit qui écrase}~: un risque fréquent et mineur
          (4~$\times$~1) et un risque rare et critique (1~$\times$~4) ont la
          même criticité, 4~;
    \item \textbf{des frontières arbitraires}~: une fréquence de 0,32 par
          an et une de 0,34 tombent dans deux niveaux différents, une
          fréquence de 0,34 et une de 0,99 dans le même.
\end{itemize}

\begin{boxalert}
\textbf{Violation typique.} Accepter un risque rare et critique parce que
sa criticité tombe en zone verte.

\smallskip
\textbf{Règle~:} la criticité classe, elle ne décide pas. Un risque
d'impact critique se présente toujours à l'instance qui répond des
risques, avec son scénario. Quand la décision dépend de l'ordre de
grandeur des pertes, on quantifie (Chapitre~\ref{chap:fair}).
\end{boxalert}

% ============================================================================
\section{Inhérent, Résiduel, Appétence}
\label{sec:c02-residuel}
% ============================================================================

\subsection{Trois niveaux}

\begin{boxdef}
\termecle{Risque inhérent}\index{risque inhérent}~: niveau du risque
avant la prise en compte des mesures de sécurité.

\medskip
\termecle{Risque résiduel}\index{risque résiduel}~: niveau du risque qui
subsiste compte tenu des mesures réellement en place \emph{et dont
l'efficacité a été établie}.

\medskip
\termecle{Appétence au risque}\index{appétence au risque}~: niveau de
risque que l'organisation accepte pour atteindre ses objectifs, fixé par
l'organe de gouvernance (Chapitre~\ref{chap:grc}).
\end{boxdef}

\subsection{Une convention de réduction}

Pour que deux analystes obtiennent le même résiduel, il faut une règle de
réduction écrite. Le cours retient la suivante, fondée sur les étapes du
scénario.

\begin{boxprocedure}
\textbf{Réduire un risque, étape par étape.}
\begin{itemize}
    \item On classe les mesures selon l'étape du scénario sur laquelle
          elles agissent~: \emph{entrée} (empêcher l'accès initial),
          \emph{propagation} (empêcher l'extension), \emph{récupération}
          (limiter les conséquences).
    \item Une étape couverte par au moins une mesure \textbf{testée et
          efficace} abaisse d'un niveau~: la vraisemblance pour l'entrée et
          la propagation, l'impact pour la récupération.
    \item Une étape ne compte qu'une fois, quel que soit le nombre de
          mesures qui la couvrent~; aucun niveau ne descend sous~1.
    \item Une mesure non testée, ou cotée \NEK{}, n'abaisse rien.
    \item Le papier de travail cite, pour chaque réduction, la mesure et la
          preuve de son test.
\end{itemize}
\end{boxprocedure}

\subsection{R-AG-01, coté}

\begin{boxscenario}{Opération SANAGA. R-AG-01, inhérent et résiduel}
\textbf{Inhérent.} Une banque de 38 agences qui reçoit chaque jour des
courriels de l'extérieur sera visée au moins une fois par an~:
vraisemblance~4. Le scénario redouté touche tout le réseau plus de trois
jours~: impact~4. Criticité inhérente~: $4 \times 4 = 16$.

\medskip
\textbf{Mesures de 2025}, avec les cotes de l'audit interne (prologue,
Tableau~P.3)~:
\begin{itemize}
    \item entrée~: filtrage des pièces jointes et antivirus, \OUI{}~; second
          facteur, \NON{} $\Rightarrow$ étape couverte, $-1$ en
          vraisemblance~;
    \item propagation~: cloisonnement, \NON{} $\Rightarrow$ aucune
          réduction~;
    \item récupération~: sauvegardes hors ligne testées, \NON{}~; plan de
          réponse testé, \NON{} $\Rightarrow$ aucune réduction.
\end{itemize}

\textbf{Résiduel.} Vraisemblance~3, impact~4~: criticité $3 \times 4 = 12$,
intolérable.

\medskip
\textbf{Analyse} (confiance moyenne). Le registre de 2025 aurait dû porter
un risque intolérable sur les agences. Il n'en portait aucun, parce que le
scénario n'y figurait pas. La cote résiduelle de 12 n'est pas une
prédiction~: elle dit que, faute de mesures sur deux étapes sur trois, la
banque dépendait d'une seule barrière, et que cette barrière arrête les
attaques connues, pas les nouvelles.
\end{boxscenario}

% ============================================================================
\section{Une Cote sans Preuve est une Incertitude}
\label{sec:c02-incertitude}
% ============================================================================

\subsection{Deux incertitudes}

Le mot « incertitude » figure dans la définition du risque. Il recouvre
deux choses que la gestion des risques confond souvent.

\begin{boxdef}
\termecle{Incertitude aléatoire}\index{incertitude aléatoire}~: part de
l'incertitude qui tient au hasard des événements~: même parfaitement
informé, on ne sait pas si l'attaque aura lieu cette année. Elle ne se
réduit pas~; elle se décrit par une probabilité.

\medskip
\termecle{Incertitude épistémique}\index{incertitude épistémique}~: part
de l'incertitude qui tient à ce que l'on ne sait pas de l'état du système
ou de la menace~: on ignore si les journaux contiennent des traces
d'intrusion. Elle se réduit en collectant des preuves.
\end{boxdef}

\begin{tableaucours}{L{2.8cm}L{4.6cm}L{4.8cm}}{%
    \tableauHeader{} &
    \tableauHeader{Aléatoire} &
    \tableauHeader{Épistémique}
}
Porte sur & Les événements futurs & L'état présent du système et de la
menace \\
Se réduit~? & Non & Oui, par des preuves \\
Se mesure par & Une probabilité, un intervalle de pertes & Le taux
d'incertitude $U$ et l'entropie $H$ des réponses \\
À la BANQUE SANAGA & Le groupe attaquera-t-il en 2027~? & Les journaux du
pare-feu contiennent-ils des traces d'intrusion~? \\
\end{tableaucours}
\vspace{-0.8em}
\begin{center}{\small\itshape Tableau~2.4. Incertitude aléatoire et
incertitude épistémique}\end{center}

Une vraisemblance cotée « 1 » sans base de preuve mélange les deux~: elle
prétend décrire le hasard (l'aléatoire), alors qu'elle traduit une
ignorance (l'épistémique). Le programme de l'auditologie fait de cette
distinction l'une de ses thèses~\cite{minka2026auditologie,minka2026nek}~:
avant de mesurer un risque, on mesure ce que l'on sait de lui.

\subsection{Coter la cote}

Chaque cote d'un registre repose sur une \termecle{base de
preuve}\index{base de preuve}~: des mesures internes (journaux, incidents
recensés, tests), des données externes (statistiques du secteur, alertes
des autorités), un avis d'expert motivé, ou rien. On cote donc la cote
elle-même~:
\begin{itemize}
    \item \OUI{}~: la base de preuve existe et soutient la cote (un avis
          d'expert motivé et écrit compte, et se signale comme tel)~;
    \item \NON{}~: la base de preuve existe et contredit la cote~;
    \item \NEK{}~: aucune base de preuve, ou une base qui ne permet pas de
          conclure (une déclaration, un silence non surveillé).
\end{itemize}

\begin{boxscenario}{Opération SANAGA. Le registre de 2025, revu}
Mme~ABENA revoit les douze cotes (six vraisemblances, six impacts) des six
scénarios du registre approuvé en novembre 2025.

\medskip
\begin{tableaucours}{cL{3.4cm}ccL{3.4cm}cc}{%
    \tableauHeader{Réf.} &
    \tableauHeader{Scénario} &
    \tableauHeader{V} &
    \tableauHeader{I} &
    \tableauHeader{Base de la vraisemblance} &
    \tableauHeader{V} &
    \tableauHeader{I}
}
DC-01 & Coupure électrique prolongée & 2 & 4 & Journal des coupures sur
cinq ans & \OUI{} & \OUI{} \\
DC-02 & Défaillance du prestataire DIGICORE & 2 & 3 & « Prestataire
fiable »~; aucune revue, pas de droit d'audit & \NEK{} & \OUI{} \\
DC-03 & Intrusion par Internet & 1 & 4 & « Jamais attaqués »~; journaux
du pare-feu d'accès non lus & \NEK{} & \OUI{} \\
DC-04 & Erreur lors d'un changement & 1 & 3 & Sept incidents liés à des
changements en 2025 & \NON{} & \OUI{} \\
DC-05 & Copie de données par un administrateur & 1 & 4 & « Confiance dans
les équipes » & \NEK{} & \NEK{} \\
DC-06 & Déni de service des services en ligne & 1 & 3 & Quatre attaques
signalées par le fournisseur d'accès en 2025 & \NON{} & \OUI{} \\
\end{tableaucours}

\smallskip
{\small Colonnes~: V et I cotés (1 à 4)~; puis cote de la base de preuve
de la vraisemblance (V) et de l'impact (I). Les impacts reposent sur
l'analyse d'impact sur l'activité de 2024, sauf pour DC-05~: personne n'a
établi quelles données un administrateur pouvait copier.}
\end{boxscenario}

\begin{boxlecture}{Base de preuve du registre de 2025}
\textbf{Constatation.} Sur douze cotes~: 6~\OUI{}, 2~\NON{}, 4~\NEK{}.
\[
\mu = 0{,}50~;\quad U = \tfrac{4}{12} \approx 0{,}33~;\quad
H = -\left(\tfrac{1}{2}\log_2\tfrac{1}{2} + \tfrac{1}{6}\log_2\tfrac{1}{6}
  + \tfrac{1}{3}\log_2\tfrac{1}{3}\right) \approx 0{,}50 + 0{,}43 + 0{,}53
  \approx 1{,}46~\text{bit}.
\]
Par colonne~: impacts 5/0/1 ($\mu \approx 0{,}83$, $U \approx 0{,}17$,
$H \approx 0{,}65$)~; vraisemblances 1/2/3 ($\mu \approx 0{,}17$,
$U = 0{,}50$, $H \approx 1{,}46$).

\medskip
\textbf{Analyse} (confiance élevée). Le registre est solide sur les
conséquences et fragile sur les fréquences~: une seule vraisemblance sur
six est fondée. Les deux \NON{} sont les plus graves~: DC-04 et DC-06
étaient contredits par des pièces que la banque détenait. DC-04 annonçait
la bascule ratée de mars 2026. Enfin, ce tableau ne mesure que ce que le
registre contient~: le scénario des agences, absent, échappe même à cette
revue. C'est le \NEK{} du prologue, au niveau du périmètre.

\medskip
\hfill\lecture{0{,}50}{0{,}33}{1{,}46}
\end{boxlecture}

\subsection{Quand la cote est un intervalle}

Une estimation honnête de la vraisemblance prend souvent la forme d'un
intervalle. L'atelier de quantification du cours estime la fréquence de
R-AG-01 entre 0,2 et 0,6 événement par an (Chapitre~\ref{chap:fair}).
Sur l'échelle du Tableau~2.3, cet intervalle chevauche deux niveaux~: de 0,2
à 0,33, niveau~2~; de 0,33 à 0,6, niveau~3. Le résiduel vaut donc 8 ou 12
selon le niveau retenu~: « à traiter dans l'année » ou « intolérable ».

\begin{boxinfo}
\textbf{Une décision qui dépend de l'incertitude.} Quand l'intervalle
d'une cote traverse un seuil d'appétence, la décision dépend de ce que
l'on ignore. Deux réponses sont légitimes~: retenir le niveau le plus
défavorable en le signalant, ou réduire l'incertitude avant de décider
(une revue des journaux, une simulation d'attaque, une quantification).
Choisir silencieusement le niveau favorable ne l'est pas.
\end{boxinfo}

\subsection{La fiche de risque, avec sa base de preuve}

\begin{boxobjet}{R-DC-03 · Intrusion par Internet (fiche révisée)}
\textbf{Valeur métier~:} les opérations du système central (comptes,
virements).

\textbf{Biens supports~:} pare-feu, serveurs exposés, comptes
d'administration à distance.

\textbf{Source de risque et objectif visé~:} groupe cybercriminel, fraude
ou extorsion.

\textbf{Événement redouté~:} prise de contrôle du système central.

\textbf{Vraisemblance~:} 1 en 2025. Base~: absence d'incident connu,
journaux du pare-feu d'accès effacés à trente jours sans lecture. Cote de
la base~: \NEK{}.

\textbf{Impact~:} 4. Base~: analyse d'impact sur l'activité de 2024. Cote
de la base~: \OUI{}.

\textbf{Pour réduire l'incertitude~:} centraliser et conserver les
journaux un an~; faire analyser trois mois de journaux par un tiers~;
commander un test d'intrusion externe (UE MSN-2S3-SEC03).

\textbf{Décision provisoire~:} faute de fréquence observée, la
vraisemblance est recalculée par la convention de réduction~: inhérent~4
(serveurs exposés à Internet), entrée couverte par des mesures testées au
centre (second facteur, correctifs~: prologue, Tableau~P.3), $-1$~;
propagation couverte (cloisonnement testé), $-1$. Vraisemblance~2,
impact~4~: criticité~8, à traiter dans l'année, et présentée au comité en
raison de l'impact~4. À confirmer par la revue des journaux.
\end{boxobjet}

% ============================================================================
\section{Le Cadre de Référence}
\label{sec:c02-cadre}
% ============================================================================

L'article~26 de la loi n°~2010/012, cité au Chapitre~\ref{chap:grc}, impose
aux exploitants de systèmes d'information des dispositifs normalisés pour
\emph{identifier} et \emph{évaluer} en permanence les risques de
sécurité~\cite{loi2010012}. Pour une banque, le règlement COBAC R-2016/04
range les risques liés aux systèmes d'information parmi les risques que le
dispositif de contrôle interne doit couvrir~\cite{cobac2016r04}.

\begin{boxinfo}
\textbf{Analyse} (lecture du cours). Ni la loi ni le règlement ne disent
comment coter. Mais une évaluation que l'on ne peut pas justifier n'en est
pas une~: devant le superviseur comme devant un juge, une cote se défend
par sa base de preuve. Le registre qui distingue ce qu'il sait de ce qu'il
suppose est plus défendable que celui qui affiche des chiffres sans
source. Le Chapitre~\ref{chap:droit} revient sur ces obligations.
\end{boxinfo}

% ============================================================================
\section{Ce que la Cote demande à chaque Métier}
\label{sec:c02-metiers}
% ============================================================================

\begin{boxmetier}{Gestionnaire des risques cyber}
\textbf{Documenter chaque cote}
\begin{enumerate}
    \item Écrire, pour chaque vraisemblance et chaque impact, la base de
          preuve et sa date.
    \item Coter la base~: \OUI{}, \NON{} ou \NEK{}~; ne jamais transformer
          un \NEK{} en « faible ».
    \item Donner un intervalle quand l'estimation le justifie, et signaler
          tout intervalle qui traverse un seuil d'appétence.
    \item Ne réduire le résiduel que par des mesures testées, selon la
          convention écrite.
\end{enumerate}
\end{boxmetier}

\begin{boxmetier}{Auditeur en cybersécurité}
\textbf{Auditer le registre, pas seulement les mesures}

L'auditeur prélève des cotes et demande leur base. Une vraisemblance sans
base est une constatation~; une vraisemblance contredite par une pièce
détenue par l'organisation est une constatation grave, parce qu'elle
montre que l'information existait et n'a pas été utilisée.
\end{boxmetier}

\begin{boxmetier}{RSSI}
\textbf{Transformer les \NEK{} en travaux}

Chaque \NEK{} du registre appelle une action qui produit une preuve~:
lecture de journaux, test de restauration, revue d'un prestataire, test
d'intrusion. Le RSSI les classe par l'importance de la décision qu'ils
débloquent~: d'abord les cotes dont l'intervalle traverse un seuil.
\end{boxmetier}

\begin{boxmetier}{Responsable GRC}
\textbf{Une échelle, une convention}

Le responsable GRC fait approuver par le conseil l'échelle, l'appétence et
la convention de réduction, puis veille à ce que tous les registres de
l'organisation les appliquent. Deux registres qui cotent différemment ne
peuvent pas être comparés, donc pas arbitrés.
\end{boxmetier}

\separateur

\begin{boxresume}
\textbf{Ce qu'il faut retenir de ce chapitre~:}
\begin{itemize}
    \item \textbf{Un risque} est l'effet de l'incertitude sur un objectif~;
          un scénario relie une source de risque, par des biens supports
          vulnérables, à un événement redouté sur une valeur métier.
    \item \textbf{Le risque cyber} dépend d'un adversaire qui s'adapte, a
          peu d'historique, peut rester invisible, frappe plusieurs sites à
          la fois et produit des pertes très dispersées.
    \item \textbf{La criticité classe, elle ne décide pas}~: les matrices
          écrasent les écarts~; tout impact critique remonte à l'instance
          qui répond des risques.
    \item \textbf{Le résiduel} ne se réduit que par des mesures testées,
          selon une convention écrite~: une étape du scénario couverte
          abaisse d'un niveau.
    \item \textbf{Deux incertitudes}~: l'aléatoire se décrit, l'épistémique
          se réduit par des preuves.
    \item \textbf{On cote la cote}~: \OUI{}, \NON{} ou \NEK{} selon sa base
          de preuve. Le registre 2025 de la SANAGA~: $\mu = 0{,}50$,
          $U \approx 0{,}33$, et une seule vraisemblance fondée sur six.
\end{itemize}
\end{boxresume}

\begin{boxinfo}
\textbf{Points de contrôle.} Vous avez compris ce chapitre si vous savez~:
\begin{itemize}
    \item décomposer un scénario en source, objectif visé, biens supports,
          vulnérabilités, événement redouté et valeur métier~;
    \item calculer une criticité inhérente et résiduelle avec la convention
          de réduction~;
    \item montrer, sur un exemple, ce qu'une matrice de risques cache~;
    \item coter la base de preuve d'une vraisemblance, et dire ce qu'il
          faudrait pour la rendre \OUI{}.
\end{itemize}
\end{boxinfo}

% ============================================================================
\section*{Exercices}
% ============================================================================

\begin{boxexo}[title={\ding{45}\quad Exercice~2.1. Le vocabulaire du scénario \niveaubase}]
Associez chacun des éléments suivants, tirés de l'incident de septembre
2026, à l'une des notions du chapitre (valeur métier, bien support,
menace, vulnérabilité, source de risque, objectif visé, événement redouté,
scénario).
\begin{enumerate}[(a)]
    \item Un groupe cybercriminel organisé.
    \item Obtenir le paiement d'une rançon.
    \item Le partage de fichiers ouvert à toutes les agences.
    \item Les postes de travail des agences.
    \item Le service au guichet.
    \item La fermeture des guichets pendant plusieurs jours.
    \item Un rançongiciel.
    \item Courriel piégé, vol d'identifiants, déplacement par le partage,
          chiffrement des postes.
\end{enumerate}
\end{boxexo}

\begin{boxexo}[title={\ding{45}\quad Exercice~2.2. Six mois après \niveaumoyen}]
En mars 2027, la situation de R-AG-01 est la suivante~: le second facteur
est déployé et testé sur tous les comptes~; le partage de fichiers est
cloisonné par agence, mais le cloisonnement n'a pas encore été testé~; les
sauvegardes hors ligne ont été restaurées avec succès lors d'un test.
\begin{enumerate}[(a)]
    \item Calculez la criticité résiduelle avec la convention du chapitre
          (inhérent~: $4 \times 4$). Quelle décision l'appétence impose-t-elle~?
    \item Que deviendrait le résiduel si le cloisonnement était testé avec
          succès~?
    \item Le second facteur n'a rien changé au calcul. Expliquez pourquoi,
          et dites si c'est un défaut de la convention ou une qualité.
\end{enumerate}
\end{boxexo}

\begin{boxexo}[title={\ding{45}\quad Exercice~2.3. Ce que cache la matrice \niveaumoyen}]
Le registre de la banque contient deux risques~:
\begin{itemize}
    \item A~: erreurs de saisie au guichet, vraisemblance~4, impact~1~;
    \item B~: fraude massive au virement par compromission du système de
          paiement, vraisemblance~1, impact~4.
\end{itemize}
\begin{enumerate}[(a)]
    \item Calculez leur criticité. Que conclurait un comité qui lit
          seulement la criticité~?
    \item Les erreurs A se produisent douze fois par an et coûtent
          1~million FCFA chacune~; la fraude B a une chance sur vingt de se
          produire chaque année et coûterait 2~milliards FCFA. Calculez la
          perte annuelle attendue de chacun.
    \item Expliquez l'écart entre (a) et (b) à l'aide des trois défauts des
          matrices vus dans le chapitre.
    \item Quelle règle de l'appétence de la banque évite la mauvaise
          décision~?
\end{enumerate}
\end{boxexo}

\begin{boxexo}[title={\ding{45}\quad Exercice~2.4. La base de preuve du registre de LiZbiZ \niveauexpert}]
Mme~NKENG reprend le registre R1 à R8 de LiZbiZ (cotes de probabilité
inchangées) et recherche la base de chaque probabilité~:
\begin{itemize}
    \item R1, ajustement frauduleux du stock~: écart de 186~millions FCFA
          constaté, ajustements identifiés~;
    \item R2, fournisseurs fictifs~: fraude TRANSLOG établie~;
    \item R3, compte d'un agent parti~: « on désactive les comptes au
          départ », déclaration du responsable informatique, sans
          extraction~;
    \item R4, arrêt prolongé du centre de Garoua~: journal des coupures sur
          trois ans, un arrêt de plus d'un jour~;
    \item R5, perte de données faute de sauvegarde restaurable~: « les
          sauvegardes tournent chaque nuit », aucun test de restauration~;
    \item R6, fuite de données des clients~: aucune information~;
    \item R7, erreurs de facturation~: trois erreurs par mois relevées par
          le contrôle de gestion~;
    \item R8, rançongiciel (probabilité~2)~: « nous n'avons jamais été
          touchés »~; l'audit des quatre centres a pourtant relevé des
          comptes par défaut, des services en clair et des flux non
          autorisés.
\end{itemize}
\begin{enumerate}[(a)]
    \item Cotez chaque base \OUI{}, \NON{} ou \NEK{}, en justifiant.
    \item Calculez $\mu$, $U$ et $H$.
    \item Proposez une nouvelle cote de probabilité pour R8, recalculez sa
          criticité résiduelle (impact~4, aucune mesure testée) et dites ce
          que l'appétence impose.
    \item Rédigez pour Mme~NKENG une note de dix lignes~: quelles cotes
          rétablir en priorité, par quelles preuves, et pourquoi dans cet
          ordre.
\end{enumerate}
\end{boxexo}

\bigskip

\begin{boxinfo}
\textbf{Pour aller plus loin.}
\begin{itemize}
    \item L'ISO~31000:2018~\cite{iso31000_2018} et l'ISO/IEC~27005:2022~\cite{iso27005_2022},
          pour le vocabulaire et le processus.
    \item Le guide EBIOS Risk Manager~\cite{anssi2018ebios}, pour les notions
          de valeur métier, bien support et source de risque.
    \item L'article de Cox sur les matrices de risque~\cite{cox2008matrices}.
    \item L'article sur l'incertitude épistémique~\cite{minka2026nek}, qui
          fonde la cote \NEK{}.
\end{itemize}
\end{boxinfo}

\begin{boxinfo}
\textbf{Transition.} La Partie~I a posé qui décide (Chapitre~\ref{chap:grc})
et ce que l'on décide (ce chapitre). La Partie~II passe aux méthodes~:
comment conduire une appréciation des risques complète, de l'établissement
du contexte à l'acceptation du résiduel. C'est l'objet du
Chapitre~\ref{chap:iso27005}.
\end{boxinfo}


    % ======================================================================
    % PARTIE 2 — MÉTHODES : GÉRER LES RISQUES, CONDUIRE L'AUDIT
    % Standards de référence : ISO/IEC 27005:2022, EBIOS Risk Manager, FAIR, ISO 19011, ISO/IEC 27007 ; articles A0 et A1.
    % Niveau : fondamental à intermédiaire
    % ======================================================================
    \part{Méthodes : Gérer les Risques, Conduire l'Audit}
    \label{part:methode}

        % ============================================================================
% Fichier  : partie2/P02_C03_iso31000_iso27005.tex
% Titre    : Gérer les Risques : ISO 31000 et ISO/IEC 27005
% Auteur   : MINKA MI NGUIDJOÏ Thierry Emmanuel
% Version  : 1.0 -- novembre 2026
% Corrigés : corriges/Corr_C03.tex
% ============================================================================
% Faits du cas (bible des cas) : BANQUE SANAGA, R-AG-01 (inhérent 16,
% résiduel 12), convention de réduction par étape, appétence provisoire de
% novembre 2026, registre 2025 DC-01 à DC-06, DIGICORE sans clause d'audit,
% pare-feu d'accès au siège (journaux 30 jours). Ex. 2.2 : en mars 2027,
% second facteur testé, cloisonnement non testé, sauvegardes restaurées.
% Faits nouveaux (reportés dans la bible) : critères d'acceptation de
% décembre 2026 ; identification par les deux approches (GAB sous système
% obsolète, droits d'administration locale des chefs d'agence) ; plan de
% traitement A1 à A10 (193 M FCFA pour R-AG-01) ; état en septembre 2027 :
% 5 OUI / 2 NON / 3 NEK ; décision d'acceptation de R-DC-02.
% Sources : ISO 31000:2018 (principes, cadre, processus, options de
% traitement : paraphrasés) ; ISO/IEC 27005:2022 (deux approches, critères :
% paraphrasés, sans numéro de paragraphe) ; ISO/IEC 27001:2022, art. 6.1.3
% (paraphrasé) ; loi n° 2010/012, art. 26 al. 1 et 29 al. 1 (texte français
% officiel, Ann. D d'ACSI).
% ============================================================================

\chapter{Gérer les Risques : ISO 31000 et ISO/IEC 27005}
\label{chap:iso27005}
\chaptermark{ISO 31000 et ISO/IEC 27005}

\epigraph{%
    « Un registre est une photographie.\\
    La gestion des risques est un processus. »%
}{Principe du cours}

\niveauChapitre{Fondamental}{%
    Maîtrise du Chapitre~\ref{chap:risque}%
}

\begin{boxobjectifs}
\begin{itemize}
    \item Situer les principes, le cadre et le processus de l'ISO~31000.
    \item Établir le contexte d'une appréciation~: périmètre, parties
          intéressées, critères de risque et critères d'acceptation.
    \item Appliquer les deux approches d'identification de
          l'ISO/IEC~27005:2022 et montrer qu'elles se complètent.
    \item Choisir une option de traitement et rédiger un plan de
          traitement.
    \item Faire accepter formellement le risque résiduel, et refuser une
          acceptation fondée sur une cote \NEK{}.
    \item Suivre l'avancement d'un plan de traitement par $\mu$, $U$ et
          $H$.
\end{itemize}
\end{boxobjectifs}

\bigskip

% ============================================================================
\section*{Avant de commencer : « Donnez-nous un plan »}
% ============================================================================

Douala, décembre 2026. Le comité des risques du conseil de la BANQUE
SANAGA\footnote{Comme dans tout le cours, la BANQUE SANAGA, ses
dirigeants, ses prestataires, les dates et les montants sont fictifs~; les
normes et les textes cités sont exacts.} a reçu en novembre la cotation de
R-AG-01~: 12, intolérable. Il ne demande plus une cote~; il demande un
plan. Que fait-on, qui le fait, pour quand, à quel coût, et que restera-t-il
ensuite~?

Mme~ABENA mesure l'écart. Le registre de 2025 était une liste de six
scénarios, revue une fois par an. Il n'avait ni critères d'acceptation, ni
plan de traitement, ni suivi. La banque avait un registre~; elle n'avait
pas de processus. La RSSI propose de reprendre la démarche depuis le
début, selon l'ISO~31000 et l'ISO/IEC~27005, sur toute la banque cette
fois.

\medskip

Ce chapitre ouvre la Partie~II, consacrée aux méthodes. Il présente le
processus général de gestion des risques, puis la manière dont l'ISO/IEC
27005 l'applique à la sécurité de l'information, jusqu'à l'acceptation
formelle de ce qui reste.

\begin{boxalert}
\textbf{Violation typique.} Tenir un registre des risques sans plan de
traitement ni décision d'acceptation~: les risques sont décrits, cotés, et
laissés en l'état.

\smallskip
\textbf{Règle~:} chaque risque d'un registre aboutit à une décision
écrite~: un traitement, avec un responsable et une échéance, ou une
acceptation, signée par l'instance compétente. Un risque coté sans
décision n'est pas géré.
\end{boxalert}

\separateur

% ============================================================================
\section{L'ISO 31000 : Principes, Cadre, Processus}
\label{sec:c03-iso31000}
% ============================================================================

\subsection{Une norme générale}

L'ISO~31000:2018 donne des lignes directrices pour le management du
risque, quel qu'il soit~: financier, industriel, juridique,
cyber~\cite{iso31000_2018}. Elle n'est pas certifiable et ne prescrit pas
de méthode. Elle s'articule en trois ensembles~: des \emph{principes}, qui
disent ce qu'est un management du risque efficace~; un \emph{cadre
organisationnel}, qui l'intègre dans la gouvernance~; un
\emph{processus}, qui en décrit les étapes.

\begin{tableaucours}{L{2.6cm}L{9.6cm}}{%
    \tableauHeader{Ensemble} &
    \tableauHeader{Contenu (d'après l'ISO~31000:2018)}
}
Principes & Intégré~; structuré et global~; adapté~; inclusif~;
dynamique~; fondé sur la meilleure information disponible~; attentif aux
facteurs humains et culturels~; en amélioration continue \\
Cadre & Leadership et engagement de la direction~; intégration dans les
activités~; conception~; mise en œuvre~; évaluation~; amélioration \\
Processus & Communication et consultation~; domaine d'application,
contexte et critères~; appréciation (identification, analyse,
évaluation)~; traitement~; suivi et revue~; enregistrement et rapports \\
\end{tableaucours}
\vspace{-0.8em}
\begin{center}{\small\itshape Tableau~3.1. Les trois ensembles de
l'ISO~31000}\end{center}

Le cadre fait le lien avec le Chapitre~\ref{chap:grc}~: sans leadership de
la direction ni intégration dans la gouvernance, le processus tourne à vide.
Le principe de « meilleure information disponible » fait le lien avec le
Chapitre~\ref{chap:risque}~: la norme demande que l'information soit
disponible, et elle ajoute qu'il faut tenir compte de ses limites et de ses
incertitudes. C'est exactement ce que fait la cote de la base de preuve.

\subsection{Le processus}

\begin{figure}[H]
\centering
\begin{tikzpicture}[
    etape/.style={draw=CouleurPrimaire, thick, rounded corners,
                  fill=CouleurPrimaire!8, text width=5.2cm,
                  minimum height=0.8cm, align=center, font=\small},
    transverse/.style={draw=CouleurAccent, thick, rounded corners,
                       fill=CouleurAccent!12, text width=2.4cm,
                       minimum height=5.6cm, align=center, font=\small},
    fleche/.style={->, thick, CouleurSecondaire}]
    \node[etape] (C) at (0,2.4) {Domaine d'application, contexte, critères};
    \node[etape, minimum height=1.6cm] (A) at (0,0.6)
        {\textbf{Appréciation}\\identification\\analyse\\évaluation};
    \node[etape] (T) at (0,-1.2) {Traitement};
    \node[etape] (E) at (0,-2.4) {Enregistrement et rapports};
    \draw[fleche] (C) -- (A);
    \draw[fleche] (A) -- (T);
    \draw[fleche] (T) -- (E);
    \node[transverse] (CC) at (-4.4,0) {Communication\\et\\consultation};
    \node[transverse] (SR) at (4.4,0) {Suivi\\et\\revue};
\end{tikzpicture}
\caption{Le processus de management du risque (d'après l'ISO~31000:2018)}
\label{fig:c03-processus}
\end{figure}

Deux activités encadrent les autres et durent tout le processus~: la
\emph{communication et consultation}, qui associe les parties intéressées
à chaque étape, et le \emph{suivi et revue}, qui vérifie que les décisions
sont appliquées et que les cotes restent valables. À la BANQUE SANAGA, ce
sont précisément les deux activités qui manquaient~: le registre était
rédigé par la DSI seule, et revu une fois par an.

\subsection{L'ISO/IEC 27005, déclinaison pour la sécurité de l'information}

L'ISO/IEC~27005:2022 applique ce processus à la sécurité de l'information
et le relie aux exigences de l'ISO/IEC~27001~\cite{iso27005_2022}. Elle
n'impose pas non plus de méthode~: EBIOS Risk Manager
(Chapitre~\ref{chap:ebios}) ou une méthode maison peuvent la mettre en
œuvre. Sa version de 2022 a refondu le texte de 2018~; elle insiste sur
l'identification des risques à partir des événements, et non plus
seulement des actifs.

% ============================================================================
\section{Établir le Contexte}
\label{sec:c03-contexte}
% ============================================================================

\subsection{Périmètre et parties intéressées}

Le premier livrable fixe le périmètre de l'appréciation et les parties
intéressées. La RSSI retient cette fois toute la banque~: le siège, le
centre de données, le site de secours de Bonabéri et les 38 agences. Elle
recense les parties intéressées et ce qu'elles attendent.

\begin{tableaucours}{L{3.4cm}L{8.8cm}}{%
    \tableauHeader{Partie intéressée} &
    \tableauHeader{Attente en matière de sécurité}
}
Clients & Guichets ouverts, comptes exacts, données protégées \\
Actionnaires et conseil & Pertes contenues dans l'appétence, réputation
préservée \\
COBAC & Dispositif de maîtrise des risques couvrant les systèmes
d'information~; comptes rendus après incident \\
DIGICORE et autres prestataires & Engagements contractuels clairs, accès
encadrés \\
Agents des agences & Outils disponibles, consignes simples \\
\end{tableaucours}
\vspace{-0.8em}
\begin{center}{\small\itshape Tableau~3.2. Les parties intéressées de
l'appréciation}\end{center}

\subsection{Critères de risque, critères d'acceptation}

L'ISO/IEC~27005 demande d'établir, avant toute appréciation, deux familles
de critères~: ceux qui servent à apprécier (les échelles de vraisemblance
et d'impact) et ceux qui servent à accepter~\cite{iso27005_2022}. La
banque dispose des premiers depuis novembre~2026 (Chapitre~\ref{chap:risque},
Tableau~2.3). Il lui manque les seconds~: la règle d'appétence dit quels
risques traiter, elle ne dit pas qui peut accepter quoi.

\begin{boxdef}
\termecle{Critères d'acceptation du risque}\index{critères
d'acceptation}~: règles, fixées à l'avance par l'organe de direction, qui
disent quel niveau de risque résiduel peut être accepté, par qui, à quelles
conditions et pour combien de temps.
\end{boxdef}

En décembre 2026, le comité des risques complète l'appétence par les
critères suivants.

\begin{tableaucours}{cL{3.6cm}L{7.4cm}}{%
    \tableauHeader{Résiduel} &
    \tableauHeader{Qui peut accepter} &
    \tableauHeader{Conditions}
}
1 à 4 & Le propriétaire du risque & Décision écrite~; revue annuelle \\
5 à 8 & Le comité de sécurité, sur proposition du propriétaire & Aucune
mesure supplémentaire proportionnée~; revue semestrielle~; sinon,
traitement dans l'année \\
9 et plus & Personne, sauf le comité des risques du conseil, à titre
temporaire & Plan de traitement daté~; durée de six mois au plus,
renouvelable une fois \\
Toute cote & -- & Pas d'acceptation si une cote de vraisemblance ou
d'impact a une base \NEK{} ou \NON{}, sauf acceptation « en l'état des
connaissances » avec une date de revue \\
\end{tableaucours}
\vspace{-0.8em}
\begin{center}{\small\itshape Tableau~3.3. Critères d'acceptation de la
BANQUE SANAGA (décembre 2026)}\end{center}

\begin{boxinfo}
\textbf{Pourquoi la dernière ligne~?} Accepter un risque, c'est décider en
connaissance de cause de ne pas le traiter davantage. Si la cote repose sur
une ignorance, la décision ne porte pas sur le risque, mais sur une
supposition. La règle n'interdit pas d'accepter~; elle oblige à dire ce que
l'on accepte d'ignorer, et jusqu'à quand.
\end{boxinfo}

% ============================================================================
\section{Apprécier : Deux Approches qui se Complètent}
\label{sec:c03-apprecier}
% ============================================================================

\subsection{Identifier}

L'ISO/IEC~27005:2022 propose deux approches d'identification~: par les
\emph{événements} et par les \emph{actifs}~\cite{iso27005_2022}.

\begin{tableaucours}{L{2.6cm}L{4.8cm}L{4.8cm}}{%
    \tableauHeader{} &
    \tableauHeader{Par les événements} &
    \tableauHeader{Par les actifs}
}
Point de départ & Les sources de risque, leurs objectifs, et les
événements redoutés sur les valeurs métier & L'inventaire des biens
supports, puis leurs menaces et vulnérabilités \\
Niveau & Stratégique~: on raisonne en scénarios & Détaillé~: on raisonne
composant par composant \\
Force & Rapide, centrée sur ce qui compte pour la direction & Exhaustive,
fait apparaître des faiblesses précises \\
Faiblesse & Peut oublier un composant que personne n'a imaginé & Lourde~;
peut noyer l'essentiel dans des centaines de lignes \\
\end{tableaucours}
\vspace{-0.8em}
\begin{center}{\small\itshape Tableau~3.4. Les deux approches
d'identification}\end{center}

\begin{boxscenario}{Opération SANAGA. Ce que chaque approche a trouvé}
\textbf{Par les événements.} Un atelier réunit la direction du réseau
d'agences, la DSI, la direction des risques et la direction financière. Il
retient quatre scénarios stratégiques~: le rançongiciel en agence
(R-AG-01)~; la fraude au virement par compromission d'un poste du
siège~; la défaillance de DIGICORE~; la fuite du fichier des clients.

\medskip
\textbf{Par les actifs.} L'inventaire des biens supports des 38 agences,
mené en parallèle, fait apparaître deux faiblesses qu'aucun scénario
n'avait imaginées~:
\begin{itemize}
    \item les automates de retrait (GAB) de 14 agences fonctionnent sous un
          système d'exploitation qui ne reçoit plus de correctifs~;
    \item dans 31 agences, le chef d'agence dispose des droits
          d'administration locale sur tous les postes, pour « dépanner ».
\end{itemize}

\textbf{Analyse} (confiance élevée). Les deux approches ne se recouvrent
pas. Les scénarios ont donné les priorités de la direction~; l'inventaire
a donné des faiblesses concrètes, dont la seconde alimente directement
R-AG-01 (elle facilite la propagation). On ne choisit pas entre les deux~:
on part des événements pour savoir ce qui compte, et l'on vérifie par les
actifs que rien d'important n'échappe.
\end{boxscenario}

\subsection{Analyser et évaluer}

L'\emph{analyse} cote chaque risque~: vraisemblance, impact, criticité,
avec leur base de preuve (Chapitre~\ref{chap:risque}). L'\emph{évaluation}
compare le résultat aux critères et en tire une priorité~: traiter sans
délai, traiter dans l'année, ou proposer à l'acceptation. Les deux
faiblesses de l'inventaire ne sont pas des risques à part entière~; ce sont
des vulnérabilités qui modifient la cote d'un scénario existant, ou qui en
appellent un nouveau (« compromission d'un GAB »).

% ============================================================================
\section{Traiter}
\label{sec:c03-traiter}
% ============================================================================

\subsection{Les options}

L'ISO~31000 énumère les options de traitement~: éviter le risque en
renonçant à l'activité~; prendre ou augmenter un risque pour saisir une
opportunité~; supprimer la source de risque~; modifier la vraisemblance~;
modifier les conséquences~; partager le risque~; maintenir le risque par
une décision argumentée~\cite{iso31000_2018}. En sécurité de
l'information, on les regroupe souvent en quatre familles.

\begin{tableaucours}{L{2.4cm}L{5.0cm}L{4.8cm}}{%
    \tableauHeader{Famille} &
    \tableauHeader{Principe} &
    \tableauHeader{Limite}
}
Réduire & Modifier la vraisemblance ou les conséquences par des mesures &
Ne réduit le résiduel que si les mesures sont testées \\
Partager & Transférer une partie des conséquences financières (assurance,
contrat) & Ne transfère ni la responsabilité, ni l'atteinte à la
réputation \\
Éviter & Renoncer à l'activité ou au composant qui crée le risque &
Supprime aussi la valeur que l'activité apportait \\
Maintenir & Accepter le résiduel, selon les critères & Exige une décision
écrite et une base de preuve \\
\end{tableaucours}
\vspace{-0.8em}
\begin{center}{\small\itshape Tableau~3.5. Les quatre familles de
traitement}\end{center}

\subsection{Les options pour R-AG-01}

\begin{boxscenario}{Opération SANAGA. Choisir un traitement pour R-AG-01}
\textbf{Maintenir~:} impossible. La criticité résiduelle est de 12~; seul
le comité des risques du conseil pourrait l'accepter, pour six mois, et
avec un plan.

\medskip
\textbf{Éviter~:} supprimer le partage de fichiers commun aux agences.
Option retenue en partie~: le partage commun est remplacé par un espace
par agence. Supprimer la messagerie externe en agence, en revanche,
priverait les agences d'un outil indispensable.

\medskip
\textbf{Partager~:} une assurance contre les risques cyber, si le marché
le permet. Ordre de grandeur obtenu d'un courtier~: prime de 35~millions
FCFA par an, franchise de 100~millions, plafond de 1~milliard. Les
assureurs exigent en général, avant de couvrir, un second facteur et des
sauvegardes hors ligne testées~: l'assurance vient \emph{après} la
réduction, pas à sa place.

\medskip
\textbf{Réduire~:} couvrir les trois étapes du scénario par des mesures
testées (plan ci-dessous).

\medskip
\textbf{Analyse} (confiance moyenne). La combinaison retenue est
«~réduire, éviter en partie, puis partager~». Le résiduel visé, selon la
convention de réduction, est de vraisemblance~2 (entrée et propagation
couvertes) et d'impact~3 (récupération couverte)~: criticité~6. La
convention ne permet pas de descendre plus bas~; un résiduel de~6 relève
du comité de sécurité.
\end{boxscenario}

\subsection{Le plan de traitement}

L'ISO/IEC~27001 demande, en substance, que l'organisation choisisse ses
options de traitement, détermine les mesures nécessaires, les compare à
celles de son annexe~A, établisse une déclaration d'applicabilité, formule
un plan de traitement, et obtienne des propriétaires des risques
l'approbation de ce plan et l'acceptation des risques
résiduels~\cite{iso27001_2022}. La déclaration d'applicabilité est étudiée
au Chapitre~\ref{chap:iso27000}~; voici le plan.

\begin{tableaucours}{cL{4.4cm}cL{2.2cm}cr}{%
    \tableauHeader{Réf.} &
    \tableauHeader{Action} &
    \tableauHeader{Étape} &
    \tableauHeader{Responsable} &
    \tableauHeader{Échéance} &
    \tableauHeader{M FCFA}
}
A1 & Second facteur sur tous les comptes & Entrée & DSI & mars 2027 & 25 \\
A2 & Cloisonnement du partage, un espace par agence & Propag. & DSI &
fév. 2027 & 40 \\
A3 & Test d'intrusion interne vérifiant le cloisonnement & Propag. & RSSI
& juin 2027 & 15 \\
A4 & Sauvegardes hors ligne des serveurs d'agence, test trimestriel &
Récup. & DSI & mars 2027 & 60 \\
A5 & Procédure de guichet en mode dégradé, exercée & Récup. & Dir. réseau
d'agences & avril 2027 & 5 \\
A6 & Sensibilisation des agents des 38 agences & Entrée & RSSI & juin 2027
& 10 \\
A7 & Plan de réponse aux incidents testé par un exercice & Récup. & RSSI &
sept. 2027 & 8 \\
A8 & Journalisation centralisée des postes et du partage & Détection &
DSI & juin 2027 & 30 \\
\textbf{Total} & & & & & \textbf{193} \\
\end{tableaucours}
\vspace{-0.8em}
\begin{center}{\small\itshape Tableau~3.6. Plan de traitement de R-AG-01
(décembre 2026)}\end{center}

Deux remarques. L'action A3 n'ajoute aucune mesure~: elle produit la
preuve sans laquelle A2 n'abaisserait rien. Et l'action A8 ne correspond à
aucune étape de la convention~: elle ne réduit pas la cote, mais elle
réduit l'incertitude, puisque la banque saura désormais si elle est
attaquée.

\begin{boxprocedure}
\textbf{Le contenu minimal d'un plan de traitement.} Pour chaque action~:
\begin{itemize}
    \item le risque traité et l'étape du scénario visée~;
    \item le responsable (une personne, pas un service) et l'échéance~;
    \item le coût estimé et la source du financement~;
    \item la \textbf{preuve attendue} qui permettra de la déclarer
          réalisée (procès-verbal de test, extraction, rapport)~;
    \item l'effet attendu sur la cote, selon la convention écrite.
\end{itemize}
Le plan entier porte l'approbation du propriétaire du risque, et sa date
de revue.
\end{boxprocedure}

% ============================================================================
\section{Accepter le Risque Résiduel}
\label{sec:c03-accepter}
% ============================================================================

\subsection{Une décision, pas un constat}

\begin{boxdef}
\termecle{Acceptation du risque}\index{acceptation du risque}~: décision
argumentée, prise par une personne ou une instance qui en a l'autorité, de
ne pas traiter davantage un risque résiduel, pour une durée donnée et à
des conditions données.
\end{boxdef}

L'acceptation est la décision que les registres oublient le plus souvent.
Un risque qui reste dans le registre d'une année sur l'autre, sans
traitement ni signature, est accepté de fait~: personne ne l'a décidé, et
personne n'en répond. C'est la forme la plus discrète du défaut de
gouvernance étudié au Chapitre~\ref{chap:grc}.

\subsection{Le cas de DIGICORE}

Le scénario DC-02, défaillance du prestataire d'hébergement DIGICORE,
était coté $2 \times 3 = 6$ en 2025, avec une vraisemblance sans base
(\NEK{}). Le contrat ne contient pas de clause de droit d'audit~; la banque
ne sait pas comment le prestataire gère ses propres risques.

\begin{boxobjet}{R-DC-02 · Décision d'acceptation (janvier 2027)}
\textbf{Risque~:} défaillance prolongée de DIGICORE, prestataire
d'hébergement du centre de données.

\textbf{Propriétaire~:} le directeur des opérations.

\textbf{Résiduel~:} vraisemblance~2, impact~3, criticité~6. Base de la
vraisemblance~: \NEK{} (aucune revue du prestataire, pas de droit
d'audit). Base de l'impact~: \OUI{} (analyse d'impact sur l'activité).

\textbf{Mesures en place~:} site de secours de Bonabéri~; plan de
continuité approuvé en 2022, jamais testé (aucune réduction retenue).

\textbf{Décision~:} accepté par le comité de sécurité \emph{en l'état des
connaissances}, pour six mois. Conditions~: avenant au contrat introduisant
un droit d'audit et un rapport annuel d'assurance du prestataire (A10)~;
test du plan de continuité avant la revue.

\textbf{Revue~:} juillet 2027. Si l'avenant est refusé, le risque est
présenté au comité des risques du conseil.

\textbf{Signatures~:} directeur des opérations (propriétaire)~; président
du comité de sécurité~; RSSI (avis).
\end{boxobjet}

\begin{boxalert}
\textbf{Violation typique.} Faire signer l'acceptation par le RSSI ou par
le DSI parce qu'ils «~connaissent le sujet~».

\smallskip
\textbf{Règle~:} accepte celui qui a l'autorité et supporte les
conséquences, selon les critères d'acceptation. Le RSSI donne un avis~; il
ne signe pas à la place du propriétaire.
\end{boxalert}

% ============================================================================
\section{Suivre et Revoir}
\label{sec:c03-suivre}
% ============================================================================

\subsection{Le programme 2027, neuf mois plus tard}

Le programme de traitement de 2027 compte les huit actions de R-AG-01 et
deux autres~: A9, conservation et lecture des journaux du pare-feu
d'accès (pour R-DC-03), et A10, avenant au contrat de DIGICORE (pour
R-DC-02). Fin septembre 2027, le comité de sécurité fait le point. Une
action est cotée \OUI{} si elle est réalisée \emph{et} que la preuve
attendue est au dossier, \NON{} si elle n'est pas réalisée, \NEK{} si elle
est déclarée réalisée sans la preuve attendue.

\begin{tableaucours}{cL{5.6cm}cL{4.2cm}}{%
    \tableauHeader{Réf.} &
    \tableauHeader{Action} &
    \tableauHeader{Cote} &
    \tableauHeader{Pièce}
}
A1 & Second facteur & \OUI{} & Extraction de la configuration \\
A2 & Cloisonnement & \OUI{} & Rapport du test A3 \\
A3 & Test d'intrusion interne & \OUI{} & Rapport du prestataire \\
A4 & Sauvegardes hors ligne testées & \OUI{} & Procès-verbaux de
restauration \\
A5 & Mode dégradé exercé & \NEK{} & Procédure rédigée~; aucun exercice
tracé \\
A6 & Sensibilisation & \NEK{} & Listes de présence de 22 agences sur 38 \\
A7 & Exercice de réponse aux incidents & \NON{} & Reporté à 2028 \\
A8 & Journalisation centralisée & \NEK{} & « 20 agences raccordées »,
sans extraction \\
A9 & Journaux du pare-feu d'accès & \OUI{} & Configuration, rapports de
revue mensuels \\
A10 & Avenant DIGICORE & \NON{} & Refus du prestataire \\
\end{tableaucours}
\vspace{-0.8em}
\begin{center}{\small\itshape Tableau~3.7. Avancement du programme de
traitement (septembre 2027)}\end{center}

\begin{boxlecture}{Programme de traitement 2027}
\textbf{Constatation.} 5~\OUI{}, 2~\NON{}, 3~\NEK{} sur 10 actions~:
\[
\mu = 0{,}50~;\quad U = 0{,}30~;\quad
H = -(0{,}5\log_2 0{,}5 + 0{,}2\log_2 0{,}2 + 0{,}3\log_2 0{,}3)
  \approx 0{,}500 + 0{,}464 + 0{,}521 \approx 1{,}49~\text{bit}.
\]

\textbf{Analyse} (confiance élevée). Les actions techniques sont faites et
prouvées~; les actions humaines et organisationnelles (mode dégradé,
sensibilisation, exercice) sont en retard ou sans preuve. Pour R-AG-01,
l'entrée et la propagation sont couvertes~: vraisemblance~2. La
récupération l'est aussi grâce à A4~: impact~3. Le résiduel visé, 6, est
atteint. Mais deux \NEK{} portent sur la récupération (A5) et la détection
(A8)~: le jour d'une attaque, la banque ne sait pas encore si ses agences
sauront servir en mode dégradé. Le refus de DIGICORE (A10) déclenche, selon
la décision de janvier, la présentation de R-DC-02 au conseil.

\medskip
\hfill\lecture{0{,}50}{0{,}30}{1{,}49}
\end{boxlecture}

\subsection{Quand revoir}

Une appréciation se revoit à date fixe (au moins une fois par an), et
aussi chaque fois que le contexte change~: incident significatif, nouveau
système, nouveau prestataire, nouvelle menace signalée par les autorités,
changement de réglementation. Le Chapitre~\ref{chap:grc-pratique} montre
comment des indicateurs de risque déclenchent ces revues sans attendre la
date anniversaire.

% ============================================================================
\section{Le Cadre Camerounais}
\label{sec:c03-cadre}
% ============================================================================

\begin{boxjuris}
\textbf{\loicam{n°~2010/012}, article~26, alinéa~1}~\cite{loi2010012}~:

\medskip
\textit{« Les exploitants des systèmes d'information prennent toutes les
mesures techniques et administratives afin de garantir la sécurité des
services offerts. A cet effet, ils se dotent de systèmes normalisés leur
permettant d'identifier, d'évaluer, de traiter et de gérer continûment les
risques liés à la sécurité des systèmes d'information dans le cadre des
services offerts directement ou indirectement. »}

\medskip
\textbf{Article~29, alinéa~1}~:

\medskip
\textit{« Les exploitants des systèmes d'information ont l'obligation de
conserver les données de connexion et de trafic de leurs systèmes
d'information pendant une période de dix (10) ans. »}

\medskip
\textbf{Portée.} Les quatre verbes de l'article~26 (identifier, évaluer,
traiter, gérer continûment) recouvrent les étapes du processus de
l'ISO~31000~: identification, analyse et évaluation, traitement, suivi et
revue. Une démarche conforme à l'ISO/IEC~27005 est donc une manière de
satisfaire cette obligation, sans être la seule.

\medskip
\textbf{Analyse} (confiance moyenne~: la portée exacte des « données de
connexion et de trafic » est étudiée au Chapitre~\ref{chap:droit}). Des
journaux de pare-feu effacés au bout de trente jours, comme ceux de la
BANQUE SANAGA en 2025, se concilient mal avec l'article~29. L'action A9 ne
traite donc pas seulement un risque~: elle répond aussi à une obligation.
\end{boxjuris}

% ============================================================================
\section{Ce que le Processus demande à chaque Métier}
\label{sec:c03-metiers}
% ============================================================================

\begin{boxmetier}{Gestionnaire des risques cyber}
\textbf{Faire tourner le processus}
\begin{enumerate}
    \item Faire approuver les critères d'acceptation avant toute
          appréciation.
    \item Combiner l'approche par les événements et l'approche par les
          actifs.
    \item Pour chaque risque, obtenir une décision écrite~: un plan, ou une
          acceptation signée.
    \item Suivre le plan par ses preuves, pas par les déclarations.
\end{enumerate}
\end{boxmetier}

\begin{boxmetier}{RSSI}
\textbf{Un plan qui produit des preuves}

Chaque action du plan prévoit la preuve qui la clôturera. Le RSSI inscrit
au plan les tests eux-mêmes (comme A3), car une mesure non testée
n'abaisse rien, et refuse de déclarer réalisée une action sans sa pièce.
\end{boxmetier}

\begin{boxmetier}{Consultant en mise en œuvre de SMSI}
\textbf{Livrer un processus, pas un registre}

Le consultant qui accompagne une organisation laisse derrière lui des
critères approuvés, une méthode écrite, un plan avec des responsables et
une première revue datée. Un registre bien présenté, livré seul, sera
obsolète à la première réorganisation.
\end{boxmetier}

\begin{boxmetier}{Auditeur en cybersécurité}
\textbf{Auditer les décisions}

L'auditeur prélève des risques du registre et demande, pour chacun, la
décision~: le plan approuvé ou l'acceptation signée. Il vérifie que le
signataire avait l'autorité selon les critères, et qu'aucune acceptation ne
repose sur une cote sans base, sauf acceptation datée en l'état des
connaissances.
\end{boxmetier}

\separateur

\begin{boxresume}
\textbf{Ce qu'il faut retenir de ce chapitre~:}
\begin{itemize}
    \item \textbf{L'ISO~31000} donne des principes, un cadre et un
          processus~; la communication et le suivi encadrent toutes les
          étapes.
    \item \textbf{Le contexte} fixe le périmètre, les parties intéressées
          et deux familles de critères~: pour apprécier, et pour accepter.
    \item \textbf{Deux approches d'identification}~: par les événements
          pour savoir ce qui compte, par les actifs pour vérifier que rien
          n'échappe.
    \item \textbf{Quatre familles de traitement}~: réduire, partager,
          éviter, maintenir~; l'assurance vient après la réduction.
    \item \textbf{Un plan de traitement} donne, pour chaque action, un
          responsable, une échéance, un coût et la preuve attendue.
    \item \textbf{L'acceptation} est une décision écrite, par l'instance
          compétente, à durée limitée~; sur une cote \NEK{}, elle se fait
          « en l'état des connaissances », avec une date de revue.
    \item \textbf{Le suivi} cote les actions par leurs preuves~: en
          septembre 2027, $\mu = 0{,}50$, $U = 0{,}30$ à la BANQUE SANAGA.
\end{itemize}
\end{boxresume}

\begin{boxinfo}
\textbf{Points de contrôle.} Vous avez compris ce chapitre si vous savez~:
\begin{itemize}
    \item placer une activité dans le processus de l'ISO~31000~;
    \item dire ce que l'approche par les actifs apporte à l'approche par les
          événements, sur un exemple~;
    \item rédiger une ligne de plan de traitement complète~;
    \item dire qui peut accepter un résiduel donné, et à quelles
          conditions.
\end{itemize}
\end{boxinfo}

% ============================================================================
\section*{Exercices}
% ============================================================================

\begin{boxexo}[title={\ding{45}\quad Exercice~3.1. Où en est-on dans le processus~? \niveaubase}]
Placez chacune des activités suivantes dans le processus de l'ISO~31000
(Figure~\ref{fig:c03-processus}).
\begin{enumerate}[(a)]
    \item Le comité des risques fixe qui peut accepter un résiduel de 5 à 8.
    \item L'inventaire des postes des 38 agences.
    \item La RSSI cote la vraisemblance de la fraude au virement.
    \item Le comité compare les criticités aux critères et classe les
          risques.
    \item La DSI déploie le second facteur.
    \item Le comité de sécurité fait le point sur le programme 2027.
    \item La direction du réseau d'agences participe à l'atelier des
          scénarios.
    \item Le rapport annuel des risques est transmis au conseil.
\end{enumerate}
\end{boxexo}

\begin{boxexo}[title={\ding{45}\quad Exercice~3.2. Choisir un traitement \niveaumoyen}]
Pour chacun des risques suivants de LiZbiZ (appétence~: 9 et plus
intolérable, 5 à 8 traité dans l'année, 1 à 4 accepté), proposez une
famille de traitement et une mesure, et justifiez.
\begin{enumerate}[(a)]
    \item R1, ajustement frauduleux du stock, résiduel 12.
    \item R4, arrêt prolongé du centre de Garoua, résiduel 4, base de
          probabilité \OUI{}.
    \item R6, fuite de données des clients, résiduel 3, base de probabilité
          \NEK{} (Exercice~2.4).
    \item R8, rançongiciel, résiduel 12 après nouvelle cote.
\end{enumerate}
Indiquez, pour (d), pourquoi une assurance seule ne suffit pas.
\end{boxexo}

\begin{boxexo}[title={\ding{45}\quad Exercice~3.3. Une acceptation contestable \niveaumoyen}]
En février 2027, le dossier de R-DC-06 (déni de service des services en
ligne, coté $1 \times 3 = 3$ en 2025) contient la décision suivante~:
« Risque accepté. Signé~: le DSI. » Aucune date de revue n'est indiquée.
\begin{enumerate}[(a)]
    \item Relevez tous les défauts de cette décision au regard des critères
          d'acceptation (Tableau~3.3).
    \item La base de la vraisemblance était cotée \NON{} (quatre attaques
          signalées en 2025). Proposez une nouvelle vraisemblance, recalculez
          la criticité et dites qui peut désormais accepter le risque.
    \item Rédigez la décision corrigée sur le modèle de R-DC-02.
\end{enumerate}
\end{boxexo}

\begin{boxexo}[title={\ding{45}\quad Exercice~3.4. Le plan de traitement de R8 \niveauexpert}]
LiZbiZ décide de traiter R8 (rançongiciel, probabilité~3, impact~4,
criticité~12, aucune mesure testée). L'audit des quatre centres a relevé
des comptes par défaut, des services en clair et des flux non autorisés.
\begin{enumerate}[(a)]
    \item Décomposez le scénario en étapes (entrée, propagation,
          récupération) et reliez chaque constat de l'audit à une étape.
    \item Rédigez un plan de traitement de six actions au plus, au format du
          Tableau~3.6, en indiquant pour chacune la preuve attendue.
    \item Calculez le résiduel visé avec la convention de réduction, en
          partant d'un inhérent de $4 \times 4$, et dites qui pourra
          l'accepter selon des critères identiques à ceux de la BANQUE
          SANAGA.
    \item Un an plus tard, quatre actions sont prouvées, une est déclarée
          faite sans preuve, une n'est pas faite. Calculez $\mu$, $U$ et
          $H$, et dites si le résiduel visé est atteint.
\end{enumerate}
\end{boxexo}

\bigskip

\begin{boxinfo}
\textbf{Pour aller plus loin.}
\begin{itemize}
    \item L'ISO~31000:2018~\cite{iso31000_2018}, pour les principes, le
          cadre et le processus.
    \item L'ISO/IEC~27005:2022~\cite{iso27005_2022}, pour l'application à
          la sécurité de l'information.
    \item L'ISO/IEC~27001:2022~\cite{iso27001_2022}, article~6.1.3, sur le
          traitement des risques et la déclaration d'applicabilité.
\end{itemize}
\end{boxinfo}

\begin{boxinfo}
\textbf{Transition.} Le processus dit quoi faire~; il ne dit pas comment
imaginer les scénarios. L'atelier de décembre 2026 a trouvé quatre
scénarios stratégiques~: comment les construire avec méthode, à partir des
sources de risque et de leurs objectifs, puis les détailler en modes
opératoires~? C'est l'objet de la méthode EBIOS Risk Manager, au
Chapitre~\ref{chap:ebios}.
\end{boxinfo}

        % ============================================================================
% Fichier  : partie2/P02_C04_ebios_risk_manager.tex
% Titre    : EBIOS Risk Manager
% Auteur   : MINKA MI NGUIDJOÏ Thierry Emmanuel
% Version  : 1.0 -- novembre 2026
% Corrigés : corriges/Corr_C04.tex
% ============================================================================
% Vérifié sur cyber.gouv.fr (consulté le 05/11/2026) : noms des cinq
% ateliers ; synthèse « conformité » et « scénarios » ; socle de sécurité ;
% scénarios centrés sur les menaces intentionnelles et ciblées ; guide
% publié le 9 octobre 2018 avec le Club EBIOS. Version 1.5 de mars 2024
% signalée par Wikipédia (lien vers le guide) : à confirmer sur le guide.
% NON RELU dans le guide (PDF non téléchargé) et signalé à l'auteur : la
% formule du niveau de menace d'une partie prenante (dépendance ×
% pénétration / maturité × confiance), la séquence « connaître, rentrer,
% trouver, exploiter », le nom « plan d'amélioration continue de la
% sécurité », les deux cycles. Les échelles de gravité et de vraisemblance
% du guide ne sont pas citées : le chapitre emploie l'échelle du cours.
% Faits du cas (bible) : LiZbiZ, quatre centres de données (Douala,
% Yaoundé, Garoua, Bafoussam), douze agences, application intégrée ;
% M. ETEME, Mme NKENG, M. KOUNOU, M. BEYALA (AUDIX CONSEIL) ; TRANSLOG
% (112,7 M FCFA, OWONA) ; ajustements de stock de FOUDA à Douala ; écart de
% 186 M FCFA ; registre R1 à R8 ; constats de l'audit des quatre centres.
% Faits nouveaux (reportés dans la bible) : demande du client PHARMACAM
% (fictif), étude EBIOS RM de janvier à mars 2027 animée par Mme NKENG ;
% ER1 à ER4 ; SR1 à SR5 ; parties prenantes P1 à P5 et leurs cotes ; base
% de preuve des 20 cotes de l'écosystème : 13 OUI / 1 NON / 6 NEK.
% ============================================================================

\chapter{EBIOS Risk Manager}
\label{chap:ebios}
\chaptermark{EBIOS Risk Manager}

\epigraph{%
    « On ne se protège pas d'une menace\\
    que l'on n'a pas pris la peine d'imaginer. »%
}{Principe du cours}

\niveauChapitre{Intermédiaire}{%
    Maîtrise du Chapitre~\ref{chap:iso27005}%
}

\begin{boxobjectifs}
\begin{itemize}
    \item Situer EBIOS Risk Manager~: son origine, sa logique « conformité
          et scénarios », ses cinq ateliers.
    \item Cadrer une étude~: valeurs métier, biens supports, événements
          redoutés, gravité, socle de sécurité.
    \item Identifier et retenir des couples source de risque et objectif
          visé.
    \item Cartographier l'écosystème et construire des scénarios
          stratégiques.
    \item Détailler un scénario opérationnel et en coter la vraisemblance.
    \item Coter la base de preuve d'une étude EBIOS, et en tirer les
          priorités de connaissance.
    \item Situer MEHARI et le NIST SP~800-30 par rapport à EBIOS.
\end{itemize}
\end{boxobjectifs}

\bigskip

% ============================================================================
\section*{Avant de commencer : le client demande une preuve}
% ============================================================================

Douala, janvier 2027. PHARMACAM\footnote{LiZbiZ, PHARMACAM, leurs
dirigeants, les dates et les montants sont fictifs~; la méthode et les
textes cités sont exacts.}, distributeur de médicaments et premier client
de LiZbiZ, prépare le renouvellement de son contrat de transport. Sa
direction des achats a ajouté une clause~: le prestataire devra remettre
« une analyse des risques de cybersécurité conduite selon une méthode
reconnue ». Les médicaments voyagent sous chaîne du froid~; un arrêt du
suivi des expéditions de plusieurs jours, et ce sont des lots entiers
qu'il faut détruire.

M.~ETEME, directeur général de LiZbiZ, appelle d'abord M.~BEYALA, d'AUDIX
CONSEIL, qui connaît l'entreprise. Celui-ci refuse d'animer l'étude~: son
cabinet a audité LiZbiZ et l'auditera encore~; il ne peut pas auditer
demain une analyse qu'il aurait conduite aujourd'hui. Il conseille de
confier l'animation à Mme~NKENG, contrôleuse interne, et de choisir EBIOS
Risk Manager.

\medskip

Ce chapitre présente cette méthode et la déroule sur LiZbiZ, atelier par
atelier. Les rançongiciels, la fraude au stock et les faux fournisseurs,
que le registre R1 à R8 décrivait séparément, y deviennent des scénarios
reliés à des attaquants, à leurs objectifs et à l'écosystème de
l'entreprise.

\begin{boxalert}
\textbf{Violation typique.} Faire animer l'analyse des risques par le
cabinet qui en auditera ensuite les résultats.

\smallskip
\textbf{Règle~:} on n'audite pas son propre travail. Un auditeur peut
conseiller sur la méthode~; il ne construit pas l'analyse qu'il devra
évaluer. Le Chapitre~\ref{chap:droit} revient sur ces règles
d'indépendance.
\end{boxalert}

\separateur

% ============================================================================
\section{La Méthode}
\label{sec:c04-methode}
% ============================================================================

\subsection{Origine et logique}

EBIOS Risk Manager est la méthode d'appréciation et de traitement du risque
numérique publiée en 2018 par l'Agence nationale de la sécurité des
systèmes d'information (\sigle{ANSSI}, France), avec le soutien du Club
EBIOS~\cite{anssi2018ebios}. Elle succède aux versions d'EBIOS de 1995,
2004 et 2010. Elle est compatible avec l'ISO~31000 et l'ISO/IEC~27005~:
elle en est une mise en œuvre, pas une alternative.

Sa logique tient en une phrase, que l'ANSSI formule comme une
\emph{synthèse entre conformité et scénarios}~:
\begin{itemize}
    \item la \textbf{conformité} construit un \emph{socle de sécurité}~:
          l'organisation applique les mesures d'un référentiel (hygiène
          informatique, ISO/IEC~27002) et mesure ses écarts~;
    \item les \textbf{scénarios} mettent ce socle à l'épreuve de menaces
          \emph{intentionnelles et ciblées}~: un attaquant précis, avec un
          objectif précis, qui passe éventuellement par l'écosystème de
          l'organisation.
\end{itemize}

\begin{boxinfo}
\textbf{Pourquoi cette division du travail~?} Les menaces courantes
(maliciel opportuniste, erreur, panne) se traitent par le socle~: il serait
inutile de construire un scénario pour chacune. L'effort d'imagination est
réservé aux attaques qui visent l'organisation et peuvent contourner le
socle. C'est l'inverse d'un registre qui énumère cent menaces au même
niveau de détail.
\end{boxinfo}

\subsection{Cinq ateliers}

\begin{tableaucours}{L{3.4cm}L{4.4cm}L{4.4cm}}{%
    \tableauHeader{Atelier} &
    \tableauHeader{Question} &
    \tableauHeader{Principaux livrables}
}
1. Cadrage et socle de sécurité & Que protège-t-on, de quoi a-t-on peur,
où en est-on de la conformité~? & Périmètre, valeurs métier, biens
supports, événements redoutés et gravité, écarts au socle \\
2. Sources de risque & Qui pourrait nous attaquer, et pour quoi faire~? &
Couples source de risque et objectif visé, retenus ou écartés \\
3. Scénarios stratégiques & Par où passerait l'attaquant, y compris par
nos partenaires~? & Cartographie de l'écosystème, scénarios stratégiques,
mesures sur l'écosystème \\
4. Scénarios opérationnels & Comment s'y prendrait-il, techniquement~? &
Modes opératoires, vraisemblance \\
5. Traitement du risque & Que décide-t-on, et comment le suit-on~? &
Stratégie de traitement, mesures, risques résiduels, cadre de suivi \\
\end{tableaucours}
\vspace{-0.8em}
\begin{center}{\small\itshape Tableau~4.1. Les cinq ateliers d'EBIOS Risk
Manager}\end{center}

Les ateliers~1 à~3 relèvent de la direction et des métiers~: ils
raisonnent en enjeux et en attaquants. L'atelier~4 relève des
techniciens~: il raisonne en vulnérabilités et en chemins d'attaque.
L'atelier~5 revient à la direction, qui décide. Les ateliers stratégiques se
revoient rarement, les ateliers opérationnels plus souvent, au rythme des
systèmes et des menaces.

\begin{boxinfo}
\textbf{Échelles.} Le guide propose des échelles de gravité et de
vraisemblance que chaque organisation adapte. Le cours emploie la sienne
(Chapitre~\ref{chap:risque}, Tableau~2.3), en appelant \emph{gravité}
l'impact d'un événement redouté, comme le fait la méthode. On note G1 à G4
et V1 à V4 les niveaux de cette échelle.
\end{boxinfo}

% ============================================================================
\section{Atelier 1 : Cadrage et Socle de Sécurité}
\label{sec:c04-atelier1}
% ============================================================================

\subsection{Périmètre, valeurs métier, biens supports}

Mme~NKENG réunit M.~ETEME, le directeur des opérations, le directeur
financier, M.~KOUNOU et un chef d'agence. Le périmètre retenu est le
système d'information logistique~: les quatre centres de données de Douala,
Yaoundé, Garoua et Bafoussam, les douze agences et l'application intégrée.
L'atelier retient trois valeurs métier~: le \emph{suivi des expéditions},
la \emph{facturation} et la \emph{gestion des stocks}, et une quatrième,
ajoutée à la demande de PHARMACAM~: les \emph{données des clients}.

\subsection{Événements redoutés et gravité}

\begin{tableaucours}{cL{5.6cm}L{4.2cm}c}{%
    \tableauHeader{Réf.} &
    \tableauHeader{Événement redouté} &
    \tableauHeader{Conséquences} &
    \tableauHeader{Grav.}
}
ER1 & Indisponibilité du suivi des expéditions pendant plus de trois jours
& Lots sous chaîne du froid perdus, pénalités, perte du client & G4 \\
ER2 & Altération des données de stock par des ajustements non autorisés &
Écarts de stock, réserve du commissaire aux comptes & G3 \\
ER3 & Paiement de factures à un fournisseur fictif & Perte financière
directe & G3 \\
ER4 & Divulgation des données des clients & Rupture de contrat,
sanction de l'autorité de protection & G3 \\
\end{tableaucours}
\vspace{-0.8em}
\begin{center}{\small\itshape Tableau~4.2. Événements redoutés de LiZbiZ}\end{center}

Les gravités s'appuient sur des faits~: l'écart de stock de 186~millions
FCFA (ER2) et la fraude TRANSLOG de 112,7~millions (ER3) se situent au
niveau majeur de l'échelle~; pour ER1, PHARMACAM a chiffré la valeur d'un
seul camion de vaccins à plus de 500~millions.

\subsection{Le socle de sécurité}

L'atelier compare ensuite l'existant à un référentiel. Faute de mieux,
LiZbiZ reprend les constats de l'audit de ses quatre centres~: des comptes
par défaut actifs, des services d'administration en clair, des flux non
autorisés entre centres. Ces écarts ne sont pas des scénarios~; ils
constituent un socle troué, que les scénarios des ateliers suivants vont
exploiter. Un écart au socle se corrige, en principe, sans attendre la fin
de l'étude.

% ============================================================================
\section{Atelier 2 : Sources de Risque}
\label{sec:c04-atelier2}
% ============================================================================

\subsection{Des couples, pas des menaces}

\begin{boxdef}
\termecle{Couple SR/OV}\index{couple SR/OV}~: association d'une source de
risque (qui~?) et d'un objectif visé (pour quoi faire~?). Un même
attaquant peut poursuivre plusieurs objectifs~; chaque couple se traite à
part.
\end{boxdef}

L'atelier recense les couples plausibles, puis les apprécie selon la
\emph{motivation} de la source à atteindre l'objectif, ses
\emph{ressources}, et son \emph{activité} connue dans le secteur. Il en
retient quelques-uns, jugés pertinents.

\begin{tableaucours}{cL{3.0cm}L{2.8cm}L{4.0cm}c}{%
    \tableauHeader{Réf.} &
    \tableauHeader{Source de risque} &
    \tableauHeader{Objectif visé} &
    \tableauHeader{Appréciation} &
    \tableauHeader{Retenu}
}
SR1 & Groupe cybercriminel & Extorsion par rançongiciel & Très motivé,
outillé, très actif contre les entreprises de transport & Oui \\
SR2 & Employé malveillant, seul ou avec un complice & Détourner du stock
et masquer l'écart & Précédent connu (ajustements de Douala) & Oui \\
SR3 & Fournisseur indélicat ou faux fournisseur & Se faire payer des
factures fictives & Précédent connu (TRANSLOG) & Oui \\
SR4 & Concurrent & Obtenir les tarifs et la liste des clients & Motivé,
mais peu d'activité connue & Non \\
SR5 & Militant & Nuire à l'image & Peu motivé par une entreprise de
logistique & Non \\
\end{tableaucours}
\vspace{-0.8em}
\begin{center}{\small\itshape Tableau~4.3. Couples source de risque et
objectif visé}\end{center}

\begin{boxscenario}{Opération LIZBIZ. Deux fraudes devenues des sources de risque}
\textbf{Constatation.} Les couples SR2 et SR3 ne sont pas imaginés~: ils
reprennent les deux fraudes établies par AUDIX CONSEIL, les ajustements de
stock de FOUDA et le fournisseur fictif TRANSLOG (112,7~millions FCFA).

\medskip
\textbf{Analyse} (confiance élevée). Le registre de 2026 traitait ces
fraudes comme des risques de contrôle interne (R1, R2), et le rançongiciel
comme un risque informatique (R8). EBIOS les met sur le même plan~: trois
attaquants, trois objectifs, qui exploitent souvent les mêmes faiblesses.
Les droits illimités de M.~KOUNOU servent aussi bien un fraudeur interne
qu'un cybercriminel qui aurait volé son compte.
\end{boxscenario}

% ============================================================================
\section{Atelier 3 : Scénarios Stratégiques}
\label{sec:c04-atelier3}
% ============================================================================

\subsection{L'écosystème}

Un attaquant ciblé passe souvent par un partenaire, moins protégé que sa
cible. L'atelier~3 cartographie donc les \emph{parties prenantes}~:
prestataires, sous-traitants, fournisseurs, clients reliés au système. Le
guide de l'ANSSI les évalue selon quatre critères~: la \emph{dépendance}
de l'organisation envers la partie prenante, la \emph{pénétration} de
celle-ci dans le système, sa \emph{maturité} en cybersécurité, et la
\emph{confiance} qu'on peut lui accorder. Il en déduit un niveau de
menace, d'autant plus élevé que la partie prenante est indispensable et
présente dans le système, et d'autant plus faible qu'elle est mûre et
fiable~:
\[
\text{Menace} = \frac{\text{Dépendance} \times \text{Pénétration}}
                     {\text{Maturité} \times \text{Confiance}}
\]
Chaque critère se cote de 1 à 4. L'atelier fixe lui-même les seuils~: une
menace de 3 ou plus place la partie prenante en zone de \emph{danger}, de
1 à 3 en zone de \emph{contrôle}, sous 1 en zone de \emph{veille}.

\begin{tableaucours}{cL{4.2cm}ccccc}{%
    \tableauHeader{Réf.} &
    \tableauHeader{Partie prenante} &
    \tableauHeader{Dép.} &
    \tableauHeader{Pén.} &
    \tableauHeader{Mat.} &
    \tableauHeader{Conf.} &
    \tableauHeader{Menace}
}
P1 & Éditeur de l'application intégrée (maintenance à distance) & 4 & 4 & 2
& 2 & 4,00 \\
P2 & Transporteurs sous-traitants (portail de suivi) & 3 & 2 & 1 & 2 &
3,00 \\
P3 & Opérateur de télécommunications & 4 & 1 & 3 & 3 & 0,44 \\
P4 & Cabinet comptable externe & 2 & 2 & 2 & 3 & 0,67 \\
P5 & Fournisseurs (portail de dépôt des factures) & 2 & 1 & 1 & 1 & 2,00 \\
\end{tableaucours}
\vspace{-0.8em}
\begin{center}{\small\itshape Tableau~4.4. Les parties prenantes de
LiZbiZ}\end{center}

L'éditeur de l'application (P1) et les transporteurs (P2) sont en zone de
danger~; les fournisseurs (P5) en zone de contrôle~; l'opérateur et le
cabinet comptable en zone de veille.

\subsection{Trois scénarios stratégiques}

\begin{boxdef}
\termecle{Scénario stratégique}\index{scénario stratégique}~: chemin
d'attaque de haut niveau, qui relie un couple SR/OV à un événement redouté,
en passant éventuellement par une ou plusieurs parties prenantes. Il se
décrit sans détail technique et prend la gravité de l'événement redouté.
\end{boxdef}

\begin{tableaucours}{cL{8.4cm}cc}{%
    \tableauHeader{Réf.} &
    \tableauHeader{Chemin d'attaque} &
    \tableauHeader{ER} &
    \tableauHeader{Grav.}
}
SS1 & SR1 compromet l'accès de maintenance de l'éditeur (P1), atteint les
centres de données et chiffre l'application et ses sauvegardes & ER1 & G4 \\
SS2 & SR2, magasinier de Douala, obtient l'aide d'un détenteur de droits
d'administration et passe des ajustements de stock sans contrôle & ER2 &
G3 \\
SS3 & SR3 s'enregistre comme fournisseur sur le portail (P5) et fait
valider des factures fictives par un complice interne & ER3 & G3 \\
\end{tableaucours}
\vspace{-0.8em}
\begin{center}{\small\itshape Tableau~4.5. Scénarios stratégiques de
LiZbiZ}\end{center}

L'atelier propose aussi des mesures sur l'écosystème, qui ne dépendent pas
du détail technique~: imposer à l'éditeur un second facteur et une
journalisation de ses connexions de maintenance~; limiter l'accès des
transporteurs aux seules expéditions qui les concernent~; vérifier
l'existence de tout nouveau fournisseur avant son premier paiement.

\subsection{Ce que valent les cotes de l'écosystème}

Les vingt cotes du Tableau~4.4 n'ont pas toutes la même solidité. La
dépendance et la pénétration se lisent dans les contrats et dans les
droits d'accès~; la maturité et la confiance se devinent.

\begin{boxlecture}{Base de preuve des cotes de l'écosystème}
\textbf{Constatation.} Mme~NKENG cote la base de chacune des vingt
cotes~:
\begin{itemize}
    \item dépendance~: 5~\OUI{} (contrats, chiffre d'affaires
          concerné)~;
    \item pénétration~: 4~\OUI{} et 1~\NON{}~: P2 était coté 2, mais les
          droits du portail montrent que chaque transporteur voit toutes
          les expéditions~;
    \item maturité~: 1~\OUI{} (l'opérateur est certifié ISO/IEC~27001), 4~\NEK{}~;
    \item confiance~: 3~\OUI{} (historique de la relation), 2~\NEK{} (P1, P2).
\end{itemize}
Sur vingt cotes~: 13~\OUI{}, 1~\NON{}, 6~\NEK{}~:
\[
\mu = 0{,}65~;\quad U = 0{,}30~;\quad
H = -(0{,}65\log_2 0{,}65 + 0{,}05\log_2 0{,}05 + 0{,}30\log_2 0{,}30)
  \approx 0{,}404 + 0{,}216 + 0{,}521 \approx 1{,}14~\text{bit}.
\]

\textbf{Analyse} (confiance élevée sur le calcul, moyenne sur la lecture).
L'incertitude se concentre au dénominateur de la formule~: la maturité et
la confiance. Or ce sont elles qui font baisser la menace. Une maturité
cotée 2 sans preuve peut cacher un 1, ce qui doublerait la menace de P1. Les
zones de danger sont donc probablement sous-estimées, et non surestimées.
Pour P2, le \NON{} se corrige tout de suite~: avec une pénétration de 3, sa
menace passe de 3,00 à 4,50.

\medskip
\hfill\lecture{0{,}65}{0{,}30}{1{,}14}
\end{boxlecture}

\begin{boxprocedure}
\textbf{Réduire l'incertitude sur l'écosystème.} Pour chaque partie
prenante en zone de danger ou de contrôle, le papier de travail prévoit~:
\begin{itemize}
    \item un questionnaire de maturité, signé par la partie prenante, et
          les attestations qu'elle peut fournir (certification, rapport
          d'assurance)~;
    \item l'extraction des droits d'accès dont elle dispose réellement~;
    \item une clause contractuelle de droit d'audit et de notification des
          incidents~;
    \item la date de la dernière revue, et la cote de base de chaque
          critère.
\end{itemize}
\end{boxprocedure}

% ============================================================================
\section{Atelier 4 : Scénarios Opérationnels}
\label{sec:c04-atelier4}
% ============================================================================

\subsection{Du chemin au mode opératoire}

\begin{boxdef}
\termecle{Scénario opérationnel}\index{scénario opérationnel}~: suite
d'actions élémentaires, sur les biens supports, par lesquelles une source
de risque réalise un scénario stratégique. On le décompose en quatre
phases~: \emph{connaître} la cible, y \emph{rentrer}, y \emph{trouver} ce
qui permet d'atteindre l'objectif, puis l'\emph{exploiter}.
\end{boxdef}

\subsection{Le mode opératoire de SS1}

M.~KOUNOU et deux techniciens de l'éditeur détaillent SS1. Pour chaque
phase, ils cotent la vraisemblance que l'attaquant y réussisse, au vu des
constats de l'audit.

\begin{tableaucours}{L{1.8cm}L{6.2cm}cL{2.6cm}}{%
    \tableauHeader{Phase} &
    \tableauHeader{Action élémentaire} &
    \tableauHeader{V} &
    \tableauHeader{Base}
}
Connaître & Repérer sur Internet l'éditeur et son outil de maintenance à
distance & V4 & Sources publiques \\
Rentrer & Voler l'identifiant d'un technicien de l'éditeur par
hameçonnage, ou utiliser un compte par défaut du serveur de maintenance &
V3 & Constat d'audit (comptes par défaut) \\
Trouver & Se déplacer d'un centre à l'autre par les flux non autorisés,
capter des mots de passe sur les services en clair & V4 & Constats
d'audit \\
Exploiter & Chiffrer l'application dans les quatre centres et effacer les
sauvegardes accessibles & V3 & Sauvegardes en ligne, jamais restaurées \\
\end{tableaucours}
\vspace{-0.8em}
\begin{center}{\small\itshape Tableau~4.6. Le scénario opérationnel SO1}\end{center}

\begin{boxinfo}
\textbf{Une règle de chaîne.} Un mode opératoire ne réussit que si toutes
ses phases réussissent. Le cours retient donc, pour vraisemblance du
scénario, celle de la phase la plus difficile pour l'attaquant~: ici V3.
C'est une convention prudente~; le guide décrit aussi des manières plus
fines de combiner les actions élémentaires.
\end{boxinfo}

SO1 obtient une vraisemblance V3 et hérite de la gravité G4 de ER1~:
criticité $3 \times 4 = 12$, intolérable. On retrouve la cote de R8 après
sa révision (Exercice~2.4)~: EBIOS ne la contredit pas, il l'explique. Et
il montre où agir~: la phase «~rentrer~» est la seule cotée V3 au lieu de
V4~; la fermer (second facteur pour l'éditeur, suppression des comptes par
défaut) ferait tomber la vraisemblance de tout le scénario.

% ============================================================================
\section{Atelier 5 : Traitement du Risque}
\label{sec:c04-atelier5}
% ============================================================================

L'atelier~5 rejoint le processus du Chapitre~\ref{chap:iso27005}~: options
de traitement, plan, risques résiduels, acceptation. EBIOS y ajoute deux
apports~: les mesures se rattachent aux \emph{phases} des scénarios et aux
\emph{parties prenantes}, et le plan s'inscrit dans un cadre de suivi qui
fait revenir l'étude à intervalles réguliers. Le guide parle d'un plan
d'amélioration continue de la sécurité.

\begin{tableaucours}{cL{5.6cm}L{2.4cm}c}{%
    \tableauHeader{Scénario} &
    \tableauHeader{Mesures principales} &
    \tableauHeader{Phase ou partie prenante} &
    \tableauHeader{Résiduel visé}
}
SO1 (SS1) & Second facteur et journalisation pour l'éditeur~; suppression
des comptes par défaut~; filtrage entre centres~; sauvegardes hors ligne
testées & Rentrer, trouver, exploiter~; P1 & $2 \times 3 = 6$ \\
SS2 & Séparation des tâches de réception et d'ajustement~; revue des
ajustements par la contrôleuse interne~; comptes d'administration
nominatifs & Exploiter & $2 \times 3 = 6$ \\
SS3 & Vérification de tout nouveau fournisseur avant paiement~; double
validation des factures & P5 & $1 \times 3 = 3$ \\
\end{tableaucours}
\vspace{-0.8em}
\begin{center}{\small\itshape Tableau~4.7. Synthèse du traitement}\end{center}

\begin{boxobjet}{SO1 · Rançongiciel par l'accès de maintenance}
\textbf{Valeur métier~:} suivi des expéditions.

\textbf{Événement redouté~:} ER1, indisponibilité de plus de trois jours
(G4).

\textbf{Source de risque et objectif visé~:} SR1, groupe cybercriminel,
extorsion.

\textbf{Partie prenante~:} P1, éditeur de l'application (menace 4,00,
maturité et confiance \NEK{}).

\textbf{Mode opératoire~:} connaître (V4), rentrer (V3), trouver (V4),
exploiter (V3)~; vraisemblance retenue V3.

\textbf{Criticité~:} 12, intolérable. Base~: constats de l'audit des quatre
centres (\OUI{}) pour les phases rentrer, trouver et exploiter.

\textbf{Traitement~:} réduire (Tableau~4.7), puis envisager une assurance.

\textbf{Propriétaire~:} le directeur des opérations.

\textbf{Revue~:} après le test des sauvegardes et la signature de
l'avenant avec l'éditeur, et au plus tard dans un an.
\end{boxobjet}

% ============================================================================
\section{Situer EBIOS parmi les Méthodes}
\label{sec:c04-methodes}
% ============================================================================

EBIOS Risk Manager n'est pas la seule méthode. Deux autres sont souvent
citées~: MEHARI, développée par le CLUSIF, l'association française des
professionnels de la sécurité, et le guide d'appréciation des risques du
NIST, la publication spéciale 800-30~\cite{nist2012sp80030}.

\begin{tableaucours}{L{2.4cm}L{3.3cm}L{3.3cm}L{3.3cm}}{%
    \tableauHeader{} &
    \tableauHeader{EBIOS RM} &
    \tableauHeader{MEHARI} &
    \tableauHeader{NIST SP~800-30}
}
Origine & ANSSI, 2018 & CLUSIF & NIST, révision~1 de 2012 \\
Point de départ & Sources de risque, objectifs, écosystème & Bases de
connaissances, diagnostic des services de sécurité & Sources de menace,
événements de menace, vulnérabilités \\
Force & Scénarios lisibles par la direction~; écosystème & Questionnaires
détaillés, utiles à l'audit & Démarche à trois niveaux (organisation,
mission, système) \\
Limite & Demande des animateurs formés & Lourde sans outil & Pensée pour
l'administration américaine~; à adapter \\
\end{tableaucours}
\vspace{-0.8em}
\begin{center}{\small\itshape Tableau~4.8. Trois méthodes d'appréciation
des risques}\end{center}

Pour une organisation camerounaise, le choix n'est pas dicté par la loi~:
l'article~26 de la loi n°~2010/012 exige des « systèmes normalisés »
d'identification et d'évaluation des risques (Chapitre~\ref{chap:iso27005}),
sans désigner de méthode~\cite{loi2010012}. EBIOS a l'avantage d'être
gratuite, publiée en français et compatible avec l'ISO/IEC~27005~; elle
reste une méthode étrangère, que l'on adapte aux échelles, au droit et aux
menaces du pays.

% ============================================================================
\section{Ce que la Méthode demande à chaque Métier}
\label{sec:c04-metiers}
% ============================================================================

\begin{boxmetier}{Gestionnaire des risques cyber}
\textbf{Animer, pas rédiger seul}
\begin{enumerate}
    \item Réunir les bons participants à chaque atelier~: la direction et
          les métiers pour les ateliers~1 à~3, les techniciens pour le~4.
    \item Partir des faits connus (fraudes, incidents, constats d'audit)
          pour les couples SR/OV.
    \item Coter la base de preuve des critères de l'écosystème, et
          programmer les questionnaires et les extractions qui manquent.
    \item Relier chaque mesure à une phase ou à une partie prenante.
\end{enumerate}
\end{boxmetier}

\begin{boxmetier}{Auditeur en cybersécurité}
\textbf{Rester de l'autre côté}

M.~BEYALA a eu raison de refuser l'animation. L'auditeur peut évaluer une
étude EBIOS~: il vérifie que les couples écartés l'ont été pour de bonnes
raisons, que les cotes de l'écosystème ont une base, et que les mesures
annoncées existent. Il ne l'écrit pas.
\end{boxmetier}

\begin{boxmetier}{Responsable GRC}
\textbf{Ramener l'étude dans le registre}

Une étude EBIOS produit des scénarios~; le registre doit les recevoir. Le
responsable GRC remplace les lignes R1, R2 et R8 par les scénarios SS2,
SS3 et SO1, avec leurs propriétaires, et veille à ce que l'étude soit
revue quand un nouveau partenaire entre dans l'écosystème.
\end{boxmetier}

\begin{boxmetier}{Consultant en mise en œuvre de SMSI}
\textbf{Transmettre la méthode}

Quand un consultant anime l'étude, il forme en même temps un animateur
interne. La demande de PHARMACAM se répétera chaque année~; LiZbiZ doit
pouvoir refaire l'étude sans lui.
\end{boxmetier}

\separateur

\begin{boxresume}
\textbf{Ce qu'il faut retenir de ce chapitre~:}
\begin{itemize}
    \item \textbf{EBIOS Risk Manager} (ANSSI, 2018) réalise une synthèse
          entre conformité (un socle de sécurité) et scénarios (des menaces
          intentionnelles et ciblées).
    \item \textbf{Cinq ateliers}~: cadrage et socle, sources de risque,
          scénarios stratégiques, scénarios opérationnels, traitement.
    \item \textbf{Un couple SR/OV} relie un attaquant à un objectif~; on
          part des faits connus pour les retenir.
    \item \textbf{L'écosystème} se cote par la dépendance, la pénétration,
          la maturité et la confiance~; les deux derniers critères sont
          souvent \NEK{}, et ce sont eux qui font baisser la menace.
    \item \textbf{Un scénario opérationnel} se décompose en connaître,
          rentrer, trouver, exploiter~; sa vraisemblance est celle de la
          phase la plus difficile.
    \item \textbf{Indépendance}~: celui qui auditera l'étude ne l'anime pas.
\end{itemize}
\end{boxresume}

\begin{boxinfo}
\textbf{Points de contrôle.} Vous avez compris ce chapitre si vous savez~:
\begin{itemize}
    \item dire à quel atelier appartient un livrable donné~;
    \item calculer le niveau de menace d'une partie prenante et la placer
          dans une zone~;
    \item décomposer un scénario opérationnel en quatre phases et en coter
          la vraisemblance~;
    \item expliquer pourquoi une maturité cotée sans preuve sous-estime
          probablement la menace.
\end{itemize}
\end{boxinfo}

% ============================================================================
\section*{Exercices}
% ============================================================================

\begin{boxexo}[title={\ding{45}\quad Exercice~4.1. À quel atelier~? \niveaubase}]
Indiquez l'atelier d'EBIOS Risk Manager auquel appartient chacun des
éléments suivants.
\begin{enumerate}[(a)]
    \item La liste des écarts de LiZbiZ au guide d'hygiène informatique.
    \item Le couple « concurrent, obtenir la liste des clients », écarté.
    \item La cote de pénétration du cabinet comptable.
    \item L'action « capter des mots de passe sur les services en clair ».
    \item La gravité G4 de l'indisponibilité du suivi des expéditions.
    \item La décision de faire signer un avenant à l'éditeur.
    \item Le chemin « faux fournisseur, portail, complice interne ».
\end{enumerate}
\end{boxexo}

\begin{boxexo}[title={\ding{45}\quad Exercice~4.2. Trois nouvelles parties prenantes \niveaumoyen}]
LiZbiZ envisage trois nouveaux partenaires~:
\begin{itemize}
    \item Q1, un hébergeur en nuage pour le portail de suivi~:
          dépendance~4, pénétration~3, maturité~3, confiance~2~;
    \item Q2, une société de maintenance des chambres froides, reliée au
          réseau pour la télésurveillance~: dépendance~3, pénétration~2,
          maturité~1, confiance~2~;
    \item Q3, un prestataire de paie~: dépendance~2, pénétration~1,
          maturité~2, confiance~3.
\end{itemize}
\begin{enumerate}[(a)]
    \item Calculez leur niveau de menace et placez-les dans les zones du
          chapitre.
    \item Les maturités ont été cotées sur la base d'une plaquette
          commerciale. Quelle cote de base leur donnez-vous~? Recalculez la
          menace de Q1 si sa maturité réelle est de~2.
    \item Proposez, pour Q2, deux mesures sur l'écosystème et la preuve
          qui permettrait de relever sa confiance.
\end{enumerate}
\end{boxexo}

\begin{boxexo}[title={\ding{45}\quad Exercice~4.3. Le mode opératoire de SS2 \niveaumoyen}]
Détaillez SS2 (le magasinier de Douala et son complice) en scénario
opérationnel.
\begin{enumerate}[(a)]
    \item Décrivez une action élémentaire pour chacune des quatre phases.
    \item Cotez la vraisemblance de chaque phase et justifiez par un fait du
          cas. Déduisez la vraisemblance du scénario selon la règle de
          chaîne.
    \item Quelle phase faut-il fermer en priorité~? Montrez l'effet sur la
          criticité.
\end{enumerate}
\end{boxexo}

\begin{boxexo}[title={\ding{45}\quad Exercice~4.4. EBIOS à la BANQUE SANAGA \niveauexpert}]
L'atelier de décembre 2026 de la BANQUE SANAGA a retenu le scénario
« fraude au virement par compromission d'un poste du siège »
(Chapitre~\ref{chap:iso27005}).
\begin{enumerate}[(a)]
    \item Atelier~1~: nommez la valeur métier, deux biens supports et
          l'événement redouté, et cotez sa gravité avec l'échelle du cours.
    \item Atelier~2~: proposez deux couples SR/OV, dont un que vous
          écarterez, en justifiant par la motivation, les ressources et
          l'activité.
    \item Atelier~3~: identifiez une partie prenante qui pourrait servir de
          point d'entrée, cotez ses quatre critères, et indiquez pour chacun
          la base de preuve dont vous disposez.
    \item Atelier~4~: rédigez le mode opératoire en quatre phases et
          déduisez la vraisemblance.
    \item Comparez la démarche avec la fiche R-AG-01~: qu'apporte EBIOS que
          la fiche de risque ne donnait pas~?
\end{enumerate}
\end{boxexo}

\bigskip

\begin{boxinfo}
\textbf{Pour aller plus loin.}
\begin{itemize}
    \item Le guide EBIOS Risk Manager de l'ANSSI~\cite{anssi2018ebios} et ses
          fiches méthodes, sur le site de l'ANSSI.
    \item Le guide NIST SP~800-30, révision~1~\cite{nist2012sp80030}.
    \item L'ISO/IEC~27005:2022~\cite{iso27005_2022}, pour l'articulation
          avec le processus.
\end{itemize}
\end{boxinfo}

\begin{boxinfo}
\textbf{Transition.} EBIOS a coté SO1 à 12~: intolérable. Mais combien
LiZbiZ ou la BANQUE SANAGA perdraient-elles, en francs CFA, si le scénario
se réalisait~? Et combien faut-il dépenser pour le traiter~? Les échelles
qualitatives ne répondent pas à ces questions. Le Chapitre~\ref{chap:fair}
apprend à quantifier.
\end{boxinfo}

        % ============================================================================
% Fichier  : partie2/P02_C05_quantifier_fair.tex
% Titre    : Quantifier le Risque Cyber : FAIR
% Auteur   : MINKA MI NGUIDJOÏ Thierry Emmanuel
% Version  : 1.0 -- novembre 2026
% Corrigés : corriges/Corr_C05.tex
% ============================================================================
% Sources : Open FAIR (The Open Group), O-RT 3.0.1 et O-RA 2.0.1 (versions
% vérifiées sur opengroup.org le 05/11/2026) ; taxonomie présentée d'après
% Freund et Jones (2015). Le renommage « vulnérabilité » → « susceptibilité »
% dans O-RT 3 est signalé à l'auteur comme À VÉRIFIER. Estimations calibrées
% : Hubbard et Seiersen (2016).
% Faits du cas (bible) : R-AG-01, fréquence 0,2 à 0,6, ampleur 150 à 900 M
% FCFA ; incident réel 410 M ; plan A1 à A8 = 193 M ; assurance (prime
% 35 M, franchise 100 M, plafond 1 Md).
% Faits nouveaux (reportés dans la bible) : valeurs les plus probables
% 0,35 et 410 M ; ventilation des 410 M (productivité 190, réponse 140,
% remplacement 80) ; après traitement : fréquence PERT(0,05 ; 0,10 ; 0,20),
% ampleur PERT(50 ; 150 ; 400) ; coût récurrent du plan 30 M par an.
% CHIFFRES CALCULÉS par simulation (scratchpad fair_sanaga.py, numpy,
% graine 2027, 100 000 itérations, fréquence PERT tirée à chaque itération,
% nombre d'événements de Poisson, ampleur PERT par événement) :
%   actuel : moyenne 164,6 ; P(L>0) 0,305 ; P50 0 ; P90 579,3 ; P95 733,6 ;
%   P99 1 169,8 ; P(L>500) 0,143 ; P(lambda>1/3) 0,640 ; P(k>=2) 0,054 ;
%   dépassements 100:0,305 250:0,288 500:0,143 750:0,047 1000:0,021.
%   après : moyenne 18,6 ; P(L>0) 0,101 ; P90 60,0 ; P95 174,4 ; P99 283,2 ;
%   P(L>500) 0,0006 ; dépassements 100:0,089 250:0,018 500:0,0006.
%   assurance : remboursement moyen 126,6 (actuel), 8,2 (après).
%   sensibilité, max 1 500 : moyenne 200,9 ; P99 1 507,0 ; P(L>500) 0,184.
%   table des 10 années : graine 7.
% ============================================================================

\chapter{Quantifier le Risque Cyber : FAIR}
\label{chap:fair}
\chaptermark{Quantifier le risque}

\epigraph{%
    « Un conseil ne vote pas une couleur.\\
    Il vote un budget. »%
}{Principe du cours}

\niveauChapitre{Intermédiaire}{%
    Maîtrise du Chapitre~\ref{chap:iso27005}%
}

\begin{boxobjectifs}
\begin{itemize}
    \item Expliquer pourquoi et quand quantifier un risque cyber.
    \item Décomposer un risque selon la taxonomie FAIR~: fréquence des
          événements de perte et ampleur des pertes.
    \item Produire des estimations calibrées (minimum, plus probable,
          maximum) et en coter la base de preuve.
    \item Conduire une simulation de Monte-Carlo pas à pas.
    \item Interpréter une distribution de pertes annuelles~: moyenne,
          quantiles, probabilité de dépassement.
    \item Comparer des traitements en francs CFA, y compris une assurance.
    \item Dire ce que la quantification apporte et ce qu'elle masque.
\end{itemize}
\end{boxobjectifs}

\bigskip

% ============================================================================
\section*{Avant de commencer : 193 millions, est-ce trop~?}
% ============================================================================

Douala, janvier 2027. Le comité des risques du conseil de la BANQUE
SANAGA\footnote{Comme dans tout le cours, la BANQUE SANAGA, ses
dirigeants, les dates et les montants sont fictifs~; la méthode et les
normes citées sont exactes.} examine le plan de traitement de R-AG-01~:
huit actions, 193~millions FCFA (Chapitre~\ref{chap:iso27005}). Le
directeur financier pose la question que les échelles ne savent pas
traiter~: « Vous me dites que le risque est coté 12. Combien nous
coûte-t-il~? Et combien nous coûtera-t-il après ces 193~millions~? »

La RSSI n'a pas de réponse en francs CFA. La cote~12 classe le risque
parmi les intolérables~; elle ne dit pas si 193~millions sont une bonne
affaire, une dépense insuffisante ou un luxe. Le comité lui demande de
revenir en mars avec une quantification.

\medskip

Ce chapitre présente la méthode qu'elle emploie, FAIR, et la déroule sur
R-AG-01. Tous les résultats numériques ont été obtenus par une simulation
réelle de 100\,000 années, dont le chapitre explique chaque étape.

\begin{boxalert}
\textbf{Violation typique.} Présenter au conseil une perte « attendue » de
143,5~millions, obtenue en multipliant une fréquence par une perte, comme
si c'était ce qui allait arriver.

\smallskip
\textbf{Règle~:} une quantification se présente comme une distribution~:
la plupart des années ne coûtent rien, quelques-unes coûtent très cher. On
donne au moins la moyenne, la probabilité d'une perte et un quantile élevé,
avec les hypothèses et leur base de preuve.
\end{boxalert}

\separateur

% ============================================================================
\section{Pourquoi Quantifier}
\label{sec:c05-pourquoi}
% ============================================================================

Les échelles qualitatives ont deux mérites~: elles sont rapides, et
chacun les comprend. Le Chapitre~\ref{chap:risque} en a montré les
limites~: niveaux ordinaux, produit qui écrase les écarts, frontières
arbitraires. Ces limites deviennent bloquantes dans trois situations~:
\begin{itemize}
    \item \textbf{comparer un risque à une dépense}~: un plan de
          193~millions, une prime d'assurance, un recrutement~;
    \item \textbf{comparer des risques de nature différente}~: un risque
          cyber et un risque de crédit, que le conseil arbitre avec le même
          budget~;
    \item \textbf{une cote qui chevauche un seuil}~: la fréquence de
          R-AG-01, estimée entre 0,2 et 0,6 par an, traverse les niveaux~2
          et~3 de l'échelle, donc la frontière entre «~à traiter dans
          l'année~» et «~intolérable~» (Chapitre~\ref{chap:risque}).
\end{itemize}

On ne quantifie pas tout. La quantification coûte du temps et des
compétences~; on la réserve aux risques qui pèsent sur une décision
importante, typiquement ceux d'impact critique ou dont la cote traverse un
seuil.

% ============================================================================
\section{La Taxonomie FAIR}
\label{sec:c05-taxonomie}
% ============================================================================

\subsection{Une décomposition}

FAIR (\emph{Factor Analysis of Information Risk}) est un modèle de
quantification du risque d'information. Il est normalisé par The Open
Group sous le nom d'Open FAIR, en deux normes~: la taxonomie du risque
(O-RT, version~3.0.1) et l'analyse du risque (O-RA,
version~2.0.1)~\cite{opengroup2026openfair}. Le cours présente la
taxonomie d'après l'ouvrage de référence de Freund et
Jones~\cite{freund2015fair}.

\begin{boxdef}
\termecle{Risque (au sens de FAIR)}\index{FAIR}~: fréquence et ampleur
probables des pertes futures. Il se décompose en deux branches~:
\begin{itemize}
    \item la \textbf{fréquence des événements de perte}~: combien de fois
          par an un événement de menace aboutit à une perte~;
    \item l'\textbf{ampleur des pertes}~: combien coûte chaque événement.
\end{itemize}
\end{boxdef}

\begin{figure}[H]
\centering
\begin{tikzpicture}[
    noeud/.style={draw=CouleurPrimaire, thick, rounded corners,
                  fill=CouleurPrimaire!8, text width=2.3cm,
                  minimum height=0.9cm, align=center, font=\scriptsize},
    racine/.style={noeud, fill=CouleurAccent!15, draw=CouleurAccent},
    lien/.style={thick, CouleurSecondaire}]
    \node[racine] (R) at (0,0) {\textbf{Risque}};
    \node[noeud] (F) at (-3.4,-1.5) {Fréquence des événements de perte};
    \node[noeud] (A) at (3.4,-1.5) {Ampleur des pertes};
    \node[noeud] (T) at (-5.2,-3.1) {Fréquence des événements de menace};
    \node[noeud] (V) at (-1.8,-3.1) {Vulnérabilité};
    \node[noeud] (C) at (-6.0,-4.7) {Fréquence de contact};
    \node[noeud] (P) at (-3.4,-4.7) {Probabilité d'action};
    \node[noeud] (Cap) at (-1.0,-4.7) {Capacité de la menace};
    \node[noeud] (Res) at (1.4,-4.7) {Force de résistance};
    \node[noeud] (PP) at (2.0,-3.1) {Pertes primaires};
    \node[noeud] (PS) at (4.8,-3.1) {Pertes secondaires};
    \draw[lien] (R) -- (F); \draw[lien] (R) -- (A);
    \draw[lien] (F) -- (T); \draw[lien] (F) -- (V);
    \draw[lien] (T) -- (C); \draw[lien] (T) -- (P);
    \draw[lien] (V) -- (Cap); \draw[lien] (V) -- (Res);
    \draw[lien] (A) -- (PP); \draw[lien] (A) -- (PS);
\end{tikzpicture}
\caption{La taxonomie FAIR (d'après Freund et Jones, 2015)}
\label{fig:c05-taxonomie}
\end{figure}

\begin{tableaucours}{L{3.6cm}L{8.6cm}}{%
    \tableauHeader{Facteur} &
    \tableauHeader{Définition}
}
Fréquence des événements de menace & Nombre de fois par an où une source
de menace agit contre l'actif~: produit de la fréquence de contact et de
la probabilité qu'elle passe à l'action \\
Vulnérabilité & Probabilité qu'un événement de menace devienne un
événement de perte~: elle dépend de la capacité de la menace comparée à la
force de résistance des mesures \\
Pertes primaires & Pertes subies directement par l'organisation du fait de
l'événement \\
Pertes secondaires & Pertes qui viennent des réactions de tiers~: clients,
régulateur, tribunaux, presse~; elles ont leur propre fréquence et leur
propre ampleur \\
\end{tableaucours}
\vspace{-0.8em}
\begin{center}{\small\itshape Tableau~5.1. Les facteurs de la taxonomie}\end{center}

\begin{boxinfo}
\textbf{À vérifier dans la norme.} La version~3 de la taxonomie Open FAIR
a révisé certains termes~; on y lit notamment \emph{susceptibilité} là où
les textes antérieurs parlent de \emph{vulnérabilité}. Le cours garde le
vocabulaire de Freund et Jones, plus répandu, et signale l'écart.
\end{boxinfo}

\subsection{Six formes de pertes}

FAIR ventile l'ampleur en six formes de pertes~: la \emph{productivité}
(activité perdue), la \emph{réponse} (gestion de l'incident,
communication, conseils), le \emph{remplacement} (actifs à reconstruire
ou racheter), les \emph{amendes et jugements}, l'\emph{avantage
concurrentiel} perdu, et la \emph{réputation}~\cite{freund2015fair}.

\begin{boxscenario}{Opération SANAGA. Les 410 millions, ventilés}
\textbf{Constatation.} La direction financière ventile le coût de
l'incident de septembre 2026~:
\begin{itemize}
    \item productivité~: 190~millions (neuf agences fermées trois jours,
          opérations perdues ou reportées)~;
    \item réponse~: 140~millions (prestataire de réponse à incident,
          cellule de crise, communication auprès de la clientèle)~;
    \item remplacement~: 80~millions (reconstruction de 180 postes)~;
    \item amendes et jugements, avantage concurrentiel~: néant~;
    \item réputation~: non chiffrée.
\end{itemize}

\textbf{Analyse} (confiance élevée sur les pertes primaires, faible sur
les secondaires). Les 410~millions sont presque entièrement des pertes
primaires. La perte de clients qui ont changé de banque n'est pas mesurée~:
elle est \NEK{}. Le coût réel de l'incident est donc d'au moins
410~millions, pas de 410~millions.
\end{boxscenario}

% ============================================================================
\section{Estimer}
\label{sec:c05-estimer}
% ============================================================================

\subsection{Des estimations calibrées}

FAIR ne demande pas des valeurs exactes, que personne n'a, mais des
\termecle{estimations calibrées}\index{estimation calibrée}~: pour chaque
facteur, un minimum, une valeur la plus probable et un maximum, choisis de
sorte que l'expert soit confiant, à 90~\% environ, que la vraie valeur se
trouve entre le minimum et le maximum. Hubbard et Seiersen montrent que
cette confiance s'apprend, par des exercices de calibrage, et que la
plupart des experts non entraînés donnent des intervalles trop
étroits~\cite{hubbard2016measure}.

\subsection{Les estimations de R-AG-01}

En février 2027, la RSSI anime un atelier avec la direction financière, la
direction du réseau d'agences et le prestataire de réponse à incident.

\begin{tableaucours}{L{3.0cm}ccc L{4.4cm}}{%
    \tableauHeader{Facteur} &
    \tableauHeader{Min.} &
    \tableauHeader{Plus prob.} &
    \tableauHeader{Max.} &
    \tableauHeader{Justification}
}
Fréquence des événements de perte (par an) & 0,2 & 0,35 & 0,6 & Un
incident réussi en 2026~; tentatives fréquentes dans le secteur~; défenses
de 2025 \\
Ampleur par événement (M FCFA) & 150 & 410 & 900 & Minimum~: trois agences
touchées~; plus probable~: l'incident vécu~; maximum~: propagation à tout
le réseau et pertes secondaires \\
\end{tableaucours}
\vspace{-0.8em}
\begin{center}{\small\itshape Tableau~5.2. Estimations calibrées pour
R-AG-01 (défenses de 2025)}\end{center}

Pour la fréquence, l'atelier estime directement la fréquence des
événements de perte, sans descendre jusqu'à la fréquence de contact et à la
vulnérabilité~: la taxonomie permet de s'arrêter au niveau où l'on dispose
de la meilleure information.

\subsection{La base de preuve des estimations}

\begin{boxlecture}{Base de preuve des six valeurs}
\textbf{Constatation.} La RSSI cote la base de chacune des six valeurs
du Tableau~5.2~:
\begin{itemize}
    \item fréquence~: minimum \NEK{} (avis non motivé), plus probable
          \OUI{} (incident de 2026, données sectorielles d'un assureur),
          maximum \NEK{}~;
    \item ampleur~: minimum \OUI{} (devis de reconstruction pour trois
          agences), plus probable \OUI{} (coût de l'incident), maximum
          \NEK{} (pertes secondaires non chiffrées).
\end{itemize}
3~\OUI{}, 0~\NON{}, 3~\NEK{}~: $\mu = 0{,}50$~; $U = 0{,}50$~;
$H = -(2 \times 0{,}5\log_2 0{,}5) = 1{,}00$~bit.

\medskip
\textbf{Analyse} (confiance élevée). La quantification est solide en son
centre et incertaine dans ses queues. Or ce sont les queues qui
intéressent le conseil~: le maximum d'ampleur décide du quantile à
99~\%, et le maximum de fréquence de la probabilité d'avoir deux incidents
la même année. La Section~\ref{sec:c05-masque} mesure ce que pèse cette
incertitude.

\medskip
\hfill\lecture{0{,}50}{0{,}50}{1{,}00}
\end{boxlecture}

% ============================================================================
\section{La Simulation de Monte-Carlo}
\label{sec:c05-montecarlo}
% ============================================================================

\subsection{Le principe}

Multiplier la fréquence la plus probable par l'ampleur la plus probable
donne $0{,}35 \times 410 = 143{,}5$~millions. Ce chiffre n'est ni ce qui
arrivera une année donnée (le plus souvent, rien), ni même la moyenne,
parce qu'il ignore la forme des distributions. La \termecle{simulation de
Monte-Carlo}\index{simulation de Monte-Carlo} remplace ce calcul par un
grand nombre d'années tirées au hasard, selon les estimations.

\begin{boxprocedure}
\textbf{Simuler une année.}
\begin{enumerate}
    \item Tirer une fréquence $\lambda$ dans la distribution de fréquence
          (ici une loi PERT, qui suit le minimum, le plus probable et le
          maximum).
    \item Tirer un nombre d'événements dans l'année, selon une loi de
          Poisson de paramètre $\lambda$~: le plus souvent 0, parfois 1,
          rarement 2 ou plus.
    \item Pour chaque événement, tirer une ampleur dans la distribution
          d'ampleur (loi PERT).
    \item La perte de l'année est la somme des ampleurs.
\end{enumerate}
On répète l'opération un grand nombre de fois (ici 100\,000), puis on lit
la distribution des pertes annuelles obtenues. Le papier de travail
conserve le programme, la graine du générateur aléatoire et les
estimations, pour que le résultat puisse être reproduit.
\end{boxprocedure}

\begin{boxinfo}
\textbf{La loi PERT.} Elle se construit à partir d'un minimum $a$, d'une
valeur la plus probable $m$ et d'un maximum $b$. Sa moyenne vaut
$\frac{a + 4m + b}{6}$~: pour la fréquence,
$\frac{0{,}2 + 1{,}4 + 0{,}6}{6} \approx 0{,}367$~; pour l'ampleur,
$\frac{150 + 1\,640 + 900}{6} \approx 448$~millions. La moyenne de
l'ampleur dépasse la valeur la plus probable, parce que le maximum est plus
éloigné du mode que le minimum.
\end{boxinfo}

\subsection{Dix années, à la main}

Pour comprendre la mécanique, voici les dix premières années d'une
simulation.

\begin{tableaucours}{cccL{3.4cm}r}{%
    \tableauHeader{Année} &
    \tableauHeader{$\lambda$ tirée} &
    \tableauHeader{Événements} &
    \tableauHeader{Ampleurs (M FCFA)} &
    \tableauHeader{Perte}
}
1 & 0,38 & 0 & -- & 0 \\
2 & 0,26 & 0 & -- & 0 \\
3 & 0,39 & 0 & -- & 0 \\
4 & 0,43 & 1 & 385 & 385 \\
5 & 0,38 & 0 & -- & 0 \\
6 & 0,26 & 0 & -- & 0 \\
7 & 0,43 & 0 & -- & 0 \\
8 & 0,35 & 0 & -- & 0 \\
9 & 0,40 & 1 & 521 & 521 \\
10 & 0,35 & 0 & -- & 0 \\
\textbf{Moyenne} & & & & \textbf{90,6} \\
\end{tableaucours}
\vspace{-0.8em}
\begin{center}{\small\itshape Tableau~5.3. Dix années simulées}\end{center}

Huit années sur dix ne coûtent rien~; deux coûtent 385 et 521~millions.
La moyenne de ces dix années, 90,6~millions, est loin de la moyenne
obtenue sur 100\,000 années (164,6~millions)~: dix tirages ne suffisent
pas. C'est la leçon de la simulation, et aussi celle de l'expérience~: une
banque qui n'a pas eu d'incident pendant huit ans n'en déduit rien sur la
neuvième.

\subsection{Cent mille années}

\begin{tableaucours}{L{6.6cm}rr}{%
    \tableauHeader{Indicateur (M FCFA, par an)} &
    \tableauHeader{Défenses 2025} &
    \tableauHeader{Après traitement}
}
Perte annuelle moyenne & 164,6 & 18,6 \\
Probabilité d'une perte dans l'année & 30,5~\% & 10,1~\% \\
Médiane (une année sur deux) & 0 & 0 \\
Quantile à 90~\% (une année sur dix au-delà) & 579 & 60 \\
Quantile à 95~\% (une année sur vingt) & 734 & 174 \\
Quantile à 99~\% (une année sur cent) & 1\,170 & 283 \\
Probabilité d'une perte supérieure à 500~M & 14,3~\% & 0,06~\% \\
\end{tableaucours}
\vspace{-0.8em}
\begin{center}{\small\itshape Tableau~5.4. Distribution des pertes annuelles
dues à R-AG-01 (100\,000 années simulées)}\end{center}

La colonne «~après traitement~» reprend la simulation avec les estimations
de l'atelier pour la situation visée par le plan~: fréquence entre 0,05 et
0,20 (plus probable 0,10), grâce au second facteur et au cloisonnement~;
ampleur entre 50 et 400~millions (plus probable 150), grâce aux
sauvegardes testées et au mode dégradé.

\begin{tableaucours}{cccccc}{%
    \tableauHeader{Perte supérieure à (M)} &
    \tableauHeader{0} &
    \tableauHeader{250} &
    \tableauHeader{500} &
    \tableauHeader{750} &
    \tableauHeader{1\,000}
}
Défenses 2025 & 30,5~\% & 28,8~\% & 14,3~\% & 4,7~\% & 2,1~\% \\
Après traitement & 10,1~\% & 1,8~\% & 0,06~\% & 0~\% & 0~\% \\
\end{tableaucours}
\vspace{-0.8em}
\begin{center}{\small\itshape Tableau~5.5. Probabilités de dépassement
(courbe de dépassement, en points)}\end{center}

% ============================================================================
\section{Interpréter et Décider}
\label{sec:c05-decider}
% ============================================================================

\subsection{Lire la distribution}

\begin{boxlecture}{Ce que la RSSI présente au comité, en mars 2027}
\textbf{Constatation.}
\begin{itemize}
    \item Avec les défenses de 2025, une année sur trois environ connaît
          un incident~; la perte moyenne est de 165~millions par an, et une
          année sur sept dépasse 500~millions.
    \item Après traitement, une année sur dix connaît un incident~; la
          perte moyenne tombe à 19~millions, et une perte de plus de
          500~millions devient exceptionnelle (moins d'une année sur mille).
\end{itemize}

\textbf{Analyse} (confiance moyenne~: voir la base de preuve des
estimations). Le plan réduit la perte moyenne d'environ 146~millions par
an. Il coûte 193~millions la première année, puis environ 30~millions par
an (licences, tests trimestriels, exercices). Le gain net annuel, de
l'ordre de 116~millions, rembourse l'investissement en moins de deux ans.
Sur ce seul risque, 193~millions ne sont pas un luxe.
\end{boxlecture}

\subsection{L'appétence en francs CFA}

La quantification permet aussi d'exprimer l'appétence autrement qu'en
couleurs. La RSSI propose au comité une formulation~: « la probabilité
qu'un même risque cyber cause, une année donnée, une perte supérieure à
500~millions FCFA ne dépasse pas 5~\% ». Avec les défenses de 2025,
R-AG-01 la viole largement (14,3~\%)~; après traitement, il la respecte
(0,06~\%).

La simulation tranche aussi la question laissée ouverte au
Chapitre~\ref{chap:risque}~: la fréquence tirée dépasse 1/3 par an dans
64~\% des années simulées. Le niveau~3 de vraisemblance est donc plus
probable que le niveau~2~; la cote prudente de 12 était justifiée.

\subsection{L'assurance, chiffrée}

L'offre du courtier (prime de 35~millions par an, franchise de
100~millions par événement, plafond de 1~milliard) se juge maintenant sur
les mêmes simulations.

\begin{tableaucours}{L{6.4cm}rr}{%
    \tableauHeader{M FCFA par an} &
    \tableauHeader{Défenses 2025} &
    \tableauHeader{Après traitement}
}
Perte moyenne sans assurance & 164,6 & 18,6 \\
Remboursement moyen de l'assureur & 126,6 & 8,2 \\
Coût moyen avec assurance (perte $-$ remboursement $+$ prime) & 73,0 &
45,4 \\
\end{tableaucours}
\vspace{-0.8em}
\begin{center}{\small\itshape Tableau~5.6. L'assurance, avant et après
traitement}\end{center}

Avant traitement, l'assurance divise le coût moyen par plus de deux. Après
traitement, elle coûte en moyenne plus qu'elle ne rapporte (45,4 contre
18,6). Deux conclusions~: l'assurance est utile comme protection
temporaire, pendant la mise en œuvre du plan~; ensuite, elle ne se justifie
que si la prime baisse, ou pour protéger la banque contre les années
extrêmes que la simulation estime mal. Et l'on ne pourra pas l'obtenir sans
les mesures, que les assureurs exigent~: la réduction vient d'abord.

% ============================================================================
\section{Ce que la Quantification Masque}
\label{sec:c05-masque}
% ============================================================================

\subsection{Des chiffres faits d'opinions}

Une simulation transforme des estimations en distribution~; elle ne les
rend pas plus vraies. Trois des six valeurs d'entrée sont \NEK{}. Pour
mesurer ce que cela pèse, la RSSI relance la simulation en portant le
maximum d'ampleur de 900 à 1\,500~millions, pour tenir compte des pertes
secondaires non chiffrées.

\begin{tableaucours}{L{6.6cm}rr}{%
    \tableauHeader{Indicateur (M FCFA)} &
    \tableauHeader{Max. 900} &
    \tableauHeader{Max. 1\,500}
}
Perte annuelle moyenne & 164,6 & 200,9 \\
Quantile à 99~\% & 1\,170 & 1\,507 \\
Probabilité d'une perte supérieure à 500~M & 14,3~\% & 18,4~\% \\
\end{tableaucours}
\vspace{-0.8em}
\begin{center}{\small\itshape Tableau~5.7. Sensibilité au maximum
d'ampleur}\end{center}

Une seule valeur \NEK{}, modifiée, déplace la moyenne de 22~\% et le
quantile à 99~\% de 29~\%. La décision de traiter ne change pas~; mais le
montant de couverture d'assurance, lui, en dépend. C'est le rôle de
l'analyse de sensibilité~: identifier les \NEK{} dont la décision dépend,
et y consacrer l'effort de collecte de preuves.

\subsection{Cinq angles morts}

\begin{itemize}
    \item \textbf{Le scénario absent.} On ne simule que ce que l'on a
          imaginé. En 2025, R-AG-01 n'existait pas dans le registre~:
          aucune simulation ne l'aurait chiffré.
    \item \textbf{La fausse précision.} « 164,6~millions » suggère une
          exactitude que les estimations n'ont pas. On présente des ordres
          de grandeur arrondis («~environ 165~millions~»).
    \item \textbf{L'indépendance supposée.} Le modèle tire les événements
          indépendamment. Une même faiblesse peut provoquer plusieurs
          incidents corrélés, et une crise du secteur toucher plusieurs
          banques à la fois.
    \item \textbf{L'adversaire qui s'adapte.} Les estimations décrivent la
          menace de 2027~; un attaquant qui change de méthode les rend
          caduques (Chapitre~\ref{chap:risque}, Tableau~2.2).
    \item \textbf{L'incertitude de connaissance.} La distribution mêle
          l'aléa (l'attaque aura-t-elle lieu~?) et l'ignorance (la vraie
          fréquence)~; elle ne dit pas quelle part vient de chacun. La cote
          de la base de preuve des estimations le dit.
\end{itemize}

\begin{boxinfo}
\textbf{La quantification et l'auditologie.} Le cours tient les deux
ensemble~: la simulation décrit l'incertitude \emph{aléatoire} sur les
pertes, sous des hypothèses~; la cote \OUI{}~/ \NON{}~/ \NEK{} des
hypothèses décrit l'incertitude \emph{épistémique} sur ces hypothèses. Une
quantification présentée sans la seconde est incomplète~; le
Chapitre~\ref{chap:auditologie} montre comment les combiner.
\end{boxinfo}

\begin{boxinfo}
\textbf{Et le cadre juridique~?} Aucun texte camerounais n'impose de
quantifier les risques cyber. L'article~26 de la loi n°~2010/012 demande
de les \emph{évaluer} (Chapitre~\ref{chap:iso27005})~\cite{loi2010012}~;
la quantification est une manière de le faire, que le conseil, le
superviseur et l'assureur comprennent, parce qu'elle parle leur langue~:
l'argent.
\end{boxinfo}

% ============================================================================
\section{Ce que la Quantification demande à chaque Métier}
\label{sec:c05-metiers}
% ============================================================================

\begin{boxmetier}{Gestionnaire des risques cyber}
\textbf{Quantifier ce qui décide}
\begin{enumerate}
    \item Réserver la quantification aux risques dont dépend une décision
          importante.
    \item Faire produire des estimations calibrées, par des experts
          entraînés, et coter leur base de preuve.
    \item Présenter une distribution (moyenne, probabilité, quantile élevé),
          jamais un seul chiffre.
    \item Tester la sensibilité aux valeurs \NEK{}.
\end{enumerate}
\end{boxmetier}

\begin{boxmetier}{RSSI}
\textbf{Défendre un budget}

La RSSI ne demande plus 193~millions pour « réduire un risque coté 12 ».
Elle montre une réduction de perte moyenne de l'ordre de 146~millions par
an, et le passage de la probabilité d'une année à plus de 500~millions de
14~\% à moins de 0,1~\%. Le conseil vote un budget, avec ses hypothèses.
\end{boxmetier}

\begin{boxmetier}{Auditeur en cybersécurité}
\textbf{Auditer les entrées, pas le logiciel}

L'auditeur ne refait pas les 100\,000 tirages~; il vérifie que le programme
et la graine sont conservés, que chaque estimation a une justification
écrite, que les auteurs ont été calibrés, et que la sensibilité aux
valeurs non fondées a été testée.
\end{boxmetier}

\begin{boxmetier}{Responsable GRC}
\textbf{Une même monnaie}

Le responsable GRC propose au conseil une appétence exprimée en francs
CFA, applicable à tous les risques quantifiés, cyber ou non. Les risques
deviennent comparables, et les arbitrages budgétaires, explicables.
\end{boxmetier}

\separateur

\begin{boxresume}
\textbf{Ce qu'il faut retenir de ce chapitre~:}
\begin{itemize}
    \item \textbf{Quantifier} sert à comparer un risque à une dépense, des
          risques entre eux, ou à trancher une cote qui chevauche un seuil.
    \item \textbf{FAIR} décompose le risque en fréquence des événements de
          perte et ampleur des pertes~; l'ampleur se ventile en six formes
          de pertes, primaires et secondaires.
    \item \textbf{Des estimations calibrées} (minimum, plus probable,
          maximum) remplacent les valeurs exactes que personne n'a~; leur
          base de preuve se cote.
    \item \textbf{La simulation de Monte-Carlo} tire des années~; on lit la
          moyenne, la probabilité d'une perte et les quantiles élevés.
    \item \textbf{R-AG-01}~: environ 165~millions de perte moyenne par an
          avant traitement, 19 après~; le plan de 193~millions se rembourse
          en moins de deux ans.
    \item \textbf{La quantification masque} le scénario absent, la fausse
          précision, les corrélations, l'adversaire qui s'adapte, et
          l'incertitude de connaissance.
\end{itemize}
\end{boxresume}

\begin{boxinfo}
\textbf{Points de contrôle.} Vous avez compris ce chapitre si vous savez~:
\begin{itemize}
    \item placer un facteur dans la taxonomie FAIR~;
    \item ventiler un coût d'incident en formes de pertes~;
    \item simuler quelques années à la main et dire pourquoi elles ne
          suffisent pas~;
    \item lire une distribution et en tirer une décision de traitement ou
          d'assurance.
\end{itemize}
\end{boxinfo}

% ============================================================================
\section*{Exercices}
% ============================================================================

\begin{boxexo}[title={\ding{45}\quad Exercice~5.1. Six formes de pertes \niveaubase}]
Classez chacun des coûts suivants dans l'une des six formes de pertes de
FAIR, et dites s'il s'agit d'une perte primaire ou secondaire.
\begin{enumerate}[(a)]
    \item Les heures supplémentaires des agents pour ressaisir les
          opérations.
    \item Le rachat de vingt postes de travail.
    \item L'amende infligée par le superviseur pour notification tardive.
    \item Les honoraires du prestataire de réponse à incident.
    \item Les clients qui ferment leur compte après l'incident.
    \item Le contrat d'une grande entreprise remporté par une banque
          concurrente pendant la fermeture des guichets.
\end{enumerate}
\end{boxexo}

\begin{boxexo}[title={\ding{45}\quad Exercice~5.2. Le chiffre trompeur \niveaumoyen}]
\begin{enumerate}[(a)]
    \item Calculez les moyennes PERT de la fréquence et de l'ampleur pour la
          situation après traitement (Section~\ref{sec:c05-montecarlo}), et
          leur produit. Comparez au résultat de la simulation (18,6).
    \item Pourquoi la médiane de la perte annuelle est-elle nulle, avant
          comme après traitement~?
    \item Un administrateur propose de provisionner chaque année
          165~millions «~pour couvrir le risque~». Expliquez ce que cette
          provision couvre et ce qu'elle ne couvre pas, à l'aide du
          Tableau~5.4.
\end{enumerate}
\end{boxexo}

\begin{boxexo}[title={\ding{45}\quad Exercice~5.3. Une autre offre d'assurance \niveaumoyen}]
Un second assureur propose, après traitement, une prime de 12~millions par
an, une franchise de 200~millions par événement et un plafond de
500~millions. On admet qu'après traitement il rembourserait en moyenne
1,5~million par an.
\begin{enumerate}[(a)]
    \item Calculez le coût moyen annuel avec cette assurance et comparez-le
          à l'absence d'assurance et à la première offre.
    \item À l'aide du Tableau~5.5, dites dans quelles années cette
          assurance joue, et ce qu'elle protège réellement.
    \item Recommandez une décision au comité, en deux phrases.
\end{enumerate}
\end{boxexo}

\begin{boxexo}[title={\ding{45}\quad Exercice~5.4. Quantifier SO1 pour PHARMACAM \niveauexpert}]
LiZbiZ veut présenter à PHARMACAM une quantification simple du scénario
SO1 (Chapitre~\ref{chap:ebios}). L'atelier retient~: un événement au plus
par an, avec une probabilité de 0,25~; une ampleur triangulaire de minimum
100, plus probable 300, maximum 800~millions FCFA (moyenne~: le tiers de la
somme des trois valeurs).
\begin{enumerate}[(a)]
    \item Calculez la perte annuelle moyenne.
    \item Simulez dix années à la main avec les nombres aléatoires suivants
          (un événement survient si le nombre est inférieur à 0,25)~: 0,81~;
          0,12~; 0,55~; 0,37~; 0,09~; 0,66~; 0,93~; 0,21~; 0,48~; 0,70.
          Les ampleurs tirées pour les événements sont, dans l'ordre, 320,
          610 et 180~millions. Calculez la moyenne des dix années et
          commentez.
    \item Cotez la base de preuve des quatre valeurs d'entrée, sachant que
          LiZbiZ n'a jamais subi de rançongiciel et que l'ampleur
          maximale correspond à l'estimation de PHARMACAM pour un camion de
          vaccins perdu.
    \item Rédigez le paragraphe destiné à PHARMACAM~: ce que la
          quantification dit, ce qu'elle suppose, ce qu'elle ne dit pas.
\end{enumerate}
\end{boxexo}

\bigskip

\begin{boxinfo}
\textbf{Pour aller plus loin.}
\begin{itemize}
    \item Freund et Jones, \emph{Measuring and Managing Information
          Risk}~\cite{freund2015fair}, l'exposé de référence de FAIR.
    \item Les normes Open FAIR de The Open Group~\cite{opengroup2026openfair}.
    \item Hubbard et Seiersen, \emph{How to Measure Anything in
          Cybersecurity Risk}~\cite{hubbard2016measure}, sur le calibrage
          des experts.
\end{itemize}
\end{boxinfo}

\begin{boxinfo}
\textbf{Transition.} La Partie~II a jusqu'ici appris à apprécier, traiter
et chiffrer les risques. Reste à vérifier que les mesures annoncées
existent et fonctionnent~: c'est le travail de l'auditeur. Le
Chapitre~\ref{chap:audit} présente la conduite d'un audit de
cybersécurité, du mandat au rapport.
\end{boxinfo}

        % ============================================================================
% Fichier  : partie2/P02_C06_conduire_audit_cyber.tex
% Titre    : Conduire un Audit de Cybersécurité
% Auteur   : MINKA MI NGUIDJOÏ Thierry Emmanuel
% Version  : 1.0 -- novembre 2026
% Corrigés : corriges/Corr_C06.tex
% ============================================================================
% Sources : ISO 19011:2018 (sept principes, programme et conduite d'audit :
% paraphrasés) ; ISO/IEC 27007:2020 ; IIA, normes mondiales d'audit interne
% (2024) ; ITAF 4e éd. (2020) ; NIST SP 800-53A rév. 5 (méthodes examiner,
% interroger, tester) ; loi n° 2010/012, art. 4 point 9, art. 13 al. 1,
% art. 32 al. 1 et 2 (texte français officiel, Ann. D d'ACSI).
% Faits du cas (bible) : BANQUE SANAGA, Mme EKAMBI (audit interne), plan
% A1 à A8, état de septembre 2027 (5/2/3), mission sur les agences
% programmée pour 2028 (Ex. 1.4), GAB de 14 agences sous un système sans
% correctifs, droits d'administration locale des chefs d'agence dans 31
% agences (C03).
% Faits nouveaux (reportés dans la bible) : mission « sécurité du réseau
% d'agences », février à avril 2028 ; échantillon de 8 agences sur 38 ;
% 1 240 comptes ; 8 tests T1 à T8 : 3 OUI / 3 NON / 2 NEK ; Kribi échoue à
% l'exercice de mode dégradé ; 31 agences sur 38 journalisées ; 3 agences
% sur 8 gardent l'administration locale ; le prestataire des GAB refuse
% l'accès ; contestation de la DSI à la réunion de clôture.
% ============================================================================

\chapter{Conduire un Audit de Cybersécurité}
\label{chap:audit}
\chaptermark{Conduire un audit}

\epigraph{%
    « L'auditeur ne croit pas.\\
    Il vérifie, ou il dit qu'il n'a pas pu. »%
}{Principe du cours}

\niveauChapitre{Fondamental}{%
    Maîtrise du Chapitre~\ref{chap:grc}%
}

\begin{boxobjectifs}
\begin{itemize}
    \item Définir l'audit, ses critères, ses preuves, ses constats et ses
          conclusions.
    \item Appliquer les sept principes de l'ISO~19011 et situer
          l'ISO/IEC~27007.
    \item Distinguer audit interne, audit de certification, audit
          contractuel et audit réglementaire.
    \item Décrire le cycle d'une mission, du mandat au suivi, et organiser
          le contradictoire.
    \item Choisir une technique de collecte et apprécier la fiabilité d'une
          preuve.
    \item Échantillonner, et dire ce qu'un échantillon permet de conclure.
    \item Coter un test \OUI{}, \NON{} ou \NEK{}, et traiter une limitation
          de l'étendue.
\end{itemize}
\end{boxobjectifs}

\bigskip

% ============================================================================
\section*{Avant de commencer : « Les actions sont faites »}
% ============================================================================

Douala, février 2028. Le plan de traitement de R-AG-01 est déclaré achevé.
Le rapport de la RSSI au comité des risques l'affirme~: second facteur,
cloisonnement, sauvegardes, mode dégradé, sensibilisation, journalisation,
tout est en place\footnote{Comme dans tout le cours, la BANQUE SANAGA, ses
dirigeants, ses prestataires, les dates et les montants sont fictifs~; les
normes et les textes cités sont exacts.}. Le comité ne doute pas de la
bonne foi de Mme~ABENA. Mais il a appris, en 2026, ce que vaut une
assurance qu'on ne vérifie pas. Il confirme donc la mission inscrite au
programme de l'audit interne~: la sécurité du réseau d'agences.

Mme~EKAMBI, directrice de l'audit interne, sait que la mission sera
regardée de près. Si elle conclut que tout va bien, le comité réduira la
surveillance~; si elle se trompe, la prochaine attaque le dira. Elle doit
donc pouvoir montrer, pour chaque conclusion, la preuve qui la fonde, ou
dire clairement qu'elle ne l'a pas obtenue.

\medskip

Ce chapitre décrit la conduite d'un audit de cybersécurité, de l'audit
interne à l'audit réglementaire. Le Chapitre~\ref{chap:conclure} apprendra
ensuite à décomposer l'objet audité et à conclure.

\begin{boxalert}
\textbf{Violation typique.} Conclure qu'une mesure est en place sur la
foi d'un rapport de l'audité, d'une capture d'écran qu'il a fournie, ou
d'un entretien.

\smallskip
\textbf{Règle~:} une déclaration n'est pas une preuve de fonctionnement.
L'auditeur obtient lui-même la preuve, ou une preuve produite par le
système sous ses yeux~; à défaut, il cote \NEK{} et le dit.
\end{boxalert}

\separateur

% ============================================================================
\section{Ce qu'est un Audit}
\label{sec:c06-definir}
% ============================================================================

\subsection{Cinq notions}

\begin{boxdef}
\termecle{Audit}\index{audit}~: processus méthodique, indépendant et
documenté, qui consiste à obtenir des preuves objectives et à les évaluer
de manière objective pour déterminer dans quelle mesure des critères sont
satisfaits~\cite{iso19011_2018}.

\medskip
\termecle{Critère d'audit}\index{critère d'audit}~: ensemble d'exigences
qui sert de référence~: une norme, une politique, une loi, un contrat, un
plan de traitement.

\medskip
\termecle{Preuve d'audit}\index{preuve d'audit}~: enregistrement,
énoncé de fait ou autre information, vérifiable et pertinent au regard des
critères.

\medskip
\termecle{Constat d'audit}\index{constat d'audit}~: résultat de
l'évaluation des preuves au regard des critères~: conformité,
non-conformité, ou impossibilité de conclure.

\medskip
\termecle{Conclusion d'audit}\index{conclusion d'audit}~: résultat de
l'audit, tiré de l'ensemble des constats, au regard des objectifs de la
mission.
\end{boxdef}
{\small\itshape Définitions reformulées d'après l'ISO~19011:2018~; la
troisième issue du constat est propre au cours.}

\medskip

L'ISO~19011 envisage deux issues pour un constat~: conforme ou non
conforme. Le cours en ajoute une troisième, \NEK{}~: la preuve n'a pas pu
être obtenue, ou ne permet pas de trancher. Ce n'est pas une
complaisance~; c'est une information, souvent la plus importante du
rapport (Chapitre~\ref{chap:conclure}).

\subsection{Sept principes}

L'ISO~19011:2018 pose sept principes, qui s'appliquent à tout audit de
système de management~\cite{iso19011_2018}.

\begin{tableaucours}{L{3.6cm}L{8.6cm}}{%
    \tableauHeader{Principe} &
    \tableauHeader{Ce qu'il signifie pour l'auditeur en cybersécurité}
}
Intégrité & Honnêteté, diligence, compétence~; refus de toute pression sur
les constats \\
Présentation impartiale & Rapporter fidèlement, y compris les obstacles
rencontrés et les divergences non résolues \\
Conscience professionnelle & Le soin qu'exige l'importance de la mission
et la confiance des commanditaires \\
Confidentialité & Protéger les informations obtenues, en particulier les
faiblesses découvertes \\
Indépendance & Ne pas auditer son propre travail~; rester objectif \\
Approche fondée sur la preuve & Conclure sur des preuves vérifiables,
obtenues sur un échantillon raisonné \\
Approche fondée sur le risque & Concentrer l'effort là où le risque, pour
la mission et pour l'audité, est le plus grand \\
\end{tableaucours}
\vspace{-0.8em}
\begin{center}{\small\itshape Tableau~6.1. Les sept principes de
l'audit (d'après l'ISO~19011:2018)}\end{center}

\begin{boxinfo}
\textbf{Les autres référentiels.} L'ISO/IEC~27007:2020 complète
l'ISO~19011 pour l'audit des systèmes de management de la sécurité de
l'information~\cite{iso27007_2020}~; le Chapitre~\ref{chap:smsi} l'emploie.
Les auditeurs internes appliquent en outre les normes mondiales d'audit
interne de l'IIA, publiées en 2024~\cite{iia2024gias}, et les auditeurs des
systèmes d'information le cadre ITAF de l'ISACA~\cite{isaca2020itaf}.
\end{boxinfo}

\subsection{Quatre familles d'audit}

\begin{tableaucours}{L{2.6cm}L{3.4cm}L{3.0cm}L{3.2cm}}{%
    \tableauHeader{Famille} &
    \tableauHeader{Commanditaire} &
    \tableauHeader{Critères} &
    \tableauHeader{Exemple du cours}
}
Interne (première partie) & L'organisation elle-même, par son audit
interne & Ses politiques, ses plans, les normes qu'elle applique &
Mission de Mme~EKAMBI sur les agences \\
Contractuel (seconde partie) & Un client ou un partenaire & Le contrat &
PHARMACAM chez LiZbiZ \\
Certification (tierce partie) & L'organisation, auprès d'un organisme
accrédité & La norme certifiable & CERTIQUAL AFRIQUE, ISO/IEC~27001 \\
Réglementaire & La loi, par une autorité & Les textes & L'audit de
sécurité obligatoire de l'ANPP \\
\end{tableaucours}
\vspace{-0.8em}
\begin{center}{\small\itshape Tableau~6.2. Les familles d'audit}\end{center}

Un test d'intrusion n'est pas un audit au sens strict~: il cherche des
chemins d'attaque, sans évaluer l'ensemble des critères. Il peut fournir
des preuves à un audit, comme le test A3 du plan de la banque. Il est
enseigné dans l'UE MSN-2S3-SEC03.

% ============================================================================
\section{Le Cycle de la Mission}
\label{sec:c06-cycle}
% ============================================================================

\subsection{Programme et mission}

L'ISO~19011 distingue le \emph{programme d'audit}, qui organise les audits
sur plusieurs années selon les risques, et la \emph{mission}, qui en
réalise un. Le programme de l'audit interne de la BANQUE SANAGA n'avait
pas couvert la sécurité en 2024 et 2025 (Chapitre~\ref{chap:grc},
Tableau~1.7)~; celui de 2028 l'inscrit en tête, parce que le risque
résiduel de R-AG-01 repose sur des mesures nouvelles.

\begin{figure}[H]
\centering
\begin{tikzpicture}[
    etape/.style={draw=CouleurPrimaire, thick, rounded corners,
                  fill=CouleurPrimaire!8, text width=2.0cm,
                  minimum height=1.2cm, align=center, font=\scriptsize},
    fleche/.style={->, thick, CouleurSecondaire}]
    \node[etape] (M) at (0,0) {\textbf{Mandat}\\objectifs, périmètre, critères};
    \node[etape] (P) at (2.5,0) {\textbf{Préparation}\\analyse de risque, programme de travail};
    \node[etape] (T) at (5.0,0) {\textbf{Terrain}\\collecte et évaluation des preuves};
    \node[etape] (C) at (7.5,0) {\textbf{Contradictoire}\\réunion de clôture, réponses};
    \node[etape] (R) at (10.0,0) {\textbf{Rapport}\\constats, conclusions};
    \node[etape] (S) at (5.0,-2.0) {\textbf{Suivi}\\vérification des actions};
    \draw[fleche] (M) -- (P); \draw[fleche] (P) -- (T);
    \draw[fleche] (T) -- (C); \draw[fleche] (C) -- (R);
    \draw[fleche] (R.south) |- (S.east);
    \draw[fleche, dashed] (S.west) -| (M.south);
\end{tikzpicture}
\caption{Le cycle d'une mission d'audit}
\label{fig:c06-cycle}
\end{figure}

\subsection{Le mandat}

\begin{boxscenario}{Opération SANAGA. La lettre de mission}
\textbf{Commanditaire~:} le comité des risques du conseil, auquel l'audit
interne rend compte.

\medskip
\textbf{Objectif~:} établir si les mesures du plan de traitement de
R-AG-01, et les mesures sur les agences issues de l'appréciation de
décembre 2026, sont en place et fonctionnent.

\medskip
\textbf{Périmètre~:} les 38 agences, les comptes et les postes des agents,
les espaces de fichiers, les sauvegardes des serveurs d'agence, les
automates de retrait (GAB)~; période~: de septembre 2027 à mars 2028.

\medskip
\textbf{Critères~:} la politique de sécurité de la banque~; le plan de
traitement (actions A1 à A8, preuves attendues)~; les mesures
correspondantes de l'ISO/IEC~27002:2022.

\medskip
\textbf{Équipe~:} Mme~EKAMBI et deux auditeurs internes, avec l'appui
d'un auditeur technique externe pour les extractions. Aucun membre de
l'équipe n'a participé à la mise en œuvre du plan.

\medskip
\textbf{Accès~:} la lettre de mission, signée par le président du comité,
donne à l'équipe l'accès à tous les systèmes, documents et personnes
nécessaires, y compris chez les prestataires dans la limite des contrats.
\end{boxscenario}

La dernière phrase compte. Sans elle, l'auditeur dépend de la bonne volonté
de chacun. Avec elle, un refus d'accès devient un constat, que le comité
doit traiter.

\subsection{La préparation}

L'équipe analyse les risques de la mission~: où une erreur de conclusion
serait la plus grave, où la preuve sera la plus difficile à obtenir. Elle
en tire un \emph{programme de travail}~: pour chaque critère, le test, la
technique, l'échantillon, la preuve attendue.

\begin{tableaucours}{cL{3.4cm}L{4.4cm}L{3.4cm}}{%
    \tableauHeader{Réf.} &
    \tableauHeader{Critère} &
    \tableauHeader{Test} &
    \tableauHeader{Preuve attendue}
}
T1 & Second facteur sur tous les comptes (A1) & Extraction exhaustive des
1\,240 comptes~; essai de connexion sans second facteur & Extraction,
procès-verbal d'essai \\
T2 & Cloisonnement par agence (A2) & Essais d'accès d'une agence vers les
espaces des autres & Procès-verbal d'essai \\
T3 & Sauvegardes hors ligne testées (A4) & Restauration observée d'un
serveur choisi par l'auditeur & Procès-verbal, empreinte des fichiers \\
T4 & Mode dégradé exercé (A5) & Exercice inopiné dans deux agences &
Compte rendu, chronométrage \\
T5 & Sensibilisation des agents (A6) & Rapprochement des listes de présence
et des effectifs & Listes, fichier du personnel \\
T6 & Journalisation centralisée (A8) & Extraction du collecteur~: agences
qui envoient des journaux & Extraction datée \\
T7 & Fin de l'administration locale des chefs d'agence & Extraction des
groupes locaux des postes & Extraction \\
T8 & Correctifs des GAB & Revue des configurations auprès du prestataire &
Configurations \\
\end{tableaucours}
\vspace{-0.8em}
\begin{center}{\small\itshape Tableau~6.3. Le programme de travail de la
mission}\end{center}

% ============================================================================
\section{Collecter et Apprécier la Preuve}
\label{sec:c06-preuve}
% ============================================================================

\subsection{Trois techniques, et leurs variantes}

Le NIST décrit trois méthodes d'évaluation des mesures de sécurité~:
\emph{examiner} (des documents, des configurations, des journaux),
\emph{interroger} (des personnes), \emph{tester} (faire fonctionner la
mesure dans des conditions définies)~\cite{nist2022sp80053a}. Les auditeurs
y ajoutent l'\emph{observation} (regarder une opération se faire) et la
\emph{ré-exécution} (refaire soi-même un contrôle pour voir s'il donne le
même résultat).

\subsection{La fiabilité de la preuve}

Toutes les preuves ne se valent pas. Une preuve est d'autant plus fiable
qu'elle est obtenue directement par l'auditeur, qu'elle est produite par
un système plutôt que rédigée par une personne, et qu'elle vient d'une
source indépendante de l'audité.

\begin{tableaucours}{cL{5.4cm}L{5.8cm}}{%
    \tableauHeader{Rang} &
    \tableauHeader{Type de preuve} &
    \tableauHeader{Exemple dans la mission}
}
1 & Obtenue directement par l'auditeur (test, ré-exécution, observation)
& Essai de connexion sans second facteur, restauration observée \\
2 & Produite par le système, extraite sous les yeux de l'auditeur &
Extraction des comptes, du collecteur de journaux \\
3 & Document de tiers indépendant & Rapport du test d'intrusion A3 \\
4 & Document produit par l'audité & Rapport de la RSSI, captures d'écran
fournies \\
5 & Déclaration orale & « Toutes les agences sont raccordées » \\
\end{tableaucours}
\vspace{-0.8em}
\begin{center}{\small\itshape Tableau~6.4. Une hiérarchie de la fiabilité
des preuves}\end{center}

\begin{boxinfo}
\textbf{Suffisante, pertinente, fiable.} Une preuve doit être
\emph{pertinente} (elle porte sur le critère), \emph{fiable} (rang élevé
du Tableau~6.4) et \emph{suffisante} (en quantité, pour la conclusion
visée). Une déclaration orale peut orienter un test~; elle ne suffit
jamais à coter \OUI{}.
\end{boxinfo}

\subsection{Échantillonner}

On ne vérifie pas toujours tout. Pour un \emph{contrôle automatisé} (le
second facteur, géré par un serveur central), on extrait l'ensemble de la
population~: 1\,240 comptes, sans échantillon. Pour ce qui se fait agence
par agence, l'équipe retient un échantillon de huit agences sur
trente-huit~:
\begin{itemize}
    \item Ngaoundéré, point d'entrée de l'attaque de 2026~;
    \item deux autres agences touchées en 2026~;
    \item les trois agences qui traitent le plus d'opérations~;
    \item deux agences tirées au hasard parmi les autres, dont Kribi.
\end{itemize}

\begin{boxdef}
\termecle{Échantillon raisonné}\index{échantillonnage}~: échantillon
choisi par l'auditeur selon les risques. Il permet de trouver des défauts
là où ils sont probables~; il ne permet pas d'extrapoler un taux à la
population. Seul un \termecle{échantillon aléatoire}, de taille suffisante,
autorise une estimation statistique.
\end{boxdef}

La distinction n'est pas un détail~: si trois agences de l'échantillon sur
huit présentent un défaut, on ne peut pas en déduire que «~37,5~\% des
agences~» le présentent. On peut seulement dire que le défaut existe, et
demander une vérification exhaustive.

% ============================================================================
\section{Coter les Tests}
\label{sec:c06-coter}
% ============================================================================

\subsection{Les résultats de terrain}

\begin{tableaucours}{cL{6.8cm}cL{2.8cm}}{%
    \tableauHeader{Réf.} &
    \tableauHeader{Résultat} &
    \tableauHeader{Cote} &
    \tableauHeader{Rang de la preuve}
}
T1 & 1\,240 comptes sur 1\,240 avec second facteur~; connexion refusée sans
lui dans les huit agences & \OUI{} & 1 et 2 \\
T2 & Tous les essais d'accès entre agences refusés & \OUI{} & 1 \\
T3 & Trois restaurations réussies, empreintes identiques & \OUI{} & 1 \\
T4 & Garoua sert en mode dégradé en 40 minutes~; Kribi ne peut pas servir
pendant plus de deux heures (procédure introuvable, codes inconnus) &
\NON{} & 1 \\
T5 & Listes de 30 agences, signatures non nominatives~; rapprochement
avec les effectifs impossible & \NEK{} & 4 \\
T6 & 31 agences sur 38 envoient des journaux au collecteur & \NON{} & 2 \\
T7 & Chef d'agence encore administrateur local dans 3 agences sur 8 &
\NON{} & 2 \\
T8 & Le prestataire des GAB refuse l'accès aux configurations, en
invoquant son contrat & \NEK{} & -- \\
\end{tableaucours}
\vspace{-0.8em}
\begin{center}{\small\itshape Tableau~6.5. Résultats des tests (mars 2028)}\end{center}

\begin{boxlecture}{La mission sur le réseau d'agences}
\textbf{Constatation.} Sur huit tests~: 3~\OUI{}, 3~\NON{}, 2~\NEK{}.
\[
\mu = 0{,}375~;\quad U = 0{,}25~;\quad
H = -(2 \times 0{,}375\log_2 0{,}375 + 0{,}25\log_2 0{,}25)
  \approx 0{,}531 + 0{,}531 + 0{,}500 \approx 1{,}56~\text{bit}.
\]

\textbf{Analyse} (confiance élevée sur les \OUI{} et les \NON{}). Les trois
mesures techniques centrales du plan (T1 à T3) fonctionnent~: le résiduel
de R-AG-01 sur l'entrée et la propagation est confirmé. Mais le rapport de
la RSSI déclarait «~tout en place~», et trois mesures ne le sont pas. Ce
n'est pas un défaut d'honnêteté~: la RSSI a reçu des déclarations des
agences et de la DSI, et les a transmises. C'est exactement la raison
d'être de l'audit~: remplacer des déclarations de rang~4 ou~5 par des
preuves de rang~1 ou~2.

\medskip
\hfill\lecture{0{,}375}{0{,}25}{1{,}56}
\end{boxlecture}

\subsection{Le cas du NEK}

Les deux \NEK{} ne sont pas de même nature.

\begin{itemize}
    \item \textbf{T5} est un \NEK{} de \emph{preuve insuffisante}~: les
          listes existent, mais ne permettent pas de conclure. L'auditeur
          recommande une preuve meilleure (émargement nominatif, quiz
          enregistré) et pourra conclure au prochain passage.
    \item \textbf{T8} est un \NEK{} de \emph{limitation de l'étendue}~: on
          a empêché l'auditeur de voir. C'est en soi un constat, grave, que
          le rapport porte au comité~: quatorze agences ont des GAB sous un
          système sans correctifs (Chapitre~\ref{chap:iso27005}), et la
          banque ne peut pas vérifier ce que fait leur prestataire.
\end{itemize}

\begin{boxalert}
\textbf{Violation typique.} Retirer du rapport un test que l'on n'a pas pu
réaliser, pour ne présenter que ce qui a été vérifié.

\smallskip
\textbf{Règle~:} une limitation de l'étendue figure dans le rapport, avec
sa cause et ce qu'elle empêche de conclure. Le principe de présentation
impartiale l'exige~: le commanditaire doit savoir ce que l'audit n'a pas
pu voir.
\end{boxalert}

\subsection{Le papier de travail}

\begin{boxprocedure}
\textbf{Le papier de travail d'un test.} Chaque test du programme laisse
un papier de travail qui permet à un autre auditeur de refaire le
raisonnement~:
\begin{itemize}
    \item la référence du test, le critère, la population et l'échantillon,
          avec la méthode de sélection~;
    \item la technique employée et la date~;
    \item les preuves obtenues, leur rang de fiabilité, et l'endroit où
          elles sont conservées (extraction, empreinte, procès-verbal)~;
    \item le résultat et la cote (\OUI{}, \NON{} ou \NEK{}), avec, pour un
          \NEK{}, sa nature (preuve insuffisante ou limitation)~;
    \item l'auteur et le réviseur.
\end{itemize}
\end{boxprocedure}

% ============================================================================
\section{Le Contradictoire et le Rapport}
\label{sec:c06-contradictoire}
% ============================================================================

\subsection{La réunion de clôture}

Avant le rapport, l'équipe présente ses constats aux audités~: la RSSI, la
DSI, la direction du réseau d'agences. C'est le
\termecle{contradictoire}\index{contradictoire}~: l'audité peut contester
un fait, apporter une preuve que l'auditeur n'avait pas, ou expliquer une
cause. Un constat qui a survécu au contradictoire est solide~; un constat
qui le contourne est fragile.

\begin{boxscenario}{Opération SANAGA. La DSI conteste T6}
\textbf{Faits.} À la réunion de clôture, le DSI affirme que les sept
agences manquantes «~ont été raccordées la semaine dernière~».

\medskip
\textbf{Réponse de l'équipe.} Mme~EKAMBI demande une extraction datée du
collecteur, faite devant un auditeur. L'extraction, obtenue le lendemain,
montre 38 agences sur 38.

\medskip
\textbf{Analyse} (confiance élevée). Le constat n'est pas supprimé~: pendant
la période auditée, sept agences n'étaient pas journalisées, et une
attaque y serait passée inaperçue. Le rapport le maintient, et ajoute
«~corrigé en cours de mission, extraction du 22~mars 2028~». La correction
est une bonne nouvelle~; elle ne réécrit pas le passé.
\end{boxscenario}

\subsection{Le rapport et le suivi}

Le rapport présente les objectifs, le périmètre, les critères, les
limitations, puis les constats et les conclusions. Chaque constat donne le
critère, le fait, la cote et la preuve~; le Chapitre~\ref{chap:conclure}
montre comment le développer en observation complète (cause, conséquence,
recommandation) et comment tirer une conclusion d'ensemble des cotes.
L'audité répond par un plan d'actions daté~; l'audit interne vérifie, à
l'échéance, que les actions sont faites, par des preuves du même rang que
celles de la mission.

\begin{boxobjet}{R-AG-01 · Après l'audit de mars 2028}
\textbf{Mesures confirmées par l'audit~:} second facteur (T1),
cloisonnement (T2), sauvegardes restaurées (T3).

\textbf{Mesures non confirmées~:} mode dégradé (T4, \NON{} à Kribi),
sensibilisation (T5, \NEK{}), journalisation (T6, corrigée en cours de
mission), administration locale (T7, \NON{} dans 3 agences sur 8).

\textbf{Résiduel~:} inchangé à $2 \times 3 = 6$ selon la convention de
réduction, car les étapes entrée, propagation et récupération restent
couvertes par des mesures prouvées. Mais la récupération repose sur les
sauvegardes seules~: dans une agence comme Kribi, les guichets resteraient
fermés le temps de la restauration.

\textbf{Nouveau risque à apprécier~:} compromission d'un GAB (T8), base
\NEK{}~; présentation au comité des risques en raison de la limitation de
l'étendue.
\end{boxobjet}

% ============================================================================
\section{Le Cadre Camerounais}
\label{sec:c06-cadre}
% ============================================================================

\begin{boxjuris}
\textbf{\loicam{n°~2010/012}, article~4, point~9}~\cite{loi2010012}~:

\medskip
\textit{« Audit de sécurité~: examen méthodique des composantes et des
acteurs de la sécurité, de la politique, des mesures, des solutions, des
procédures et des moyens mis en oeuvre par une organisation, pour
sécuriser son environnement, effectuer des contrôles de conformité, des
contrôles d'évaluation de l'adéquation des moyens (organisationnels,
techniques, humains, financiers) investis au regard des risques encourus,
d'optimisation, de rationalité et de performance. »}

\medskip
\textbf{Article~13, alinéa~1}~:

\medskip
\textit{« Sont soumis à un audit de sécurité obligatoire, les réseaux de
communications électroniques et les systèmes d'information, des
opérateurs, des autorités de certification et des fournisseurs de services
de communications électroniques. »}

\medskip
\textbf{Article~32, alinéas~1 et~2}~:

\medskip
\textit{« (1) Les réseaux de communications électroniques et les systèmes
d'information sont soumis à un régime d'audit de sécurité obligatoire et
périodique de leurs systèmes de sécurité par l'Agence. (2) L'audit de
sécurité et les mesures d'impact de gravité sont effectués chaque année ou
lorsque les circonstances l'exigent. »}

\medskip
\textbf{Portée.} La définition légale de l'article~4 couvre à la fois la
conformité et l'adéquation des moyens «~au regard des risques encourus~»~:
elle correspond à l'approche fondée sur le risque de l'ISO~19011. Les
articles~13 et~32 instituent un audit réglementaire annuel, conduit sous
l'autorité de l'Agence nationale des technologies de l'information et de la
communication. Leur champ exact (quels systèmes, quelles organisations)
et leur application à une agence publique comme l'ANPP sont étudiés au
Chapitre~\ref{chap:droit}.
\end{boxjuris}

% ============================================================================
\section{Ce que l'Audit demande à chaque Métier}
\label{sec:c06-metiers}
% ============================================================================

\begin{boxmetier}{Auditeur en cybersécurité}
\textbf{Préférer la preuve que l'on obtient soi-même}
\begin{enumerate}
    \item Obtenir un mandat écrit qui garantit l'accès, y compris chez les
          prestataires.
    \item Écrire, avant le terrain, la preuve attendue de chaque test.
    \item Viser les preuves de rang~1 et~2~; ne jamais coter \OUI{} sur une
          déclaration.
    \item Distinguer, pour chaque \NEK{}, preuve insuffisante et limitation
          de l'étendue, et porter la seconde au commanditaire.
    \item Maintenir un constat corrigé en cours de mission, en signalant la
          correction.
\end{enumerate}
\end{boxmetier}

\begin{boxmetier}{RSSI}
\textbf{Préparer l'audit, pas le subir}

La RSSI aurait évité l'écart entre son rapport et l'audit en demandant,
pour chaque action, la preuve prévue au plan (Chapitre~\ref{chap:iso27005})
plutôt qu'une déclaration. Un RSSI dont les rapports reposent sur des
preuves de rang~2 voit l'audit confirmer ses chiffres.
\end{boxmetier}

\begin{boxmetier}{Responsable GRC}
\textbf{Faire entrer l'audit dans les contrats}

La limitation de T8 vient d'un contrat sans droit d'audit, comme celui de
DIGICORE. Le responsable GRC fait inscrire, dans chaque contrat de
prestataire qui touche au système d'information, un droit d'accès de
l'audit interne et des auditeurs mandatés.
\end{boxmetier}

\separateur

\begin{boxresume}
\textbf{Ce qu'il faut retenir de ce chapitre~:}
\begin{itemize}
    \item \textbf{L'audit} obtient et évalue des preuves au regard de
          critères~; le constat peut être conforme, non conforme, ou
          \NEK{}.
    \item \textbf{Sept principes} (ISO~19011)~: intégrité, présentation
          impartiale, conscience professionnelle, confidentialité,
          indépendance, approche fondée sur la preuve, approche fondée sur
          le risque.
    \item \textbf{Quatre familles}~: interne, contractuel, certification,
          réglementaire.
    \item \textbf{Le cycle}~: mandat, préparation, terrain, contradictoire,
          rapport, suivi.
    \item \textbf{La preuve} vaut par sa pertinence, sa fiabilité et sa
          suffisance~; une déclaration n'est jamais suffisante.
    \item \textbf{Un échantillon raisonné} trouve des défauts~; il ne
          permet pas d'extrapoler un taux.
    \item \textbf{Une limitation de l'étendue} est un constat, porté au
          commanditaire.
\end{itemize}
\end{boxresume}

\begin{boxinfo}
\textbf{Points de contrôle.} Vous avez compris ce chapitre si vous savez~:
\begin{itemize}
    \item classer une preuve dans la hiérarchie de fiabilité~;
    \item rédiger une ligne de programme de travail avec sa preuve
          attendue~;
    \item dire ce qu'un échantillon raisonné permet de conclure~;
    \item distinguer les deux sortes de \NEK{} et traiter chacune.
\end{itemize}
\end{boxinfo}

% ============================================================================
\section*{Exercices}
% ============================================================================

\begin{boxexo}[title={\ding{45}\quad Exercice~6.1. Techniques et preuves \niveaubase}]
Pour chacun des éléments suivants, indiquez la technique (examiner,
interroger, tester, observer, ré-exécuter) et le rang de fiabilité de la
preuve (Tableau~6.4).
\begin{enumerate}[(a)]
    \item Le DSI explique que les correctifs sont appliqués chaque mois.
    \item L'auditeur lance lui-même la restauration d'un serveur.
    \item La RSSI remet une capture d'écran de la console de sauvegarde.
    \item L'auditeur assiste à l'ouverture d'une agence en mode dégradé.
    \item L'auditeur extrait la liste des comptes d'administration, devant
          lui, depuis l'annuaire.
    \item L'auditeur recalcule le délai d'application des correctifs à
          partir du journal de déploiement.
\end{enumerate}
\end{boxexo}

\begin{boxexo}[title={\ding{45}\quad Exercice~6.2. Une ligne de programme \niveaumoyen}]
Le comité demande d'ajouter à la mission un test sur les comptes des agents
partis~: aucun compte d'un agent ayant quitté la banque ne doit rester
actif plus de cinq jours après son départ.
\begin{enumerate}[(a)]
    \item Rédigez la ligne du programme de travail~: critère, test,
          population, échantillon éventuel, preuve attendue.
    \item La DRH fournit une liste de 46 départs en 2027~; l'extraction de
          l'annuaire montre 4 comptes de partants actifs plus de cinq jours.
          Cotez le test et justifiez.
    \item La DRH refuse de communiquer les départs du dernier trimestre,
          «~pour des raisons de confidentialité~». Quelle est la nature de
          ce \NEK{}, et que fait l'auditeur~?
\end{enumerate}
\end{boxexo}

\begin{boxexo}[title={\ding{45}\quad Exercice~6.3. Ce que dit l'échantillon \niveaumoyen}]
\begin{enumerate}[(a)]
    \item Justifiez chacun des quatre critères de sélection de l'échantillon
          de huit agences.
    \item Un administrateur lit le résultat de T7 et conclut que «~14
          agences environ~» gardent l'administration locale. Que
          répondez-vous~?
    \item Si les huit agences avaient été conformes pour T7, l'auditeur
          aurait-il pu coter \OUI{} pour les 38 agences~? Proposez la
          meilleure preuve possible pour ce test.
\end{enumerate}
\end{boxexo}

\begin{boxexo}[title={\ding{45}\quad Exercice~6.4. L'audit contractuel de PHARMACAM \niveauexpert}]
PHARMACAM exerce, en 2028, son droit d'audit sur LiZbiZ, pour vérifier le
plan de traitement du scénario SO1 (Chapitres~\ref{chap:ebios}
et~\ref{chap:fair}).
\begin{enumerate}[(a)]
    \item De quelle famille d'audit s'agit-il~? Qui peut le conduire~?
          AUDIX CONSEIL peut-il s'en charger~? Justifiez.
    \item Rédigez la lettre de mission~: objectif, périmètre, critères,
          accès.
    \item Rédigez un programme de travail de cinq tests, avec la preuve
          attendue et son rang.
    \item L'éditeur de l'application refuse de montrer la configuration de
          son accès de maintenance. Comment le rapport le traite-t-il, et
          quelle conséquence PHARMACAM peut-elle en tirer~?
\end{enumerate}
\end{boxexo}

\bigskip

\begin{boxinfo}
\textbf{Pour aller plus loin.}
\begin{itemize}
    \item L'ISO~19011:2018~\cite{iso19011_2018}, sur les principes, le
          programme et la conduite des audits.
    \item L'ISO/IEC~27007:2020~\cite{iso27007_2020}, pour l'audit d'un SMSI.
    \item Le NIST SP~800-53A~\cite{nist2022sp80053a}, sur les méthodes
          d'évaluation des mesures de sécurité.
    \item Les normes mondiales d'audit interne de l'IIA~\cite{iia2024gias}.
\end{itemize}
\end{boxinfo}

\begin{boxinfo}
\textbf{Transition.} La mission a produit huit cotes~: 3~\OUI{}, 3~\NON{},
2~\NEK{}. Que conclure~? Le réseau d'agences est-il « sécurisé »~? Pour
répondre, il faut décomposer l'objet audité, coter chaque assertion selon
plusieurs critères, et appliquer une règle de décision qui tienne compte de
l'incertitude. C'est l'objet du Chapitre~\ref{chap:conclure}.
\end{boxinfo}

        % ============================================================================
% Fichier  : partie2/P02_C07_decomposer_coter_conclure.tex
% Titre    : Décomposer, Coter, Conclure
% Auteur   : MINKA MI NGUIDJOÏ Thierry Emmanuel
% Version  : 1.0 -- novembre 2026
% Corrigés : corriges/Corr_C07.tex
% ============================================================================
% Méthode reprise sans modification du corpus (ACSI, chapitres 7 et 9) :
% gabarit canonique à huit niveaux (article A0) ; fiche d'assertion à cinq
% attributs ; règle d'atomicité ; trois conditions (exhaustivité,
% non-redondance, borne cognitive) ; cote d'un contrôle à plusieurs critères
% (un NON suffit ; à défaut un NEK ; OUI si tous OUI) ; mesures mu, U, H
% (article A1) ; quadrant aux seuils du cours (mu 0,5 ; U 0,20) ; règle de
% décision « U d'abord, mu ensuite, H pour contrôler, les cotes de
% l'auditeur prévalent » ; conversion en maturité (convention ACSI) ;
% observation en sept éléments. Domaine D de LiZbiZ : valeurs d'ACSI.
% Faits du cas (bible) : mission de mars 2028 sur le réseau d'agences, T1 à
% T8 = 3/3/2. Faits nouveaux : décomposition en six sous-domaines S1 à S6 ;
% exercice inopiné de Kribi le 14 mars 2028 (2 h 15), contradictoire du
% 20 mars ; risque R-AG-02 (compromission d'un GAB).
% ISO/IEC 27002:2022, mesure 5.30 (préparation des TIC pour la continuité
% d'activité) : citée par son numéro et son intitulé.
% ============================================================================

\chapter{Décomposer, Coter, Conclure}
\label{chap:conclure}
\chaptermark{Décomposer et conclure}

\epigraph{%
    « Conclure, ce n'est pas compter les \emph{oui}.\\
    C'est savoir ce qu'ils valent. »%
}{Principe de la méthode}

\niveauChapitre{Intermédiaire}{%
    Maîtrise du Chapitre~\ref{chap:audit}%
}

\begin{boxobjectifs}
\begin{itemize}
    \item Décomposer un domaine de sécurité selon le gabarit canonique à
          huit niveaux, et vérifier la qualité de la décomposition.
    \item Rédiger une fiche d'assertion à plusieurs critères, en respectant
          l'atomicité.
    \item Déduire la cote d'un contrôle de celles de ses critères.
    \item Calculer $\mu$, $U$ et $H$ par sous-domaine, et localiser
          l'incertitude.
    \item Situer un domaine dans le quadrant conformité~$\times$ incertitude
          et appliquer la règle de décision.
    \item Rédiger une observation en sept éléments et la conclusion d'un
          rapport.
\end{itemize}
\end{boxobjectifs}

\bigskip

% ============================================================================
\section*{Avant de commencer : « Alors, c'est sécurisé~? »}
% ============================================================================

Douala, avril 2028. Mme~EKAMBI présente les résultats de la mission sur le
réseau d'agences au comité des risques~: huit tests, 3~\OUI{}, 3~\NON{},
2~\NEK{} (Chapitre~\ref{chap:audit}). Un administrateur l'interrompt~:
« Trois sur huit, c'est 37~\%. Les agences sont-elles sécurisées, oui ou
non~? »

La question est légitime, la réponse ne peut pas être un pourcentage. Les
huit cotes ne portent pas sur la même chose~: trois concernent les
barrières techniques du scénario R-AG-01, une la capacité à continuer de
servir, une le personnel, une un prestataire qui a refusé l'accès. Les
additionner, c'est mélanger un pare-feu et une limitation de l'étendue.

\medskip

Ce chapitre apprend à structurer ce que l'on audite, à coter chaque
contrôle selon ses critères, puis à conclure, en disant à la fois ce qui va
mal et ce que l'on ne sait pas. La méthode est celle de l'auditologie,
commune aux cours du corpus~\cite{minka2026auditologie,minka2026nek}.

\begin{boxalert}
\textbf{Violation typique.} Résumer un audit par le pourcentage de
contrôles conformes.

\smallskip
\textbf{Règle~:} une conclusion se donne par sous-domaine, avec trois
mesures (conformité, incertitude, dispersion), selon une règle de décision
écrite à l'avance. Un taux unique cache où sont les défauts et où est
l'ignorance.
\end{boxalert}

\separateur

% ============================================================================
\section{Décomposer}
\label{sec:c07-decomposer}
% ============================================================================

\subsection{Pourquoi décomposer}

Un domaine de sécurité compte des dizaines de contrôles. La mémoire de
travail humaine traite environ sept éléments à la fois~\cite{miller1956}~:
au-delà, l'auditeur oublie, recouvre ou mélange. La décomposition découpe
l'objet audité en niveaux emboîtés, chacun de taille maîtrisable. Elle a
trois usages~: vérifier que rien n'échappe au programme, situer chaque
cote, et localiser l'incertitude.

\subsection{Le gabarit canonique}

Le programme de l'auditologie propose un \termecle{gabarit
canonique}\index{gabarit canonique} à huit niveaux~\cite{minka2026auditologie}.

\begin{tableaucours}{cL{2.2cm}L{4.4cm}L{4.6cm}}{%
    \tableauHeader{N°} &
    \tableauHeader{Niveau} &
    \tableauHeader{Question} &
    \tableauHeader{Exemple (réseau d'agences)}
}
1 & Domaine & Que soumet-on à l'audit~? & Sécurité du réseau d'agences \\
2 & Sous-domaine & Quels grands ensembles le composent~? & Récupération \\
3 & Processus & Quelles suites d'activités~? & Continuité du service au
guichet \\
4 & Activité & Quelles actions élémentaires~? & Servir en mode dégradé \\
5 & Contrôle & Quel dispositif maîtrise l'activité~? & Procédure de mode
dégradé, exercée \\
6 & Test & Comment l'auditeur le vérifie-t-il~? & Exercice inopiné,
chronométré \\
7 & Preuve & Quel élément le test produit-il~? & Compte rendu signé,
chronométrage \\
8 & Trace & Quel enregistrement l'atteste~? & Journal des opérations en
mode dégradé \\
\end{tableaucours}
\vspace{-0.8em}
\begin{center}{\small\itshape Tableau~7.1. Le gabarit canonique
(d'après~\cite{minka2026auditologie})}\end{center}

Les quatre premiers niveaux disent \emph{ce qui est audité}~; les quatre
derniers, \emph{comment on le vérifie}. Le contrôle fait la charnière~:
c'est l'objet sur lequel l'auditeur se prononce.

\subsection{Trois conditions de qualité}

\begin{tableaucours}{L{3.0cm}L{5.2cm}L{4.0cm}}{%
    \tableauHeader{Condition} &
    \tableauHeader{Sens} &
    \tableauHeader{Défaut si elle manque}
}
Exhaustivité & Les éléments d'un niveau couvrent entièrement leur parent &
Une partie du domaine n'est pas auditée \\
Non-redondance & Deux éléments d'un même niveau ne se recouvrent pas~;
chacun a un seul parent & Double comptage, conclusions contradictoires \\
Borne cognitive & Aucun élément n'a plus de sept enfants environ & Grille
illisible, éléments oubliés \\
\end{tableaucours}
\vspace{-0.8em}
\begin{center}{\small\itshape Tableau~7.2. Les trois conditions d'une bonne
décomposition}\end{center}

\subsection{Le réseau d'agences, décomposé}

\begin{boxscenario}{Opération SANAGA. Six sous-domaines}
Mme~EKAMBI décompose le domaine « sécurité du réseau d'agences » en six
sous-domaines. Les quatre premiers suivent les étapes du scénario R-AG-01
(Chapitre~\ref{chap:risque})~; les deux derniers couvrent ce que le
scénario ne prévoyait pas.

\medskip
\begin{tableaucours}{cL{3.4cm}L{4.8cm}c}{%
    \tableauHeader{Code} &
    \tableauHeader{Sous-domaine} &
    \tableauHeader{Contrôles testés} &
    \tableauHeader{Tests}
}
S1 & Entrée (accès, authentification) & Second facteur sur tous les
comptes & T1 \\
S2 & Propagation (cloisonnement, privilèges) & Cloisonnement par agence~;
fin de l'administration locale & T2, T7 \\
S3 & Détection & Journalisation centralisée & T6 \\
S4 & Récupération & Sauvegardes hors ligne testées~; mode dégradé & T3,
T4 \\
S5 & Facteur humain & Sensibilisation des agents & T5 \\
S6 & Équipements de tiers & Correctifs des GAB & T8 \\
\end{tableaucours}

\medskip
\textbf{Vérification.} Six sous-domaines, sous la borne de sept~;
chaque test a un seul parent~; l'ensemble couvre le périmètre de la lettre
de mission. L'exhaustivité appelle une réserve~: d'autres contrôles
existent dans chaque sous-domaine (par exemple les correctifs des postes en
S1)~; la mission n'a testé que ceux du plan de traitement. Le rapport doit
le dire.
\end{boxscenario}

% ============================================================================
\section{Coter un Contrôle}
\label{sec:c07-coter}
% ============================================================================

\subsection{La fiche d'assertion}

Le gabarit dit \emph{où} se trouve un contrôle. La \termecle{fiche
d'assertion}\index{fiche d'assertion} dit ce que l'auditeur affirme à son
sujet, et ce qui arrive si l'affirmation est fausse.

\begin{tableaucours}{L{2.4cm}L{5.0cm}L{4.8cm}}{%
    \tableauHeader{Attribut} &
    \tableauHeader{Rôle} &
    \tableauHeader{Forme}
}
Objectif & Ce que le contrôle doit garantir & À l'infinitif~: « S'assurer
que\ldots » \\
Critère & La condition vérifiable qui définit la conformité & Au présent~:
une condition unique, observable \\
Question & La question posée pour coter & Interrogative~; réponse \OUI{},
\NON{} ou \NEK{} \\
Risque & L'état de non-conformité si la réponse est \NON{} & L'antithèse du
critère \\
Conséquence & L'effet sur l'organisation si le risque se réalise &
Financier, opérationnel, réglementaire \\
\end{tableaucours}
\vspace{-0.8em}
\begin{center}{\small\itshape Tableau~7.3. Les attributs de la fiche
d'assertion}\end{center}

\begin{boxinfo}
\textbf{Règle d'atomicité.} L'objectif, le critère et la question ne
contiennent aucun «~et~» ni «~ou~». Toute assertion qui en contient est
scindée en autant d'assertions qu'il y a de conditions indépendantes. Ainsi
une réponse \NON{} désigne un seul défaut.
\end{boxinfo}

\subsection{Un contrôle, trois critères}

\begin{boxscenario}{Opération SANAGA. Fiche d'assertion du contrôle S4.2}
\textbf{Position~:} domaine «~réseau d'agences~» $\rightarrow$ S4
(récupération) $\rightarrow$ continuité du service au guichet
$\rightarrow$ servir en mode dégradé $\rightarrow$ contrôle S4.2.

\medskip
\textbf{Objectif~:} s'assurer que chaque agence peut servir ses clients
sans le système d'information habituel.

\medskip
\textbf{Critère~1~:} la procédure de mode dégradé est disponible sur place,
sous forme papier. \emph{Question~:} la procédure est-elle présente dans
l'agence~?

\textbf{Critère~2~:} les codes d'accès au mode dégradé sont connus de la
personne qui dirige l'agence le jour du test. \emph{Question~:} le
responsable présent connaît-il les codes~?

\textbf{Critère~3~:} l'agence sert un premier client en moins d'une heure.
\emph{Question~:} le premier client a-t-il été servi en moins d'une
heure~?

\medskip
\textbf{Test, preuve, trace~:} exercice inopiné dans deux agences de
l'échantillon~; compte rendu signé et chronométrage~; journal des
opérations en mode dégradé.

\medskip
\textbf{Risque~:} une agence touchée par un incident reste fermée.

\textbf{Conséquence~:} perte d'exploitation (190~millions pour neuf agences
en 2026, Chapitre~\ref{chap:fair}), atteinte à la réputation.
\end{boxscenario}

\subsection{De la cote des critères à celle du contrôle}

Chaque critère reçoit sa propre cote. La cote du contrôle s'en déduit par
une règle simple~:
\begin{itemize}
    \item un seul critère \NON{} suffit pour coter le contrôle \NON{}~;
    \item à défaut, un seul critère \NEK{} donne \NEK{}~;
    \item le contrôle n'est coté \OUI{} que si tous ses critères le sont.
\end{itemize}

À Garoua, les trois critères sont \OUI{}. À Kribi, le critère~1 est \NON{}
(procédure introuvable), le critère~2 \NON{} (codes inconnus du chef
d'agence remplaçant), donc le critère~3 aussi (2~h~15). Un contrôle
d'agence se teste agence par agence~: S4.2 est \NON{} pour l'échantillon,
puisqu'une agence sur deux ne le satisfait pas.

\begin{boxalert}
\textbf{Violation typique.} Coter « \OUI{} à 50~\% » un contrôle conforme
dans une agence et défaillant dans l'autre.

\smallskip
\textbf{Règle~:} la cote porte sur le critère, qui est binaire. On cote
\NON{}, et le papier de travail précise où le contrôle fonctionne et où il
échoue.
\end{boxalert}

% ============================================================================
\section{Mesurer et Localiser l'Incertitude}
\label{sec:c07-mesurer}
% ============================================================================

\subsection{Trois mesures}

\begin{boxdef}
Pour un ensemble de cotes, on note $p(\OUI)$, $p(\NON)$ et $p(\NEK)$ les
parts de chaque réponse~\cite{minka2026nek}.

\medskip
\termecle{Taux de conformité}\index{taux de conformité}~: $\mu = p(\OUI)$.

\medskip
\termecle{Taux d'incertitude}\index{taux d'incertitude}~: $U = p(\NEK)$.

\medskip
\termecle{Entropie}\index{entropie}~: dispersion des réponses, en bits,
\[
    H = -\sum_{m \in \{\OUI, \NON, \NEK\}} p(m) \log_2 p(m),
    \qquad 0 \leq H \leq \log_2 3 \approx 1{,}585,
\]
avec la convention $0 \log_2 0 = 0$.
\end{boxdef}

\begin{tableaucours}{L{4.6cm}cccL{3.4cm}}{%
    \tableauHeader{Situation} &
    \tableauHeader{$\mu$} & \tableauHeader{$U$} & \tableauHeader{$H$} &
    \tableauHeader{Lecture}
}
Tout est \OUI{} & 1 & 0 & 0 & Conformité certaine \\
Tout est \NON{} & 0 & 0 & 0 & Non-conformité certaine \\
Tout est \NEK{} & 0 & 1 & 0 & \textbf{Ignorance unanime} \\
Réponses également réparties & $1/3$ & $1/3$ & 1,585 & Dispersion maximale \\
\end{tableaucours}
\vspace{-0.8em}
\begin{center}{\small\itshape Tableau~7.4. Les cas limites des trois
mesures (d'après~\cite{minka2026nek})}\end{center}

\begin{boxalert}
\textbf{Violation typique.} Conclure « certitude » parce que l'entropie
est nulle.

\smallskip
\textbf{Règle~:} une entropie nulle dit seulement que toutes les cotes sont
identiques. Si elles sont toutes \NEK{}, on ne sait rien~: $H = 0$ et
$U = 1$. On ne lit jamais $H$ sans $U$ ni $\mu$.
\end{boxalert}

\subsection{Par sous-domaine}

\begin{boxlecture}{Le réseau d'agences, sous-domaine par sous-domaine}
\begin{tableaucours}{cL{3.4cm}cccc}{%
    \tableauHeader{Code} &
    \tableauHeader{Sous-domaine} &
    \tableauHeader{\OUI{}/\NON{}/\NEK{}} &
    \tableauHeader{$\mu$} & \tableauHeader{$U$} & \tableauHeader{$H$}
}
S1 & Entrée & 1 / 0 / 0 & 1,00 & 0,00 & 0,00 \\
S2 & Propagation & 1 / 1 / 0 & 0,50 & 0,00 & 1,00 \\
S3 & Détection & 0 / 1 / 0 & 0,00 & 0,00 & 0,00 \\
S4 & Récupération & 1 / 1 / 0 & 0,50 & 0,00 & 1,00 \\
S5 & Facteur humain & 0 / 0 / 1 & 0,00 & 1,00 & 0,00 \\
S6 & Équipements de tiers & 0 / 0 / 1 & 0,00 & 1,00 & 0,00 \\
\textbf{Domaine} & & \textbf{3 / 3 / 2} & \textbf{0,375} & \textbf{0,25} &
\textbf{1,56} \\
\end{tableaucours}

\medskip
\textbf{Constatation.} Toute l'incertitude du domaine se concentre dans S5
et S6~: deux cas d'ignorance unanime ($H = 0$, $U = 1$). Les quatre autres
sous-domaines sont connus ($U = 0$)~: S1 est conforme, S3 non conforme, S2
et S4 partagés.

\medskip
\textbf{Analyse} (confiance élevée). Lu d'un bloc, le domaine paraît
« à 37~\% » et en partie inconnu. Décomposé, il dit trois choses
distinctes~: les barrières techniques du scénario R-AG-01 tiennent (S1,
le cloisonnement de S2, les sauvegardes de S4)~; des défauts connus restent
à corriger (administration locale, journalisation, mode dégradé)~; et la
banque ignore tout de deux sujets, le personnel et les GAB. Ce ne sont pas
les mêmes décisions.

\medskip
\hfill\lecture{0{,}375}{0{,}25}{1{,}56}
\end{boxlecture}

\begin{boxinfo}
\textbf{Peu de cotes, mesures fragiles.} S2 et S4 ne comptent que deux
contrôles chacun~: $\mu = 0{,}50$ y signifie « un sur deux », pas une
proportion stable. À ce niveau, on conclut contrôle par contrôle~; les
mesures servent surtout à comparer les sous-domaines et à repérer les
\NEK{}.
\end{boxinfo}

% ============================================================================
\section{Conclure}
\label{sec:c07-conclure}
% ============================================================================

\subsection{Le quadrant conformité~$\times$ incertitude}

L'article sur l'incertitude épistémique propose de lire ensemble la
conformité et l'incertitude, dans un quadrant à quatre
régions~\cite{minka2026nek}. Le cours retient, par convention, un seuil de
0,5 pour $\mu$ et de 0,20 pour $U$~; chaque mission peut fixer les siens, à
condition de les écrire dans le plan d'audit.

\begin{figure}[H]
\centering
\begin{tikzpicture}[x=6cm, y=4.5cm]
    \fill[red!15] (0,0) rectangle (0.5,0.2);
    \fill[green!18] (0.5,0) rectangle (1,0.2);
    \fill[orange!15] (0,0.2) rectangle (0.5,1);
    \fill[orange!25] (0.5,0.2) rectangle (1,1);
    \draw[thick, color=CouleurPrimaire] (0,0) rectangle (1,1);
    \draw[dashed, color=CouleurPrimaire] (0.5,0) -- (0.5,1);
    \draw[dashed, color=CouleurPrimaire] (0,0.2) -- (1,0.2);
    \node[align=center, font=\scriptsize\bfseries] at (0.25,0.1)
        {Non-conformité claire};
    \node[align=center, font=\scriptsize\bfseries] at (0.75,0.1)
        {Conformité claire};
    \node[align=center, font=\scriptsize\bfseries] at (0.25,0.6)
        {Non-conformité\\opaque};
    \node[align=center, font=\scriptsize\bfseries] at (0.75,0.6)
        {Conformité\\opaque};
    \fill[CouleurAccent] (0.375,0.25) circle (0.012);
    \node[font=\scriptsize, anchor=west] at (0.38,0.27) {Domaine};
    \fill[CouleurPrimaire] (0.02,0.98) circle (0.010);
    \node[font=\scriptsize, anchor=west] at (0.03,0.94) {S5, S6};
    \fill[CouleurPrimaire] (0.98,0.02) circle (0.010);
    \node[font=\scriptsize, anchor=east] at (0.97,0.045) {S1};
    \node[font=\scriptsize] at (0.5,-0.06) {0,5};
    \node[font=\scriptsize] at (-0.06,0.2) {0,20};
    \node[font=\small] at (0.5,-0.16) {Taux de conformité $\mu$};
    \node[font=\small, rotate=90] at (-0.15,0.5) {Taux d'incertitude $U$};
\end{tikzpicture}
\caption{Le quadrant diagnostique (d'après l'article sur l'incertitude
épistémique~; seuils du cours)}
\label{fig:c07-quadrant}
\end{figure}

\begin{tableaucours}{L{3.4cm}L{4.0cm}L{4.8cm}}{%
    \tableauHeader{Région} &
    \tableauHeader{Situation} &
    \tableauHeader{Action}
}
Conformité claire & Conformité établie, peu d'inconnu & Validation
courante \\
Non-conformité claire & Non-conformité établie, peu d'inconnu &
Correctifs \\
\textbf{Conformité opaque} & Beaucoup de \OUI{}, mais une part importante
de \NEK{} & \textbf{Alerte prioritaire}~: rechercher la preuve avant toute
conclusion \\
Non-conformité opaque & Beaucoup de \NON{} et une part importante de
\NEK{} & Rechercher la preuve et engager des premiers correctifs \\
\end{tableaucours}
\vspace{-0.8em}
\begin{center}{\small\itshape Tableau~7.5. Les quatre régions et leurs
actions (d'après~\cite{minka2026nek})}\end{center}

La conformité opaque est la région la plus dangereuse~: une grille binaire
la confond avec la conformité claire. Le rapport de la RSSI de février
2028 en donne l'exemple~: huit mesures déclarées « en place », soit une
conformité apparente de 100~\%~; mais ces \OUI{} reposaient sur des
déclarations, qui ne valaient que des \NEK{} tant qu'on ne les avait pas
vérifiées.

\subsection{La règle de décision}

\begin{boxprocedure}
\textbf{Conclure sur un domaine.} L'auditeur ne conclut qu'après avoir
consigné, pour chaque domaine et chaque sous-domaine~:
\begin{enumerate}
    \item \textbf{$U$ d'abord.} Si $U \geq 0{,}20$, l'ensemble est opaque~:
          aucune conclusion de conformité n'est possible. L'auditeur
          recherche la preuve, ou conclut \emph{sur l'opacité elle-même}.
    \item \textbf{$\mu$ ensuite.} Si $U < 0{,}20$, l'ensemble est clair~:
          conforme si $\mu \geq 0{,}5$, non conforme sinon.
    \item \textbf{$H$ pour contrôler.} Une entropie élevée avec $U$ faible
          signale un désaccord entre \OUI{} et \NON{}~: l'ensemble n'est pas
          homogène. On le décompose avant de conclure.
    \item \textbf{Les cotes de l'auditeur prévalent.} Les mesures orientent~;
          la conclusion repose sur les contrôles cotés, preuve à l'appui.
\end{enumerate}
\end{boxprocedure}

\begin{boxscenario}{Opération SANAGA. La conclusion de Mme EKAMBI}
\textbf{Domaine.} $U = 0{,}25 \geq 0{,}20$~: le domaine est opaque, en
région de non-conformité opaque ($\mu = 0{,}375$). Aucune conclusion de
conformité d'ensemble n'est possible~; la mission recherche les preuves
manquantes et engage des premiers correctifs.

\medskip
\textbf{Sous-domaines.}
\begin{itemize}
    \item S1~: clair et conforme.
    \item S2 et S4~: clairs, $H = 1$~: non homogènes~; conclusion par
          contrôle~: cloisonnement et sauvegardes efficaces, administration
          locale et mode dégradé non efficaces.
    \item S3~: clair, non conforme (corrigé en cours de mission).
    \item S5 et S6~: ignorance unanime~; l'auditeur conclut sur l'opacité.
          Pour S6, l'opacité vient d'une limitation de l'étendue, portée au
          comité.
\end{itemize}

\textbf{Phrase pour le comité.} « Les barrières techniques qui ont été
défaillantes en 2026 sont désormais en place et fonctionnent. Trois
défauts connus restent à corriger. Et la banque ne sait pas, à ce jour, si
ses agents sont formés ni si ses GAB sont protégés~: sur ces deux points,
nous ne pouvons pas conclure, et nous en donnons les raisons. »
\end{boxscenario}

\subsection{Pas de niveau de maturité pour un domaine opaque}

Les directions aiment résumer un domaine par un niveau de 1 à 5. Le corpus
propose une conversion du taux de conformité en niveau de maturité~:
moins de 0,35, niveau~1~; de 0,35 à 0,55, niveau~2~; de 0,55 à 0,75,
niveau~3~; de 0,75 à 0,90, niveau~4~; au-delà, niveau~5. C'est une
\emph{convention pédagogique déclarée}, non une norme, et elle a une règle~:
si $U \geq 0{,}20$, le niveau est \emph{indéterminé}. Le réseau d'agences
n'a donc pas de niveau en 2028~; lui en attribuer un, sur la foi de
$\mu = 0{,}375$, ferait croire au comité qu'il connaît un état qu'en partie
il ignore.

% ============================================================================
\section{Rédiger l'Observation}
\label{sec:c07-observation}
% ============================================================================

\subsection{Sept éléments}

Chaque non-conformité significative donne lieu à une observation, de
structure invariante.

\begin{tableaucours}{cL{3.0cm}L{8.4cm}}{%
    \tableauHeader{N°} &
    \tableauHeader{Élément} &
    \tableauHeader{Contenu}
}
1 & Norme & Le critère, et le texte ou le référentiel d'où il vient \\
2 & Constatation & Le fait établi, chiffré, avec sa preuve et sa cote \\
3 & Risques & Ce qui peut arriver~: le risque de la fiche d'assertion et sa
conséquence \\
4 & Contradictoire & La réponse de l'audité, datée \\
5 & Analyse & La cause de l'écart et l'interprétation, avec un niveau de
confiance \\
6 & Conclusion & La position finale de l'auditeur \\
7 & Recommandations & Ce qu'il faut faire, en traitant la cause \\
\end{tableaucours}
\vspace{-0.8em}
\begin{center}{\small\itshape Tableau~7.6. Les sept éléments de
l'observation}\end{center}

\begin{boxscenario}{Opération SANAGA. Observation S4.2 · Mode dégradé}
\textbf{1. Norme~:} plan de traitement de R-AG-01, action A5 («~procédure
de guichet en mode dégradé, exercée~»)~; ISO/IEC~27002:2022, mesure~5.30
(préparation des TIC pour la continuité d'activité)~\cite{iso27002_2022}.

\medskip
\textbf{2. Constatation~:} exercice inopiné du 14~mars 2028. Garoua sert
un premier client en 40~minutes. À Kribi, la procédure est introuvable, le
chef d'agence remplaçant ignore les codes du mode dégradé, et aucun client
n'est servi pendant 2~h~15 (compte rendu signé, chronométrage). Cote~:
\NON{}.

\medskip
\textbf{3. Risques~:} une agence touchée par un incident reste fermée~;
perte d'exploitation et atteinte à la réputation.

\medskip
\textbf{4. Contradictoire~:} la direction du réseau d'agences indique, le
20~mars, que la procédure a été diffusée par courriel en mai~2027 et que le
chef d'agence titulaire de Kribi était en congé.

\medskip
\textbf{5. Analyse} (confiance élevée)~: la procédure a été diffusée par
la messagerie, c'est-à-dire par l'outil qui manquera le jour d'un incident.
Elle n'a jamais été intégrée au dossier papier de l'agence, et sa
connaissance repose sur une seule personne.

\medskip
\textbf{6. Conclusion~:} le contrôle S4.2 n'est pas efficace pour
l'échantillon~; son efficacité dans les 36 autres agences n'est pas
établie.

\medskip
\textbf{7. Recommandations~:} placer la procédure et les codes, sous
enveloppe scellée, dans le coffre de chaque agence~; former au moins deux
personnes par agence~; exercer chaque agence une fois par an, avec
chronométrage~; vérification par l'audit interne lors du suivi.
\end{boxscenario}

\subsection{Le nouveau risque}

La limitation de l'étendue sur les GAB n'appelle pas d'observation au sens
strict~: on ne peut pas décrire un écart que l'on n'a pas vu. Elle appelle
une appréciation de risque, que le rapport demande au gestionnaire des
risques.

\begin{boxobjet}{R-AG-02 · Compromission d'un GAB}
\textbf{Valeur métier~:} le retrait d'espèces dans les agences.

\textbf{Biens supports~:} automates de retrait de 14 agences, sous un
système d'exploitation qui ne reçoit plus de correctifs~; contrat de
maintenance sans droit d'audit.

\textbf{Événement redouté~:} retraits frauduleux, ou indisponibilité des
automates.

\textbf{Vraisemblance~:} \NEK{} (configuration non vérifiable).

\textbf{Impact~:} à apprécier (montant des encaisses, nombre de clients
concernés).

\textbf{Action immédiate~:} négociation d'un avenant de droit d'audit avec
le prestataire~; à défaut, présentation au comité des risques, comme pour
DIGICORE (Chapitre~\ref{chap:iso27005}).
\end{boxobjet}

% ============================================================================
\section{Ce que la Conclusion demande à chaque Métier}
\label{sec:c07-metiers}
% ============================================================================

\begin{boxmetier}{Auditeur en cybersécurité}
\textbf{Conclure par sous-domaine}
\begin{enumerate}
    \item Décomposer le domaine avant le terrain, et faire revoir la
          décomposition.
    \item Rédiger une fiche d'assertion atomique pour chaque contrôle.
    \item Calculer $\mu$, $U$ et $H$ par sous-domaine, et appliquer la règle
          de décision écrite dans le plan d'audit.
    \item Conclure sur l'opacité quand $U \geq 0{,}20$, au lieu de la taire.
    \item Développer chaque non-conformité significative en observation à
          sept éléments.
\end{enumerate}
\end{boxmetier}

\begin{boxmetier}{RSSI}
\textbf{Lire le rapport par sous-domaine}

La RSSI ne retient pas « 37~\% ». Elle retient que ses barrières
techniques tiennent, qu'elle a trois défauts à corriger, et que deux
sujets lui échappent~: le personnel et les GAB. Son prochain rapport au
comité suivra la même décomposition.
\end{boxmetier}

\begin{boxmetier}{Responsable GRC}
\textbf{Une même décomposition partout}

Si le registre des risques, le plan de traitement et le programme d'audit
emploient la même décomposition (ici, les étapes du scénario plus le
facteur humain et les tiers), les cotes de l'audit alimentent directement le
registre. C'est le dispositif unique du Chapitre~\ref{chap:grc-pratique}.
\end{boxmetier}

\separateur

\begin{boxresume}
\textbf{Ce qu'il faut retenir de ce chapitre~:}
\begin{itemize}
    \item \textbf{Le gabarit canonique} a huit niveaux~: domaine,
          sous-domaine, processus, activité, contrôle, test, preuve, trace.
    \item \textbf{Une bonne décomposition} est exhaustive, non redondante et
          respecte la borne de sept éléments environ.
    \item \textbf{La fiche d'assertion} donne objectif, critère, question,
          risque, conséquence~; elle est atomique.
    \item \textbf{La cote d'un contrôle}~: un \NON{} suffit~; à défaut un
          \NEK{}~; \OUI{} seulement si tous les critères le sont.
    \item \textbf{La règle de décision}~: $U$ d'abord, $\mu$ ensuite, $H$
          pour contrôler, les cotes de l'auditeur prévalent.
    \item \textbf{Une entropie nulle} n'est pas une certitude si $U = 1$.
    \item \textbf{L'observation} a sept éléments, dont le contradictoire et
          une conclusion explicite.
\end{itemize}
\end{boxresume}

\begin{boxinfo}
\textbf{Points de contrôle.} Vous avez compris ce chapitre si vous savez~:
\begin{itemize}
    \item placer un élément au bon niveau du gabarit~;
    \item scinder une assertion qui n'est pas atomique~;
    \item calculer les trois mesures d'un sous-domaine et le situer dans le
          quadrant~;
    \item rédiger une conclusion qui dit ce qui va mal et ce que l'on ne
          sait pas.
\end{itemize}
\end{boxinfo}

% ============================================================================
\section*{Exercices}
% ============================================================================

\begin{boxexo}[title={\ding{45}\quad Exercice~7.1. À quel niveau~? \niveaubase}]
Placez chaque élément au niveau du gabarit canonique qui lui correspond.
\begin{enumerate}[(a)]
    \item L'extraction des groupes d'administration locale, avec son
          empreinte.
    \item La propagation.
    \item Le retrait des droits d'administration aux chefs d'agence.
    \item L'entrée du journal qui enregistre la modification d'un groupe
          local.
    \item Comparer les membres des groupes locaux à la liste des
          administrateurs autorisés.
    \item La gestion des droits sur les postes de travail.
    \item Aucun poste d'agence ne compte d'administrateur local autre que
          le compte central de la DSI.
\end{enumerate}
\end{boxexo}

\begin{boxexo}[title={\ding{45}\quad Exercice~7.2. Rendre atomique \niveaumoyen}]
La première version du programme portait la question suivante~: « Les
chefs d'agence ont-ils perdu leurs droits d'administration et les postes
sont-ils gérés par l'outil central ou par un prestataire agréé~? »
\begin{enumerate}[(a)]
    \item Expliquez pourquoi cette question n'admet pas de réponse honnête
          dans une agence où les droits ont été retirés mais où les postes
          ne sont gérés par personne.
    \item Scindez-la en questions atomiques.
    \item Rédigez la fiche d'assertion complète du contrôle «~fin de
          l'administration locale~» (S2.2), avec deux critères.
\end{enumerate}
\end{boxexo}

\begin{boxexo}[title={\ding{45}\quad Exercice~7.3. Le domaine D de LiZbiZ \niveaumoyen}]
AUDIX CONSEIL a décomposé le domaine~D (contrôle d'accès à l'application
intégrée de LiZbiZ) en cinq sous-domaines~:
\begin{itemize}
    \item D.1, politique et organisation~: 4 / 0 / 0~;
    \item D.2, cycle de vie des comptes~: 1 / 1 / 2~;
    \item D.3, authentification~: 4 / 0 / 0~;
    \item D.4, comptes à privilèges~: 2 / 0 / 2~;
    \item D.5, journalisation et revue~: 2 / 1 / 1.
\end{itemize}
(nombre de \OUI{}, \NON{} et \NEK{}).
\begin{enumerate}[(a)]
    \item Calculez $\mu$, $U$ et $H$ pour chaque sous-domaine et pour le
          domaine.
    \item Situez chacun dans le quadrant et appliquez la règle de décision.
    \item Quel niveau de maturité peut-on attribuer au domaine~?
    \item Où portez-vous les tests du second passage, et pourquoi~?
\end{enumerate}
\end{boxexo}

\begin{boxexo}[title={\ding{45}\quad Exercice~7.4. Observation et conclusion \niveauexpert}]
\begin{enumerate}[(a)]
    \item Rédigez l'observation à sept éléments du test T6 (journalisation~:
          31 agences sur 38, corrigée le 22~mars 2028, Chapitre~\ref{chap:audit}).
    \item Rédigez le paragraphe de conclusion du rapport consacré à S5
          (facteur humain), en concluant sur l'opacité.
    \item Un administrateur demande que le rapport «~retire les deux
          \NEK{}, qui affaiblissent le message~». Répondez-lui en citant
          les principes de l'audit (Chapitre~\ref{chap:audit}).
\end{enumerate}
\end{boxexo}

\bigskip

\begin{boxinfo}
\textbf{Pour aller plus loin.}
\begin{itemize}
    \item L'article \emph{Vers l'auditologie}~\cite{minka2026auditologie},
          pour le gabarit canonique et l'observabilité partielle.
    \item L'article sur l'incertitude épistémique~\cite{minka2026nek}, pour
          les trois mesures et le quadrant.
    \item Miller~\cite{miller1956}, sur la borne de la mémoire de travail.
\end{itemize}
\end{boxinfo}

\begin{boxinfo}
\textbf{Transition.} La Partie~II est achevée~: on sait apprécier,
traiter, chiffrer, auditer et conclure. Il reste à connaître les textes
auxquels se mesurent risques et audits~: la famille des normes ISO/IEC
27000, les autres référentiels et le droit. C'est l'objet de la
Partie~III, qui s'ouvre avec le Chapitre~\ref{chap:iso27000}.
\end{boxinfo}


    % ======================================================================
    % PARTIE 3 — CADRE NORMATIF ET JURIDIQUE
    % Standards de référence : ISO/IEC 27001 et 27002:2022, NIST CSF 2.0, NIST SP 800-53, CIS Controls, COBIT 2019 ; loi n° 2010/012, COBAC R-2016/04, Convention de Malabo.
    % Niveau : fondamental
    % ======================================================================
    \part{Cadre Normatif et Juridique}
    \label{part:normatif}

        % ============================================================================
% Fichier  : partie3/P03_C08_famille_iso27000.tex
% Titre    : La Famille ISO/IEC 27000
% Auteur   : MINKA MI NGUIDJOÏ Thierry Emmanuel
% Version  : 1.0 -- novembre 2026
% Corrigés : corriges/Corr_C08.tex
% ============================================================================
% Vérifié sur iso.org le 06/11/2026 : ISO/IEC 27002:2022, édition 3, publiée
% en février 2022 (remplace l'édition 2013 et ses rectificatifs) ; ISO/IEC
% 27001:2022 et son amendement 1:2024 ; ISO/IEC 27000:2018 en vigueur.
% Textes sous licence : articles et mesures PARAPHRASÉS, jamais cités ;
% intitulés des mesures traduits par le cours. Le dossier Normes ne contient
% que l'édition 2013 : écart signalé dans le texte.
% Structure de l'édition 2022 (93 mesures ; 37/8/14/34 ; 11 mesures
% nouvelles ; 5 attributs), date de fin de transition (31 octobre 2025),
% objet de l'amendement 2024 (changement climatique) : d'après des sources
% publiques concordantes, à confirmer sur le texte par l'auteur.
% Faits du cas (bible) : certificat de juin 2025 (CERTIQUAL AFRIQUE, 3 ans,
% centre de données seul) ; périmètre étendu approuvé par le conseil
% (Ex. 1.4) ; résultats de l'audit de mars 2028 (T1 à T8).
% Faits nouveaux (reportés dans la bible) : renouvellement du certificat en
% juin 2028 sur le périmètre étendu ; extrait de la déclaration
% d'applicabilité (10 lignes) ; exclusion de 8.28 contredite par les
% interfaces développées par la DSI ; revue croisée 3/3/3.
% ============================================================================

\chapter{La Famille ISO/IEC 27000}
\label{chap:iso27000}
\chaptermark{La famille ISO/IEC 27000}

\epigraph{%
    « Une norme ne dit pas ce que vous faites.\\
    Elle dit ce que vous devrez prouver. »%
}{Principe du cours}

\niveauChapitre{Fondamental}{%
    Maîtrise du Chapitre~\ref{chap:iso27005}%
}

\begin{boxobjectifs}
\begin{itemize}
    \item Situer les principales normes de la famille ISO/IEC~27000 et leur
          rôle.
    \item Présenter les exigences de l'ISO/IEC~27001:2022, articles~4 à~10,
          et le cycle de certification.
    \item Utiliser l'ISO/IEC~27002:2022~: 93 mesures, quatre thèmes,
          attributs, et l'écart avec l'édition de 2013.
    \item Rédiger une déclaration d'applicabilité et justifier une
          exclusion.
    \item Confronter une déclaration d'applicabilité aux preuves d'audit.
    \item Situer les ISO/IEC~27004, 27006 et 27007.
\end{itemize}
\end{boxobjectifs}

\bigskip

% ============================================================================
\section*{Avant de commencer : le renouvellement}
% ============================================================================

Douala, avril 2028. Le certificat ISO/IEC~27001 de la BANQUE
SANAGA\footnote{Comme dans tout le cours, la BANQUE SANAGA, l'organisme
CERTIQUAL AFRIQUE, les dates et les montants sont fictifs~; les normes
citées sont exactes et paraphrasées, leur texte étant protégé.}, obtenu en
juin 2025 pour trois ans, arrive à échéance. Le conseil a approuvé
l'extension du système de management aux 38 agences~; l'audit de
renouvellement de CERTIQUAL AFRIQUE portera donc, pour la première fois,
sur toute la banque.

Mme~ABENA doit remettre à l'organisme certificateur une pièce maîtresse~:
la \emph{déclaration d'applicabilité}, qui dit, mesure par mesure, ce que
la banque applique et pourquoi. La version de 2025 tenait en quelques pages,
sur le seul centre de données. Celle de 2028 doit couvrir les agences, au
moment même où l'audit interne vient d'y relever trois non-conformités et
deux inconnues (Chapitre~\ref{chap:conclure}).

\medskip

Ce chapitre ouvre la Partie~III, consacrée au cadre normatif et
juridique. Il présente la famille de normes la plus utilisée en sécurité
de l'information, puis la manière de rédiger, et de vérifier, une
déclaration d'applicabilité.

\begin{boxalert}
\textbf{Violation typique.} Rédiger la déclaration d'applicabilité en
cochant « mise en œuvre » pour toutes les mesures, à partir des
déclarations des services.

\smallskip
\textbf{Règle~:} la déclaration d'applicabilité est un document
\emph{opposable}~: l'organisme certificateur et les auditeurs la
confrontent aux preuves. Chaque « mise en œuvre » doit pouvoir être
démontrée~; chaque exclusion doit être justifiée par l'appréciation des
risques.
\end{boxalert}

\separateur

% ============================================================================
\section{Une Famille de Normes}
\label{sec:c08-famille}
% ============================================================================

La famille ISO/IEC~27000 est publiée par le sous-comité~27 du comité
technique commun de l'ISO et de la CEI (JTC~1/SC~27). Elle compte des
dizaines de documents~; une seule norme est certifiable, l'ISO/IEC~27001.

\begin{tableaucours}{L{2.6cm}L{5.4cm}L{4.2cm}}{%
    \tableauHeader{Norme} &
    \tableauHeader{Objet} &
    \tableauHeader{Dans ce cours}
}
ISO/IEC~27000 & Vue d'ensemble et vocabulaire~\cite{iso27000_2018} &
Définitions (propriétaire de risque) \\
ISO/IEC~27001 & Exigences du SMSI~; seule norme certifiable de la
famille~\cite{iso27001_2022} & Ce chapitre, Chapitre~\ref{chap:smsi} \\
ISO/IEC~27002 & Recommandations pour les mesures de
sécurité~\cite{iso27002_2022} & Ce chapitre \\
ISO/IEC~27004 & Surveillance, mesure, analyse et évaluation~\cite{iso27004_2016}
& Chapitre~\ref{chap:smsi} \\
ISO/IEC~27005 & Gestion des risques~\cite{iso27005_2022} &
Chapitre~\ref{chap:iso27005} \\
ISO/IEC~27006 & Exigences pour les organismes qui auditent et certifient
les SMSI & Ce chapitre \\
ISO/IEC~27007 & Lignes directrices pour l'audit des SMSI~\cite{iso27007_2020}
& Chapitres~\ref{chap:audit} et~\ref{chap:smsi} \\
ISO/IEC~27014 & Gouvernance de la sécurité~\cite{iso27014_2020} &
Chapitre~\ref{chap:grc} \\
\end{tableaucours}
\vspace{-0.8em}
\begin{center}{\small\itshape Tableau~8.1. Les principales normes de la
famille ISO/IEC~27000}\end{center}

D'autres normes de la famille déclinent les mesures pour un secteur ou un
sujet~: les services en nuage, la gestion des incidents, la protection de
la vie privée. Le Chapitre~\ref{chap:emergents} y revient pour le nuage.

% ============================================================================
\section{L'ISO/IEC 27001:2022}
\label{sec:c08-27001}
% ============================================================================

\subsection{Une structure commune}

L'ISO/IEC~27001 suit la structure commune à toutes les normes de systèmes
de management de l'ISO (qualité, environnement, continuité)~: dix
articles, dont les sept derniers contiennent les exigences. Une
organisation qui a déjà un système de management de la qualité retrouve
donc le même plan.

\begin{tableaucours}{cL{3.0cm}L{8.4cm}}{%
    \tableauHeader{Art.} &
    \tableauHeader{Titre} &
    \tableauHeader{Ce que la norme exige, en substance}
}
4 & Contexte de l'organisation & Comprendre les enjeux internes et
externes, les parties intéressées et leurs exigences~; fixer le
périmètre~; établir le SMSI et ses processus \\
5 & Leadership & Engagement de la direction~; politique de sécurité~;
rôles, responsabilités et autorités \\
6 & Planification & Actions face aux risques et opportunités~:
appréciation et traitement des risques, déclaration d'applicabilité~;
objectifs de sécurité~; planification des modifications \\
7 & Support & Ressources~; compétences~; sensibilisation~;
communication~; informations documentées \\
8 & Fonctionnement & Planifier et maîtriser les opérations~; réaliser les
appréciations et les traitements des risques prévus \\
9 & Évaluation des performances & Surveillance et mesure~; audit
interne~; revue de direction \\
10 & Amélioration & Amélioration continue~; non-conformités et actions
correctives \\
\end{tableaucours}
\vspace{-0.8em}
\begin{center}{\small\itshape Tableau~8.2. Les articles~4 à~10 de
l'ISO/IEC~27001:2022 (paraphrasés)}\end{center}

Les Chapitres~\ref{chap:grc} à~\ref{chap:conclure} ont déjà parcouru
l'essentiel de ces exigences sans les nommer~: le périmètre et les parties
intéressées (article~4), le rôle de la direction (article~5),
l'appréciation et le traitement des risques (articles~6 et~8), l'audit
interne (article~9). La norme ne prescrit pas \emph{comment} faire~; elle
dit ce que l'organisation doit faire et pouvoir prouver.

\subsection{Ce qui a changé en 2022}

L'édition de 2022 remplace celle de 2013. Les principaux changements~:
\begin{itemize}
    \item l'annexe~A, liste des mesures de référence, est alignée sur la
          nouvelle ISO/IEC~27002:2022 (93 mesures au lieu de 114)~;
    \item une exigence nouvelle impose de planifier les modifications du
          SMSI (article~6.3)~;
    \item le SMSI doit identifier ses processus et leurs interactions
          (article~4.4).
\end{itemize}
Un amendement publié en 2024 demande en outre à l'organisation de
déterminer si le changement climatique est un enjeu pertinent pour son
système de management. Les certificats délivrés sur l'édition de 2013
devaient avoir migré au plus tard le 31~octobre 2025. Le certificat de la
BANQUE SANAGA, délivré en juin 2025 sur l'édition de 2022, n'était pas
concerné.

\begin{boxinfo}
\textbf{Le dossier du cours et l'édition 2013.} Les normes archivées dans
les ressources du cours sont celles de 2013. Le cours enseigne l'édition
de 2022~; les numéros de mesures cités ici sont ceux de 2022. Un étudiant
qui consulte un document de 2013 trouvera d'autres numéros (de A.5 à A.18)
et quatorze domaines au lieu de quatre thèmes.
\end{boxinfo}

\subsection{Le cycle de certification}

\begin{tableaucours}{L{3.4cm}L{8.8cm}}{%
    \tableauHeader{Étape} &
    \tableauHeader{Contenu}
}
Audit initial, étape~1 & Revue documentaire et de l'état de préparation~:
périmètre, politique, appréciation des risques, déclaration
d'applicabilité \\
Audit initial, étape~2 & Évaluation de la mise en œuvre et de
l'efficacité du SMSI sur le terrain~; décision de certification \\
Audits de surveillance & Au moins un par an, pendant les trois ans du
certificat, sur une partie du SMSI \\
Audit de renouvellement & Avant l'échéance des trois ans, sur l'ensemble
du SMSI \\
\end{tableaucours}
\vspace{-0.8em}
\begin{center}{\small\itshape Tableau~8.3. Le cycle de certification}\end{center}

L'organisme certificateur doit lui-même satisfaire aux exigences de
l'ISO/IEC~27006, et être accrédité par un organisme d'accréditation. La
certification est volontaire~: elle ne remplace ni l'audit interne, ni
l'audit de sécurité réglementaire (Chapitre~\ref{chap:audit}).

% ============================================================================
\section{L'ISO/IEC 27002:2022}
\label{sec:c08-27002}
% ============================================================================

\subsection{93 mesures, quatre thèmes}

L'ISO/IEC~27002 n'est pas certifiable~: elle donne, pour chaque mesure, un
objectif et des recommandations de mise en œuvre~\cite{iso27002_2022}.
L'édition de 2022 réorganise les mesures en quatre thèmes.

\begin{tableaucours}{L{3.4cm}cL{7.0cm}}{%
    \tableauHeader{Thème} &
    \tableauHeader{Mesures} &
    \tableauHeader{Exemples (intitulés traduits par le cours)}
}
Organisationnelles (5.1 à 5.37) & 37 & Politiques, rôles, renseignement
sur les menaces (5.7), relations avec les fournisseurs (5.19, 5.20),
préparation des TIC pour la continuité (5.30) \\
Liées aux personnes (6.1 à 6.8) & 8 & Sélection, sensibilisation et
formation (6.3), signalement des événements \\
Physiques (7.1 à 7.14) & 14 & Périmètres, surveillance de la sécurité
physique (7.4), matériels \\
Technologiques (8.1 à 8.34) & 34 & Droits d'accès privilégiés (8.2),
authentification sécurisée (8.5), vulnérabilités techniques (8.8),
sauvegarde (8.13), journalisation (8.15), cloisonnement des réseaux
(8.22), codage sécurisé (8.28) \\
\end{tableaucours}
\vspace{-0.8em}
\begin{center}{\small\itshape Tableau~8.4. Les quatre thèmes de
l'ISO/IEC~27002:2022}\end{center}

Onze mesures sont nouvelles, dont le renseignement sur les menaces (5.7),
la sécurité des services en nuage (5.23), la préparation des TIC pour la
continuité (5.30), la surveillance de la sécurité physique (7.4), la gestion
de la configuration (8.9), la prévention des fuites de données (8.12), les
activités de surveillance (8.16), le filtrage web (8.23) et le codage
sécurisé (8.28). Les autres mesures de 2013 ont été regroupées ou
reformulées.

\subsection{Les attributs}

L'édition de 2022 associe à chaque mesure des \termecle{attributs}\index{attributs
de mesure}, qui permettent de trier les mesures selon plusieurs vues~:
\begin{itemize}
    \item le \emph{type de mesure}~: préventive, détective, corrective~;
    \item les \emph{propriétés de sécurité}~: confidentialité, intégrité,
          disponibilité~;
    \item les \emph{concepts de cybersécurité}~: identifier, protéger,
          détecter, répondre, rétablir (les fonctions historiques du cadre
          du NIST, Chapitre~\ref{chap:grc})~;
    \item les \emph{capacités opérationnelles} (gestion des identités,
          continuité, gestion des fournisseurs, etc.)~;
    \item les \emph{domaines de sécurité}~: gouvernance et écosystème,
          protection, défense, résilience.
\end{itemize}

\begin{boxinfo}
\textbf{À quoi servent les attributs~?} À vérifier l'équilibre d'un
dispositif. Les mesures de la BANQUE SANAGA en 2025 étaient presque toutes
\emph{préventives}~; les mesures \emph{détectives} (journalisation,
surveillance) et \emph{correctives} (sauvegardes, mode dégradé) manquaient
en agence. Trier la déclaration d'applicabilité par type de mesure l'aurait
montré. C'est aussi ce que faisaient apparaître les étapes du scénario
(Chapitre~\ref{chap:risque}).
\end{boxinfo}

% ============================================================================
\section{La Déclaration d'Applicabilité}
\label{sec:c08-dda}
% ============================================================================

\subsection{Ce qu'elle contient}

\begin{boxdef}
\termecle{Déclaration d'applicabilité}\index{déclaration d'applicabilité}
(DdA)~: document exigé par l'ISO/IEC~27001 (article~6.1.3) qui recense les
mesures nécessaires au traitement des risques, justifie leur inclusion,
indique si elles sont mises en œuvre, et justifie l'exclusion de toute
mesure de l'annexe~A~\cite{iso27001_2022}.
\end{boxdef}

La déclaration fait le pont entre l'appréciation des risques
(Chapitre~\ref{chap:iso27005}) et la liste des mesures. Une mesure y figure
parce qu'un risque l'exige, ou parce qu'une obligation légale ou
contractuelle l'impose. Une mesure en est exclue parce qu'aucun risque ni
aucune obligation ne la rend nécessaire, et non parce qu'elle est difficile.

\subsection{Un extrait}

\begin{boxscenario}{Opération SANAGA. La déclaration d'applicabilité de 2028 (extrait)}
\begin{tableaucours}{cL{3.6cm}cL{4.4cm}c}{%
    \tableauHeader{Mesure} &
    \tableauHeader{Intitulé (traduit)} &
    \tableauHeader{Incl.} &
    \tableauHeader{Justification} &
    \tableauHeader{Mise en œuvre}
}
5.20 & Sécurité dans les accords avec les fournisseurs & Oui & R-DC-02,
R-AG-02 & Oui \\
5.30 & Préparation des TIC pour la continuité & Oui & R-AG-01, récupération
& Oui \\
6.3 & Sensibilisation et formation & Oui & R-AG-01, entrée & Oui \\
8.2 & Droits d'accès privilégiés & Oui & R-AG-01, propagation & Oui \\
8.5 & Authentification sécurisée & Oui & R-AG-01, entrée & Oui \\
8.8 & Gestion des vulnérabilités techniques & Oui & R-AG-01, R-AG-02 & Oui \\
8.13 & Sauvegarde des informations & Oui & R-AG-01, récupération & Oui \\
8.15 & Journalisation & Oui & R-DC-03, loi n°~2010/012, art.~29 & Oui \\
8.22 & Cloisonnement des réseaux & Oui & R-AG-01, propagation & Oui \\
8.28 & Codage sécurisé & Non & « Pas de développement interne~: le
système central est un progiciel » & -- \\
\end{tableaucours}

\smallskip
{\small La première version remise par les directions, en avril 2028,
porte « Oui » dans la colonne « mise en œuvre » pour les neuf mesures
incluses.}
\end{boxscenario}

\subsection{Confronter la déclaration aux preuves}

La RSSI a appris la leçon de février~: avant d'envoyer la déclaration,
elle la confronte aux résultats de l'audit de mars~2028
(Chapitres~\ref{chap:audit} et~\ref{chap:conclure}).

\begin{boxlecture}{Les neuf « mise en œuvre » face aux preuves}
\textbf{Constatation.}
\begin{itemize}
    \item \OUI{}~: 8.5 (second facteur, T1), 8.13 (sauvegardes, T3), 8.22
          (cloisonnement, T2)~;
    \item \NON{}~: 5.30 (mode dégradé, T4), 8.2 (administration locale,
          T7), 8.15 (journalisation incomplète pendant la période, T6)~;
    \item \NEK{}~: 6.3 (sensibilisation, T5), 5.20 (le contrat des GAB n'a
          pas de droit d'audit, T8), 8.8 (vulnérabilités~: GAB non vérifiés,
          postes non testés).
\end{itemize}
3~\OUI{}, 3~\NON{}, 3~\NEK{}~:
\[
\mu = U = \tfrac{1}{3}~;\quad
H = \log_2 3 \approx 1{,}585~\text{bit}.
\]

\textbf{Analyse} (confiance élevée). La déclaration affirmait neuf mesures
sur neuf~; les preuves en confirment trois. La dispersion est maximale~:
tout ce que l'on peut dire de cette déclaration, c'est qu'elle ne
renseigne pas sur l'état réel. La version envoyée à CERTIQUAL AFRIQUE est
corrigée~: « partielle » pour 5.30, 8.2 et 8.15, avec le plan d'actions~;
« à vérifier » pour 6.3, 5.20 et 8.8, avec la date de la vérification
prévue.

\medskip
\hfill\lecture{0{,}33}{0{,}33}{1{,}58}
\end{boxlecture}

\subsection{L'exclusion de 8.28}

L'exclusion du codage sécurisé paraît solide~: la banque n'écrit pas son
système central. Mais la DSI développe, depuis 2026, les interfaces entre
le système central, l'application de paiement et le portail en ligne.
L'exclusion est donc contredite par un fait~: sa base est \NON{}. La mesure
est réintégrée, limitée aux développements d'interfaces.

\begin{boxprocedure}
\textbf{Vérifier une exclusion.} Le papier de travail qui examine une
exclusion de la déclaration d'applicabilité vérifie~:
\begin{itemize}
    \item que la justification renvoie à l'appréciation des risques (aucun
          risque ne rend la mesure nécessaire), et non à une difficulté de
          mise en œuvre~;
    \item qu'aucune obligation légale ou contractuelle n'impose la
          mesure~;
    \item que les faits invoqués sont exacts (ici~: existe-t-il vraiment
          aucun développement interne~?)~;
    \item que l'exclusion est approuvée par la direction, et datée.
\end{itemize}
\end{boxprocedure}

% ============================================================================
\section{Mesurer et Auditer le SMSI}
\label{sec:c08-autres}
% ============================================================================

Deux normes de la famille accompagnent le fonctionnement du SMSI.
L'ISO/IEC~27004 aide à construire des indicateurs qui répondent à
l'exigence de surveillance et de mesure de l'article~9.1~\cite{iso27004_2016}.
L'ISO/IEC~27007 adapte à l'audit d'un SMSI les lignes directrices générales
de l'ISO~19011~\cite{iso27007_2020}. Le Chapitre~\ref{chap:smsi} met en
œuvre l'une et l'autre~: les indicateurs de la RSSI, puis l'audit interne du
SMSI élargi et la préparation de la revue de direction.

\begin{boxinfo}
\textbf{La norme et la loi.} L'article~26 de la loi n°~2010/012 impose aux
exploitants des systèmes d'information de se doter de « systèmes
normalisés » de gestion des risques (Chapitre~\ref{chap:iso27005})~\cite{loi2010012}.
Un SMSI conforme à l'ISO/IEC~27001 en est une forme reconnue
internationalement~; la loi n'en fait pas une obligation, et le certificat
ne dispense pas de l'audit de sécurité réglementaire des articles~13 et~32.
Le Chapitre~\ref{chap:droit} précise l'articulation.
\end{boxinfo}

% ============================================================================
\section{Ce que la Norme demande à chaque Métier}
\label{sec:c08-metiers}
% ============================================================================

\begin{boxmetier}{Consultant en mise en œuvre de SMSI}
\textbf{Une déclaration défendable}
\begin{enumerate}
    \item Partir de l'appréciation des risques, pas de la liste de
          l'annexe~A.
    \item Pour chaque mesure incluse, citer le risque ou l'obligation qui
          l'exige.
    \item Pour chaque « mise en œuvre », désigner la preuve qui le
          démontre.
    \item Justifier chaque exclusion par le risque, et en vérifier les
          faits.
\end{enumerate}
\end{boxmetier}

\begin{boxmetier}{Auditeur en cybersécurité}
\textbf{Lire la déclaration comme un ensemble d'assertions}

Chaque ligne « mise en œuvre » est une assertion de l'audité~; l'auditeur
de certification la teste, comme Mme~EKAMBI a testé le plan de traitement.
Il prélève en priorité les mesures nouvelles de 2022, les exclusions, et
les mesures qui portent les risques les plus élevés.
\end{boxmetier}

\begin{boxmetier}{RSSI}
\textbf{Ne rien déclarer que l'on ne peut prouver}

La RSSI préfère une déclaration honnête, avec des « partielle » et des
« à vérifier », à une déclaration parfaite qui ne résisterait pas à
l'étape~2 de l'audit. Un organisme certificateur pardonne une mesure en
cours~; il ne pardonne pas une déclaration fausse.
\end{boxmetier}

\begin{boxmetier}{Responsable GRC}
\textbf{Relier la déclaration au registre}

La colonne « justification » de la déclaration renvoie aux références du
registre des risques (R-AG-01, R-DC-02\ldots). Le responsable GRC vérifie
que tout risque intolérable a ses mesures dans la déclaration, et que toute
mesure de la déclaration a un risque ou une obligation.
\end{boxmetier}

\separateur

\begin{boxresume}
\textbf{Ce qu'il faut retenir de ce chapitre~:}
\begin{itemize}
    \item \textbf{La famille ISO/IEC~27000} compte une seule norme
          certifiable, l'ISO/IEC~27001~; les autres l'outillent.
    \item \textbf{L'ISO/IEC~27001:2022} suit la structure commune des
          systèmes de management~: articles~4 à~10, du contexte à
          l'amélioration.
    \item \textbf{L'ISO/IEC~27002:2022} compte 93 mesures en quatre thèmes
          (37 organisationnelles, 8 liées aux personnes, 14 physiques, 34
          technologiques), avec des attributs.
    \item \textbf{La déclaration d'applicabilité} relie les risques aux
          mesures~; chaque « mise en œuvre » se prouve, chaque exclusion se
          justifie par le risque.
    \item \textbf{Confrontée aux preuves}, la déclaration de la BANQUE
          SANAGA passe de 9 sur 9 à 3 \OUI{}, 3 \NON{}, 3 \NEK{}.
    \item \textbf{La certification} est volontaire et triennale~; elle ne
          remplace ni l'audit interne, ni l'audit réglementaire.
\end{itemize}
\end{boxresume}

\begin{boxinfo}
\textbf{Points de contrôle.} Vous avez compris ce chapitre si vous savez~:
\begin{itemize}
    \item dire dans quel article de l'ISO/IEC~27001 se trouve une exigence~;
    \item placer une mesure dans son thème et lui associer des attributs~;
    \item rédiger une ligne de déclaration d'applicabilité avec sa
          justification~;
    \item contester une exclusion par les faits.
\end{itemize}
\end{boxinfo}

% ============================================================================
\section*{Exercices}
% ============================================================================

\begin{boxexo}[title={\ding{45}\quad Exercice~8.1. Quel article~? \niveaubase}]
Indiquez l'article de l'ISO/IEC~27001:2022 (4 à 10) dont relève chacune des
activités suivantes.
\begin{enumerate}[(a)]
    \item Le conseil approuve l'extension du périmètre aux agences.
    \item La RSSI rédige la déclaration d'applicabilité.
    \item L'audit interne conduit la mission de mars 2028.
    \item La DSI planifie le remplacement du collecteur de journaux sans
          interrompre la journalisation.
    \item La direction du réseau d'agences corrige le mode dégradé de Kribi
          après l'observation S4.2.
    \item La RSSI organise la formation de deux personnes par agence.
\end{enumerate}
\end{boxexo}

\begin{boxexo}[title={\ding{45}\quad Exercice~8.2. Thèmes et attributs \niveaumoyen}]
Pour chacune des mesures suivantes, donnez son thème et son type
(préventive, détective, corrective), puis dites à quelle étape du scénario
R-AG-01 elle se rattache.
\begin{enumerate}[(a)]
    \item 8.5, authentification sécurisée.
    \item 8.15, journalisation.
    \item 8.13, sauvegarde des informations.
    \item 6.3, sensibilisation et formation.
    \item 8.22, cloisonnement des réseaux.
    \item 5.30, préparation des TIC pour la continuité.
\end{enumerate}
Que montre la répartition par type de mesure~?
\end{boxexo}

\begin{boxexo}[title={\ding{45}\quad Exercice~8.3. Quatre exclusions \niveaumoyen}]
Une banque de la sous-région exclut de sa déclaration d'applicabilité les
mesures suivantes. Pour chacune, dites si l'exclusion est recevable, et
justifiez.
\begin{enumerate}[(a)]
    \item 5.23 (services en nuage)~: « Nous n'utilisons aucun service en
          nuage. » La messagerie est hébergée chez un fournisseur en nuage.
    \item 7.4 (surveillance de la sécurité physique)~: « Trop coûteux pour
          les petites agences. »
    \item 8.28 (codage sécurisé)~: « Aucun développement~; tous les logiciels
          sont des progiciels, paramétrés par l'éditeur. » L'inventaire le
          confirme.
    \item 8.12 (prévention des fuites de données)~: « Mesure couverte par la
          politique de classification. »
\end{enumerate}
\end{boxexo}

\begin{boxexo}[title={\ding{45}\quad Exercice~8.4. La déclaration de LiZbiZ \niveauexpert}]
LiZbiZ veut s'engager vers la certification ISO/IEC~27001 pour rassurer
PHARMACAM (Chapitres~\ref{chap:ebios} à~\ref{chap:audit}).
\begin{enumerate}[(a)]
    \item Proposez un périmètre défendable, en tenant compte des
          dépendances (éditeur, transporteurs, centres de données).
    \item Rédigez six lignes de déclaration d'applicabilité pour les mesures
          qui traitent le scénario SO1, avec leur justification.
    \item L'audit contractuel de PHARMACAM a donné~: comptes par défaut
          supprimés (\OUI{}), services chiffrés (\OUI{}), filtrage entre
          centres (\NON{}), second facteur de l'éditeur (\NEK{}, accès
          refusé), restauration observée (\OUI{}), plan de réponse
          (\NON{}). Remplissez la colonne «~mise en œuvre~», calculez $\mu$,
          $U$ et $H$, et dites si LiZbiZ peut demander l'audit initial.
    \item Quelle mesure de l'ISO/IEC~27002 traite la cause du \NEK{}, et que
          doit faire LiZbiZ avant l'étape~2~?
\end{enumerate}
\end{boxexo}

\bigskip

\begin{boxinfo}
\textbf{Pour aller plus loin.}
\begin{itemize}
    \item L'ISO/IEC~27001:2022~\cite{iso27001_2022}, articles~4 à~10 et
          annexe~A.
    \item L'ISO/IEC~27002:2022~\cite{iso27002_2022}, pour les recommandations
          de mise en œuvre de chaque mesure.
    \item L'ISO/IEC~27000~\cite{iso27000_2018}, pour le vocabulaire.
\end{itemize}
\end{boxinfo}

\begin{boxinfo}
\textbf{Transition.} L'ISO/IEC~27001 n'est pas le seul référentiel. La
banque doit aussi répondre au superviseur, et ses partenaires parlent
parfois NIST ou CIS. Comment passer d'un référentiel à l'autre sans tout
refaire~? C'est l'objet du Chapitre~\ref{chap:referentiels}.
\end{boxinfo}

        % ============================================================================
% Fichier  : partie3/P03_C09_referentiels_correspondances.tex
% Titre    : Les Autres Référentiels et leurs Correspondances
% Auteur   : MINKA MI NGUIDJOÏ Thierry Emmanuel
% Version  : 1.0 -- novembre 2026
% Corrigés : corriges/Corr_C09.tex
% ============================================================================
% Sources : NIST CSF 2.0 (2024 : 6 fonctions, 22 catégories, 106
% sous-catégories, quatre niveaux de mise en œuvre, profils) ; NIST SP
% 800-53 rév. 5 (20 familles) ; CIS Controls v8 (2021, 18 contrôles, 153
% mesures, 3 groupes de mise en œuvre, IG1 = 56 mesures ; révision 8.1 de
% 2024) ; PCI DSS v4.0 (2022) ; critères de services de confiance de l'AICPA
% (SOC 2) ; ISAE 3402 (types 1 et 2) ; COBIT 2019 (40 objectifs). Chiffres
% d'après les publications des éditeurs, à confirmer par l'auteur ;
% intitulés traduits par le cours.
% Faits du cas (bible et corpus) : ANPP (Yaoundé), agence publique de
% paiement contrôlée par la Chambre des comptes ; M. ESSOMBA, directeur
% général ; M. NJOYA, magistrat vérificateur ; 4 agents payés après leur
% départ ; aucune procédure écrite de désactivation des comptes ; pièces
% « aux archives ».
% Faits nouveaux (reportés dans la bible) : en 2027, un bailleur de fonds
% exige un profil NIST CSF 2.0 ; base commune de 10 contrôles B1 à B10 ;
% résultats 4/3/3 ; projection par fonction du CSF ; niveau actuel 1
% « partiel », cible 3 ; DIGICORE pourrait remettre un rapport ISAE 3402 de
% type 2.
% ============================================================================

\chapter{Les Autres Référentiels et leurs Correspondances}
\label{chap:referentiels}
\chaptermark{Référentiels et correspondances}

\epigraph{%
    « Tester une fois,\\
    rendre compte autant de fois qu'il le faut. »%
}{Principe du cours}

\niveauChapitre{Fondamental}{%
    Maîtrise du Chapitre~\ref{chap:iso27000}%
}

\begin{boxobjectifs}
\begin{itemize}
    \item Présenter le cadre de cybersécurité du NIST (CSF~2.0), le
          catalogue NIST SP~800-53, les CIS Controls et PCI DSS.
    \item Situer les rapports d'assurance SOC~2 et ISAE~3402, et COBIT~2019.
    \item Construire une table de correspondance entre référentiels, et en
          connaître les pièges.
    \item Bâtir une base commune de contrôles, testée une fois et projetée
          sur plusieurs référentiels.
    \item Choisir les référentiels d'une organisation.
\end{itemize}
\end{boxobjectifs}

\bigskip

% ============================================================================
\section*{Avant de commencer : trois audits pour un seul système}
% ============================================================================

Yaoundé, mars 2027. L'ANPP\footnote{L'ANPP, ses dirigeants, le bailleur,
les dates et les montants sont fictifs~; les référentiels cités sont
exacts.}, agence publique de paiement, finance la modernisation de ses
systèmes avec l'appui d'un bailleur de fonds international. La convention
de financement exige un « profil de cybersécurité selon le cadre du NIST,
version~2.0 », actuel et cible. La Chambre des comptes, par la voix du
magistrat vérificateur M.~NJOYA, prépare son contrôle des paiements et de
leurs systèmes. Et l'audit de sécurité prévu par la loi n°~2010/012 pourrait
s'appliquer à l'agence (Chapitre~\ref{chap:droit}).

M.~ESSOMBA, directeur général, refuse d'organiser trois chantiers~: «~Nos
serveurs sont les mêmes pour tout le monde. Je veux un seul travail, et
trois rapports.~»

\medskip

Ce chapitre présente les principaux référentiels qui coexistent avec
l'ISO/IEC~27001, puis la technique qui permet d'en satisfaire plusieurs à
la fois~: la table de correspondance et la base commune de contrôles.

\begin{boxalert}
\textbf{Violation typique.} Conclure qu'une organisation est conforme au
cadre du NIST parce qu'elle est certifiée ISO/IEC~27001, au motif qu'une
table de correspondance relie les deux.

\smallskip
\textbf{Règle~:} une correspondance dit que deux exigences portent sur le
même sujet, pas qu'elles disent la même chose. La conformité se démontre
référentiel par référentiel, à partir des preuves.
\end{boxalert}

\separateur

% ============================================================================
\section{Le Cadre de Cybersécurité du NIST}
\label{sec:c09-csf}
% ============================================================================

\subsection{Six fonctions, vingt-deux catégories}

Le cadre de cybersécurité du NIST (\sigle{CSF}), dans sa version~2.0 de
2024, organise la gestion des risques cyber en six fonctions, déclinées en
22 catégories et 106 sous-catégories~\cite{nist2024csf}. Le
Chapitre~\ref{chap:grc} a présenté la fonction \emph{Gouverner}~; voici
l'ensemble.

\begin{tableaucours}{L{2.4cm}L{9.8cm}}{%
    \tableauHeader{Fonction} &
    \tableauHeader{Catégories (intitulés traduits par le cours)}
}
Gouverner (GV) & Contexte organisationnel (OC)~; stratégie de gestion des
risques (RM)~; rôles, responsabilités et autorités (RR)~; politique (PO)~;
supervision (OV)~; risques de la chaîne d'approvisionnement (SC) \\
Identifier (ID) & Gestion des actifs (AM)~; appréciation des risques
(RA)~; amélioration (IM) \\
Protéger (PR) & Identités, authentification et contrôle d'accès (AA)~;
sensibilisation et formation (AT)~; sécurité des données (DS)~; sécurité
des plateformes (PS)~; résilience de l'infrastructure (IR) \\
Détecter (DE) & Surveillance continue (CM)~; analyse des événements
indésirables (AE) \\
Répondre (RS) & Gestion des incidents (MA)~; analyse (AN)~; communication
(CO)~; atténuation (MI) \\
Rétablir (RC) & Exécution du plan de rétablissement (RP)~; communication
(CO) \\
\end{tableaucours}
\vspace{-0.8em}
\begin{center}{\small\itshape Tableau~9.1. Les fonctions et catégories du
NIST CSF~2.0}\end{center}

\subsection{Profils et niveaux}

Le CSF ne fixe pas d'exigences~; il donne un langage. Une organisation
décrit son \emph{profil actuel} (ce qu'elle fait, sous-catégorie par
sous-catégorie) et son \emph{profil cible} (ce qu'elle vise), puis comble
l'écart. Elle peut aussi caractériser la rigueur de sa gestion des risques
par l'un des quatre \emph{niveaux de mise en œuvre}~: partiel, informé par
le risque, reproductible, adaptatif.

\begin{boxinfo}
\textbf{Pourquoi les bailleurs aiment le CSF.} Il est gratuit, lisible par
une direction, et ne suppose pas de certification. Un profil actuel et un
profil cible donnent au bailleur une feuille de route qu'il peut suivre
d'une année sur l'autre. C'est aussi sa limite~: un profil est une
auto-description~; il ne vaut que par les preuves qui l'appuient.
\end{boxinfo}

% ============================================================================
\section{Catalogues de Mesures}
\label{sec:c09-catalogues}
% ============================================================================

\subsection{NIST SP~800-53}

Le catalogue NIST SP~800-53, révision~5, compte plus d'un millier de
mesures et d'améliorations, réparties en vingt familles (contrôle d'accès,
audit et responsabilité, réponse aux incidents, etc.)~\cite{nist2020sp80053}.
Il est obligatoire pour les systèmes fédéraux américains et sert ailleurs
de réservoir de mesures détaillées. Son compagnon, le SP~800-53A, décrit la
manière d'évaluer chaque mesure (Chapitre~\ref{chap:audit}).

\subsection{Les CIS Controls}

Les CIS Controls, version~8, publiés par le Center for Internet Security,
retiennent 18 contrôles et 153 mesures, classés par ordre de
priorité~\cite{cis2021controls}. Leur originalité est de définir trois
\emph{groupes de mise en œuvre}~: le premier (IG1) réunit 56 mesures,
présentées comme l'hygiène de base que toute organisation devrait appliquer~;
les groupes~2 et~3 ajoutent des mesures pour les organisations plus
exposées ou plus outillées.

\begin{tableaucours}{cL{5.0cm}cL{5.0cm}}{%
    \tableauHeader{N°} &
    \tableauHeader{Contrôle (traduit)} &
    \tableauHeader{N°} &
    \tableauHeader{Contrôle (traduit)}
}
1 & Inventaire des actifs matériels & 10 & Défenses contre les maliciels \\
2 & Inventaire des logiciels & 11 & Restauration des données \\
3 & Protection des données & 12 & Gestion de l'infrastructure réseau \\
4 & Configuration sécurisée & 13 & Surveillance et défense du réseau \\
5 & Gestion des comptes & 14 & Sensibilisation et formation \\
6 & Gestion du contrôle d'accès & 15 & Gestion des prestataires \\
7 & Gestion continue des vulnérabilités & 16 & Sécurité des applications \\
8 & Gestion des journaux & 17 & Réponse aux incidents \\
9 & Protection de la messagerie et du navigateur & 18 & Tests d'intrusion \\
\end{tableaucours}
\vspace{-0.8em}
\begin{center}{\small\itshape Tableau~9.2. Les 18 CIS Controls (version~8)}\end{center}

\begin{boxinfo}
\textbf{Un point de départ pour les petites organisations.} Pour une
entreprise comme LiZbiZ, qui n'a ni RSSI ni audit interne, le groupe IG1
des CIS Controls est un objectif réaliste et vérifiable avant toute
certification. Le Chapitre~\ref{chap:audit-technique} s'appuie sur les
guides de configuration du même organisme.
\end{boxinfo}

\subsection{PCI DSS}

La norme PCI DSS, version~4.0, est imposée par les réseaux de cartes de
paiement à toute organisation qui stocke, traite ou transmet des données
de cartes~\cite{pcissc2022dss}. Elle est contractuelle, détaillée et
vérifiée par des évaluateurs agréés. Elle ne concerne pas l'ANPP, qui paie
par virement~; elle concerne la BANQUE SANAGA pour son activité monétique,
sur un périmètre (l'environnement des données de cartes) distinct de celui
de son SMSI.

% ============================================================================
\section{Rapports d'Assurance et Gouvernance}
\label{sec:c09-assurance}
% ============================================================================

\subsection{SOC~2 et ISAE~3402}

Quand une organisation confie une activité à un prestataire, elle ne peut
pas toujours l'auditer elle-même. Elle peut alors demander un
\termecle{rapport d'assurance}\index{rapport d'assurance} établi par un
auditeur indépendant sur les contrôles du prestataire.

\begin{tableaucours}{L{2.4cm}L{4.6cm}L{5.2cm}}{%
    \tableauHeader{Rapport} &
    \tableauHeader{Cadre} &
    \tableauHeader{Contenu}
}
ISAE~3402 & Norme internationale d'assurance~\cite{iaasb2011isae3402} &
Contrôles d'un prestataire qui affectent l'information financière de ses
clients \\
SOC~2 & Critères de services de confiance de l'AICPA~\cite{aicpa2017tsc} &
Sécurité, disponibilité, intégrité des traitements, confidentialité, vie
privée \\
\end{tableaucours}
\vspace{-0.8em}
\begin{center}{\small\itshape Tableau~9.3. Deux rapports d'assurance sur les
prestataires}\end{center}

Chacun existe en deux types~: le \emph{type~1} porte sur la conception des
contrôles à une date~; le \emph{type~2} sur leur efficacité pendant une
période, généralement de six à douze mois. Seul le type~2 dit si les
contrôles ont fonctionné.

\begin{boxscenario}{Opération SANAGA. Une issue pour DIGICORE}
\textbf{Faits.} DIGICORE refuse l'avenant qui donnerait à la banque un droit
d'audit (Chapitre~\ref{chap:iso27005}).

\medskip
\textbf{Analyse} (confiance moyenne). Un prestataire qui sert plusieurs
banques refuse souvent d'être audité par chacune. La banque peut exiger, à
la place, un rapport d'assurance de type~2 par an, établi par un auditeur
indépendant, sur les contrôles qui la concernent. La base \NEK{} de
R-DC-02 deviendrait alors une base \OUI{} ou \NON{}, sans que la banque
entre chez le prestataire.
\end{boxscenario}

\subsection{COBIT~2019}

COBIT~2019, de l'ISACA, est un cadre de gouvernance et de gestion de
l'informatique, qui compte quarante objectifs~\cite{isaca2018cobitintro}.
Il n'est pas spécifique à la sécurité~: il la situe dans l'ensemble du
pilotage du système d'information. Le cours GCI (MSN-2S3-MSI03) l'étudie en
détail~; NRAC n'en retient que la distinction entre gouvernance et gestion
(Chapitre~\ref{chap:grc}).

% ============================================================================
\section{Les Tables de Correspondance}
\label{sec:c09-correspondances}
% ============================================================================

\subsection{Ce qu'est une correspondance}

\begin{boxdef}
\termecle{Table de correspondance}\index{table de correspondance}~:
tableau qui relie les exigences d'un référentiel à celles d'un autre, en
indiquant, pour chaque lien, s'il est complet ou partiel. Elle sert à
réutiliser un travail fait pour un référentiel au profit d'un autre.
\end{boxdef}

\begin{boxprocedure}
\textbf{Construire une correspondance fiable.} Le papier de travail qui
établit ou utilise une table de correspondance vérifie~:
\begin{itemize}
    \item la \textbf{source} de chaque lien~: publication officielle (le
          NIST et le CIS publient des correspondances), travail interne,
          ou fournisseur~;
    \item la \textbf{nature} du lien~: complet (les deux exigences ont le
          même contenu) ou partiel (l'une couvre une partie de l'autre)~;
    \item le \textbf{sens} du lien~: satisfaire A peut suffire pour B sans
          que l'inverse soit vrai~;
    \item la \textbf{multiplicité}~: une exigence peut correspondre à
          plusieurs autres, et réciproquement~;
    \item la \textbf{version} de chaque référentiel~: une table ISO/IEC
          27002:2013 ne vaut pas pour l'édition de 2022.
\end{itemize}
\end{boxprocedure}

\subsection{La base commune de l'ANPP}

Le responsable de la sécurité de l'ANPP construit une \termecle{base commune
de contrôles}\index{base commune de contrôles}~: une liste de contrôles
propres à l'agence, chacun relié aux référentiels qui l'exigent.

\begin{tableaucours}{cL{4.2cm}ccc}{%
    \tableauHeader{Réf.} &
    \tableauHeader{Contrôle de l'ANPP} &
    \tableauHeader{ISO/IEC 27002} &
    \tableauHeader{CSF~2.0} &
    \tableauHeader{CIS v8}
}
B1 & Inventaire des postes et serveurs & 5.9 & ID.AM & 1 \\
B2 & Désactivation des comptes au départ & 5.18 & PR.AA & 5 \\
B3 & Authentification multifacteur & 8.5 & PR.AA & 6 \\
B4 & Journalisation et revue & 8.15 & DE.CM & 8 \\
B5 & Sauvegardes testées & 8.13 & PR.DS, RC.RP & 11 \\
B6 & Application des correctifs & 8.8 & ID.RA & 7 \\
B7 & Sensibilisation des agents & 6.3 & PR.AT & 14 \\
B8 & Sécurité des prestataires & 5.19, 5.20 & GV.SC & 15 \\
B9 & Réponse aux incidents & 5.24, 5.26 & RS.MA & 17 \\
B10 & Test d'intrusion annuel & (partiel) & ID.IM & 18 \\
\end{tableaucours}
\vspace{-0.8em}
\begin{center}{\small\itshape Tableau~9.4. La base commune de l'ANPP
(correspondances établies par le cours)}\end{center}

Trois remarques. B5 correspond à deux catégories du CSF~: la sauvegarde
protège les données et sert au rétablissement. B10 n'a pas de mesure
dédiée dans l'ISO/IEC~27002~; le test d'intrusion n'y apparaît que comme un
moyen parmi d'autres, d'où la mention «~partiel~». Enfin, la base ne couvre
pas tout~: plusieurs catégories du CSF (le contexte organisationnel, la
politique, la communication) n'ont pas de contrôle dans cette première
version.

\subsection{Tester une fois, rendre compte trois fois}

\begin{boxlecture}{La base commune de l'ANPP, évaluée en 2027}
\textbf{Constatation.} L'équipe de l'ANPP, appuyée par un cabinet, teste
une fois les dix contrôles~:
\begin{itemize}
    \item \OUI{}~: B1, B5, B9 (exercice de 2026), B10 (rapport de 2027)~;
    \item \NON{}~: B2 (quatre agents payés après leur départ, comptes
          encore actifs~; aucune procédure écrite), B3 (pas de second
          facteur), B7 (aucune sensibilisation depuis 2023)~;
    \item \NEK{}~: B4 (journaux «~aux archives~», non produits), B6
          (aucun relevé des correctifs), B8 (contrats sans clause de
          sécurité, aucun rapport des prestataires).
\end{itemize}
4~\OUI{}, 3~\NON{}, 3~\NEK{}~:
\[
\mu = 0{,}40~;\quad U = 0{,}30~;\quad
H = -(0{,}4\log_2 0{,}4 + 2 \times 0{,}3\log_2 0{,}3)
  \approx 0{,}529 + 0{,}521 + 0{,}521 \approx 1{,}57~\text{bit}.
\]

\textbf{Projection sur le CSF} (par fonction)~: Gouverner 0/0/1 (B8)~;
Identifier 2/0/1 (B1, B10~; B6)~; Protéger 1/3/0 (B5~; B2, B3, B7)~;
Détecter 0/0/1 (B4)~; Répondre 1/0/0 (B9)~; Rétablir 1/0/0 (B5).

\medskip
\textbf{Analyse} (confiance moyenne). La même série de tests donne trois
lectures. Pour le bailleur, le profil actuel est faible en
\emph{Protéger} (des défauts établis) et inconnu en \emph{Gouverner} et
\emph{Détecter}~: le niveau de mise en œuvre est au mieux «~partiel~», et
l'ANPP vise «~reproductible~». Pour la Chambre des comptes, B2 rejoint le
constat des quatre agents payés après leur départ~: le défaut de
désactivation est aussi un défaut de contrôle des paiements. Pour
l'audit réglementaire, les mêmes preuves serviront, référentiel par
référentiel.

\medskip
\hfill\lecture{0{,}40}{0{,}30}{1{,}57}
\end{boxlecture}

\begin{boxalert}
\textbf{Violation typique.} Reporter une cote sur un référentiel par la
seule table, alors que le lien est partiel.

\smallskip
\textbf{Règle~:} un \OUI{} ne se transporte que par un lien complet. Par un
lien partiel, le \OUI{} devient au mieux un \NEK{} sur l'autre référentiel,
jusqu'à ce que la partie non couverte soit testée.
\end{boxalert}

% ============================================================================
\section{Choisir ses Référentiels}
\label{sec:c09-choisir}
% ============================================================================

On ne choisit pas un référentiel pour sa réputation, mais pour ce qu'il
oblige ou permet.

\begin{tableaucours}{L{3.4cm}L{8.8cm}}{%
    \tableauHeader{Critère} &
    \tableauHeader{Question à se poser}
}
Obligation & La loi, le superviseur ou un contrat impose-t-il un
référentiel (PCI DSS, exigence d'un bailleur, audit réglementaire)~? \\
Reconnaissance & Les clients ou partenaires attendent-ils une certification
(ISO/IEC~27001) ou un rapport (SOC~2, ISAE~3402)~? \\
Taille et maturité & L'organisation peut-elle faire vivre un SMSI, ou
doit-elle commencer par une hygiène de base (CIS IG1)~? \\
Langage & Faut-il parler à une direction et à un bailleur (CSF) ou à des
techniciens (SP~800-53, guides de configuration)~? \\
Coût de la preuve & Combien coûte la démonstration de conformité, et
peut-elle être mutualisée (base commune)~? \\
\end{tableaucours}
\vspace{-0.8em}
\begin{center}{\small\itshape Tableau~9.5. Critères de choix d'un référentiel}\end{center}

\begin{boxscenario}{Trois organisations, trois choix}
\textbf{BANQUE SANAGA.} ISO/IEC~27001 certifiée, sur toute la banque~;
PCI DSS pour la monétique~; exigences du superviseur
(Chapitre~\ref{chap:droit})~; rapports ISAE~3402 demandés aux prestataires.

\medskip
\textbf{LiZbiZ.} CIS Controls IG1 d'abord, comme socle vérifiable~;
ISO/IEC~27001 ensuite, si PHARMACAM l'exige au renouvellement suivant.

\medskip
\textbf{ANPP.} Une base commune ISO/IEC~27002, projetée sur le CSF pour le
bailleur et sur les textes pour la Chambre des comptes et l'audit
réglementaire.
\end{boxscenario}

% ============================================================================
\section{Ce que les Référentiels demandent à chaque Métier}
\label{sec:c09-metiers}
% ============================================================================

\begin{boxmetier}{Responsable GRC}
\textbf{Une base, plusieurs rapports}
\begin{enumerate}
    \item Recenser les référentiels imposés et attendus.
    \item Construire une base commune de contrôles, chacun relié aux
          référentiels par des liens qualifiés (complet, partiel).
    \item Tester chaque contrôle une fois, avec des preuves réutilisables.
    \item Produire les rapports par projection, en signalant les liens
          partiels.
\end{enumerate}
\end{boxmetier}

\begin{boxmetier}{Auditeur en cybersécurité}
\textbf{Ne pas auditer la table}

L'auditeur à qui l'on présente une table de correspondance vérifie sa
source, sa version et la nature des liens. Il refuse de reporter une
conformité par un lien partiel, et prélève des contrôles dans les zones
que la base ne couvre pas.
\end{boxmetier}

\begin{boxmetier}{Gestionnaire des risques cyber}
\textbf{Le référentiel ne remplace pas le risque}

Un référentiel donne des mesures~; il ne dit pas lesquelles comptent le
plus pour l'organisation. Le gestionnaire des risques vérifie que les
contrôles de la base traitent les scénarios du registre, et pas seulement
les cases des référentiels.
\end{boxmetier}

\separateur

\begin{boxresume}
\textbf{Ce qu'il faut retenir de ce chapitre~:}
\begin{itemize}
    \item \textbf{NIST CSF~2.0}~: six fonctions, 22 catégories~; un langage,
          des profils et des niveaux, pas d'exigences.
    \item \textbf{Catalogues}~: NIST SP~800-53 (détaillé), CIS Controls (18
          contrôles priorisés, groupe IG1 pour l'hygiène de base), PCI DSS
          (cartes de paiement, contractuel).
    \item \textbf{Rapports d'assurance}~: ISAE~3402 et SOC~2~; seul le
          type~2 dit si les contrôles ont fonctionné.
    \item \textbf{Une correspondance} relie des sujets, pas des
          contenus~; elle a une source, une nature, un sens et une
          version.
    \item \textbf{Une base commune} se teste une fois et se projette sur
          plusieurs référentiels~; un \OUI{} ne se transporte que par un
          lien complet.
\end{itemize}
\end{boxresume}

\begin{boxinfo}
\textbf{Points de contrôle.} Vous avez compris ce chapitre si vous savez~:
\begin{itemize}
    \item placer une pratique dans une fonction du CSF~;
    \item dire ce qu'apporte un rapport de type~2 par rapport au type~1~;
    \item construire une ligne de base commune avec trois référentiels~;
    \item refuser une conformité déduite d'un lien partiel.
\end{itemize}
\end{boxinfo}

% ============================================================================
\section*{Exercices}
% ============================================================================

\begin{boxexo}[title={\ding{45}\quad Exercice~9.1. Quelle fonction du CSF~? \niveaubase}]
Placez chaque pratique dans une fonction du NIST CSF~2.0, et, si possible,
dans une catégorie.
\begin{enumerate}[(a)]
    \item Le conseil fixe l'appétence au risque cyber.
    \item La DSI tient l'inventaire des postes des agences.
    \item Les agents suivent une formation à l'hameçonnage.
    \item Le collecteur de journaux déclenche une alerte sur une connexion
          inhabituelle.
    \item La cellule de crise décide d'isoler les agences touchées.
    \item Les agences rouvrent après restauration des postes.
    \item La banque exige un rapport de type~2 de DIGICORE.
\end{enumerate}
\end{boxexo}

\begin{boxexo}[title={\ding{45}\quad Exercice~9.2. Correspondances partielles \niveaumoyen}]
\begin{enumerate}[(a)]
    \item Pour chacun des contrôles B2, B4 et B10 de l'ANPP, dites si le lien
          avec la mesure ISO/IEC~27002 indiquée est complet ou partiel, et
          pourquoi.
    \item Un consultant affirme que l'ANPP, étant «~conforme à 60~\% de
          l'ISO/IEC~27002~», est «~au niveau~2 du CSF~». Relevez les deux
          erreurs de raisonnement.
    \item B10 est \OUI{}. Quelle cote reporter sur l'ISO/IEC~27002, et
          pourquoi~?
\end{enumerate}
\end{boxexo}

\begin{boxexo}[title={\ding{45}\quad Exercice~9.3. L'hygiène de base de LiZbiZ \niveaumoyen}]
\begin{enumerate}[(a)]
    \item Justifiez le choix du groupe IG1 des CIS Controls pour LiZbiZ, à
          l'aide des critères du Tableau~9.5.
    \item Reliez chacun des constats de l'audit des quatre centres (comptes
          par défaut, services en clair, flux non autorisés, sauvegardes non
          restaurées) à un contrôle CIS.
    \item PHARMACAM demande «~une certification~». Que répondez-vous à
          M.~ETEME~: CIS, ISO/IEC~27001 ou rapport SOC~2~? Justifiez.
\end{enumerate}
\end{boxexo}

\begin{boxexo}[title={\ding{45}\quad Exercice~9.4. Trois rapports pour l'ANPP \niveauexpert}]
À partir des résultats de la base commune (4/3/3)~:
\begin{enumerate}[(a)]
    \item Calculez $\mu$, $U$ et $H$ pour la fonction \emph{Protéger} et
          pour la fonction \emph{Identifier}, et appliquez la règle de
          décision du Chapitre~\ref{chap:conclure}.
    \item Rédigez, en dix lignes, la partie «~profil actuel et profil
          cible~» destinée au bailleur.
    \item Rédigez, en cinq lignes, la note destinée à M.~NJOYA sur le
          contrôle B2, en reliant le défaut de sécurité au défaut de
          paiement.
    \item Proposez trois contrôles à ajouter à la base pour couvrir les
          catégories du CSF qu'elle ignore, avec leur correspondance
          ISO/IEC~27002.
\end{enumerate}
\end{boxexo}

\bigskip

\begin{boxinfo}
\textbf{Pour aller plus loin.}
\begin{itemize}
    \item Le NIST CSF~2.0~\cite{nist2024csf} et ses références informatives.
    \item Les CIS Controls, version~8~\cite{cis2021controls}.
    \item La norme ISAE~3402~\cite{iaasb2011isae3402} et les critères de
          services de confiance de l'AICPA~\cite{aicpa2017tsc}.
\end{itemize}
\end{boxinfo}

\begin{boxinfo}
\textbf{Transition.} Les référentiels sont volontaires ou contractuels. La
loi, elle, s'impose~: quelles organisations doivent être auditées, par
qui, avec quelles garanties pour l'auditeur et l'audité~? Le
Chapitre~\ref{chap:droit} présente le cadre juridique camerounais et la
déontologie de l'auditeur.
\end{boxinfo}

        % ============================================================================
% Fichier  : partie3/P03_C10_cadre_juridique_deontologie.tex
% Titre    : Le Cadre Camerounais et Africain, la Déontologie
% Auteur   : MINKA MI NGUIDJOÏ Thierry Emmanuel
% Version  : 1.0 -- novembre 2026
% Corrigés : corriges/Corr_C10.tex
% ============================================================================
% Sources et degré de vérification (repris d'ACSI, chapitre 12, et Ann. D) :
%  - Loi n° 2010/012 : texte français officiel lu ; cités textuellement ici :
%    art. 7 al. 2 et 14 ; art. 61 al. 1 ; art. 4 (9), 13 al. 1 et 32 renvoyés
%    au chapitre 6 ; art. 26, 29 et 30 résumés ; art. 61 al. 3 paraphrasé.
%  - Règlement COBAC R-2016/04 : art. 1 et 2 relevés sur une base en ligne
%    (cours GCI), À COLLATIONNER avec la version du SG-COBAC.
%  - Convention de Malabo (UA, 2014) : entrée en vigueur en 2023 ; statut
%    camerounais À VÉRIFIER.
%  - Loi n° 2024/017 : présentée sans citation, À VÉRIFIER au Journal
%    officiel ; étudiée dans l'UE MSN-2S3-JUR01.
%  - Déontologie : principes éthiques des normes mondiales d'audit interne
%    de l'IIA (2024) ; menaces sur l'indépendance d'après le cadre
%    conceptuel du code de l'IESBA.
% Faits du cas : ANPP (M. ESSOMBA), BANQUE SANAGA (COBAC, rapport après
% l'incident), LiZbiZ (AUDIX CONSEIL, M. BEYALA, PHARMACAM).
% Faits nouveaux (reportés dans la bible) : en 2028, l'ANTIC annonce un
% audit de sécurité à l'ANPP ; cinq obligations légales cotées 0/3/2 ; la
% COBAC demande en 2026 un rapport sur l'incident de septembre ; M. ETEME
% demande à AUDIX d'adoucir la synthèse destinée à PHARMACAM ; la DSI de la
% SANAGA propose un « petit-déjeuner » au chef de mission de CERTIQUAL.
% ============================================================================

\chapter{Le Cadre Camerounais et Africain, la Déontologie}
\label{chap:droit}
\chaptermark{Cadre juridique et déontologie}

\epigraph{%
    « L'auditeur constate un écart au critère.\\
    Le juge, seul, qualifie l'infraction. »%
}{Règle de prudence}

\niveauChapitre{Fondamental}{%
    Maîtrise du Chapitre~\ref{chap:audit}%
}

\begin{boxobjectifs}
\begin{itemize}
    \item Distinguer, dans un texte juridique, ce qui est vérifié de ce qui
          reste à vérifier, et ne citer que le premier.
    \item Identifier les obligations d'audit et de sécurité de la loi
          n°~2010/012, et le rôle de l'ANTIC.
    \item Situer le règlement COBAC R-2016/04, la Convention de Malabo et
          la loi sur les données personnelles.
    \item Faire de la loi une source de critères d'audit.
    \item Appliquer les principes déontologiques de l'auditeur et
          reconnaître les menaces sur son indépendance.
    \item Rédiger un constat sans qualifier juridiquement les faits.
\end{itemize}
\end{boxobjectifs}

\bigskip

% ============================================================================
\section*{Avant de commencer : trois lettres}
% ============================================================================

Trois lettres, la même semaine de 2028\footnote{Les organisations, les
personnes, les dates et les montants des cas sont fictifs~; les textes
cités sont exacts, avec le degré de vérification indiqué.}.

À Yaoundé, M.~ESSOMBA reçoit de l'Agence nationale des technologies de
l'information et de la communication (\sigle{ANTIC}) l'annonce d'un audit
de sécurité de l'ANPP. Il se demande si son agence y est vraiment soumise,
et ce qu'il doit remettre.

À Douala, M.~ETEME écrit à AUDIX CONSEIL~: la synthèse destinée à
PHARMACAM « pourrait être formulée plus positivement », d'autant que le
renouvellement du contrat de LiZbiZ, et donc la suite des missions
d'AUDIX, en dépend.

À la BANQUE SANAGA enfin, la DSI propose au chef de mission de CERTIQUAL
AFRIQUE, la veille de l'étape~2, un « petit-déjeuner de travail » dans un
grand hôtel.

\medskip

Ce chapitre présente les textes qui encadrent la sécurité et son audit au
Cameroun et en Afrique, puis les règles de déontologie qui protègent
l'auditeur et la valeur de son travail.

\begin{boxalert}
\textbf{Règle du chapitre.} Un texte de loi n'est cité entre guillemets
que s'il a été lu dans sa version officielle. Les autres dispositions sont
présentées sans citation, avec la mention « à vérifier ». Un auditeur
applique la même règle dans ses rapports~: une citation inexacte ruine la
crédibilité de tout le reste.
\end{boxalert}

\separateur

% ============================================================================
\section{Panorama des Textes}
\label{sec:c10-panorama}
% ============================================================================

\begin{tableaucours}{L{3.2cm}L{3.0cm}L{3.8cm}L{2.2cm}}{%
    \tableauHeader{Texte} &
    \tableauHeader{Assujettis} &
    \tableauHeader{Ce qu'il impose} &
    \tableauHeader{Vérification}
}
Loi n°~2010/012 (cybersécurité et cybercriminalité) & Opérateurs,
exploitants de SI, autorités de certification, organismes publics & Mesures
de sécurité, gestion des risques, conservation des journaux, audit de
sécurité & Texte officiel lu \\
Règlement COBAC R-2016/04 & Établissements de crédit de la CEMAC &
Système de contrôle interne, dont le risque opérationnel & Art.~1 et~2
relevés, à collationner \\
Convention de Malabo (Union africaine, 2014) & États parties & Cadres
nationaux de cybersécurité et de protection des données & Statut
camerounais à vérifier \\
Loi n°~2024/017 (données personnelles) & Responsables de traitement &
Protection des données, y compris pendant l'audit & À vérifier au Journal
officiel \\
Droit OHADA des sociétés & Sociétés commerciales & Commissaire aux
comptes, comité d'audit & Présenté, à vérifier \\
\end{tableaucours}
\vspace{-0.8em}
\begin{center}{\small\itshape Tableau~10.1. Les textes et leur degré de
vérification dans ce cours}\end{center}

% ============================================================================
\section{La Loi sur la Cybersécurité}
\label{sec:c10-2010}
% ============================================================================

\subsection{L'audit de sécurité et l'ANTIC}

Le Chapitre~\ref{chap:audit} a cité la définition légale de l'audit de
sécurité (article~4, point~9) et le régime d'audit obligatoire et
périodique (articles~13 et~32)~\cite{loi2010012}. La loi confie à l'ANTIC,
désignée comme «~l'Agence~», une mission de contrôle de ces audits.

\begin{boxjuris}
\textbf{\loicam{n°~2010/012}, article~7, alinéa~2} (extrait)~: l'Agence a
notamment pour mission \textit{« de s'assurer de la régularité, de
l'effectivité des audits de sécurité des systèmes d'information suivant
les normes en la matière, des organismes publics et des autorités de
certification »}.

\medskip
\textbf{Portée.} Deux points. La loi renvoie aux \emph{normes}~: un audit
de sécurité légal se conduit selon les normes et référentiels des
Chapitres~\ref{chap:audit} et~\ref{chap:iso27000}. Et elle vise
expressément les \emph{organismes publics}, ce qui concerne l'ANPP.
\end{boxjuris}

\subsection{L'ANPP est-elle soumise~?}

\begin{boxscenario}{Opération ANPP. Lire le texte avant de conclure}
\textbf{Constatations.}
\begin{itemize}
    \item L'article~7, alinéa~2, charge l'Agence de s'assurer de la
          régularité des audits de sécurité des organismes publics.
    \item L'article~13, alinéa~1, soumet à l'audit obligatoire «~les
          réseaux de communications électroniques et les systèmes
          d'information, des opérateurs, des autorités de certification et
          des fournisseurs de services de communications électroniques~».
    \item Les conditions et modalités sont renvoyées à des textes
          d'application.
\end{itemize}

\textbf{Analyse} (confiance faible sur la portée de l'article~13, élevée
sur l'article~7). La ponctuation de l'article~13 admet deux lectures~:
soit tous les systèmes d'information sont visés, soit seulement ceux des
opérateurs, autorités et fournisseurs. L'article~7 tranche en partie pour
l'ANPP~: ses audits de sécurité relèvent du contrôle de l'Agence en tant
qu'organisme public. La réponse sûre à M.~ESSOMBA est donc~: l'audit
annoncé s'inscrit dans la mission légale de l'ANTIC~; ses modalités se
lisent dans les textes d'application en vigueur, que l'ANPP doit demander.
Un auditeur ne tranche pas seul une ambiguïté de texte~: il la signale.
\end{boxscenario}

\subsection{Des obligations qui deviennent des critères}

Plusieurs articles imposent aux exploitants des systèmes d'information des
obligations qui deviennent, pour l'auditeur, des critères opposables. Le
tableau les résume d'après le texte officiel~; seul l'article~29, alinéa~1,
a été cité textuellement (Chapitre~\ref{chap:iso27005}).

\begin{tableaucours}{L{1.6cm}L{6.2cm}L{4.4cm}}{%
    \tableauHeader{Art.} &
    \tableauHeader{Obligation (résumé du texte officiel)} &
    \tableauHeader{Critère d'audit}
}
26 (1) & Prendre les mesures techniques et administratives de sécurité~;
se doter de systèmes normalisés de gestion continue des risques & Un
processus de gestion des risques existe et fonctionne \\
26 (2) et (3) & Mettre en place des mécanismes de sécurité (disponibilité,
intégrité, authentification, non-répudiation, confidentialité, sécurité
physique), approuvés et visés par l'Agence & Les mécanismes existent~; le
visa de l'Agence est au dossier \\
29 (1) & Conserver les données de connexion et de trafic pendant dix ans &
La politique de conservation couvre dix ans \\
29 (2) & Installer des mécanismes de surveillance du contrôle d'accès aux
données & Les accès aux données sont journalisés et surveillés \\
30 (1) & Évaluer et réviser les systèmes de sécurité selon l'évolution des
technologies & Une revue périodique est tracée \\
\end{tableaucours}
\vspace{-0.8em}
\begin{center}{\small\itshape Tableau~10.2. Obligations de la loi
n°~2010/012 et critères d'audit}\end{center}

\begin{boxlecture}{L'ANPP face à la loi}
\textbf{Constatation.} À partir de la base commune du
Chapitre~\ref{chap:referentiels}, l'équipe de l'ANPP cote les cinq
obligations~:
\begin{itemize}
    \item 26~(1)~: \NON{}~: aucun registre des risques, aucune appréciation
          formalisée~;
    \item 26~(2) et~(3)~: \NEK{}~: des mécanismes existent, mais personne ne
          sait s'ils ont été soumis au visa de l'Agence~;
    \item 29~(1)~: \NEK{}~: les journaux sont «~aux archives~», leur durée
          de conservation n'est pas connue~;
    \item 29~(2)~: \NON{}~: les comptes des agents partis restent actifs, et
          aucune surveillance des accès n'est en place~;
    \item 30~(1)~: \NON{}~: aucune revue de sécurité depuis 2023.
\end{itemize}
0~\OUI{}, 3~\NON{}, 2~\NEK{}~: $\mu = 0$~; $U = 0{,}40$~;
$H = -(0{,}6\log_2 0{,}6 + 0{,}4\log_2 0{,}4) \approx 0{,}97$~bit.

\medskip
\textbf{Analyse} (confiance moyenne). Non-conformité opaque~: aucune
obligation n'est démontrée. Les deux \NEK{} se lèvent à peu de frais (une
demande à l'ANTIC sur le visa, une extraction des archives de journaux)~;
l'ANPP a intérêt à le faire \emph{avant} l'audit, pour que l'Agence
constate des écarts connus et traités plutôt qu'une ignorance.

\medskip
\hfill\lecture{0{,}00}{0{,}40}{0{,}97}
\end{boxlecture}

\subsection{Le secret, et l'obstruction}

\begin{boxjuris}
\textbf{\loicam{n°~2010/012}, article~14}~: \textit{« Le personnel de
l'Agence et les experts commis en vue d'accomplir des opérations d'audit
sont astreints au secret professionnel. »}

\medskip
\textbf{Article~61, alinéa~1}~: \textit{« Sont punis d'un emprisonnement de
trois (03) mois à trois (03) ans et d'une amende de 20.000 (vingt mille) à
100.000 (cent mille) F CFA, les personnels de l'Agence et les experts des
personnes morales chargés des audits qui révèlent sans autorisation, des
informations confidentielles dont ils ont eu connaissance à l'occasion
d'un audit de sécurité. »}

\medskip
\textbf{Article~61, alinéa~3} (paraphrasé)~: la loi punit aussi celui qui
fait obstacle aux audits de sécurité ou refuse de fournir les informations
ou documents qui s'y rapportent.

\medskip
\textbf{Portée.} La loi protège l'audit dans les deux sens~: elle punit
l'auditeur qui divulgue, et l'audité qui fait obstruction. Dans un audit de
sécurité légal, le refus du prestataire des GAB (Chapitre~\ref{chap:audit})
ne serait pas seulement une limitation de l'étendue~: ce serait un fait à
consigner, que l'Agence apprécierait.
\end{boxjuris}

% ============================================================================
\section{La Réglementation Bancaire}
\label{sec:c10-cobac}
% ============================================================================

Les établissements de crédit de la CEMAC sont soumis au règlement COBAC
R-2016/04 relatif au contrôle interne~\cite{cobac2016r04}.

\begin{boxjuris}
\textbf{Règlement COBAC R-2016/04, articles~1 et~2} (texte relevé sur une
base réglementaire en ligne, à collationner avec la version du
Secrétariat général de la COBAC)~:

\medskip
Article~1~: \textit{« Le présent règlement fixe les exigences minimales
relatives au système de contrôle interne que les établissements assujettis
doivent mettre en œuvre. »}

\medskip
Article~2, risque opérationnel~: \textit{« risque de pertes résultant de
l'inadéquation ou de la défaillance des processus internes, des personnes
et des systèmes, ou d'évènements externes. »}

\medskip
\textbf{Portée.} Une attaque par rançongiciel est un risque opérationnel
au sens du règlement~: défaillance des systèmes et des personnes,
provoquée par un événement externe. Le rapport demandé par la COBAC après
l'incident de septembre 2026 relève donc du contrôle interne de la banque,
et non d'une simple information technique.
\end{boxjuris}

Le règlement donne le cadre~; il ne dit pas comment évaluer une mesure de
sécurité. Pour le rapport au superviseur, la banque s'appuie sur ses
propres critères (ISO/IEC~27002, plan de traitement) et sur les résultats
de son audit interne. Les exigences précises du règlement sur les systèmes
d'information, au-delà des deux articles relevés, sont à lire dans le
texte collationné.

% ============================================================================
\section{Le Cadre Africain et les Données Personnelles}
\label{sec:c10-afrique}
% ============================================================================

\subsection{La Convention de Malabo}

La Convention de l'Union africaine sur la cybersécurité et la protection
des données à caractère personnel a été adoptée à Malabo en
2014~\cite{ua2014malabo}. Elle invite les États parties à se doter d'un
cadre national de cybersécurité, d'une protection des données personnelles
et d'une répression de la cybercriminalité. Elle est entrée en vigueur en
2023, après avoir réuni le nombre de ratifications requis. L'état de sa
signature et de sa ratification par le Cameroun est \emph{à vérifier}~;
le cours ne tire donc d'elle aucune obligation directe pour les
organisations camerounaises. Elle reste une référence pour lire la loi
nationale et pour les coopérations régionales.

\subsection{Les données personnelles pendant l'audit}

La loi n°~2024/017 relative à la protection des données à caractère
personnel est étudiée dans l'UE MSN-2S3-JUR01~; son texte reste à vérifier
au Journal officiel. Pour l'auditeur, une règle de prudence suffit~:
lorsqu'il extrait des données personnelles (fichiers du personnel,
bénéficiaires de paiements, clients), il ne prélève que ce que le test
exige, les protège pendant la mission, et les détruit au terme de la
conservation du dossier. L'audit de protection des données lui-même relève
de l'UE MSN-2S4-SEC02.

% ============================================================================
\section{La Déontologie de l'Auditeur}
\label{sec:c10-deontologie}
% ============================================================================

\subsection{Cinq principes}

Les normes mondiales d'audit interne de l'IIA regroupent l'éthique de
l'auditeur en cinq principes~\cite{iia2024gias}~; les codes de l'ISACA et
des organismes de certification en reprennent l'esprit.

\begin{tableaucours}{L{3.2cm}L{9.0cm}}{%
    \tableauHeader{Principe} &
    \tableauHeader{En audit de cybersécurité}
}
Intégrité & Ne rien affirmer que la preuve ne soutient pas~; ne rien
retirer sous la pression \\
Objectivité & Juger sans conflit d'intérêts ni préjugé~; ne pas auditer son
propre travail \\
Compétence & N'accepter que les missions que l'on sait conduire, ou se
faire assister (auditeur technique) \\
Conscience professionnelle & Le soin qu'exigent les enjeux~: tests
suffisants, papiers de travail complets \\
Confidentialité & Protéger les failles constatées et les données
consultées~; ne rien utiliser à des fins personnelles \\
\end{tableaucours}
\vspace{-0.8em}
\begin{center}{\small\itshape Tableau~10.3. Les principes déontologiques
(d'après les normes de l'IIA, 2024)}\end{center}

\subsection{Les menaces sur l'indépendance}

Le cadre conceptuel du code de déontologie de l'IESBA distingue cinq
menaces qui peuvent compromettre l'objectivité d'un
professionnel~\cite{iesba2023code}. Les cas du cours en offrent une
illustration chacun.

\begin{tableaucours}{L{2.6cm}L{4.6cm}L{5.0cm}}{%
    \tableauHeader{Menace} &
    \tableauHeader{Exemple du cours} &
    \tableauHeader{Sauvegarde}
}
Intérêt personnel & AUDIX dépend du contrat de LiZbiZ avec PHARMACAM pour
ses missions futures & Refuser toute modification non justifiée~;
déclarer la dépendance~; revue par un associé \\
Autorévision & Animer l'étude EBIOS que l'on auditera ensuite
(Chapitre~\ref{chap:ebios}) & Refuser l'une des deux missions \\
Représentation & Défendre l'audité devant son client ou son superviseur &
Ne s'exprimer que sur les constats~; renvoyer à l'audité \\
Familiarité & Auditer pendant des années la même équipe, devenue proche &
Rotation des chefs de mission \\
Intimidation & Une direction qui menace de retirer la mission si un
constat est maintenu & Remonter au commanditaire~; maintenir le constat \\
\end{tableaucours}
\vspace{-0.8em}
\begin{center}{\small\itshape Tableau~10.4. Les cinq menaces et leurs
sauvegardes (d'après le code de l'IESBA)}\end{center}

\begin{boxscenario}{Trois lettres, trois réponses}
\textbf{À M.~ETEME.} M.~BEYALA répond que la synthèse reprend les constats
tels qu'ils ont survécu au contradictoire~; il peut en améliorer la
clarté, mais non le sens. Il signale à son associé la dépendance de son
cabinet envers ce contrat (intérêt personnel, intimidation).

\medskip
\textbf{À la DSI de la BANQUE SANAGA.} Le chef de mission de CERTIQUAL
AFRIQUE décline l'invitation, et les règles de son organisme l'y obligent~:
un avantage offert par l'audité, même modeste, crée une apparence de
dépendance. Le refus figure au dossier.

\medskip
\textbf{À M.~ESSOMBA.} Les auditeurs de l'ANTIC sont astreints au secret
(article~14)~; l'ANPP, de son côté, doit fournir les documents demandés
(article~61, alinéa~3). La relation n'est ni hostile ni complaisante~:
elle est encadrée.
\end{boxscenario}

\subsection{Constater n'est pas qualifier}

\begin{boxalert}
\textbf{Violation typique.} Écrire dans un rapport « la DSI a dissimulé
l'absence de journalisation » ou « l'ANPP viole la loi ».

\smallskip
\textbf{Règle~:} l'auditeur constate un écart à un critère, preuve à
l'appui~: « pendant la période auditée, 7 agences sur 38 n'envoyaient pas
de journaux au collecteur~; le rapport de la RSSI de février 2028 les
déclarait raccordées ». L'intention (dissimuler) et la qualification
juridique (violation, infraction) ne se constatent pas~: la première
relève d'une enquête, la seconde du juge ou de l'autorité compétente.
\end{boxalert}

\subsection{Une obligation de moyens}

L'auditeur contractuel s'engage à conduire ses travaux selon les normes
qu'il déclare dans sa lettre de mission. Il n'est pas garant de l'absence
de faille~: son engagement est, en règle générale, une obligation de
moyens. Sa responsabilité peut être recherchée s'il n'a pas appliqué ces
normes. D'où l'importance des papiers de travail (Chapitre~\ref{chap:audit}),
qui montrent ce qu'il a fait, et des \NEK{} documentés, qui montrent ce
qu'il n'a pas pu établir~: ils le protègent autant qu'ils informent le
commanditaire.

% ============================================================================
\section{Ce que la Loi et la Déontologie demandent à chaque Métier}
\label{sec:c10-metiers}
% ============================================================================

\begin{boxmetier}{Auditeur en cybersécurité}
\textbf{Protéger la mission}
\begin{enumerate}
    \item Établir si la mission relève d'un régime légal, et en tirer les
          obligations de forme, de secret et de destinataire.
    \item Ne citer que des textes lus dans leur version officielle.
    \item Identifier, avant la mission, les menaces sur son indépendance,
          et les sauvegardes.
    \item Constater sans qualifier~; maintenir les constats sous la
          pression, et consigner la pression.
\end{enumerate}
\end{boxmetier}

\begin{boxmetier}{RSSI}
\textbf{Préparer l'audit légal}

Établir si l'organisation est assujettie, rassembler les pièces des
articles~26, 29 et~30 (appréciation des risques, mécanismes et visa,
conservation, revues), lever à l'avance les \NEK{} qui peuvent l'être, et
désigner l'interlocuteur des auditeurs~: l'obstruction est une infraction.
\end{boxmetier}

\begin{boxmetier}{Responsable GRC}
\textbf{Un registre des obligations}

Chaque article qui crée une obligation devient une ligne d'un registre des
obligations, reliée aux contrôles de la base commune
(Chapitre~\ref{chap:referentiels}) et à sa preuve. Le
Chapitre~\ref{chap:grc-pratique} montre comment le tenir.
\end{boxmetier}

\begin{boxmetier}{Consultant en mise en œuvre de SMSI}
\textbf{Conseiller sans auditer}

Le consultant qui a construit le SMSI ne participe pas à son audit, ni à
celui de l'organisme certificateur pour le compte de ce dernier. S'il est
sollicité, il le refuse ou le déclare.
\end{boxmetier}

\separateur

\begin{boxresume}
\textbf{Ce qu'il faut retenir de ce chapitre~:}
\begin{itemize}
    \item \textbf{Ne citer que le vérifié}~: chaque texte a son degré de
          vérification, et le rapport le dit.
    \item \textbf{La loi n°~2010/012} institue l'audit de sécurité, confie
          à l'ANTIC le contrôle des audits des organismes publics et des
          autorités de certification, et fournit des critères (articles~26,
          29, 30).
    \item \textbf{Elle protège l'audit dans les deux sens}~: secret de
          l'auditeur (articles~14 et~61), sanction de l'obstruction.
    \item \textbf{Le règlement COBAC R-2016/04} fait du risque informatique
          un risque opérationnel du contrôle interne.
    \item \textbf{La Convention de Malabo} encadre l'action des États~; son
          statut au Cameroun est à vérifier.
    \item \textbf{Cinq principes, cinq menaces}~: intégrité, objectivité,
          compétence, conscience professionnelle, confidentialité~; intérêt
          personnel, autorévision, représentation, familiarité,
          intimidation.
    \item \textbf{Constater n'est pas qualifier.}
\end{itemize}
\end{boxresume}

\begin{boxinfo}
\textbf{Points de contrôle.} Vous avez compris ce chapitre si vous savez~:
\begin{itemize}
    \item dire quel texte s'applique à une situation, et avec quel degré de
          certitude~;
    \item transformer un article de loi en critère d'audit~;
    \item nommer la menace sur l'indépendance dans une situation donnée, et
          la sauvegarde~;
    \item récrire une phrase de rapport qui qualifie au lieu de constater.
\end{itemize}
\end{boxinfo}

% ============================================================================
\section*{Exercices}
% ============================================================================

\begin{boxexo}[title={\ding{45}\quad Exercice~10.1. Quel texte~? \niveaubase}]
Pour chaque situation, indiquez le texte principal applicable et son degré
de vérification dans le cours.
\begin{enumerate}[(a)]
    \item La BANQUE SANAGA doit rendre compte au superviseur de l'incident
          de septembre 2026.
    \item Les journaux du pare-feu sont effacés au bout de trente jours.
    \item Un auditeur de l'ANTIC parle à un journaliste d'une faille de
          l'ANPP.
    \item Le prestataire des GAB refuse de remettre ses configurations lors
          d'un audit de sécurité légal.
    \item L'auditeur extrait le fichier du personnel de l'ANPP pour un
          rapprochement.
    \item Un étudiant affirme que la Convention de Malabo oblige LiZbiZ à
          notifier ses incidents.
\end{enumerate}
\end{boxexo}

\begin{boxexo}[title={\ding{45}\quad Exercice~10.2. Menaces sur l'indépendance \niveaumoyen}]
Nommez la menace et proposez une sauvegarde.
\begin{enumerate}[(a)]
    \item Mme~EKAMBI audite le plan de traitement qu'elle a aidé à rédiger
          en 2026.
    \item Le même chef de mission de CERTIQUAL AFRIQUE audite la banque
          depuis six ans.
    \item Un cabinet d'audit propose à LiZbiZ, pendant sa mission, de
          vendre l'outil qui corrigerait les constats.
    \item Le directeur général d'une agence publique fait savoir au
          magistrat qu'il «~en parlera au ministre~» si un constat est
          maintenu.
    \item Un auditeur témoigne devant la COBAC en faveur de la banque qu'il
          a auditée.
\end{enumerate}
\end{boxexo}

\begin{boxexo}[title={\ding{45}\quad Exercice~10.3. Constater, pas qualifier \niveaumoyen}]
Récrivez chacune des phrases suivantes sous la forme d'un constat,
preuve à l'appui, sans qualification juridique ni jugement d'intention.
\begin{enumerate}[(a)]
    \item « L'ANPP viole l'article~29 de la loi. »
    \item « Le responsable informatique de LiZbiZ a laissé des portes
          ouvertes aux fraudeurs. »
    \item « La banque a menti au certificateur dans sa déclaration
          d'applicabilité. »
    \item « Les quatre agents payés après leur départ ont détourné des
          fonds publics. »
\end{enumerate}
\end{boxexo}

\begin{boxexo}[title={\ding{45}\quad Exercice~10.4. Préparer l'audit de l'ANTIC \niveauexpert}]
M.~ESSOMBA vous confie la préparation de l'audit annoncé.
\begin{enumerate}[(a)]
    \item Rédigez la note qui répond à sa question «~sommes-nous soumis~?~»,
          en distinguant ce que le texte établit et ce qui reste à vérifier.
    \item Proposez, pour chacun des deux \NEK{} du Tableau~10.2 appliqué à
          l'ANPP, l'action qui le lève avant l'audit, et recalculez $\mu$,
          $U$ et $H$ si les deux deviennent \NON{}. La situation est-elle
          meilleure~? Justifiez.
    \item Listez les pièces à remettre aux auditeurs, et les précautions à
          prendre pour les données personnelles qu'elles contiennent.
    \item Que doit faire l'ANPP si un prestataire refuse de communiquer un
          document demandé par l'Agence~?
\end{enumerate}
\end{boxexo}

\bigskip

\begin{boxinfo}
\textbf{Pour aller plus loin.}
\begin{itemize}
    \item La loi n°~2010/012~\cite{loi2010012}, dans son texte officiel, et
          ses textes d'application.
    \item Le règlement COBAC R-2016/04~\cite{cobac2016r04}, dans la version
          du Secrétariat général de la COBAC.
    \item La Convention de Malabo~\cite{ua2014malabo}.
    \item Les normes mondiales d'audit interne~\cite{iia2024gias} et le code
          de l'IESBA~\cite{iesba2023code}.
\end{itemize}
\end{boxinfo}

\begin{boxinfo}
\textbf{Transition.} La Partie~III a posé les normes, les référentiels et
le droit. La Partie~IV passe à la pratique~: mettre en œuvre et auditer un
SMSI sur toute la BANQUE SANAGA, auditer la technique, faire vivre la GRC.
Elle s'ouvre avec le Chapitre~\ref{chap:smsi}.
\end{boxinfo}


    % ======================================================================
    % PARTIE 4 — PRATIQUE : SMSI, AUDIT TECHNIQUE ET GRC
    % Standards de référence : ISO/IEC 27001, 27004, 27007 ; CIS Benchmarks.
    % Niveau : intermédiaire à expert
    % ======================================================================
    \part{Pratique : SMSI, Audit Technique et GRC}
    \label{part:pratique}

        % ============================================================================
% Fichier  : partie4/P04_C11_smsi_mettre_oeuvre_auditer.tex
% Titre    : Mettre en Œuvre et Auditer un SMSI
% Auteur   : MINKA MI NGUIDJOÏ Thierry Emmanuel
% Version  : 1.0 -- novembre 2026
% Corrigés : corriges/Corr_C11.tex
% ============================================================================
% Sources : ISO/IEC 27001:2022 (art. 4.3, 6.2, 7.5, 9.1 à 9.3, 10.2 ;
% informations documentées exigées ; éléments d'entrée de la revue de
% direction : paraphrasés) ; ISO/IEC 27004:2016 (mesures de mise en œuvre et
% d'efficacité) ; ISO/IEC 27007:2020 ; ISO/IEC 17021-1:2015 (non-conformités
% majeure et mineure, définitions paraphrasées). Loi n° 2010/012, art. 26 et
% 30 (résumés, Chapitre 10).
% Faits du cas (bible) : BANQUE SANAGA, extension du SMSI aux agences,
% renouvellement du certificat en juin 2028, audit interne de mars 2028,
% déclaration d'applicabilité corrigée (C08), R-AG-02 (GAB) à apprécier.
% Faits nouveaux (reportés dans la bible) : projet d'extension de mars 2027
% à mai 2028 ; périmètre 2028 (monétique exclue, interfaces incluses) ; six
% indicateurs I1 à I6 (base 4/1/1) ; audit interne du SMSI de mai 2028 ;
% revue de direction du 30 mai 2028 ; étape 2 de CERTIQUAL du 12 au
% 16 juin 2028 : 1 non-conformité majeure (R-AG-02 non apprécié), 3
% mineures, 2 pistes de progrès ; correction acceptée le 25 août 2028 ;
% certificat renouvelé en septembre 2028, avec un trou de trois mois.
% ============================================================================

\chapter{Mettre en Œuvre et Auditer un SMSI}
\label{chap:smsi}
\chaptermark{Mettre en œuvre et auditer un SMSI}

\epigraph{%
    « Un système de management qui ne se mesure pas\\
    ne se gère pas~; qui ne s'audite pas ne se croit pas. »%
}{Principe du cours}

\niveauChapitre{Intermédiaire}{%
    Maîtrise du Chapitre~\ref{chap:iso27000}, du Chapitre~\ref{chap:conclure}%
}

\begin{boxobjectifs}
\begin{itemize}
    \item Conduire la mise en œuvre d'un SMSI comme un projet, de la
          décision de la direction à l'audit de certification.
    \item Définir un périmètre défendable et en justifier les exclusions.
    \item Identifier les informations documentées qu'exige
          l'ISO/IEC~27001.
    \item Construire des indicateurs selon l'ISO/IEC~27004, et coter leur
          fiabilité.
    \item Conduire l'audit interne du SMSI et préparer la revue de
          direction.
    \item Qualifier une non-conformité majeure ou mineure, et en tirer les
          conséquences pour la certification.
\end{itemize}
\end{boxobjectifs}

\bigskip

% ============================================================================
\section*{Avant de commencer : « Majeure ou mineure~? »}
% ============================================================================

Douala, 16~juin 2028, réunion de clôture de l'étape~2 de l'audit de
renouvellement\footnote{Comme dans tout le cours, la BANQUE SANAGA,
CERTIQUAL AFRIQUE, les dates et les montants sont fictifs~; les normes
citées sont exactes et paraphrasées.}. Le chef de mission de CERTIQUAL
AFRIQUE lit ses constats. Au troisième, la RSSI se raidit~: « La banque
n'a pas apprécié le risque des automates de retrait de quatorze agences,
alors que l'audit interne l'avait signalé en avril. Je qualifie cette
non-conformité de \emph{majeure}. »

Le directeur général demande ce que cela change. Tout~: le certificat
de 2025 expire dans dix jours, et l'organisme ne peut pas le renouveler
tant que la non-conformité majeure n'est pas corrigée et sa correction
vérifiée. Pendant ce temps, la banque n'aura plus de certificat.

\medskip

Ce chapitre ouvre la Partie~IV, consacrée à la pratique. Il suit
l'extension du SMSI de la BANQUE SANAGA à toute la banque, de la décision du
conseil à l'audit de certification, et explique pourquoi un seul constat a
suffi à interrompre la certification.

\begin{boxalert}
\textbf{Violation typique.} Traiter l'audit de certification comme un
examen à préparer, en mettant à jour les documents dans les semaines qui
précèdent, au lieu de faire fonctionner le système toute l'année.

\smallskip
\textbf{Règle~:} l'organisme certificateur n'évalue pas des documents~; il
évalue un système qui fonctionne, par ses traces. Une exigence qui n'a pas
de trace sur la période n'est pas satisfaite, même si le document existe.
\end{boxalert}

\separateur

% ============================================================================
\section{Mettre en Œuvre : un Projet, puis un Système}
\label{sec:c11-projet}
% ============================================================================

\subsection{Les étapes}

L'ISO/IEC~27001 décrit des exigences, pas un projet. Dans la pratique, la
mise en œuvre d'un SMSI suit pourtant un ordre presque toujours le même,
parce que chaque étape a besoin de la précédente.

\begin{tableaucours}{cL{4.4cm}L{3.2cm}L{3.6cm}}{%
    \tableauHeader{N°} &
    \tableauHeader{Étape} &
    \tableauHeader{Article} &
    \tableauHeader{BANQUE SANAGA}
}
1 & Décision de la direction, moyens & 5 & Conseil, octobre 2027 \\
2 & Contexte, parties intéressées, périmètre & 4 & Décembre 2026, revu en
2027 \\
3 & Politique et objectifs de sécurité & 5.2, 6.2 & Politique révisée
novembre 2027 \\
4 & Appréciation et traitement des risques, déclaration d'applicabilité &
6.1, 8.2, 8.3 & Décembre 2026 à avril 2028 \\
5 & Compétences, sensibilisation, documentation & 7 & 2027--2028 \\
6 & Exploitation des mesures & 8.1 & 2027--2028 \\
7 & Surveillance, audit interne, revue de direction & 9 & Mars à mai
2028 \\
8 & Non-conformités et améliorations & 10 & Continu \\
\end{tableaucours}
\vspace{-0.8em}
\begin{center}{\small\itshape Tableau~11.1. Les étapes de la mise en œuvre
d'un SMSI}\end{center}

Une fois les étapes franchies, le projet devient un système~: les
étapes~4 à~8 se répètent chaque année, et à chaque changement important.
C'est ce que l'organisme certificateur vient vérifier.

\subsection{Un périmètre défendable}

Le prologue a montré un périmètre de complaisance. L'ISO/IEC~27001 demande
de fixer le périmètre en tenant compte des enjeux, des exigences des
parties intéressées, et des interfaces et dépendances avec ce qui reste en
dehors~\cite{iso27001_2022}.

\begin{boxscenario}{Opération SANAGA. Le périmètre de 2028}
\textbf{Rédaction approuvée par le conseil.} «~Le système de management de
la sécurité de l'information couvre les activités bancaires de la BANQUE
SANAGA~: siège, centre de données principal de Douala, site de secours de
Bonabéri, et 38 agences, avec leurs postes, leurs réseaux, la messagerie et
les accès des prestataires. L'activité monétique est exclue~: elle relève
d'un programme de conformité PCI DSS distinct~; ses interfaces avec le
système central sont incluses.~»

\medskip
\textbf{Analyse} (confiance élevée). L'exclusion de la monétique est
défendable pour trois raisons~: elle est couverte par un autre dispositif,
nommé~; les interfaces, par lesquelles une attaque passerait de l'une à
l'autre, sont dans le périmètre~; et la décision est approuvée par
l'instance qui répond des risques. C'est l'inverse de 2025, où les agences
avaient été exclues sans dispositif de remplacement, sans maîtrise des
interfaces et sans décision du conseil.
\end{boxscenario}

\subsection{Les informations documentées}

\begin{boxprocedure}
\textbf{Ce que la norme exige d'écrire.} Le dossier du SMSI contient au
moins les informations documentées que l'ISO/IEC~27001 exige
expressément~:
\begin{itemize}
    \item le périmètre (4.3) et la politique de sécurité (5.2)~;
    \item le processus d'appréciation des risques (6.1.2) et ses résultats
          (8.2)~; le processus de traitement (6.1.3), la déclaration
          d'applicabilité et le plan de traitement, et ses résultats (8.3)~;
    \item les objectifs de sécurité (6.2)~;
    \item les preuves de compétence (7.2)~;
    \item les informations nécessaires à la maîtrise des processus (8.1)~;
    \item les résultats de la surveillance et de la mesure (9.1)~;
    \item le programme et les résultats des audits internes (9.2)~;
    \item les résultats des revues de direction (9.3)~;
    \item les non-conformités, les actions menées et leurs résultats
          (10.2).
\end{itemize}
Tout le reste est laissé au jugement de l'organisation (7.5)~: écrire ce
qui est nécessaire à l'efficacité, pas davantage.
\end{boxprocedure}

% ============================================================================
\section{Mesurer : les Indicateurs du SMSI}
\label{sec:c11-indicateurs}
% ============================================================================

\subsection{Deux familles de mesures}

L'article~9.1 demande de déterminer ce qui doit être surveillé et mesuré,
comment, quand, par qui, et comment les résultats sont analysés.
L'ISO/IEC~27004 aide à construire ces mesures~\cite{iso27004_2016}. Elle
distingue deux familles~:
\begin{itemize}
    \item les \textbf{mesures de mise en œuvre}~: la mesure est-elle en
          place, et à quel degré~? (part des comptes avec second facteur)~;
    \item les \textbf{mesures d'efficacité}~: la mesure produit-elle
          l'effet attendu~? (part des agents qui cliquent lors d'une
          campagne d'hameçonnage simulé).
\end{itemize}

Un tableau de bord fait seulement de mesures de mise en œuvre dit que l'on
travaille~; il ne dit pas si l'on est protégé.

\subsection{Les six indicateurs de la RSSI}

\begin{tableaucours}{cL{3.8cm}cL{2.8cm}cc}{%
    \tableauHeader{Réf.} &
    \tableauHeader{Indicateur} &
    \tableauHeader{Famille} &
    \tableauHeader{Source} &
    \tableauHeader{Cible} &
    \tableauHeader{Mai 2028}
}
I1 & Comptes avec second facteur & M.~œ. & Extraction de l'annuaire &
100~\% & 100~\% \\
I2 & Délai médian d'application des correctifs critiques & Eff. & Outil
de déploiement & 15~j & 9~j \\
I3 & Restaurations trimestrielles réussies & Eff. & Procès-verbaux & 100~\%
& 100~\% \\
I4 & Agences qui envoient des journaux & M.~œ. & Collecteur, quotidien &
100~\% & 100~\% \\
I5 & Agences ayant exercé le mode dégradé dans l'année & M.~œ. &
Déclarations des agences & 100~\% & 79~\% \\
I6 & Agents qui cliquent lors d'une campagne simulée & Eff. & Plateforme de
campagne & $<$~5~\% & 11~\% \\
\end{tableaucours}
\vspace{-0.8em}
\begin{center}{\small\itshape Tableau~11.2. Le tableau de bord du SMSI
(M.~œ.~: mise en œuvre~; Eff.~: efficacité)}\end{center}

\subsection{Un indicateur vaut ce que vaut sa source}

\begin{boxlecture}{La fiabilité du tableau de bord}
\textbf{Constatation.} Mme~EKAMBI cote la fiabilité de chaque indicateur,
c'est-à-dire la capacité de sa source à mesurer ce que dit sa définition~:
\begin{itemize}
    \item \OUI{}~: I1, I3, I4, I6 (sources automatiques ou procès-verbaux
          signés)~;
    \item \NON{}~: I2~: l'outil de déploiement ne couvre que les serveurs~;
          les postes des agences, mis à jour à part, n'entrent pas dans le
          calcul, alors que la définition parle de «~tous les systèmes~»~;
    \item \NEK{}~: I5~: les agences déclarent leurs exercices, sans compte
          rendu ni chronométrage (cf. Kribi, Chapitre~\ref{chap:conclure}).
\end{itemize}
4~\OUI{}, 1~\NON{}, 1~\NEK{}~: $\mu \approx 0{,}67$~; $U \approx 0{,}17$~;
\[
H = -\left(\tfrac{2}{3}\log_2\tfrac{2}{3} + 2 \times
  \tfrac{1}{6}\log_2\tfrac{1}{6}\right) \approx 0{,}390 + 0{,}862
  \approx 1{,}25~\text{bit}.
\]

\textbf{Analyse} (confiance élevée). Les deux indicateurs les plus
rassurants en apparence posent problème~: I2 est bon parce qu'il mesure un
sous-ensemble~; I5, à 79~\%, repose sur des déclarations. Un indicateur
dont la source est déclarative est un \NEK{}, quelle que soit sa valeur. Le
plus utile du tableau est I6, à 11~\%~: c'est une mauvaise nouvelle, mais
une mesure fiable d'efficacité.

\medskip
\hfill\lecture{0{,}67}{0{,}17}{1{,}25}
\end{boxlecture}

\begin{boxinfo}
\textbf{La fiche d'un indicateur.} Chaque indicateur se décrit par~:
l'objectif de sécurité qu'il sert~; sa définition et sa formule~; sa source
et son mode de collecte~; sa fréquence~; sa cible et ses seuils d'alerte~;
son responsable~; et la cote de fiabilité de sa source. C'est la fiche
d'assertion (Chapitre~\ref{chap:conclure}) appliquée à une mesure.
\end{boxinfo}

% ============================================================================
\section{Auditer le SMSI}
\label{sec:c11-audit}
% ============================================================================

\subsection{Deux objets}

L'audit interne exigé par l'article~9.2 porte sur deux objets~: le
\emph{système de management} (les articles~4 à~10~: le périmètre est-il
tenu, les risques appréciés, les objectifs suivis, les non-conformités
traitées~?) et les \emph{mesures} de la déclaration d'applicabilité.
L'ISO/IEC~27007 donne les lignes directrices de cet audit~\cite{iso27007_2020}.
La mission de mars~2028 sur les agences (Chapitre~\ref{chap:audit}) portait
sur des mesures~; il manquait l'audit du système lui-même.

\begin{boxscenario}{Opération SANAGA. L'audit interne du SMSI, mai 2028}
\textbf{Programme.} L'audit interne examine, pour chaque article de 4 à 10,
une ou deux exigences, avec la preuve de leur fonctionnement sur la
période~: approbation du périmètre (4.3), revue de la politique (5.2),
appréciation des risques à jour (6.1.2, 8.2), suivi des objectifs (6.2),
preuves de compétence (7.2), indicateurs (9.1), traitement des
non-conformités de l'audit de mars (10.2).

\medskip
\textbf{Constat principal.} L'appréciation des risques n'a pas été mise à
jour après l'audit de mars~: le risque R-AG-02 (compromission d'un GAB),
signalé le 20~avril, n'a ni cote, ni propriétaire, ni décision. La RSSI
l'attendait «~pour la prochaine revue annuelle~».

\medskip
\textbf{Analyse} (confiance élevée). Le défaut est de même nature qu'en
2025~: une partie du périmètre comporte un risque connu qui n'a pas été
apprécié. Le rapport de l'audit interne le porte en tête, avec une
échéance avant l'étape~2.
\end{boxscenario}

\subsection{La revue de direction}

L'article~9.3 fait de la revue de direction le moment où la direction
regarde le SMSI dans son ensemble et décide. La norme en fixe les éléments
d'entrée~: l'état des actions décidées lors des revues précédentes~; les
changements des enjeux et des attentes des parties intéressées~; les
retours sur la performance (non-conformités et actions correctives,
résultats de la surveillance, des audits, atteinte des objectifs)~; les
retours des parties intéressées~; les résultats de l'appréciation des
risques et l'état du plan de traitement~; les opportunités
d'amélioration~\cite{iso27001_2022}.

\begin{boxinfo}
\textbf{La revue de direction de la BANQUE SANAGA.} Elle se tient le
30~mai 2028, présidée par M.~FOTSO. Le dossier comprend le tableau de bord
(Tableau~11.2) et le rapport de l'audit de mars. Il ne comprend pas le
rapport de l'audit interne du SMSI, remis le 2~juin. La revue décide des
moyens pour le mode dégradé et la sensibilisation, mais ne dit rien de
R-AG-02~: personne ne le lui a présenté.
\end{boxinfo}

% ============================================================================
\section{L'Audit de Certification et ses Constats}
\label{sec:c11-certification}
% ============================================================================

\subsection{Majeure, mineure}

\begin{boxdef}
\termecle{Non-conformité majeure}\index{non-conformité majeure}~:
non-conformité qui affecte la capacité du système de management à
atteindre les résultats escomptés~; par exemple une exigence qui n'est pas
appliquée du tout, ou un défaut répété, ou un doute sérieux sur la
maîtrise d'un risque~\cite{iso17021_2015}.

\medskip
\termecle{Non-conformité mineure}\index{non-conformité mineure}~:
non-conformité qui n'affecte pas cette capacité~: un défaut ponctuel, isolé,
dans un système qui fonctionne par ailleurs.

\medskip
\termecle{Piste de progrès}\index{piste de progrès}~: constat qui n'est pas
une non-conformité, mais une occasion d'améliorer.
\end{boxdef}

Les conséquences diffèrent. Une non-conformité \emph{mineure} donne lieu à
un plan d'actions que l'organisme accepte, et dont il vérifie l'effet au
prochain audit. Une non-conformité \emph{majeure} suspend la décision de
certification~: l'organisme attend la correction et la preuve de son
efficacité, dans le délai fixé par ses règles, et peut organiser un audit
complémentaire.

\subsection{Les constats de CERTIQUAL AFRIQUE}

\begin{tableaucours}{cL{5.6cm}cL{3.6cm}}{%
    \tableauHeader{N°} &
    \tableauHeader{Constat} &
    \tableauHeader{Exig.} &
    \tableauHeader{Qualification}
}
1 & Le risque des GAB de 14 agences, connu depuis avril, n'a pas été
apprécié ni traité & 6.1.2, 8.2 & \textbf{Majeure}~: exigence non appliquée
à une partie du périmètre, défaut déjà connu \\
2 & Mode dégradé non exercé de façon vérifiable dans 8 agences sur 38 &
Annexe~A, 5.30 & Mineure~: plan d'actions engagé \\
3 & La revue de direction n'a pas examiné les résultats de l'audit
interne & 9.3 & Mineure~: revue tenue, élément d'entrée manquant \\
4 & L'indicateur des correctifs ne couvre pas les postes des agences & 9.1
& Mineure \\
5 & Exercice de réponse aux incidents non réalisé depuis 2026 & Annexe~A,
5.24 & Piste de progrès~: planifié pour juillet \\
6 & La mesure 8.28 réintégrée n'a pas encore de procédure & Annexe~A, 8.28
& Piste de progrès \\
\end{tableaucours}
\vspace{-0.8em}
\begin{center}{\small\itshape Tableau~11.3. Les constats de l'étape~2
(juin 2028)}\end{center}

\begin{boxscenario}{Opération SANAGA. Pourquoi « majeure »}
\textbf{Analyse} (confiance élevée). Le constat~1 n'est pas grave par son
objet seul (quatorze agences, un équipement)~: il l'est parce qu'il montre
que le \emph{processus} d'appréciation des risques ne fonctionne pas comme
la norme l'exige. Un risque connu a attendu deux mois sans décision, et ni
l'audit interne, ni la revue de direction n'ont pu le faire remonter à
temps. C'est exactement le mécanisme de 2025~: le système laisse une zone
sans appréciation. Les constats~2 à~4, eux, sont des défauts dans un
système qui, par ailleurs, les a détectés.
\end{boxscenario}

\subsection{Ce qui suit}

La banque apprécie R-AG-02 en juillet, obtient du prestataire des GAB un
avenant de droit d'audit, et fait vérifier la correction par l'organisme.
CERTIQUAL AFRIQUE accepte la correction le 25~août et renouvelle le
certificat en septembre 2028. Entre la fin du certificat de 2025 et le
nouveau, la banque est restée trois mois sans certificat~: il a fallu
l'expliquer à deux correspondants étrangers et à un client institutionnel.

\begin{boxobjet}{R-AG-02 · Compromission d'un GAB (appréciée en juillet 2028)}
\textbf{Propriétaire~:} le directeur du réseau d'agences.

\textbf{Vraisemblance~:} 3 (système sans correctifs, exposé dans les
agences). Base~: configurations obtenues grâce à l'avenant (\OUI{}).

\textbf{Impact~:} 3 (retraits frauduleux plafonnés par les encaisses~;
indisponibilité). Base~: \OUI{}.

\textbf{Criticité~:} 9, intolérable~: remplacement du système des 14 GAB
programmé avant la fin de 2028~; en attendant, cloisonnement des GAB sur un
réseau dédié et surveillance des retraits atypiques.
\end{boxobjet}

\begin{boxinfo}
\textbf{La norme et la loi.} L'article~30 de la loi n°~2010/012 demande
aux exploitants d'évaluer et de réviser leurs systèmes de sécurité selon
l'évolution des technologies (Chapitre~\ref{chap:droit})~\cite{loi2010012}.
La revue de direction et l'audit interne du SMSI en sont la forme
normalisée~: leurs comptes rendus sont les pièces qu'un audit réglementaire
demandera.
\end{boxinfo}

% ============================================================================
\section{Ce que le SMSI demande à chaque Métier}
\label{sec:c11-metiers}
% ============================================================================

\begin{boxmetier}{Consultant en mise en œuvre de SMSI}
\textbf{Faire fonctionner, pas seulement documenter}
\begin{enumerate}
    \item Conduire la mise en œuvre dans l'ordre des étapes, en commençant
          par la décision de la direction et le périmètre.
    \item Rédiger un périmètre qui nomme ses exclusions et inclut ses
          interfaces.
    \item Mettre en place, dès le début, l'appréciation des risques en
          continu~: un risque nouveau est apprécié quand il apparaît, pas à
          la revue annuelle.
    \item Laisser au moins un cycle complet (audit interne, revue de
          direction, actions) avant l'audit de certification.
\end{enumerate}
\end{boxmetier}

\begin{boxmetier}{Auditeur en cybersécurité}
\textbf{Qualifier par le système, pas par l'objet}

Une non-conformité est majeure quand elle révèle que le système ne
fonctionne pas, pas quand son objet impressionne. L'auditeur justifie sa
qualification par l'exigence concernée et par ce que le défaut dit du
processus.
\end{boxmetier}

\begin{boxmetier}{RSSI}
\textbf{Un indicateur fiable vaut mieux qu'un bon indicateur}

La RSSI corrige I2 (inclure les postes), remplace la source de I5 par des
comptes rendus chronométrés, et présente I6 à la direction sans
l'enjoliver~: 11~\% de clics est un chiffre utile parce qu'il est vrai.
\end{boxmetier}

\begin{boxmetier}{Responsable GRC}
\textbf{Relier les constats d'audit au registre}

Tout constat d'audit qui révèle un risque nouveau entre dans le registre
dans la semaine, avec un propriétaire et une échéance d'appréciation. Le
défaut de la banque en 2028 tenait à un délai, pas à une ignorance.
\end{boxmetier}

\separateur

\begin{boxresume}
\textbf{Ce qu'il faut retenir de ce chapitre~:}
\begin{itemize}
    \item \textbf{Un SMSI} se met en œuvre comme un projet, puis vit comme
          un système dont les étapes se répètent.
    \item \textbf{Un périmètre défendable} nomme ses exclusions, inclut ses
          interfaces et est approuvé par l'instance qui répond des risques.
    \item \textbf{Les indicateurs} sont de mise en œuvre ou d'efficacité~;
          un indicateur vaut ce que vaut sa source.
    \item \textbf{L'audit interne} porte sur le système (articles~4 à~10) et
          sur les mesures~; la revue de direction décide à partir de ses
          résultats.
    \item \textbf{Majeure ou mineure}~: la qualification dépend de ce que
          le défaut dit de la capacité du système, non de son objet.
    \item \textbf{Une majeure} suspend la certification~: à la BANQUE
          SANAGA, trois mois sans certificat.
\end{itemize}
\end{boxresume}

\begin{boxinfo}
\textbf{Points de contrôle.} Vous avez compris ce chapitre si vous savez~:
\begin{itemize}
    \item rédiger un périmètre avec une exclusion défendable~;
    \item distinguer un indicateur de mise en œuvre d'un indicateur
          d'efficacité, et coter sa fiabilité~;
    \item composer l'ordre du jour d'une revue de direction~;
    \item qualifier un constat et justifier la qualification.
\end{itemize}
\end{boxinfo}

% ============================================================================
\section*{Exercices}
% ============================================================================

\begin{boxexo}[title={\ding{45}\quad Exercice~11.1. Exigé ou non~? \niveaubase}]
Dites si chacun des documents suivants est une information documentée
expressément exigée par l'ISO/IEC~27001, et citez l'article le cas échéant.
\begin{enumerate}[(a)]
    \item La déclaration d'applicabilité.
    \item Une charte d'usage de la messagerie.
    \item Le compte rendu de la revue de direction.
    \item Les attestations de formation des administrateurs.
    \item Un schéma de l'architecture réseau.
    \item La liste des non-conformités et des actions menées.
\end{enumerate}
\end{boxexo}

\begin{boxexo}[title={\ding{45}\quad Exercice~11.2. Construire un indicateur \niveaumoyen}]
\begin{enumerate}[(a)]
    \item Classez en mise en œuvre ou efficacité~: part des agences dont le
          partage est cloisonné~; nombre de tentatives d'accès entre agences
          bloquées par mois~; délai moyen de détection d'un incident~; part
          des contrats de prestataires comportant une clause de sécurité.
    \item Rédigez la fiche complète d'un indicateur d'efficacité pour la
          mesure 5.30 (mode dégradé), avec sa source et sa cote de
          fiabilité.
    \item Corrigez la définition ou la source de I2 pour qu'il devienne
          \OUI{}.
\end{enumerate}
\end{boxexo}

\begin{boxexo}[title={\ding{45}\quad Exercice~11.3. Qualifier \niveaumoyen}]
Qualifiez chacun des constats suivants (majeure, mineure, piste de
progrès) et justifiez par ce qu'il dit du système.
\begin{enumerate}[(a)]
    \item Aucun audit interne du SMSI n'a été conduit depuis la
          certification initiale.
    \item Sur 25 comptes de partants prélevés, un est resté actif huit jours
          au lieu de cinq~; les 24 autres sont conformes.
    \item La politique de sécurité n'a pas été revue depuis trois ans,
          alors que la procédure prévoit une revue annuelle~; elle reste
          pertinente.
    \item Les objectifs de sécurité de 2028 ne sont pas mesurables
          («~renforcer la sécurité~»).
    \item Le tableau de bord pourrait comporter un indicateur sur les
          fournisseurs.
\end{enumerate}
\end{boxexo}

\begin{boxexo}[title={\ding{45}\quad Exercice~11.4. Le SMSI de LiZbiZ \niveauexpert}]
LiZbiZ décide d'engager un SMSI en vue d'une certification fin 2029
(Chapitre~\ref{chap:iso27000}, Exercice~8.4).
\begin{enumerate}[(a)]
    \item Établissez le calendrier des huit étapes du Tableau~11.1, en
          laissant un cycle complet avant l'audit initial.
    \item Proposez quatre indicateurs (deux de mise en œuvre, deux
          d'efficacité) qui couvrent le scénario SO1, avec leur source.
    \item LiZbiZ n'a pas d'audit interne. Qui peut conduire l'audit interne
          du SMSI, et à quelles conditions d'indépendance
          (Chapitre~\ref{chap:droit})~?
    \item Rédigez l'ordre du jour de la première revue de direction, en
          couvrant tous les éléments d'entrée exigés.
\end{enumerate}
\end{boxexo}

\bigskip

\begin{boxinfo}
\textbf{Pour aller plus loin.}
\begin{itemize}
    \item L'ISO/IEC~27001:2022~\cite{iso27001_2022}, articles~4.3, 9 et~10.
    \item L'ISO/IEC~27004:2016~\cite{iso27004_2016}, sur les mesures.
    \item L'ISO/IEC~27007:2020~\cite{iso27007_2020}, sur l'audit d'un SMSI.
    \item L'ISO/IEC~17021-1:2015~\cite{iso17021_2015}, sur les exigences
          applicables aux organismes de certification.
\end{itemize}
\end{boxinfo}

\begin{boxinfo}
\textbf{Transition.} L'audit d'un SMSI vérifie le système et ses mesures,
souvent à partir de documents et d'extractions. Il faut parfois aller plus
loin, jusqu'à la configuration des serveurs, la segmentation du réseau,
les comptes et les vulnérabilités. C'est l'audit technique, objet du
Chapitre~\ref{chap:audit-technique}.
\end{boxinfo}

        % ============================================================================
% Fichier  : partie4/P04_C12_audit_technique.tex
% Titre    : L'Audit Technique de Sécurité
% Auteur   : MINKA MI NGUIDJOÏ Thierry Emmanuel
% Version  : 1.0 -- novembre 2026
% Corrigés : corriges/Corr_C12.tex
% ============================================================================
% Faits du corpus repris SANS MODIFICATION (ACSI, chapitre 16) : quatre
% centres de LiZbiZ (Douala, Yaoundé, Garoua, Bafoussam) ; sous-domaines
% IN.1 à IN.5 et leurs contrôles clés ; cotes de 2026 par centre (Douala
% 4/1/0, Yaoundé 4/1/0, Garoua 0/4/1, Bafoussam 0/3/2 ; ensemble 8/9/3,
% mu 0,40, U 0,15, H 1,46) ; Garoua et Bafoussam installés par un
% prestataire ; Garoua hébergé par un prestataire local ; loi n° 2010/012,
% art. 69 (cité, texte officiel) et 70 (résumé).
% Sources : CIS Benchmarks (profils de niveau 1 et 2) ; CVSS v4.0 (FIRST,
% 2023) ; catalogue des vulnérabilités exploitées de la CISA (présenté sans
% citation) ; NIST SP 800-53A (méthodes).
% Faits nouveaux (reportés dans la bible) : second passage d'AUDIX en
% octobre 2028, sur exports figés remis le 2 octobre (empreintes SHA-256),
% après le plan B1 à B6 ; cotes 2028 : Douala 5/0/0, Yaoundé 4/1/0, Garoua
% 3/1/1, Bafoussam 2/2/1 ; ensemble 14/4/2 ; Telnet actif sur deux
% commutateurs à Bafoussam ; règle « tout vers tout » à Garoua ; six
% vulnérabilités V1 à V6.
% ============================================================================

\chapter{L'Audit Technique de Sécurité}
\label{chap:audit-technique}
\chaptermark{L'audit technique}

\epigraph{%
    « La politique dit ce qui est permis.\\
    La configuration dit ce qui est possible. »%
}{Principe du cours}

\niveauChapitre{Intermédiaire}{%
    Maîtrise du Chapitre~\ref{chap:smsi}%
}

\begin{boxobjectifs}
\begin{itemize}
    \item Distinguer audit de configuration, scan de vulnérabilités et test
          d'intrusion, et encadrer juridiquement un test technique.
    \item Travailler sur des exports figés et en garantir l'intégrité et
          l'exhaustivité.
    \item Auditer une configuration au regard d'une référence (CIS
          Benchmarks, configuration interne).
    \item Auditer la segmentation par la comparaison des flux autorisés et
          observés.
    \item Auditer les comptes et la journalisation.
    \item Auditer la gestion des vulnérabilités et prioriser leur
          correction.
    \item Comparer deux passages d'audit par centre, sans moyenne
          trompeuse.
\end{itemize}
\end{boxobjectifs}

\bigskip

% ============================================================================
\section*{Avant de commencer : ce qui est écrit est-il configuré~?}
% ============================================================================

Douala, octobre 2028. Deux ans plus tôt, AUDIX CONSEIL avait audité
l'infrastructure des quatre centres de données de LiZbiZ\footnote{LiZbiZ,
AUDIX CONSEIL, PHARMACAM, les dates et les montants sont fictifs~; les
référentiels et les textes cités sont exacts.}~: la politique interdisait
les protocoles d'administration en clair, imposait le changement des mots de
passe par défaut et séparait les postes des serveurs~; deux centres sur
quatre ne l'appliquaient pratiquement pas. Depuis, LiZbiZ a conduit le plan
de traitement de son scénario de rançongiciel (Chapitre~\ref{chap:iso27005}),
et l'audit contractuel de PHARMACAM a encore relevé un filtrage défaillant
entre centres (Chapitre~\ref{chap:audit}).

M.~ETEME demande un second passage. M.~BEYALA accepte~: auditer n'est pas
conseiller, et son cabinet n'a pas conduit le plan. Mais il pose deux
conditions~: travailler sur des exports figés des configurations, et
obtenir une autorisation écrite avant tout test sur les réseaux.

\medskip

Ce chapitre présente l'audit technique de sécurité~: ce qu'il vérifie, ce
qui le distingue d'un test d'intrusion, et comment il produit des preuves
qui tiennent.

\begin{boxalert}
\textbf{Violation typique.} Conclure qu'un centre est « durci » parce que
la politique de durcissement existe et que l'équipe déclare l'avoir
appliquée.

\smallskip
\textbf{Règle~:} seule la configuration réelle, extraite et figée, établit
l'état d'un équipement. L'auditeur la compare à une référence écrite~;
l'écart est le constat.
\end{boxalert}

\separateur

% ============================================================================
\section{Trois Activités Techniques}
\label{sec:c12-activites}
% ============================================================================

\begin{tableaucours}{L{2.6cm}L{3.4cm}L{3.0cm}L{3.2cm}}{%
    \tableauHeader{Activité} &
    \tableauHeader{Question} &
    \tableauHeader{Méthode} &
    \tableauHeader{Limite}
}
Audit de configuration & Les équipements sont-ils configurés selon la
référence~? & Lecture d'exports, comparaison & Ne dit pas ce qu'un
attaquant ferait des écarts \\
Scan de vulnérabilités & Quelles faiblesses connues présentent les
systèmes~? & Outil, de préférence authentifié & Faux positifs~; ne voit
que le connu \\
Test d'intrusion & Un attaquant peut-il atteindre tel objectif~? &
Exploitation contrôlée & Un chemin trouvé ne dit rien des autres \\
\end{tableaucours}
\vspace{-0.8em}
\begin{center}{\small\itshape Tableau~12.1. Trois activités techniques}\end{center}

L'audit technique combine en général les deux premières~; il peut
s'appuyer sur un test d'intrusion, conduit par d'autres (UE MSN-2S3-SEC03),
comme preuve d'efficacité. Ce qui le distingue d'un test d'intrusion, c'est
son objet~: l'audit vérifie l'\emph{ensemble} des critères d'un périmètre et
conclut sur leur conformité~; le test d'intrusion cherche \emph{un} chemin
vers un objectif et conclut sur son existence.

\subsection{Le droit de tester}

\begin{boxjuris}
\textbf{\loicam{n°~2010/012}, article~69}~\cite{loi2010012}~: \textit{« Est
puni d'un emprisonnement de cinq (05) à dix (10) ans et d'une amende de
10.000.000 (dix millions) à 100.000.000 (cent millions) F CFA ou de l'une
de ces deux peines seulement, celui qui accède sans droit, et en violation
des mesures de sécurité, à l'ensemble ou à une partie d'un réseau de
communications électroniques, d'un système d'information ou d'un équipement
terminal, afin d'obtenir des informations ou des données, en relation avec
un système d'information connecté à un autre système d'information. »}

\medskip
\textbf{Article~70} (résumé)~: est puni celui qui provoque, par saturation,
l'attaque d'une ressource de réseau ou d'un système d'information, dans le
but d'empêcher la réalisation des services attendus.

\medskip
\textbf{Portée.} L'autorisation écrite du responsable du système donne à
l'auditeur le droit d'accéder~: sans elle, un scan ou une tentative de
connexion peut tomber sous l'article~69. Les tests de charge et de déni de
service, qui pourraient produire l'effet visé par l'article~70, sont exclus,
sauf autorisation expresse et environnement dédié. Pour le centre de
Garoua, hébergé par un prestataire local, l'accord écrit de l'hébergeur
s'ajoute à celui de LiZbiZ.
\end{boxjuris}

\begin{boxprocedure}
\textbf{La note d'autorisation d'un test technique.} Elle précise~: le
signataire habilité (et l'hébergeur, pour un système hébergé)~; le
périmètre (adresses, systèmes inclus et exclus)~; les méthodes permises et
interdites~; les fenêtres de test~; un contact joignable et une procédure
d'arrêt immédiat~; la confidentialité, la conservation et la destruction
des résultats (articles~14 et~61, Chapitre~\ref{chap:droit}).
\end{boxprocedure}

% ============================================================================
\section{Travailler sur des Exports Figés}
\label{sec:c12-exports}
% ============================================================================

\subsection{Pourquoi figer}

L'auditeur ne travaille pas sur les équipements en production~: il
travaille sur des \termecle{exports figés}\index{export figé}, c'est-à-dire
des copies des configurations, des listes de comptes et des règles de
filtrage, extraites à une date donnée, sous ses yeux ou selon une procédure
écrite, et scellées par une empreinte.

\begin{itemize}
    \item \textbf{Sécurité}~: l'audit ne risque pas de perturber la
          production.
    \item \textbf{Reproductibilité}~: un autre auditeur peut refaire
          l'analyse sur les mêmes fichiers et vérifier les constats.
    \item \textbf{Opposabilité}~: l'empreinte établit que le fichier analysé
          est celui qui a été extrait, à la date dite.
\end{itemize}

\begin{boxprocedure}
\textbf{Le dossier d'exports.} Pour chaque export, le papier de travail
indique~: l'équipement (nom, adresse, rôle, centre)~; la commande ou la
fonction d'export employée~; la date et l'heure~; la personne qui a extrait
et le témoin~; l'empreinte cryptographique (par exemple SHA-256) calculée
au moment de la remise~; le lieu de conservation. L'auditeur recalcule
l'empreinte avant chaque analyse.
\end{boxprocedure}

\subsection{L'exhaustivité}

Un export figé est fiable sur ce qu'il contient~; il ne dit rien de ce qui
manque. L'auditeur rapproche donc la liste des exports de l'inventaire des
équipements. Un équipement présent dans l'inventaire et absent des
exports est un \NEK{}~; un équipement présent dans les exports et absent
de l'inventaire est un constat sur l'inventaire.

\begin{boxscenario}{Opération LIZBIZ. La remise du 2 octobre 2028}
\textbf{Faits.} La DSI remet 61 exports (pare-feu, commutateurs, serveurs,
annuaire), avec leurs empreintes SHA-256. L'inventaire compte 64
équipements. Les trois manquants sont les serveurs de Garoua gérés par
l'hébergeur, qui n'a pas encore fourni ses exports, et un commutateur de
Bafoussam «~en panne~».

\medskip
\textbf{Analyse} (confiance élevée). Les constats sur 61 équipements
peuvent être établis~; pour les trois autres, l'audit cotera \NEK{} et
le dira. Le commutateur «~en panne~» mérite une vérification sur place~: un
équipement que l'on ne peut pas exporter mais qui reste branché est un
équipement que personne ne surveille.
\end{boxscenario}

% ============================================================================
\section{Auditer la Configuration}
\label{sec:c12-configuration}
% ============================================================================

\subsection{Une référence écrite}

Auditer une configuration, c'est la comparer à une référence. Deux sources
sont courantes~:
\begin{itemize}
    \item les \textbf{CIS Benchmarks}, guides de configuration élaborés par
          consensus pour des centaines de produits (systèmes
          d'exploitation, équipements réseau, bases de données), avec deux
          profils~: le niveau~1, applicable presque partout, et le
          niveau~2, plus strict~\cite{cis2026benchmarks}~;
    \item la \textbf{configuration de référence interne}, qui adapte ces
          guides au contexte et que la politique de l'organisation rend
          obligatoire.
\end{itemize}
La seconde est le critère opposable~; la première sert à la construire, et
à juger si elle est suffisante.

\subsection{Lire un export}

\begin{tableaucours}{L{4.0cm}L{4.0cm}L{2.2cm}c}{%
    \tableauHeader{Ligne de l'export (simplifiée)} &
    \tableauHeader{Référence} &
    \tableauHeader{Sous-domaine} &
    \tableauHeader{Cote}
}
\texttt{line vty 0 4 / transport input telnet ssh} & Administration en SSH
seulement & IN.3 & \NON{} \\
\texttt{snmp-server community public RO} & Aucune communauté par défaut &
IN.2 & \NON{} \\
\texttt{username admin secret 9 \$9\$\ldots} & Compte nominatif, mot de
passe haché & IN.2 & \OUI{} \\
\texttt{logging host 10.20.0.15} & Journaux vers le collecteur central &
IN.5 & \OUI{} \\
\texttt{ntp server 10.20.0.1} & Horloge synchronisée & IN.5 & \OUI{} \\
\end{tableaucours}
\vspace{-0.8em}
\begin{center}{\small\itshape Tableau~12.2. Extrait de l'export d'un
commutateur de Bafoussam}\end{center}

Deux lignes suffisent à établir deux non-conformités~: Telnet reste
accepté pour l'administration, et une communauté SNMP par défaut permet de
lire la configuration à quiconque atteint l'équipement. La synchronisation
de l'horloge, souvent négligée, est une condition de la journalisation~: des
journaux mal datés ne permettent pas de reconstituer un incident.

% ============================================================================
\section{Auditer la Segmentation}
\label{sec:c12-segmentation}
% ============================================================================

La segmentation se vérifie en comparant deux matrices~: celle des flux que
la politique autorise entre zones, et celle des flux que les règles des
pare-feu laissent réellement passer. L'écart est le constat.

\begin{tableaucours}{L{3.0cm}cccc}{%
    \tableauHeader{De \textbackslash{} vers} &
    \tableauHeader{DMZ} & \tableauHeader{Applicatif} &
    \tableauHeader{Utilisateurs} & \tableauHeader{Autres centres}
}
Internet & 443 & -- & -- & -- \\
DMZ & -- & 8443 & -- & -- \\
Utilisateurs & 443 & 8443 & -- & -- \\
Administration & 22 & 22 & 22 & 22 \\
Applicatif & -- & -- & -- & 1433 (réplication) \\
\textbf{Maintenance éditeur} & -- & \textbf{toutes (observé)} & -- &
\textbf{toutes (observé)} \\
\end{tableaucours}
\vspace{-0.8em}
\begin{center}{\small\itshape Tableau~12.3. Matrice de flux du centre de
Garoua~: ports autorisés, et en gras les flux observés non autorisés}\end{center}

\begin{boxscenario}{Opération LIZBIZ. Une règle de trop}
\textbf{Constatation.} L'export du pare-feu de Garoua contient une règle
qui autorise tout trafic depuis l'adresse de maintenance de l'éditeur vers
toutes les zones du centre et vers les autres centres. Elle a été créée en
2027, «~temporairement~», lors d'une mise à jour de l'application.

\medskip
\textbf{Analyse} (confiance élevée). La règle annule la segmentation pour
l'accès qui sert de point d'entrée au scénario SO1 (Chapitre~\ref{chap:ebios})~:
un attaquant qui compromettrait l'accès de maintenance atteindrait d'emblée
les quatre centres. C'est la cause du \NON{} relevé par PHARMACAM sur le
filtrage entre centres. Une règle «~temporaire~» sans date d'expiration est
une règle permanente.
\end{boxscenario}

% ============================================================================
\section{Auditer les Comptes et la Journalisation}
\label{sec:c12-comptes}
% ============================================================================

\subsection{Les comptes}

L'audit des comptes rapproche trois listes~: les comptes présents dans les
exports (annuaire, équipements, bases de données), les personnes et les
services autorisés, et les dernières connexions. Il recherche quatre
catégories~:
\begin{itemize}
    \item les \textbf{comptes par défaut} encore actifs (IN.2)~;
    \item les \textbf{comptes orphelins}, sans titulaire (départs)~;
    \item les \textbf{comptes dormants}, sans connexion depuis longtemps~;
    \item les \textbf{comptes partagés} ou génériques, qui empêchent
          d'imputer une action.
\end{itemize}
Les comptes de service, utilisés par les applications, méritent une
attention particulière~: ils ont souvent des droits étendus, des mots de
passe qui ne changent jamais, et personne n'en est propriétaire.

\subsection{La journalisation}

Pour la journalisation (IN.5), l'auditeur vérifie quatre choses~: que
chaque équipement envoie ses journaux au collecteur (rapprochement avec
l'inventaire)~; que les horloges sont synchronisées~; que la durée de
conservation respecte la politique et la loi (article~29,
Chapitre~\ref{chap:droit})~; et que quelqu'un lit ou exploite les journaux,
par des alertes ou une revue. Des journaux collectés et jamais lus ne
détectent rien.

% ============================================================================
\section{Auditer la Gestion des Vulnérabilités}
\label{sec:c12-vulnerabilites}
% ============================================================================

\subsection{Le processus avant le scan}

L'auditeur n'audite pas seulement le résultat d'un scan~; il audite le
processus de gestion des vulnérabilités (ISO/IEC~27002, mesure~8.8)~: les
sources de veille, la fréquence des scans, le délai de correction fixé par
la politique selon la gravité, et le suivi des exceptions. Un scan
\emph{authentifié}, qui se connecte aux systèmes avec un compte de lecture,
voit les versions réellement installées~; un scan non authentifié devine à
partir des services exposés, avec plus de faux positifs.

\subsection{Prioriser}

Le système de cotation CVSS, dans sa version~4.0, attribue à chaque
vulnérabilité un score de gravité de 0 à~10~\cite{first2023cvss4}. Le score
seul ne suffit pas à prioriser~: il décrit la vulnérabilité, pas
l'organisation. Trois questions s'ajoutent~: la vulnérabilité est-elle
exploitée activement (l'agence américaine CISA publie un catalogue des
vulnérabilités exploitées connues)~? Le système est-il exposé (Internet,
accès de maintenance)~? Que porte-t-il (valeur métier, Chapitre~\ref{chap:risque})~?

\begin{tableaucours}{cL{3.6cm}ccL{3.4cm}}{%
    \tableauHeader{Réf.} &
    \tableauHeader{Système} &
    \tableauHeader{CVSS} &
    \tableauHeader{Exploitée~?} &
    \tableauHeader{Exposition}
}
V1 & Serveur de maintenance à distance (Garoua) & 8,1 & Oui & Accès de
l'éditeur \\
V2 & Base de données de l'application (Douala) & 9,0 & Non & Interne,
zone applicative \\
V3 & Portail de suivi des transporteurs & 7,5 & Oui & Internet \\
V4 & Imprimantes des agences & 9,8 & Non & Réseau des agences \\
V5 & Serveur de fichiers de Yaoundé & 6,5 & Oui & Interne \\
V6 & Commutateurs de Bafoussam (Telnet) & 7,0 & Non & Réseau
d'administration \\
\end{tableaucours}
\vspace{-0.8em}
\begin{center}{\small\itshape Tableau~12.4. Six vulnérabilités relevées
(scores et indications d'exploitation fictifs)}\end{center}

Trié par score CVSS, le tableau place V4 (des imprimantes) en tête. Trié
par risque, il place V1 et V3 en tête~: exploitées activement, exposées,
et sur le chemin du scénario SO1 pour V1. V4 vient bien plus loin~: une
imprimante d'agence n'est pas une valeur métier, et elle n'est pas
exposée à Internet.

% ============================================================================
\section{Deux Passages, Quatre Centres}
\label{sec:c12-passages}
% ============================================================================

\begin{boxlecture}{L'infrastructure de LiZbiZ, 2026 et 2028}
\begin{tableaucours}{L{2.4cm}ccccccc}{%
    \tableauHeader{Centre} &
    \tableauHeader{IN.1} & \tableauHeader{IN.2} & \tableauHeader{IN.3} &
    \tableauHeader{IN.4} & \tableauHeader{IN.5} &
    \tableauHeader{2028} & \tableauHeader{2026}
}
Douala & \OUI{} & \OUI{} & \OUI{} & \OUI{} & \OUI{} & 5/0/0 & 4/1/0 \\
Yaoundé & \OUI{} & \OUI{} & \OUI{} & \NON{} & \OUI{} & 4/1/0 & 4/1/0 \\
Garoua & \NON{} & \OUI{} & \OUI{} & \NEK{} & \OUI{} & 3/1/1 & 0/4/1 \\
Bafoussam & \NON{} & \OUI{} & \NON{} & \OUI{} & \NEK{} & 2/2/1 & 0/3/2 \\
\textbf{Ensemble} & & & & & & \textbf{14/4/2} & \textbf{8/9/3} \\
\end{tableaucours}

\medskip
\textbf{Constatation.} En 2028~: $\mu = 0{,}70$~; $U = 0{,}10$~;
\[
H = -(0{,}7\log_2 0{,}7 + 0{,}2\log_2 0{,}2 + 0{,}1\log_2 0{,}1)
  \approx 0{,}360 + 0{,}464 + 0{,}332 \approx 1{,}16~\text{bit}.
\]
En 2026~: $\mu = 0{,}40$~; $U = 0{,}15$~; $H \approx 1{,}46$~bit.

\medskip
\textbf{Analyse} (confiance élevée). L'ensemble passe de la non-conformité
claire à la conformité claire, et l'incertitude baisse. Mais la moyenne
cache deux choses~: Yaoundé régresse sur les correctifs (deux correctifs
critiques en retard), et Garoua comme Bafoussam gardent un défaut de
segmentation, le même que celui qui ouvre le scénario SO1. Les \NEK{}
viennent des exports manquants (hébergeur de Garoua, commutateur de
Bafoussam). Le rapport conclut centre par centre, et place la règle de
maintenance de Garoua en tête des recommandations.

\medskip
\hfill\lecture{0{,}70}{0{,}10}{1{,}16}
\end{boxlecture}

\begin{boxobjet}{SO1 · Rançongiciel par l'accès de maintenance (octobre 2028)}
\textbf{Phase « rentrer »~:} second facteur de l'éditeur prouvé (avenant
signé en 2028)~; mais V1, vulnérabilité exploitée, sur le serveur de
maintenance de Garoua.

\textbf{Phase « trouver »~:} la règle «~tout vers tout~» depuis l'adresse de
maintenance annule la segmentation~; Telnet persiste sur deux commutateurs
de Bafoussam.

\textbf{Phase « exploiter »~:} sauvegardes hors ligne restaurées
(Chapitre~\ref{chap:audit}).

\textbf{Cotation~:} vraisemblance V3 maintenue, tant que V1 et la règle de
Garoua ne sont pas corrigées~; gravité ramenée de G4 à G3 par la
récupération prouvée~; criticité $3 \times 3 = 9$, toujours intolérable.

\textbf{Priorité~:} corriger V1 et supprimer la règle de maintenance de
Garoua sous huit jours, puis vérifier par un nouvel export.
\end{boxobjet}

% ============================================================================
\section{Ce que l'Audit Technique demande à chaque Métier}
\label{sec:c12-metiers}
% ============================================================================

\begin{boxmetier}{Auditeur en cybersécurité}
\textbf{Lire un export}
\begin{enumerate}
    \item Vérifier l'empreinte, la date et l'équipement de chaque export.
    \item Rapprocher les exports de l'inventaire, et coter \NEK{} ce qui
          manque.
    \item Comparer chaque configuration à la référence écrite, ligne par
          ligne.
    \item Comparer les flux observés aux flux autorisés.
    \item Prioriser les vulnérabilités par le risque, pas par le seul score.
    \item Conclure par centre, avant toute agrégation.
\end{enumerate}
\end{boxmetier}

\begin{boxmetier}{RSSI}
\textbf{Préparer l'audit technique}

Fournir la matrice des flux autorisés et la configuration de référence~;
faire signer l'autorisation, avec l'accord des hébergeurs~; prévenir les
équipes de surveillance~; produire les exports avec leur empreinte~;
supprimer les règles « temporaires » à leur échéance.
\end{boxmetier}

\begin{boxmetier}{Gestionnaire des risques cyber}
\textbf{Relier les constats aux scénarios}

Chaque constat technique se rattache à une phase d'un scénario. La règle
de Garoua n'est pas un défaut de configuration parmi d'autres~: elle annule
la phase « trouver » de SO1. C'est ce rattachement qui fixe la priorité.
\end{boxmetier}

\separateur

\begin{boxresume}
\textbf{Ce qu'il faut retenir de ce chapitre~:}
\begin{itemize}
    \item \textbf{Audit de configuration, scan, test d'intrusion}~: trois
          activités, trois questions~; l'audit conclut sur un ensemble de
          critères.
    \item \textbf{Le droit de tester} vient d'une autorisation écrite~; sans
          elle, l'article~69 de la loi n°~2010/012 menace l'auditeur.
    \item \textbf{Les exports figés} rendent l'audit sûr, reproductible et
          opposable~; leur exhaustivité se vérifie contre l'inventaire.
    \item \textbf{Une configuration} se compare à une référence écrite~; une
          segmentation, à la matrice des flux autorisés.
    \item \textbf{Les vulnérabilités} se priorisent par l'exploitation,
          l'exposition et la valeur métier, pas par le seul score CVSS.
    \item \textbf{Quatre centres}~: on conclut centre par centre~; la
          moyenne de LiZbiZ progresse, mais le défaut qui ouvre SO1 demeure.
\end{itemize}
\end{boxresume}

\begin{boxinfo}
\textbf{Points de contrôle.} Vous avez compris ce chapitre si vous savez~:
\begin{itemize}
    \item dire ce qui distingue un audit technique d'un test d'intrusion~;
    \item constituer un dossier d'exports figés~;
    \item lire une ligne d'export et la coter au regard d'une référence~;
    \item prioriser des vulnérabilités en justifiant l'ordre.
\end{itemize}
\end{boxinfo}

% ============================================================================
\section*{Exercices}
% ============================================================================

\begin{boxexo}[title={\ding{45}\quad Exercice~12.1. Quelle activité, quelle autorisation~? \niveaubase}]
Pour chacune des demandes suivantes, dites s'il s'agit d'un audit de
configuration, d'un scan de vulnérabilités ou d'un test d'intrusion, et
qui doit signer l'autorisation.
\begin{enumerate}[(a)]
    \item Vérifier que les pare-feu des quatre centres appliquent la
          matrice de flux.
    \item Savoir si un attaquant parti du portail des transporteurs peut
          atteindre la base de facturation.
    \item Recenser les versions vulnérables des serveurs de Garoua,
          hébergés par un prestataire.
    \item Vérifier qu'aucun commutateur n'accepte Telnet.
    \item Mesurer la résistance du portail à un afflux massif de requêtes.
\end{enumerate}
\end{boxexo}

\begin{boxexo}[title={\ding{45}\quad Exercice~12.2. Lire un export \niveaumoyen}]
L'export d'un serveur Linux de Yaoundé contient, entre autres, les lignes
suivantes (simplifiées)~: \texttt{PermitRootLogin yes}~;
\texttt{PasswordAuthentication yes}~; \texttt{Protocol 2}~; un compte
\texttt{backup} sans connexion depuis 400 jours~; aucune ligne d'envoi des
journaux vers le collecteur.
\begin{enumerate}[(a)]
    \item Pour chaque élément, donnez le sous-domaine (IN.1 à IN.5), la
          référence attendue et la cote.
    \item Lequel de ces écarts touche directement la phase « rentrer » du
          scénario SO1~?
    \item La DSI affirme que le serveur a été corrigé la veille. Que faites-vous~?
\end{enumerate}
\end{boxexo}

\begin{boxexo}[title={\ding{45}\quad Exercice~12.3. Comparer les flux \niveaumoyen}]
La politique de la BANQUE SANAGA n'autorise, depuis le réseau des agences
vers le centre de données, que le port 8443 (application de guichet) et le
port 53 (résolution de noms) vers deux serveurs désignés. L'export du
pare-feu montre en outre~: le port 445 vers le serveur de fichiers du
siège~; le port 3389 vers tous les serveurs~; le port 8443 vers
l'application de guichet.
\begin{enumerate}[(a)]
    \item Construisez la matrice et relevez les flux non autorisés.
    \item Pour chacun, décrivez le risque en vous appuyant sur le scénario
          R-AG-01.
    \item Proposez la recommandation et la preuve de sa mise en œuvre.
\end{enumerate}
\end{boxexo}

\begin{boxexo}[title={\ding{45}\quad Exercice~12.4. Prioriser six vulnérabilités \niveauexpert}]
\begin{enumerate}[(a)]
    \item Classez V1 à V6 (Tableau~12.4) par ordre de priorité de
          correction, et justifiez chaque rang par l'exploitation,
          l'exposition et la valeur métier.
    \item La politique de LiZbiZ fixe un délai de correction de 15 jours pour
          un score supérieur à 9, et de 30 jours au-delà de 7. Montrez
          qu'elle conduit à un mauvais ordre, et proposez une règle
          meilleure.
    \item Le second passage cote IN.4 \NON{} à Yaoundé. Rédigez
          l'observation en sept éléments (Chapitre~\ref{chap:conclure}).
    \item Calculez $\mu$, $U$ et $H$ de Garoua et de Bafoussam en 2028, et
          dites ce que la décomposition par centre apporte par rapport à
          l'ensemble.
\end{enumerate}
\end{boxexo}

\bigskip

\begin{boxinfo}
\textbf{Pour aller plus loin.}
\begin{itemize}
    \item Les CIS Benchmarks~\cite{cis2026benchmarks}, pour les références
          de configuration.
    \item La spécification CVSS~4.0~\cite{first2023cvss4}.
    \item Le NIST SP~800-53A~\cite{nist2022sp80053a}, pour les méthodes
          d'évaluation.
\end{itemize}
\end{boxinfo}

\begin{boxinfo}
\textbf{Transition.} Les audits produisent des constats~; la gestion des
risques produit des scénarios~; la conformité produit des obligations. Il
reste à les faire vivre ensemble, au quotidien, dans un dispositif unique.
C'est la GRC en pratique, objet du Chapitre~\ref{chap:grc-pratique}.
\end{boxinfo}

        % ============================================================================
% Fichier  : partie4/P04_C13_grc_pratique.tex
% Titre    : La GRC en Pratique
% Auteur   : MINKA MI NGUIDJOÏ Thierry Emmanuel
% Version  : 1.0 -- novembre 2026
% Corrigés : corriges/Corr_C13.tex
% ============================================================================
% Faits du cas (bible) : BANQUE SANAGA, registre (R-AG-01, R-AG-02, R-DC-02,
% R-DC-03, R-DC-06), plan de traitement, audits de 2028, CERTIQUAL,
% indicateurs I1 à I6 (I6 = 11 %), DIGICORE (rapport ISAE 3402 de type 2),
% avenant GAB, monétique sous PCI DSS, règle « temporaire » de Garoua chez
% LiZbiZ (leçon reprise).
% Faits nouveaux (reportés dans la bible) : M. ATANGANA, responsable GRC
% nommé en octobre 2028, rattaché au directeur des risques ; comité des
% risques de janvier 2029 ; registre des obligations O1 à O12 (7/2/3) ;
% matrice risques-contrôles C1 à C9 ; indicateurs de risque K1 à K5 ;
% contrôles automatisés quotidiens ; DIGICORE n'a pas remis son rapport à
% l'échéance de décembre 2028.
% Loi n° 2010/012 : art. 26, 29, 30 résumés (Chapitre 10).
% ============================================================================

\chapter{La GRC en Pratique}
\label{chap:grc-pratique}
\chaptermark{La GRC en pratique}

\epigraph{%
    « Quatre rapports qui ne se parlent pas\\
    valent moins qu'un tableau qui les relie. »%
}{Principe du cours}

\niveauChapitre{Intermédiaire}{%
    Maîtrise du Chapitre~\ref{chap:smsi}%
}

\begin{boxobjectifs}
\begin{itemize}
    \item Relier risques, obligations, contrôles et preuves dans un
          dispositif unique.
    \item Tenir un registre des risques et un registre des obligations.
    \item Construire une matrice risques-contrôles, et y repérer les risques
          sans contrôle et les contrôles sans risque.
    \item Distinguer indicateurs de risque et indicateurs de performance, et
          fixer des seuils.
    \item Concevoir le tableau de bord GRC d'une instance de gouvernance, y
          compris ce que l'on ignore.
    \item Organiser le contrôle continu de la conformité, et en connaître
          les limites.
\end{itemize}
\end{boxobjectifs}

\bigskip

% ============================================================================
\section*{Avant de commencer : quatre rapports}
% ============================================================================

Douala, janvier 2029. Le comité des risques du conseil de la BANQUE
SANAGA\footnote{Comme dans tout le cours, la BANQUE SANAGA, ses dirigeants,
ses prestataires, les dates et les montants sont fictifs~; les textes cités
sont exacts.} reçoit, pour la même séance, quatre documents~: le tableau de
bord de la RSSI, le rapport annuel de l'audit interne, le rapport de
CERTIQUAL AFRIQUE et une lettre de la COBAC. Chacun est juste, et aucun ne
parle la langue des autres. Le tableau de bord dit «~100~\% des comptes avec
second facteur~»~; l'audit interne, «~trois non-conformités sur les
agences~»~; CERTIQUAL, «~une majeure levée~»~; la COBAC demande «~l'état des
mesures prises depuis l'incident de 2026~».

Le président du comité pose une seule question~: «~Où en sommes-nous~?~»
M.~ATANGANA, responsable GRC nommé en octobre~2028 et rattaché au directeur
des risques, a trois mois pour qu'une page y réponde.

\medskip

Ce chapitre montre comment faire vivre la gouvernance, le risque et la
conformité dans un seul dispositif~: des registres reliés, une matrice, des
indicateurs, un tableau de bord, et des contrôles qui se vérifient seuls.

\begin{boxalert}
\textbf{Violation typique.} Présenter à la gouvernance des rapports
juxtaposés (sécurité, audit, conformité, certification), chacun avec ses
références et ses échelles.

\smallskip
\textbf{Règle~:} un même référentiel relie risques, obligations, contrôles
et preuves. Chaque rapport est une vue de ce référentiel, et le tableau de
bord de la gouvernance les réunit.
\end{boxalert}

\separateur

% ============================================================================
\section{Un Dispositif Unique}
\label{sec:c13-dispositif}
% ============================================================================

Le Chapitre~\ref{chap:grc} a défini la GRC comme un seul dispositif, avec
un référentiel commun. En pratique, ce référentiel tient en quatre objets
et leurs liens.

\begin{figure}[H]
\centering
\begin{tikzpicture}[
    objet/.style={draw=CouleurPrimaire, thick, rounded corners,
                  fill=CouleurPrimaire!8, text width=2.6cm,
                  minimum height=1.1cm, align=center, font=\small},
    lien/.style={<->, thick, CouleurSecondaire},
    etiq/.style={font=\scriptsize, fill=white, inner sep=1pt}]
    \node[objet] (R) at (0,0) {\textbf{Risques}\\registre};
    \node[objet] (O) at (0,-3) {\textbf{Obligations}\\registre};
    \node[objet] (C) at (4.5,-1.5) {\textbf{Contrôles}\\base commune};
    \node[objet] (P) at (9,-1.5) {\textbf{Preuves}\\tests, extractions};
    \draw[lien] (R) -- node[etiq]{traité par} (C);
    \draw[lien] (O) -- node[etiq]{satisfaite par} (C);
    \draw[lien] (C) -- node[etiq]{démontré par} (P);
\end{tikzpicture}
\caption{Le référentiel commun de la GRC}
\label{fig:c13-referentiel}
\end{figure}

Chaque objet a son propriétaire~: le risque, un directeur métier~;
l'obligation, la fonction conformité ou juridique~; le contrôle, son
exploitant~; la preuve, celui qui l'a produite et celui qui l'a évaluée.
Les liens sont le travail propre du responsable GRC~: c'est par eux qu'un
constat d'audit change la cote d'un risque, ou qu'une nouvelle loi crée un
contrôle.

% ============================================================================
\section{Les Deux Registres}
\label{sec:c13-registres}
% ============================================================================

\subsection{Le registre des risques}

Le registre des risques reprend, pour chaque risque, les rubriques des
chapitres précédents~: scénario, propriétaire, cotes inhérente et
résiduelle avec leur base de preuve (Chapitre~\ref{chap:risque}), décision
de traitement ou d'acceptation (Chapitre~\ref{chap:iso27005}), date de
revue.

\begin{tableaucours}{L{1.5cm}L{3.4cm}L{2.1cm}cL{3.0cm}}{%
    \tableauHeader{Réf.} &
    \tableauHeader{Risque} &
    \tableauHeader{Propriétaire} &
    \tableauHeader{Rés.} &
    \tableauHeader{Décision, revue}
}
R-AG-01 & Rançongiciel en agence & Dir. réseau d'agences & 6 & Traité~;
revue juin 2029 \\
R-AG-02 & Compromission d'un GAB & Dir. réseau d'agences & 9 &
Remplacement en cours~; revue mars 2029 \\
R-DC-02 & Défaillance de DIGICORE & Dir. des opérations & 6 (\NEK{}) &
Accepté en l'état des connaissances~; rapport attendu \\
R-DC-03 & Intrusion par Internet & DSI & 8 & Traité~; revue juin 2029 \\
R-DC-06 & Déni de service des services en ligne & Dir. canaux numériques &
9 & Protection contractualisée~; test en mars \\
\end{tableaucours}
\vspace{-0.8em}
\begin{center}{\small\itshape Tableau~13.1. Extrait du registre des risques
(janvier 2029)}\end{center}

\subsection{Le registre des obligations}

\begin{boxdef}
\termecle{Registre des obligations}\index{registre des obligations}~: liste
des exigences qui s'imposent à l'organisation (lois, règlements, contrats,
normes auxquelles elle s'engage), avec, pour chacune, sa source, son
contenu résumé, les contrôles qui la satisfont, la preuve et son statut.
\end{boxdef}

\begin{tableaucours}{cL{3.6cm}L{4.4cm}c}{%
    \tableauHeader{Réf.} &
    \tableauHeader{Source} &
    \tableauHeader{Exigence (résumée)} &
    \tableauHeader{Statut}
}
O1 & Loi n°~2010/012, art.~26~(1) & Gestion continue des risques & \OUI{} \\
O2 & Loi n°~2010/012, art.~26~(2)--(3) & Mécanismes de sécurité visés par
l'Agence & \NEK{} \\
O3 & Loi n°~2010/012, art.~29~(1) & Conservation dix ans des données de
connexion & \NON{} \\
O4 & Loi n°~2010/012, art.~29~(2) & Surveillance des accès aux données &
\OUI{} \\
O5 & Loi n°~2010/012, art.~30~(1) & Revue périodique des systèmes de
sécurité & \OUI{} \\
O6 & Règlement COBAC R-2016/04 & Contrôle interne couvrant le risque
opérationnel & \OUI{} \\
O7 & PCI DSS (monétique) & Évaluation annuelle de conformité & \NEK{} \\
O8 & ISO/IEC~27001 (certificat) & Maintien du SMSI certifié & \OUI{} \\
O9 & Contrat DIGICORE & Rapport d'assurance annuel de type~2 & \NON{} \\
O10 & Avenant du prestataire des GAB & Droit d'audit de la banque & \OUI{} \\
O11 & Loi n°~2024/017 (à vérifier) & Obligations sur les données
personnelles & \NEK{} \\
O12 & Contrats des correspondants bancaires & Questionnaire de sécurité
annuel & \OUI{} \\
\end{tableaucours}
\vspace{-0.8em}
\begin{center}{\small\itshape Tableau~13.2. Le registre des obligations
(janvier 2029)}\end{center}

\begin{boxlecture}{L'état des obligations}
\textbf{Constatation.} 7~\OUI{}, 2~\NON{}, 3~\NEK{} sur 12~:
\[
\mu \approx 0{,}58~;\quad U = 0{,}25~;\quad
H = -\left(\tfrac{7}{12}\log_2\tfrac{7}{12} + \tfrac{2}{12}\log_2\tfrac{2}{12}
  + \tfrac{3}{12}\log_2\tfrac{3}{12}\right) \approx 0{,}454 + 0{,}431 + 0{,}500
  \approx 1{,}38~\text{bit}.
\]

\textbf{Analyse} (confiance élevée). Conformité opaque~: la banque satisfait
la majorité de ses obligations, mais un quart d'entre elles est dans
l'inconnu. Les trois \NEK{} n'ont pas la même cause~: O2 est une pièce à
retrouver~; O7 attend le rapport de l'évaluateur de la monétique~; O11
dépend d'un texte que la banque n'a pas encore analysé. Les deux \NON{}
sont connus et datés~: O3 tient à la capacité de stockage (deux ans au lieu
de dix), O9 au retard de DIGICORE, qui n'a pas remis son rapport à
l'échéance de décembre 2028.

\medskip
\hfill\lecture{0{,}58}{0{,}25}{1{,}38}
\end{boxlecture}

% ============================================================================
\section{La Matrice Risques-Contrôles}
\label{sec:c13-matrice}
% ============================================================================

\subsection{Croiser}

\begin{boxdef}
\termecle{Matrice risques-contrôles}\index{matrice risques-contrôles}~:
tableau qui croise les contrôles de la base commune et les risques du
registre, en indiquant quel contrôle traite quel risque, et avec la dernière
cote de chaque contrôle.
\end{boxdef}

\begin{tableaucours}{cL{3.6cm}cccccc}{%
    \tableauHeader{Réf.} &
    \tableauHeader{Contrôle} &
    \tableauHeader{AG-01} & \tableauHeader{AG-02} &
    \tableauHeader{DC-02} & \tableauHeader{DC-03} &
    \tableauHeader{DC-06} & \tableauHeader{Cote}
}
C1 & Second facteur & $\times$ & & & $\times$ & & \OUI{} \\
C2 & Cloisonnement & $\times$ & $\times$ & & & & \OUI{} \\
C3 & Sauvegardes testées & $\times$ & & $\times$ & & & \OUI{} \\
C4 & Mode dégradé & $\times$ & $\times$ & $\times$ & & $\times$ & \NON{} \\
C5 & Journalisation & $\times$ & $\times$ & & $\times$ & & \OUI{} \\
C6 & Correctifs & $\times$ & $\times$ & & $\times$ & & \NEK{} \\
C7 & Sensibilisation & $\times$ & & & $\times$ & & \NON{} \\
C8 & Assurance des prestataires & & $\times$ & $\times$ & & & \NON{} \\
C9 & Revue mensuelle des règles de pare-feu & & & & & & \OUI{} \\
\end{tableaucours}
\vspace{-0.8em}
\begin{center}{\small\itshape Tableau~13.3. Matrice risques-contrôles
(extrait)}\end{center}

\subsection{Lire la matrice}

La matrice se lit dans trois directions.
\begin{itemize}
    \item \textbf{Par ligne}~: les contrôles qui traitent beaucoup de risques
          sont des \emph{contrôles clés}~; C4 (mode dégradé) en traite
          quatre, et il est \NON{}~: c'est la priorité du plan.
    \item \textbf{Par colonne}~: un risque dont tous les contrôles sont
          \NON{} ou \NEK{} n'est pas maîtrisé. R-DC-06 n'a qu'un contrôle,
          C4, coté \NON{}~: la protection contre le déni de service,
          contractualisée avec le fournisseur d'accès, n'est pas encore un
          contrôle de la base~; il faut l'y inscrire et la tester.
    \item \textbf{Les lignes vides}~: C9 ne traite aucun risque du registre.
          Soit il est inutile, soit un risque manque. C'est le second cas~:
          la règle «~temporaire~» de LiZbiZ (Chapitre~\ref{chap:audit-technique})
          a montré le risque d'une règle de filtrage qui survit à son besoin~;
          le registre de la banque doit le porter.
\end{itemize}

\begin{boxprocedure}
\textbf{Tenir la matrice.} Le responsable GRC vérifie, à chaque mise à
jour~:
\begin{itemize}
    \item que chaque risque au-delà de l'appétence a au moins un contrôle
          coté \OUI{} par étape de son scénario~;
    \item que chaque contrôle est relié à au moins un risque ou une
          obligation~;
    \item que la cote de chaque contrôle vient du dernier test, avec sa
          date et sa source (audit, contrôle automatisé, autoévaluation)~;
    \item qu'un constat d'audit modifie la cote dans la semaine, et la cote
          du risque à la revue suivante.
\end{itemize}
\end{boxprocedure}

% ============================================================================
\section{Les Indicateurs de Risque}
\label{sec:c13-kri}
% ============================================================================

\subsection{Risque ou performance}

Les indicateurs du SMSI (Chapitre~\ref{chap:smsi}) mesurent la mise en
œuvre et l'efficacité des mesures. Un \termecle{indicateur de
risque}\index{indicateur de risque} mesure autre chose~: un signal qui
annonce que l'exposition change, souvent avant l'incident. Il se lit avec
des seuils, qui déclenchent une action.

\begin{tableaucours}{cL{4.2cm}cccL{2.6cm}}{%
    \tableauHeader{Réf.} &
    \tableauHeader{Indicateur de risque} &
    \tableauHeader{Vert} &
    \tableauHeader{Orange} &
    \tableauHeader{Rouge} &
    \tableauHeader{Risque suivi}
}
K1 & Agents qui cliquent lors d'une campagne simulée & $<$~5~\% & 5 à
10~\% & $>$~10~\% & R-AG-01 \\
K2 & Vulnérabilités exploitées non corrigées après 7~jours & 0 & 1 à 2 &
$\geq$~3 & R-AG-01, R-DC-03 \\
K3 & Règles de filtrage temporaires échues & 0 & 1 & $\geq$~2 & R-AG-01,
R-DC-03 \\
K4 & Prestataires critiques sans rapport d'assurance à jour & 0 & 1 &
$\geq$~2 & R-DC-02, R-AG-02 \\
K5 & Tentatives de connexion refusées sur l'accès distant (par semaine,
écart à la moyenne) & $<$~2 fois & 2 à 5 fois & $>$~5 fois & R-DC-03 \\
\end{tableaucours}
\vspace{-0.8em}
\begin{center}{\small\itshape Tableau~13.4. Les indicateurs de risque et
leurs seuils}\end{center}

\begin{boxinfo}
\textbf{Ce que déclenche un seuil.} Un indicateur au rouge ne change pas la
cote du risque à lui seul~: il déclenche la revue du risque par son
propriétaire, sans attendre la date prévue (Chapitre~\ref{chap:iso27005}).
En janvier 2029, K1 est au rouge (11~\%) et K4 à l'orange (DIGICORE)~: les
revues de R-AG-01 et de R-DC-02 sont avancées.
\end{boxinfo}

% ============================================================================
\section{Le Tableau de Bord GRC}
\label{sec:c13-tableau}
% ============================================================================

\subsection{Une page pour la gouvernance}

Le tableau de bord du comité des risques n'est pas celui de la RSSI. Il ne
détaille pas les mesures~; il répond à trois questions~: où sont les risques
qui dépassent l'appétence~? Respectons-nous nos obligations~? Que
savons-nous, et que ne savons-nous pas~?

\begin{boxscenario}{Opération SANAGA. La page de M. ATANGANA, avril 2029}
\textbf{1. Risques au-delà de l'appétence~:} R-AG-02 (GAB, 9~; remplacement
à 60~\%, fin prévue en juin)~; R-DC-06 (déni de service, 9~; test de la
protection en mars, réussi~: proposition de passer à 6).

\medskip
\textbf{2. Obligations~:} 7 satisfaites sur 12~; deux non satisfaites (O3,
conservation, budget de stockage demandé~; O9, DIGICORE, relance et
pénalité contractuelle)~; trois inconnues (O2, O7, O11), avec une date pour
chacune.

\medskip
\textbf{3. Contrôles clés~:} C4 (mode dégradé) non conforme, plan à fin
juin~; C6 (correctifs) inconnu, test de l'audit interne en mai.

\medskip
\textbf{4. Indicateurs de risque~:} K1 au rouge (11~\%)~; K4 à l'orange.

\medskip
\textbf{5. Ce que nous ne savons pas~:} $U = 0{,}25$ sur les obligations~;
contrôles C6 et sensibilisation à vérifier~; base \NEK{} de la vraisemblance
de R-DC-02.

\medskip
\textbf{6. Décisions demandées au comité~:} budget de stockage (O3)~;
position sur DIGICORE si le rapport n'est pas remis en juin (changement de
prestataire)~; validation de la nouvelle cote de R-DC-06.
\end{boxscenario}

La rubrique~5 est la plus rare dans les tableaux de bord, et la plus
utile. Un comité qui ne voit que des chiffres rassurants ne sait pas sur
quoi ils reposent~; un comité qui voit l'incertitude sait où porter
l'effort et ce qu'il accepte d'ignorer.

% ============================================================================
\section{Le Contrôle Continu de la Conformité}
\label{sec:c13-continu}
% ============================================================================

\subsection{Des contrôles qui se vérifient seuls}

Certains contrôles peuvent être testés automatiquement, chaque jour, par
une extraction comparée à une règle~: la part des comptes avec second
facteur (C1), la part des agences qui envoient des journaux (C5), les
règles de filtrage échues (C9, K3). C'est le \termecle{contrôle continu de
la conformité}\index{contrôle continu de la conformité}~: une alerte part
dès qu'une extraction s'écarte de la règle, au lieu d'attendre l'audit
suivant.

\begin{boxscenario}{Opération SANAGA. Trois tests quotidiens}
Depuis février 2029, trois scripts tournent chaque nuit~: extraction de
l'annuaire (comptes sans second facteur), extraction du collecteur
(agences silencieuses depuis plus de 24~heures), extraction des pare-feu
(règles dont la date d'expiration est passée). Leurs résultats alimentent
directement la cote de C1, C5 et C9 dans la matrice.

\medskip
\textbf{Le 14~mars}, le test de C5 signale une agence silencieuse depuis
trente heures~: un commutateur remplacé n'avait pas reçu la configuration
de référence. L'écart est corrigé le jour même. En 2028, le même défaut
avait attendu l'audit pour être découvert (Chapitre~\ref{chap:audit}).
\end{boxscenario}

\subsection{Les limites}

\begin{boxalert}
\textbf{Violation typique.} Tenir pour \OUI{} un contrôle parce que le test
automatique ne signale rien.

\smallskip
\textbf{Règle~:} un test automatique est lui-même un contrôle, qui doit
être validé~: voit-il toute la population (exhaustivité)~? Sa règle
est-elle juste (exactitude)~? Fonctionne-t-il encore (un script arrêté ne
signale rien)~? L'audit interne teste les tests, au moins une fois par an.
Un script qui ne tourne plus transforme un \OUI{} en \NEK{} sans que
personne ne le voie.
\end{boxalert}

Le contrôle continu de la conformité porte sur les contrôles~; l'audit
continu, qui porte sur les données et les transactions pour y détecter les
anomalies, est présenté au Chapitre~\ref{chap:emergents}.

\begin{boxinfo}
\textbf{La loi et le continu.} L'article~26 de la loi n°~2010/012 demande
aux exploitants de gérer \emph{continûment} les risques de sécurité
(Chapitre~\ref{chap:iso27005})~\cite{loi2010012}. Le dispositif décrit dans
ce chapitre (registres reliés, indicateurs à seuils, tests quotidiens) est
une manière concrète de donner un contenu à cet adverbe.
\end{boxinfo}

% ============================================================================
\section{Ce que la GRC demande à chaque Métier}
\label{sec:c13-metiers}
% ============================================================================

\begin{boxmetier}{Responsable GRC}
\textbf{Tenir les liens}
\begin{enumerate}
    \item Tenir le registre des obligations, avec la fonction juridique, et
          le relier aux contrôles.
    \item Tenir la matrice risques-contrôles à jour de chaque test et de
          chaque constat.
    \item Repérer les risques sans contrôle et les contrôles sans risque.
    \item Fixer avec les propriétaires les indicateurs de risque et leurs
          seuils.
    \item Produire une page pour la gouvernance, avec une rubrique «~ce que
          nous ne savons pas~».
\end{enumerate}
\end{boxmetier}

\begin{boxmetier}{RSSI}
\textbf{Alimenter le référentiel}

La RSSI ne tient plus son propre tableau à part~: ses indicateurs et ses
tests alimentent la matrice. Elle automatise les tests qui peuvent l'être,
et demande à l'audit interne de valider ses scripts.
\end{boxmetier}

\begin{boxmetier}{Auditeur en cybersécurité}
\textbf{Auditer le dispositif, et s'en servir}

L'auditeur peut s'appuyer sur la matrice et sur les tests automatiques pour
cibler sa mission, à condition de les avoir validés. Il audite aussi le
dispositif lui-même~: les liens sont-ils justes, les cotes à jour, les
seuils suivis~?
\end{boxmetier}

\begin{boxmetier}{Gestionnaire des risques cyber}
\textbf{Réagir aux seuils}

Un indicateur au rouge déclenche la revue du risque par son propriétaire.
Le gestionnaire des risques organise cette revue dans le mois, et en
inscrit la décision au registre.
\end{boxmetier}

\separateur

\begin{boxresume}
\textbf{Ce qu'il faut retenir de ce chapitre~:}
\begin{itemize}
    \item \textbf{Un référentiel commun}~: risques, obligations, contrôles,
          preuves, et leurs liens.
    \item \textbf{Deux registres}~: les risques (avec propriétaires et
          décisions) et les obligations (avec source et statut).
    \item \textbf{La matrice risques-contrôles} révèle les contrôles clés,
          les risques sans contrôle et les contrôles sans risque.
    \item \textbf{Les indicateurs de risque} annoncent un changement
          d'exposition~; leurs seuils déclenchent une revue.
    \item \textbf{Le tableau de bord de la gouvernance} tient en une page,
          et dit ce que l'on ne sait pas.
    \item \textbf{Le contrôle continu} détecte les écarts au jour le jour~;
          les tests automatiques sont eux-mêmes des contrôles à valider.
\end{itemize}
\end{boxresume}

\begin{boxinfo}
\textbf{Points de contrôle.} Vous avez compris ce chapitre si vous savez~:
\begin{itemize}
    \item rédiger une ligne du registre des obligations à partir d'un
          article de loi~;
    \item lire une matrice risques-contrôles dans ses trois directions~;
    \item définir un indicateur de risque et ses seuils~;
    \item dire pourquoi un test automatique silencieux n'est pas une
          preuve.
\end{itemize}
\end{boxinfo}

% ============================================================================
\section*{Exercices}
% ============================================================================

\begin{boxexo}[title={\ding{45}\quad Exercice~13.1. Classer \niveaubase}]
Dites si chacun des éléments suivants est un risque, une obligation, un
contrôle, une preuve, un indicateur de performance ou un indicateur de
risque.
\begin{enumerate}[(a)]
    \item Conserver dix ans les données de connexion.
    \item La part des comptes avec second facteur.
    \item Le procès-verbal de restauration du 12~février.
    \item Une défaillance prolongée de DIGICORE.
    \item Le nombre de vulnérabilités exploitées non corrigées après sept
          jours.
    \item La revue mensuelle des règles de pare-feu.
\end{enumerate}
\end{boxexo}

\begin{boxexo}[title={\ding{45}\quad Exercice~13.2. Compléter la matrice \niveaumoyen}]
\begin{enumerate}[(a)]
    \item Ajoutez à la matrice (Tableau~13.3) le contrôle «~protection
          contre le déni de service chez le fournisseur d'accès~» et le
          risque «~règle de filtrage temporaire maintenue~». Reliez-les.
    \item Pour chaque risque, dites s'il a au moins un contrôle \OUI{}.
          Lesquels ne sont pas maîtrisés~?
    \item Calculez $\mu$, $U$ et $H$ des neuf contrôles du tableau, puis
          celles des seuls contrôles qui traitent R-AG-01. Que montre la
          comparaison~?
\end{enumerate}
\end{boxexo}

\begin{boxexo}[title={\ding{45}\quad Exercice~13.3. Des indicateurs pour les prestataires \niveaumoyen}]
\begin{enumerate}[(a)]
    \item Proposez deux indicateurs de risque pour R-DC-02 (DIGICORE), avec
          leurs seuils et leur source.
    \item Proposez un indicateur de risque pour R-AG-02 (GAB) une fois le
          remplacement achevé.
    \item Dites, pour chacun, ce que déclenche le passage au rouge, et qui
          décide.
\end{enumerate}
\end{boxexo}

\begin{boxexo}[title={\ding{45}\quad Exercice~13.4. La page de l'ANPP \niveauexpert}]
À partir des Chapitres~\ref{chap:referentiels} et~\ref{chap:droit},
préparez le tableau de bord GRC d'une page destiné au conseil
d'administration de l'ANPP.
\begin{enumerate}[(a)]
    \item Construisez le registre des obligations (au moins six lignes),
          avec leur statut.
    \item Calculez $\mu$, $U$ et $H$ de ce registre et situez-le dans le
          quadrant.
    \item Rédigez les six rubriques de la page, sur le modèle de celle de
          M.~ATANGANA, dont la rubrique «~ce que nous ne savons pas~».
    \item Proposez deux tests automatiques quotidiens réalistes pour
          l'ANPP, et dites comment l'audit les validera.
\end{enumerate}
\end{boxexo}

\bigskip

\begin{boxinfo}
\textbf{Pour aller plus loin.}
\begin{itemize}
    \item L'ISO~31000~\cite{iso31000_2018}, pour l'intégration de la gestion
          des risques dans la gouvernance.
    \item Le cadre COSO de gestion des risques~\cite{coso2017erm}, pour
          l'appétence et les indicateurs.
    \item Le NIST CSF~2.0~\cite{nist2024csf}, fonction \emph{Gouverner},
          catégorie de supervision.
\end{itemize}
\end{boxinfo}

\begin{boxinfo}
\textbf{Transition.} Les chapitres de la Partie~IV ont traité la pratique
par morceaux~: le SMSI, l'audit technique, la GRC. Le
Chapitre~\ref{chap:cas-integrateurs} les assemble dans trois cas complets,
de l'incident au rapport.
\end{boxinfo}

        % ============================================================================
% Fichier  : partie4/P04_C14_cas_integrateurs.tex
% Titre    : Études de Cas Intégratrices
% Auteur   : MINKA MI NGUIDJOÏ Thierry Emmanuel
% Version  : 1.0 -- novembre 2026
% Corrigés : corriges/Corr_C14.tex
% ============================================================================
% Chapitre de synthèse : aucun fait nouveau majeur ; reprise des faits de la
% bible (P00 à C13) pour la BANQUE SANAGA, LiZbiZ et l'ANPP. Faits du
% corpus pour l'ANPP : 40 paiements contrôlés par la Chambre des comptes,
% 29 OUI, 6 NON, 5 NEK (4 agents payés après leur départ ; pièces « aux
% archives »). Faits nouveaux mineurs (reportés dans la bible) : lettre de
% la COBAC de janvier 2029 demandant l'état des mesures ; demande de
% PHARMACAM d'un dossier de renouvellement pour janvier 2030 ; audit de
% l'ANTIC à l'ANPP programmé en septembre 2028 (Chapitre 10).
% ============================================================================

\chapter{Études de Cas Intégratrices}
\label{chap:cas-integrateurs}
\chaptermark{Cas intégrateurs}

\epigraph{%
    « Une méthode se juge à la fin d'une histoire,\\
    pas au début d'un chapitre. »%
}{Principe du cours}

\niveauChapitre{Expert}{%
    Maîtrise du Chapitre~\ref{chap:audit-technique}, du Chapitre~\ref{chap:grc-pratique}%
}

\begin{boxobjectifs}
\begin{itemize}
    \item Retracer, pour chacun des trois cas, le chemin de l'incident ou du
          constat initial jusqu'au rapport final.
    \item Relier gouvernance, appréciation et quantification des risques,
          audit, référentiels et droit dans une même mission.
    \item Lire l'évolution d'un dispositif dans le temps par les trois
          mesures $\mu$, $U$ et $H$.
    \item Reconnaître et éviter les erreurs typiques de l'auditeur et du
          gestionnaire des risques.
    \item Préparer l'épreuve de synthèse du cours.
\end{itemize}
\end{boxobjectifs}

\bigskip

% ============================================================================
\section*{Avant de commencer : trois histoires, une méthode}
% ============================================================================

Les treize chapitres précédents ont suivi trois organisations, chapitre
après chapitre\footnote{Comme dans tout le cours, les organisations, les
personnes, les dates et les montants sont fictifs~; les méthodes, les
normes et les textes cités sont exacts.}. La BANQUE SANAGA a été compromise
malgré son certificat, puis a reconstruit sa sécurité et sa gouvernance.
LiZbiZ a découvert des fraudes, puis a dû convaincre un client exigeant.
L'ANPP, organisme public, a dû répondre à la fois à un bailleur, à la
Chambre des comptes et à l'ANTIC.

\medskip

Ce chapitre ne présente aucune notion nouvelle. Il rassemble chaque
histoire en un seul fil, montre ce que chaque méthode y a apporté, en tire
les erreurs à ne pas commettre, puis propose trois études de cas complètes.
C'est le chapitre à relire avant l'épreuve de synthèse.

\separateur

% ============================================================================
\section{Une Méthode en Dix Questions}
\label{sec:c14-methode}
% ============================================================================

Toutes les missions du cours, quel que soit leur objet, répondent aux mêmes
dix questions, dans cet ordre.

\begin{tableaucours}{cL{6.2cm}L{3.4cm}c}{%
    \tableauHeader{N°} &
    \tableauHeader{Question} &
    \tableauHeader{Outil} &
    \tableauHeader{Chap.}
}
1 & Qui décide, qui répond, qui vérifie~? & RACI, trois lignes & \ref{chap:grc} \\
2 & Quel est le périmètre, et qu'en exclut-on~? & Périmètre défendable &
\ref{chap:smsi} \\
3 & Quels scénarios, avec quelle base de preuve~? & Fiche de risque, EBIOS &
\ref{chap:risque}, \ref{chap:ebios} \\
4 & Combien cela coûterait-il~? & FAIR, Monte-Carlo & \ref{chap:fair} \\
5 & Que décide-t-on, et qui accepte le reste~? & Plan, critères
d'acceptation & \ref{chap:iso27005} \\
6 & Quels textes s'imposent, lesquels sont choisis~? & Registre des
obligations & \ref{chap:referentiels}, \ref{chap:droit} \\
7 & Les mesures sont-elles en place et efficaces~? & Audit, exports figés &
\ref{chap:audit}, \ref{chap:audit-technique} \\
8 & Que peut-on conclure, et que ne sait-on pas~? & Décomposition, $\mu$,
$U$, $H$ & \ref{chap:conclure} \\
9 & Le système se corrige-t-il seul~? & Audit interne, revue de direction &
\ref{chap:smsi} \\
10 & Une seule vue pour la gouvernance~? & Matrice, tableau de bord &
\ref{chap:grc-pratique} \\
\end{tableaucours}
\vspace{-0.8em}
\begin{center}{\small\itshape Tableau~14.1. La méthode du cours en dix
questions}\end{center}

% ============================================================================
\section{BANQUE SANAGA : du Rançongiciel au Certificat}
\label{sec:c14-sanaga}
% ============================================================================

\subsection{Le fil}

\begin{tableaucours}{L{2.2cm}L{6.6cm}c}{%
    \tableauHeader{Date} &
    \tableauHeader{Événement} &
    \tableauHeader{Chap.}
}
Juin 2025 & Certificat sur le seul centre de données & \ref{chap:prologue} \\
Sept. 2026 & Rançongiciel par l'agence de Ngaoundéré~: 180 postes, 9
agences, 410~M FCFA & \ref{chap:prologue} \\
Oct. 2026 & « Qui a décidé du périmètre~? »~: gouvernance 2/4/2 &
\ref{chap:grc} \\
Nov. 2026 & R-AG-01 coté 16 puis 12~; registre 2025 revu~: 6/2/4 &
\ref{chap:risque} \\
Déc. 2026 & Critères d'acceptation, plan A1--A8 (193~M) & \ref{chap:iso27005} \\
Mars 2027 & Quantification~: 165~M par an, puis 19 après traitement &
\ref{chap:fair} \\
Sept. 2027 & Suivi du plan~: 5/2/3 & \ref{chap:iso27005} \\
Mars 2028 & Audit des agences~: 3/3/2~; Kribi, GAB & \ref{chap:audit},
\ref{chap:conclure} \\
Avr. 2028 & Déclaration d'applicabilité confrontée aux preuves~: 3/3/3 &
\ref{chap:iso27000} \\
Juin 2028 & Étape~2~: une majeure (GAB non appréciés) & \ref{chap:smsi} \\
Sept. 2028 & Certificat renouvelé sur toute la banque & \ref{chap:smsi} \\
Avr. 2029 & Page GRC du comité~: obligations 7/2/3 & \ref{chap:grc-pratique} \\
\end{tableaucours}
\vspace{-0.8em}
\begin{center}{\small\itshape Tableau~14.2. La BANQUE SANAGA, de 2025 à
2029}\end{center}

\subsection{Ce que les mesures racontent}

\begin{boxlecture}{Le risque R-AG-01 et la connaissance qu'on en a}
\textbf{Constatation.} La criticité résiduelle de R-AG-01 passe de 12
(2025) à 9 (mars 2027), puis à 6 (septembre 2027), confirmé par l'audit de
mars 2028. Sur la même période, les évaluations successives donnent~:
\begin{itemize}
    \item agences, prologue~: $\mu = 0{,}20$, $U = 0{,}40$~;
    \item agences, audit de mars 2028~: $\mu = 0{,}375$, $U = 0{,}25$~;
    \item obligations, avril 2029~: $\mu \approx 0{,}58$, $U = 0{,}25$.
\end{itemize}

\textbf{Analyse} (confiance moyenne~: les ensembles mesurés ne sont pas
identiques). Le risque technique a baissé plus vite que l'incertitude. Les
barrières techniques ont été posées et prouvées en dix-huit mois~; la part
d'inconnu, elle, est restée autour d'un quart, parce qu'elle s'est déplacée~:
des agences vers le personnel, puis vers les prestataires et les textes. La
banque n'a pas cessé d'ignorer~; elle a appris \emph{où} elle ignorait.
\end{boxlecture}

\subsection{Ce que le cas enseigne}

\begin{itemize}
    \item Un certificat ne protège que son périmètre~; le périmètre est une
          décision du conseil.
    \item Les défaillances de 2025 et de 2028 ont la même forme~: un risque
          connu, ou connaissable, laissé sans appréciation sur une partie du
          périmètre.
    \item Les mesures techniques se corrigent vite~; les mesures humaines et
          les relations avec les prestataires, lentement.
\end{itemize}

% ============================================================================
\section{LiZbiZ : de la Fraude au Client Exigeant}
\label{sec:c14-lizbiz}
% ============================================================================

\subsection{Le fil}

\begin{tableaucours}{L{2.2cm}L{6.6cm}c}{%
    \tableauHeader{Date} &
    \tableauHeader{Événement} &
    \tableauHeader{Chap.}
}
2026 & Écart de stock de 186~M FCFA~; fraude TRANSLOG (112,7~M)~; registre
R1--R8 (cours ACSI) & \ref{chap:risque} \\
2026 & Audit des quatre centres~: 8/9/3 & \ref{chap:audit-technique} \\
2026--2027 & Base de preuve du registre~: 4/1/3~; R8 recoté à 12 &
\ref{chap:risque} \\
Janv. 2027 & PHARMACAM exige une analyse~; étude EBIOS animée par
Mme~NKENG~; SO1 à 12 & \ref{chap:ebios} \\
2027 & Quantification simple de SO1~: 100~M par an & \ref{chap:fair} \\
2028 & Audit contractuel de PHARMACAM~; déclaration d'applicabilité 3/2/1 &
\ref{chap:audit}, \ref{chap:iso27000} \\
Oct. 2028 & Second passage technique~: 14/4/2~; règle de maintenance de
Garoua & \ref{chap:audit-technique} \\
2029 & Engagement d'un SMSI, certification visée fin 2029 & \ref{chap:smsi} \\
\end{tableaucours}
\vspace{-0.8em}
\begin{center}{\small\itshape Tableau~14.3. LiZbiZ, de 2026 à 2029}\end{center}

\subsection{Ce que le cas enseigne}

\begin{itemize}
    \item Une petite organisation n'a ni RSSI ni audit interne~: elle
          compense par la séparation des rôles qu'elle peut tenir
          (contrôleuse interne, auditeur externe indépendant) et par une
          hygiène de base vérifiable.
    \item Les fraudes internes et les attaques externes exploitent les
          mêmes faiblesses~; EBIOS les met sur le même plan.
    \item Le client est un régulateur de fait~: ses clauses (analyse,
          droit d'audit) font avancer la sécurité plus vite que la loi.
    \item Le défaut qui ouvre le scénario principal peut survivre à tous les
          plans s'il est «~temporaire~»~: la règle de Garoua l'a montré.
\end{itemize}

% ============================================================================
\section{ANPP : un Organisme Public sous Trois Regards}
\label{sec:c14-anpp}
% ============================================================================

\subsection{Le fil}

\begin{tableaucours}{L{2.2cm}L{6.6cm}c}{%
    \tableauHeader{Date} &
    \tableauHeader{Événement} &
    \tableauHeader{Chap.}
}
2026 & Chambre des comptes~: 40 paiements contrôlés, 29~\OUI{}, 6~\NON{},
5~\NEK{}~; quatre agents payés après leur départ (cours ACSI) &
\ref{chap:referentiels} \\
Mars 2027 & Un bailleur exige un profil NIST CSF~; base commune B1--B10~:
4/3/3 & \ref{chap:referentiels} \\
2028 & L'ANTIC annonce un audit de sécurité~; obligations légales~: 0/3/2 &
\ref{chap:droit} \\
2028--2029 & Tableau de bord GRC pour le conseil d'administration &
\ref{chap:grc-pratique} \\
\end{tableaucours}
\vspace{-0.8em}
\begin{center}{\small\itshape Tableau~14.4. L'ANPP, de 2026 à 2029}\end{center}

\begin{boxlecture}{Les paiements de l'ANPP}
\textbf{Constatation.} 29~\OUI{}, 6~\NON{}, 5~\NEK{} sur 40 paiements~:
$\mu \approx 0{,}73$~; $U = 0{,}125$~;
\[
H = -(0{,}725\log_2 0{,}725 + 0{,}15\log_2 0{,}15 + 0{,}125\log_2 0{,}125)
  \approx 0{,}336 + 0{,}411 + 0{,}375 \approx 1{,}12~\text{bit}.
\]

\textbf{Analyse} (confiance moyenne). Lu seul, le contrôle des paiements
paraît plutôt conforme. Rapproché de la base commune (B2~: comptes des
partants actifs, \NON{}) et des obligations légales (0/3/2), il change de
sens~: les quatre agents payés après leur départ sont le symptôme
financier d'un défaut de sécurité. Le même défaut apparaît dans trois
rapports, à trois destinataires~; il n'appelle qu'une correction.

\medskip
\hfill\lecture{0{,}725}{0{,}125}{1{,}12}
\end{boxlecture}

\subsection{Ce que le cas enseigne}

\begin{itemize}
    \item Un organisme public répond à plusieurs autorités, avec des textes
          parfois ambigus~; l'auditeur signale l'ambiguïté au lieu de la
          trancher.
    \item Une base commune de contrôles, testée une fois, évite trois
          chantiers et révèle les défauts communs.
    \item La sécurité et la régularité financière se rejoignent dans les
          contrôles d'accès.
\end{itemize}

% ============================================================================
\section{Dix Erreurs à ne pas Commettre}
\label{sec:c14-erreurs}
% ============================================================================

\begin{tableaucours}{cL{6.0cm}L{3.8cm}c}{%
    \tableauHeader{N°} &
    \tableauHeader{Erreur} &
    \tableauHeader{Où le cours l'a montrée} &
    \tableauHeader{Chap.}
}
1 & Prendre un certificat pour une preuve de sécurité & Hall de la BANQUE
SANAGA & \ref{chap:prologue} \\
2 & Lire un risque absent du registre comme un risque faible & Registre
2025 & \ref{chap:prologue} \\
3 & Laisser une décision de risque sous l'instance qui en répond &
Périmètre signé par le seul DG & \ref{chap:grc} \\
4 & Coter une vraisemblance sur un silence non surveillé & « Jamais
attaqués » & \ref{chap:risque} \\
5 & Réduire un risque par une mesure non testée & Cloisonnement de mars
2027 & \ref{chap:risque} \\
6 & Présenter un chiffre unique au lieu d'une distribution & 143,5~M &
\ref{chap:fair} \\
7 & Coter \OUI{} sur une déclaration & Rapport « tout est en place » &
\ref{chap:audit} \\
8 & Taire une limitation de l'étendue & GAB, prestataire & \ref{chap:audit} \\
9 & Résumer un audit par un pourcentage & « 37~\% » & \ref{chap:conclure} \\
10 & Reporter une conformité par un lien partiel & « Conforme ISO, donc
NIST » & \ref{chap:referentiels} \\
\end{tableaucours}
\vspace{-0.8em}
\begin{center}{\small\itshape Tableau~14.5. Les dix erreurs typiques}\end{center}

Toutes ont la même racine~: confondre ce que l'on sait et ce que l'on
suppose. Le cours n'a enseigné qu'une seule discipline, sous plusieurs
noms~: dire, à chaque étape, ce qui est établi, ce qui est contredit, et ce
que l'on ignore.

% ============================================================================
\section{Préparer l'Épreuve de Synthèse}
\label{sec:c14-epreuve}
% ============================================================================

L'épreuve de synthèse prend la forme d'une étude de cas, sur le modèle des
trois études ci-dessous~: un dossier (faits, extraits de documents,
résultats de tests) et des questions qui suivent la méthode en dix
questions. Le barème détaillé, la durée et les documents autorisés sont
fixés par l'enseignant et communiqués avec le calendrier de l'UE.

\begin{boxprocedure}
\textbf{Répondre à une étude de cas.}
\begin{itemize}
    \item Lire tout le dossier avant de répondre~; repérer les acteurs, le
          périmètre et ce qui est demandé.
    \item Séparer, pour chaque réponse, la constatation (le fait, la preuve)
          de l'analyse (l'interprétation, avec un niveau de confiance).
    \item Poser les calculs~: cotes, criticités, $\mu$, $U$ et $H$, avec les
          valeurs intermédiaires.
    \item Nommer ce que l'on ne sait pas, et dire comment le savoir.
    \item Citer les normes et les textes avec leur degré de vérification~;
          ne jamais inventer un numéro d'article.
    \item Conclure par une décision, et dire qui la prend.
\end{itemize}
\end{boxprocedure}

% ============================================================================
\section{Le Parcours de chaque Métier}
\label{sec:c14-metiers}
% ============================================================================

\begin{boxmetier}{RSSI}
\textbf{De Mme ABENA en 2026 à Mme ABENA en 2029}

Elle a obtenu le droit d'être entendue par le comité, cessé d'administrer ce
qu'elle surveille, remplacé ses déclarations par des preuves, et présenté
un indicateur défavorable (11~\% de clics) sans l'enjoliver.
\end{boxmetier}

\begin{boxmetier}{Auditeur en cybersécurité}
\textbf{Mme EKAMBI, M. BEYALA, CERTIQUAL AFRIQUE}

Trois auditeurs, trois postures, une même règle~: la preuve que l'on
obtient soi-même, la limitation que l'on déclare, le constat que l'on
maintient sous la pression, l'indépendance que l'on protège.
\end{boxmetier}

\begin{boxmetier}{Gestionnaire des risques cyber}
\textbf{Du registre au scénario}

Il est passé d'une liste cotée sans preuve à des scénarios décomposés,
reliés aux attaquants et à l'écosystème, chiffrés en francs CFA, avec une
base de preuve pour chaque cote.
\end{boxmetier}

\begin{boxmetier}{Responsable GRC}
\textbf{M. ATANGANA}

Il a relié quatre rapports en une page, et fait apparaître ce que le comité
ignorait. Son outil principal n'est pas un logiciel~: ce sont les liens
entre risques, obligations, contrôles et preuves.
\end{boxmetier}

\begin{boxmetier}{Consultant en mise en œuvre de SMSI}
\textbf{Le consultant de 2025, et celui de LiZbiZ}

Le premier a livré un certificat sur un périmètre commode~; celui que
LiZbiZ retiendra devra livrer un processus qui fonctionne sans lui, et un
périmètre qui inclut ses interfaces.
\end{boxmetier}

\separateur

\begin{boxresume}
\textbf{Ce qu'il faut retenir de ce chapitre~:}
\begin{itemize}
    \item \textbf{Dix questions} structurent toute mission~: rôles,
          périmètre, scénarios, chiffrage, décisions, textes, preuves,
          conclusion, amélioration, vue d'ensemble.
    \item \textbf{BANQUE SANAGA}~: le risque baisse vite, l'inconnu se
          déplace~; les échecs de 2025 et de 2028 ont la même forme.
    \item \textbf{LiZbiZ}~: fraude et attaque exploitent les mêmes
          faiblesses~; le client fait avancer la sécurité.
    \item \textbf{ANPP}~: un même défaut d'accès apparaît dans trois
          rapports~; une base commune le corrige une fois.
    \item \textbf{Dix erreurs}, une racine~: confondre ce que l'on sait et
          ce que l'on suppose.
\end{itemize}
\end{boxresume}

% ============================================================================
\section*{Exercices}
% ============================================================================

\begin{boxexo}[title={\ding{45}\quad Étude de cas~14.1. La réponse à la COBAC \niveauexpert}]
En janvier 2029, la COBAC demande à la BANQUE SANAGA «~l'état des mesures
prises depuis l'incident de septembre 2026~».
\begin{enumerate}[(a)]
    \item Identifiez le cadre de la demande et le degré de vérification du
          texte (Chapitre~\ref{chap:droit}).
    \item Rédigez le plan de la réponse, en suivant la méthode en dix
          questions.
    \item Rédigez la partie «~risque~» de la réponse~: évolution de la
          criticité de R-AG-01, avec la base de preuve de chaque étape, et
          la quantification avant et après traitement.
    \item Rédigez la partie «~ce que la banque ne sait pas encore~», et le
          calendrier pour le savoir.
    \item Qui signe la réponse, et pourquoi~?
\end{enumerate}
\end{boxexo}

\begin{boxexo}[title={\ding{45}\quad Étude de cas~14.2. Le dossier de renouvellement de PHARMACAM \niveauexpert}]
PHARMACAM demande à LiZbiZ, pour le renouvellement de janvier 2030, un
dossier de sécurité.
\begin{enumerate}[(a)]
    \item Listez les pièces du dossier, en distinguant ce qui est prouvé et
          ce qui est déclaré.
    \item Retracez l'évolution du scénario SO1 de 2027 à 2028 (criticité,
          phases couvertes, preuves).
    \item Comparez les deux passages de l'audit technique (8/9/3 et 14/4/2)
          centre par centre, et rédigez le paragraphe qui les présente à
          PHARMACAM sans moyenne trompeuse.
    \item Quelle décision recommandez-vous à M.~ETEME sur la certification
          ISO/IEC~27001, et à quelle échéance~?
\end{enumerate}
\end{boxexo}

\begin{boxexo}[title={\ding{45}\quad Étude de cas~14.3. Un défaut, trois rapports \niveauexpert}]
À l'ANPP, le défaut de désactivation des comptes des agents partis apparaît
dans trois travaux~: le contrôle des paiements de la Chambre des comptes,
la base commune (B2), et les obligations légales (article~29, alinéa~2).
\begin{enumerate}[(a)]
    \item Montrez, pour chacun des trois destinataires (M.~NJOYA, le
          bailleur, l'ANTIC), ce que le défaut signifie dans sa langue.
    \item Rédigez une seule recommandation, avec sa preuve de mise en œuvre,
          qui réponde aux trois.
    \item Proposez le test automatique qui permettra de vérifier chaque
          nuit que le défaut ne revient pas, et dites comment l'audit le
          validera.
    \item Recalculez $\mu$, $U$ et $H$ de la base commune si B2 devient \OUI{}
          et si B4 et B6 sont établis (\OUI{} et \NON{}).
\end{enumerate}
\end{boxexo}

\bigskip

\begin{boxinfo}
\textbf{Transition.} Les trois cas reposent sur des systèmes que leurs
propriétaires détiennent. Les risques qui viennent sont ailleurs~: chez les
prestataires et dans le nuage, dans l'intelligence artificielle, dans
l'audit qui devient continu. La Partie~V les aborde au
Chapitre~\ref{chap:emergents}.
\end{boxinfo}


    % ======================================================================
    % PARTIE 5 — DÉFIS AVANCÉS
    % Standards de référence : ISO/IEC 42001, ISAE 3402 / SOC 2.
    % Niveau : expert
    % ======================================================================
    \part{Défis Avancés}
    \label{part:avance}

        % ============================================================================
% Fichier  : partie5/P05_C15_risques_emergents.tex
% Titre    : Risques Émergents et Nouveaux Objets d'Audit
% Auteur   : MINKA MI NGUIDJOÏ Thierry Emmanuel
% Version  : 1.0 -- novembre 2026
% Corrigés : corriges/Corr_C15.tex
% ============================================================================
% Sources : ISO/IEC 27002:2022 (mesures 5.19 à 5.23, intitulés traduits) ;
% ISO/IEC 27017:2015 et 27018:2019 ; ISO/IEC 42001:2023 (système de
% management de l'IA, certifiable) ; NIST AI RMF 1.0 (2023 : fonctions
% Gouverner, Cartographier, Mesurer, Gérer ; caractéristiques de
% confiance) ; ISAE 3402 / SOC 2 (Chapitre 9). Loi n° 2010/012, art. 26
% al. 1 (cité au Chapitre 3) ; loi n° 2024/017 à vérifier.
% Faits du corpus : ANPP (quatre agents payés après leur départ) ; LiZbiZ
% (148 M FCFA d'ajustements de stock, dont 71 % en fin de mois, sous le
% compte de M. FOUDA ; cours ACSI).
% Faits nouveaux (reportés dans la bible) : en 2029, la SANAGA migre sa
% messagerie vers un service en nuage (NUAGEPRO, fictif) ; registre des
% 9 prestataires critiques (4/1/4) ; modèle de détection de fraude aux
% virements acheté à un éditeur (FRAUDSCAN, fictif) ; usage d'un assistant
% d'IA générative public par des agents (copie de données de clients) ;
% audit continu des paiements de l'ANPP : 37 exceptions en juin 2029
% (29 justifiées, 6 indues, 2 inconnues).
% ============================================================================

\chapter{Risques Émergents et Nouveaux Objets d'Audit}
\label{chap:emergents}
\chaptermark{Risques émergents}

\epigraph{%
    « On peut externaliser un service.\\
    On n'externalise jamais la responsabilité. »%
}{Principe du cours}

\niveauChapitre{Expert}{%
    Maîtrise du Chapitre~\ref{chap:cas-integrateurs}%
}

\begin{boxobjectifs}
\begin{itemize}
    \item Gérer et auditer les risques liés aux tiers et à la chaîne
          d'approvisionnement.
    \item Auditer un service en nuage selon le partage des responsabilités.
    \item Situer les risques des systèmes d'intelligence artificielle, et les
          cadres qui les encadrent (ISO/IEC~42001, NIST AI RMF).
    \item Concevoir un dispositif d'audit continu, et en connaître les
          limites.
    \item Appliquer à ces nouveaux objets la cote \OUI{}~/ \NON{}~/ \NEK{} et
          les trois mesures.
\end{itemize}
\end{boxobjectifs}

\bigskip

% ============================================================================
\section*{Avant de commencer : la sécurité qui n'est plus chez soi}
% ============================================================================

Douala, 2029. La BANQUE SANAGA\footnote{Comme dans tout le cours, les
organisations, les prestataires (DIGICORE, NUAGEPRO, FRAUDSCAN), les
personnes, les dates et les montants sont fictifs~; les normes et les
cadres cités sont exacts.} a reconstruit sa sécurité sur ce qu'elle
détient~: ses postes, ses réseaux, ses agences. Mais une part croissante de
son système d'information n'est plus chez elle. Le centre de données est
hébergé par DIGICORE~; la messagerie migre en 2029 vers un service en nuage,
NUAGEPRO~; un modèle d'intelligence artificielle, acheté à l'éditeur
FRAUDSCAN, analyse désormais chaque virement. Et des agents ont pris
l'habitude de coller des extraits de dossiers de clients dans un assistant
d'IA générative public, «~pour rédiger plus vite~».

La RSSI résume la situation au comité~: «~Nous savons auditer nos systèmes.
Nous ne savons pas encore auditer ceux des autres, ni ceux qui
apprennent.~»

\medskip

Ce chapitre ouvre la Partie~V. Il présente trois objets d'audit qui
transforment le métier~: les tiers et le nuage, l'intelligence
artificielle, et l'audit lui-même, qui devient continu.

\begin{boxalert}
\textbf{Violation typique.} Considérer qu'un risque confié à un prestataire
sort du registre de l'organisation.

\smallskip
\textbf{Règle~:} le risque reste celui de l'organisation~; seule
l'exécution de certaines mesures change de main. Le registre porte le
risque, le contrat porte les obligations du prestataire, et la preuve doit
revenir à l'organisation.
\end{boxalert}

\separateur

% ============================================================================
\section{Les Tiers et la Chaîne d'Approvisionnement}
\label{sec:c15-tiers}
% ============================================================================

\subsection{Un risque qui vient de l'extérieur}

Le Chapitre~\ref{chap:ebios} a montré qu'un attaquant ciblé passe souvent
par un partenaire moins protégé. Les tiers posent trois questions
distinctes~: un prestataire peut \emph{servir de porte d'entrée} (l'accès de
maintenance de LiZbiZ)~; il peut \emph{défaillir} (DIGICORE)~; et il peut
lui-même dépendre d'autres prestataires, que l'organisation ne connaît pas
(les sous-traitants de rang deux).

L'ISO/IEC~27002:2022 consacre cinq mesures à ces sujets~: la politique de
sécurité dans les relations avec les fournisseurs (5.19), la sécurité dans
les accords (5.20), la chaîne d'approvisionnement des TIC (5.21), la
surveillance et la revue des services des fournisseurs (5.22), et la
sécurité des services en nuage (5.23)~\cite{iso27002_2022}.

\subsection{Le cycle de vie d'un tiers}

\begin{tableaucours}{L{2.6cm}L{9.6cm}}{%
    \tableauHeader{Étape} &
    \tableauHeader{Ce que l'organisation fait, et ce qu'elle garde comme preuve}
}
Sélection & Classe le tiers selon sa criticité (dépendance, pénétration~:
Chapitre~\ref{chap:ebios})~; évalue sa maturité par un questionnaire et des
attestations \\
Contrat & Inscrit les exigences de sécurité, la notification des
incidents, le droit d'audit ou le rapport d'assurance, la réversibilité,
l'accord préalable sur la sous-traitance \\
Exploitation & Suit les indicateurs de service et de sécurité, reçoit les
rapports d'assurance, exerce son droit d'audit~; met à jour la criticité \\
Sortie & Récupère les données, révoque les accès, obtient la preuve de
destruction \\
\end{tableaucours}
\vspace{-0.8em}
\begin{center}{\small\itshape Tableau~15.1. Le cycle de vie d'un tiers}\end{center}

\begin{boxlecture}{Les neuf prestataires critiques de la BANQUE SANAGA}
\textbf{Constatation.} En mars 2029, M.~ATANGANA cote, pour chacun des neuf
prestataires classés critiques, l'existence d'une preuve récente que ses
mesures fonctionnent (rapport d'assurance de type~2, audit exercé, ou
certification couvrant le service)~:
\begin{itemize}
    \item \OUI{}~: le prestataire des GAB (audit exercé grâce à
          l'avenant), le fournisseur d'accès (test de la protection contre
          le déni de service), l'éditeur du système central (rapport de
          type~2), NUAGEPRO (certification ISO/IEC~27001 et rapport de
          type~2 couvrant la messagerie)~;
    \item \NON{}~: DIGICORE (rapport annuel non remis)~;
    \item \NEK{}~: FRAUDSCAN, le prestataire de la monétique, l'imprimeur
          des chéquiers et le transporteur de fonds (aucune preuve demandée
          à ce jour).
\end{itemize}
4~\OUI{}, 1~\NON{}, 4~\NEK{}~:
\[
\mu = U \approx 0{,}44~;\quad
H = -\left(2 \times \tfrac{4}{9}\log_2\tfrac{4}{9} + \tfrac{1}{9}\log_2\tfrac{1}{9}\right)
  \approx 1{,}040 + 0{,}352 \approx 1{,}39~\text{bit}.
\]

\textbf{Analyse} (confiance élevée). La banque connaît bien ses prestataires
«~historiques~», ceux de l'incident de 2026 et de l'audit de 2028. Les
quatre \NEK{} sont des prestataires que personne n'a regardés sous l'angle
de la sécurité, parce qu'ils ne sont pas «~informatiques~». Le transporteur
de fonds et l'imprimeur des chéquiers ont pourtant accès à des
informations sensibles. C'est le périmètre de complaisance du prologue,
déplacé vers l'écosystème.

\medskip
\hfill\lecture{0{,}44}{0{,}44}{1{,}39}
\end{boxlecture}

% ============================================================================
\section{Auditer un Service en Nuage}
\label{sec:c15-nuage}
% ============================================================================

\subsection{La responsabilité partagée}

Dans un service en nuage, les mesures de sécurité se répartissent entre le
fournisseur et le client, selon le modèle de service. Le fournisseur répond
de ce qu'il exploite~; le client, de ce qu'il configure et de ce qu'il y
met.

\begin{tableaucours}{L{4.0cm}ccc}{%
    \tableauHeader{Couche} &
    \tableauHeader{Infrastructure (IaaS)} &
    \tableauHeader{Plateforme (PaaS)} &
    \tableauHeader{Logiciel (SaaS)}
}
Centre de données, matériel & Fournisseur & Fournisseur & Fournisseur \\
Système d'exploitation & Client & Fournisseur & Fournisseur \\
Application & Client & Client & Fournisseur \\
Configuration du service & Client & Client & Client \\
Identités et accès des utilisateurs & Client & Client & Client \\
Données & Client & Client & Client \\
\end{tableaucours}
\vspace{-0.8em}
\begin{center}{\small\itshape Tableau~15.2. Le partage des responsabilités
(schéma général~; chaque contrat le précise)}\end{center}

Deux lignes restent toujours au client~: les identités et les données. Une
messagerie en nuage parfaitement sécurisée par son fournisseur reste
vulnérable si les comptes des agents n'ont pas de second facteur, ou si un
agent partage un dossier de clients avec tout Internet.

\subsection{Les référentiels du nuage}

L'ISO/IEC~27017 complète l'ISO/IEC~27002 par des recommandations propres
aux services en nuage, pour le fournisseur comme pour le
client~\cite{iso27017_2015}~; l'ISO/IEC~27018 traite de la protection des
données personnelles dans le nuage public~\cite{iso27018_2019}. Pour
l'auditeur, la preuve vient le plus souvent des rapports d'assurance du
fournisseur (type~2) et de ses certifications~: il vérifie qu'ils couvrent
bien le service utilisé, la période, et le centre de données qui héberge
les données.

\begin{boxscenario}{Opération SANAGA. La messagerie dans le nuage}
\textbf{Faits.} La messagerie migre vers NUAGEPRO en mai 2029. NUAGEPRO
fournit une certification ISO/IEC~27001 et un rapport de type~2.

\medskip
\textbf{Analyse} (confiance moyenne). Ces pièces couvrent la part du
fournisseur. La part de la banque se vérifie autrement~: second facteur sur
tous les comptes (une extraction de la console d'administration du
service)~; règles de partage externe désactivées~; journaux d'accès
récupérés dans le collecteur de la banque~; localisation des données et
conditions de leur transfert, à examiner au regard de la loi
n°~2024/017 (à vérifier). Le «~rapport du fournisseur~» ne remplace aucune de
ces quatre preuves.
\end{boxscenario}

\begin{boxinfo}
\textbf{La loi et les tiers.} L'article~26, alinéa~1, de la loi
n°~2010/012 vise les risques liés aux services offerts « directement ou
indirectement » (Chapitre~\ref{chap:iso27005})~\cite{loi2010012}.
Analyse (lecture du cours)~: confier un service à un prestataire ne fait
pas sortir ses risques du champ de l'obligation de l'exploitant.
\end{boxinfo}

% ============================================================================
\section{L'Intelligence Artificielle}
\label{sec:c15-ia}
% ============================================================================

\subsection{Des risques nouveaux}

Un système d'intelligence artificielle est un logiciel dont le comportement
est appris à partir de données plutôt qu'entièrement programmé. Il ajoute
aux risques classiques des risques propres~:
\begin{itemize}
    \item \textbf{la fuite de données}~: ce que l'on confie à un assistant
          public peut être conservé, réutilisé, ou exposé~;
    \item \textbf{la manipulation}~: des données d'entraînement empoisonnées,
          ou des instructions cachées dans un document soumis à un
          assistant, peuvent détourner le système~;
    \item \textbf{l'erreur opaque}~: un modèle peut se tromper sans que l'on
          sache pourquoi, et ses performances peuvent se dégrader quand les
          données changent~;
    \item \textbf{le biais}~: un modèle peut traiter différemment des
          catégories de personnes, ce qui pose des questions de droit et de
          réputation.
\end{itemize}

\subsection{Deux cadres}

\begin{tableaucours}{L{2.6cm}L{4.8cm}L{4.8cm}}{%
    \tableauHeader{} &
    \tableauHeader{ISO/IEC~42001:2023} &
    \tableauHeader{NIST AI RMF~1.0 (2023)}
}
Nature & Système de management de l'IA, certifiable~\cite{iso42001_2023} &
Cadre volontaire de gestion des risques~\cite{nist2023airmf} \\
Structure & Même structure que l'ISO/IEC~27001 (articles~4 à~10), avec une
annexe de mesures & Quatre fonctions~: gouverner, cartographier, mesurer,
gérer \\
Apport & S'intègre au SMSI existant & Décrit les caractéristiques d'une IA
digne de confiance~: valide et fiable, sûre, sécurisée et résiliente,
responsable et transparente, explicable, respectueuse de la vie privée,
équitable \\
\end{tableaucours}
\vspace{-0.8em}
\begin{center}{\small\itshape Tableau~15.3. Deux cadres pour l'IA}\end{center}

\subsection{Deux usages à la BANQUE SANAGA}

\begin{boxscenario}{Opération SANAGA. Un modèle acheté, un assistant toléré}
\textbf{Le modèle de FRAUDSCAN.} Il attribue à chaque virement un score de
risque~; au-delà d'un seuil, le virement est bloqué pour vérification.
Questions de l'auditeur~: sur quelles données le modèle a-t-il été
entraîné~? Qui a validé le seuil, et sur quels essais~? Comment la banque
sait-elle, mois après mois, que le taux de fraudes manquées ne se dégrade
pas~? Que peut-elle expliquer à un client dont le virement est bloqué, ou
au superviseur~? À ce jour, la banque ne peut répondre qu'à la deuxième
question~: les trois autres sont \NEK{}.

\medskip
\textbf{L'assistant public.} Une revue des journaux de navigation montre
que 47 agents ont soumis des textes à un assistant d'IA générative public
en un mois, dont certains contenaient des noms et des numéros de compte de
clients. Aucune règle ne l'interdisait, aucune ne l'encadrait.

\medskip
\textbf{Analyse} (confiance élevée). Le premier usage est un risque de
\emph{tiers} et de \emph{modèle}~: il relève à la fois du cycle de vie des
prestataires et d'une surveillance des performances. Le second est un
risque de \emph{fuite}, créé par un usage non déclaré~: il appelle une
règle d'usage, une solution autorisée, et une sensibilisation. Ni l'un ni
l'autre ne figurait au registre.
\end{boxscenario}

\begin{boxobjet}{R-IA-01 · Fuite de données de clients par un assistant public}
\textbf{Valeur métier~:} la confidentialité des données des clients.

\textbf{Bien support~:} navigateurs des agents, assistant public non
contractualisé.

\textbf{Événement redouté~:} divulgation de données de clients à un tiers
sans contrat~; manquement possible aux obligations de protection des
données (loi n°~2024/017, à vérifier).

\textbf{Vraisemblance~:} 4 (usage constaté, 47 agents en un mois). Base~:
journaux de navigation (\OUI{}).

\textbf{Impact~:} 3 (atteinte à la confidentialité, réputation, risque
réglementaire). Base~: appréciation du DPD (\OUI{}).

\textbf{Traitement~:} règle d'usage de l'IA approuvée par le comité de
sécurité~; blocage de l'assistant public~; mise à disposition d'un assistant
contractualisé, sans conservation des données~; sensibilisation ciblée.

\textbf{Propriétaire~:} le directeur des opérations, avec le DPD.
\end{boxobjet}

% ============================================================================
\section{L'Audit Continu}
\label{sec:c15-continu}
% ============================================================================

\subsection{De l'échantillon à la population}

Le Chapitre~\ref{chap:grc-pratique} a présenté le contrôle continu de la
\emph{conformité}~: des tests automatiques sur les contrôles. L'\termecle{audit
continu}\index{audit continu} va plus loin~: il analyse les \emph{données}
et les \emph{transactions} elles-mêmes, sur toute la population et de
manière répétée, pour y détecter les anomalies. Là où la Chambre des
comptes examinait 40 paiements de l'ANPP, l'audit continu les examine tous,
chaque mois.

\begin{boxscenario}{Opération ANPP. Tous les paiements, chaque mois}
\textbf{Le test.} Chaque mois, un programme rapproche les paiements du
fichier des agents en activité et signale tout paiement à une personne
absente de ce fichier à la date du paiement. Le programme appartient à
l'audit interne de l'ANPP, qui le fait tourner sur une copie des données.

\medskip
\textbf{Juin 2029.} Le test signale 37 exceptions. Après examen~: 29 sont
justifiées (départs enregistrés avec retard, rappels de salaire dus), 6 sont
des paiements indus, et 2 restent sans explication, les dossiers n'ayant pas
été retrouvés.
\end{boxscenario}

\begin{boxlecture}{Les 37 exceptions de juin 2029}
\textbf{Constatation.} 29 justifiées (\OUI{}), 6 indues (\NON{}), 2 sans
explication (\NEK{})~:
\[
\mu \approx 0{,}78~;\quad U \approx 0{,}05~;\quad
H \approx 0{,}276 + 0{,}426 + 0{,}228 \approx 0{,}93~\text{bit}.
\]

\textbf{Analyse} (confiance moyenne). La plupart des exceptions sont des
faux positifs, et c'est normal~: un test exhaustif signale tout ce qui
s'écarte de la règle, y compris ce qui s'explique. Les six paiements indus
sont des constats établis sur toute la population, plus solides que ceux
qu'un échantillon de 40 aurait pu fournir. Les deux \NEK{} sont l'objet
même de l'audit~: ce sont les cas où l'explication manque.

\medskip
\hfill\lecture{0{,}78}{0{,}05}{0{,}93}
\end{boxlecture}

\subsection{Les conditions et les limites}

\begin{boxprocedure}
\textbf{Mettre en place un test d'audit continu.} Le papier de travail
décrit~:
\begin{itemize}
    \item la règle testée et son critère (texte, procédure)~;
    \item les sources de données, leur exhaustivité et leur qualité (le
          fichier des agents est-il à jour~? les paiements sont-ils tous
          extraits~?)~;
    \item le programme, sa version, et sa validation (essai sur des cas
          connus)~;
    \item la fréquence, le seuil d'exception et le circuit de traitement
          des exceptions~;
    \item qui détient le programme~: l'audit interne, et non la direction
          auditée, pour préserver l'indépendance.
\end{itemize}
\end{boxprocedure}

\begin{boxalert}
\textbf{Violation typique.} Confier le test d'audit continu à la direction
qui exécute les paiements, puis lire l'absence d'exception comme une preuve.

\smallskip
\textbf{Règle~:} un test d'audit continu ne vaut que par l'exhaustivité de
ses données et l'indépendance de celui qui le fait tourner. Le même
programme, géré par la direction, est un contrôle de premier ou de deuxième
niveau (Chapitre~\ref{chap:grc}), utile, mais qui ne remplace pas l'audit.
Et trop d'exceptions non traitées finissent par ne plus être lues.
\end{boxalert}

% ============================================================================
\section{Ce que les Nouveaux Objets demandent à chaque Métier}
\label{sec:c15-metiers}
% ============================================================================

\begin{boxmetier}{Responsable GRC}
\textbf{Un registre des tiers}

Tenir la liste des prestataires critiques, leur criticité, la preuve
attendue de chacun et son statut. Y faire entrer ceux qui ne sont pas
«~informatiques~», mais touchent à l'information.
\end{boxmetier}

\begin{boxmetier}{Auditeur en cybersécurité}
\textbf{Lire un rapport d'assurance}

Vérifier que le rapport de type~2 couvre le service utilisé, la période et
le lieu des données~; repérer les mesures que le fournisseur attend du
client (souvent listées dans le rapport) et les tester chez le client.
Pour l'IA, demander les preuves d'entraînement, de validation et de
surveillance du modèle.
\end{boxmetier}

\begin{boxmetier}{RSSI}
\textbf{Encadrer les usages avant de les interdire}

Un usage non déclaré de l'IA révèle un besoin. La RSSI propose une règle
d'usage, une solution autorisée et une sensibilisation~; le blocage seul
déplace l'usage ailleurs.
\end{boxmetier}

\begin{boxmetier}{Gestionnaire des risques cyber}
\textbf{Faire entrer l'IA dans le registre}

Chaque système d'IA en usage est inventorié, avec son propriétaire, ses
données, son fournisseur et ses scénarios de risque~: fuite, manipulation,
erreur, biais.
\end{boxmetier}

\separateur

\begin{boxresume}
\textbf{Ce qu'il faut retenir de ce chapitre~:}
\begin{itemize}
    \item \textbf{Les tiers} sont une porte d'entrée, un risque de
          défaillance et une chaîne~; le risque reste celui de
          l'organisation.
    \item \textbf{Le cycle de vie d'un tiers}~: sélection, contrat,
          exploitation, sortie, avec une preuve à chaque étape.
    \item \textbf{Dans le nuage}, les identités et les données restent
          toujours au client~; le rapport du fournisseur ne couvre que sa
          part.
    \item \textbf{L'IA} ajoute fuite, manipulation, erreur opaque et biais~;
          ISO/IEC~42001 et le NIST AI RMF l'encadrent.
    \item \textbf{L'audit continu} examine toute la population~; il vaut par
          l'exhaustivité des données et l'indépendance de son exploitant.
\end{itemize}
\end{boxresume}

\begin{boxinfo}
\textbf{Points de contrôle.} Vous avez compris ce chapitre si vous savez~:
\begin{itemize}
    \item répartir une mesure entre fournisseur et client pour un service
          donné~;
    \item dire ce qu'un rapport d'assurance prouve et ce qu'il laisse au
          client~;
    \item apprécier un usage de l'IA sous forme de fiche de risque~;
    \item concevoir un test d'audit continu et ses conditions de validité.
\end{itemize}
\end{boxinfo}

% ============================================================================
\section*{Exercices}
% ============================================================================

\begin{boxexo}[title={\ding{45}\quad Exercice~15.1. Qui en répond~? \niveaubase}]
La messagerie de la BANQUE SANAGA est un service en nuage de type logiciel
(SaaS). Pour chacune des tâches suivantes, dites si elle revient à NUAGEPRO
ou à la banque.
\begin{enumerate}[(a)]
    \item Appliquer les correctifs du système d'exploitation des serveurs.
    \item Activer le second facteur pour les agents.
    \item Désactiver le partage externe des documents.
    \item Protéger physiquement le centre de données.
    \item Désactiver le compte d'un agent parti.
    \item Conserver les journaux d'accès selon la loi.
\end{enumerate}
\end{boxexo}

\begin{boxexo}[title={\ding{45}\quad Exercice~15.2. Le contrat de FRAUDSCAN \niveaumoyen}]
\begin{enumerate}[(a)]
    \item Rédigez cinq clauses de sécurité pour le contrat du modèle de
          détection de fraude, en couvrant les quatre étapes du cycle de vie
          d'un tiers.
    \item Ajoutez deux clauses propres à l'IA (données d'entraînement,
          surveillance des performances).
    \item Quelles preuves la banque devra-t-elle obtenir chaque année pour
          lever les trois \NEK{} du modèle~?
\end{enumerate}
\end{boxexo}

\begin{boxexo}[title={\ding{45}\quad Exercice~15.3. L'IA chez LiZbiZ \niveaumoyen}]
LiZbiZ envisage un modèle qui optimise les tournées de livraison à partir
des données de suivi des expéditions et des adresses des clients.
\begin{enumerate}[(a)]
    \item Identifiez deux scénarios de risque propres à l'IA, et rédigez la
          fiche de l'un d'eux.
    \item Placez les actions de LiZbiZ dans les quatre fonctions du NIST AI
          RMF.
    \item PHARMACAM demande si ce modèle « est conforme à
          l'ISO/IEC~42001 ». Que répondez-vous~?
\end{enumerate}
\end{boxexo}

\begin{boxexo}[title={\ding{45}\quad Exercice~15.4. Un audit continu des ajustements de stock \niveauexpert}]
Le cours ACSI a établi que 148~millions FCFA d'ajustements de stock, dont
71~\% en fin de mois, avaient été enregistrés sous le compte de M.~FOUDA à
Douala.
\begin{enumerate}[(a)]
    \item Concevez un test d'audit continu qui aurait signalé ces
          ajustements~: règle, données, seuils, fréquence.
    \item Rédigez le papier de travail de mise en place (exhaustivité,
          validation, indépendance).
    \item Lors du premier mois, le test signale 120 exceptions. Comment
          éviter que les exceptions ne soient plus lues, sans perdre les
          vraies anomalies~?
    \item Un mois, 80 exceptions sont justifiées, 15 confirmées comme
          anormales, 5 sans explication. Calculez $\mu$, $U$ et $H$, et dites
          ce que chaque catégorie appelle.
\end{enumerate}
\end{boxexo}

\bigskip

\begin{boxinfo}
\textbf{Pour aller plus loin.}
\begin{itemize}
    \item L'ISO/IEC~27002:2022~\cite{iso27002_2022}, mesures 5.19 à 5.23.
    \item L'ISO/IEC~27017~\cite{iso27017_2015} et l'ISO/IEC~27018~\cite{iso27018_2019}.
    \item L'ISO/IEC~42001:2023~\cite{iso42001_2023} et le NIST AI
          RMF~\cite{nist2023airmf}.
\end{itemize}
\end{boxinfo}

\begin{boxinfo}
\textbf{Transition.} Tiers opaques, modèles qui apprennent, populations
entières analysées~: à chaque fois, la question de l'auditeur est la même~:
que sait-on, et que ne sait-on pas~? La Partie~VI, consacrée à la
recherche, montre comment l'auditologie fait de cette question un objet de
mesure. C'est le Chapitre~\ref{chap:auditologie}.
\end{boxinfo}


    % ======================================================================
    % PARTIE 6 — RECHERCHE : MESURER CE QUE L'ON NE SAIT PAS
    % Standards de référence : Articles A0 et A1 (HAL), A2 à A5 en préparation.
    % Niveau : recherche
    % ======================================================================
    \part{Recherche : Mesurer ce que l'on ne Sait pas}
    \label{part:recherche}

        % ============================================================================
% Fichier  : partie6/P06_C16_mesurer_inconnu.tex
% Titre    : Mesurer ce que l'on ne Sait pas
% Auteur   : MINKA MI NGUIDJOÏ Thierry Emmanuel
% Version  : 1.0 -- novembre 2026
% Corrigés : corriges/Corr_C16.tex
% ============================================================================
% Sources et statut (règles du corpus, mémoire « auditologie ») :
%  - A0 (hal-05701007) : mission <S, A, I, G, T>, six axiomes (énoncés repris
%    d'ACSI ch. 21, eux-mêmes repris d'A0), quatre propositions, cinq
%    piliers, falsifiabilité à la Lakatos.
%  - A1 (hal-05701270) : espace Omega_q, conditions C1 à C4 de NEK, mu, U,
%    H, D_KL, SUI et ses trois réserves.
%  - Décomposition de l'entropie : identité EXACTE de groupement (Shannon),
%    comme dans ACSI ch. 22 ; la prépublication A1 porte un terme « + U_q »
%    erroné (erratum à décider par l'auteur) : ne pas le reproduire.
%  - A2 à A5 : « en préparation », sans entrée bibliographique.
%  - Pistes propres à NRAC (incertitude épistémique et quantification FAIR,
%    base de preuve des registres) : présentées comme QUESTIONS OUVERTES,
%    non comme résultats.
% Valeurs calculées sur les cas du cours (aucun fait nouveau) : D_KL
% infrastructure LiZbiZ 2026 -> 2028 ≈ 0,27 bit ; agences SANAGA 2025 -> 2028
% ≈ 0,14 bit (ensembles non identiques) ; décompositions de H.
% ============================================================================

\chapter{Mesurer ce que l'on ne Sait pas}
\label{chap:auditologie}
\chaptermark{Mesurer l'inconnu}

\epigraph{%
    « Distinguer ce que l'on sait,\\
    ce que l'on a déposé, et ce qui reste à prouver. »%
}{Règle du chapitre}

\niveauChapitre{Recherche}{%
    Maîtrise du Chapitre~\ref{chap:conclure}%
}

\begin{boxobjectifs}
\begin{itemize}
    \item Présenter l'auditologie~: son objet, ses six axiomes, ses cinq
          piliers et sa méthode (article A0).
    \item Approfondir l'instrument \OUI{}~/ \NON{}~/ \NEK{} et ses métriques
          (article A1), et les appliquer au risque et à l'audit cyber.
    \item Décomposer l'entropie en ignorance et en désaccord, et comparer
          deux cycles par la divergence de Kullback-Leibler.
    \item Comprendre l'usage stratégique de \NEK{} et ses limites.
    \item Distinguer ce qui est établi, ce qui est déposé et ce qui est en
          préparation.
    \item Formuler un sujet de projet ou de mémoire falsifiable.
\end{itemize}
\end{boxobjectifs}

\bigskip

% ============================================================================
\section*{Avant de commencer : que vaut ce que nous avons mesuré~?}
% ============================================================================

Depuis le prologue, le cours a calculé des dizaines de valeurs de $\mu$,
$U$ et $H$~: sur les contrôles d'une banque, la base de preuve d'un
registre, les critères d'un écosystème, les obligations d'un organisme
public. Il a répété une règle~: on ne lit jamais une mesure seule, et l'on
dit toujours ce que l'on ignore.

Ce dernier chapitre change de point de vue. Il ne demande plus «~que
conclure sur la BANQUE SANAGA~?~», mais «~d'où viennent ces mesures, que
valent-elles, et que reste-t-il à établir~?~». Il présente le programme de
recherche dont le cours est l'application, l'\emph{auditologie}, en
distinguant soigneusement ce qui est publié, ce qui est en préparation, et
ce qui n'est encore qu'une question.

\begin{boxalert}
\textbf{Règle du chapitre.} Une affirmation de recherche se cite avec son
statut~: établie (théorème, résultat reproduit), déposée (prépublication
non encore évaluée par les pairs), en préparation (non publiée), ou
question ouverte. Le cours applique à son propre fondement l'exigence qu'il
fait à l'auditeur~: séparer ce que l'on sait de ce que l'on suppose.
\end{boxalert}

\separateur

% ============================================================================
\section{L'Auditologie}
\label{sec:c16-auditologie}
% ============================================================================

\subsection{L'objet}

L'article A0 part d'un constat~: l'audit des systèmes d'information dispose
de méthodes et de référentiels, mais d'aucun socle axiomatique ni d'aucun
espace de mesure unifié pour son objet central, la réduction de
l'incertitude~\cite{minka2026auditologie}.

\begin{boxdef}
\termecle{Mission auditologique}\index{mission auditologique} (A0)~:
quintuplet
$\mathcal{M} = \langle \mathcal{S}, \mathcal{A}, \mathcal{I},
\mathcal{G}, \mathcal{T} \rangle$, où $\mathcal{S}$ est le système audité,
$\mathcal{A}$ l'architecture cognitive de l'auditeur, $\mathcal{I}$ le
canal d'information par lequel les preuves parviennent à l'auditeur,
$\mathcal{G}$ la structure stratégique de l'interaction entre auditeur,
audité et tiers, et $\mathcal{T}$ la trajectoire du système au fil des
missions.

\medskip
\termecle{Auditologie}\index{auditologie}~: science appliquée de la
réduction de l'incertitude sur les missions auditologiques.
\end{boxdef}

\begin{tableaucours}{cL{3.4cm}L{8.6cm}}{%
    \tableauHeader{} &
    \tableauHeader{Composante} &
    \tableauHeader{Dans la mission de Mme~EKAMBI (mars 2028)}
}
$\mathcal{S}$ & Système audité & Les 38 agences, leurs postes, réseaux,
agents, et les prestataires qui y interviennent \\
$\mathcal{A}$ & Architecture cognitive & L'équipe d'audit, sa décomposition
en six sous-domaines, ses huit tests \\
$\mathcal{I}$ & Canal d'information & Extractions, essais, exercices~; et ce
qui ne passe pas~: les configurations des GAB \\
$\mathcal{G}$ & Structure stratégique & Comité des risques, RSSI, DSI,
direction du réseau d'agences, prestataire des GAB \\
$\mathcal{T}$ & Trajectoire & Prologue (2026), mission de 2028, suivi \\
\end{tableaucours}
\vspace{-0.8em}
\begin{center}{\small\itshape Tableau~16.1. Une mission du cours comme
mission auditologique}\end{center}

\subsection{Six axiomes}

L'article A0 pose six axiomes, dont le cours a appliqué les conséquences
depuis le prologue. Les énoncés sont ceux de l'article.

\begin{axiome}[Observabilité partielle]
Pour tout système audité $\mathcal{S}$ non dégénéré, l'auditeur n'a accès
qu'à un sous-ensemble strict des états pertinents de $\mathcal{S}$, et ce
sous-ensemble est lui-même partiellement opaque.
\end{axiome}

\begin{axiome}[Cognition limitée des agents]
L'auditeur et l'audité sont des agents cognitivement limités, au sens de
Simon, de Miller et de Kahneman~\cite{simon1957,miller1956,kahneman2011}.
\end{axiome}

\begin{axiome}[Mesurabilité de l'incertitude épistémique]
L'incertitude épistémique de $\mathcal{A}$ concernant des propositions
spécifiques sur $\mathcal{S}$ admet une mesure quantitative, définissable
par l'entropie informationnelle et ses généralisations.
\end{axiome}

\begin{axiome}[Interaction stratégique]
L'auditeur et l'audité ont des fonctions d'utilité distinctes, et leurs
actions sont, en général, mieux modélisées comme des mouvements dans un jeu
à information imparfaite et asymétrique.
\end{axiome}

\begin{axiome}[Individuation temporelle]
Le système audité $\mathcal{S}$ évolue dans le temps, et sa trajectoire est
dépendante du chemin~: l'état à la mission $n+1$ dépend de l'état à $n$ et
des interventions déclenchées par cette mission antérieure.
\end{axiome}

\begin{axiome}[Opposabilité]
Toute conclusion produite par une mission auditologique doit, en principe,
pouvoir être défendue devant un tiers disposant de l'autorité légitime pour
la contester.
\end{axiome}

\begin{tableaucours}{L{3.6cm}L{8.6cm}}{%
    \tableauHeader{Axiome} &
    \tableauHeader{Où le cours l'a rencontré}
}
Observabilité partielle & Le certificat ne voit que son périmètre
(prologue)~; les GAB que l'on ne peut pas voir (Chapitre~\ref{chap:audit}) \\
Cognition limitée & La borne de sept et la décomposition
(Chapitre~\ref{chap:conclure}) \\
Mesurabilité & $\mu$, $U$, $H$ dans chaque chapitre \\
Interaction stratégique & Les pressions sur l'auditeur
(Chapitre~\ref{chap:droit})~; l'attaquant qui s'adapte
(Chapitre~\ref{chap:risque}) \\
Individuation temporelle & Les passages successifs de LiZbiZ
(Chapitre~\ref{chap:audit-technique})~; la correction qui ne réécrit pas le
passé (Chapitre~\ref{chap:audit}) \\
Opposabilité & Le papier de travail, la preuve de rang~1, la citation
vérifiée \\
\end{tableaucours}
\vspace{-0.8em}
\begin{center}{\small\itshape Tableau~16.2. Les axiomes dans le cours}\end{center}

A0 en tire quatre conséquences, que le cours a mises en pratique~: un
instrument d'audit doit comporter une modalité explicite de connaissance
insuffisante (\NEK{})~; un espace d'audit non trivial doit être décomposé
en respectant la borne cognitive~; les conclusions ne se testent que par
l'observation longitudinale~; et l'audit fait lui-même partie de la
trajectoire du système qu'il observe.

\subsection{Cinq piliers et une méthode}

Les axiomes disent ce qu'est une mission~; cinq piliers théoriques
permettent de la modéliser~: la cognition limitée, la théorie de
l'information~\cite{shannon1948,kullback1951}, la théorie des jeux, les
systèmes dynamiques~\cite{norris1998} et l'ontologie
critique-réaliste~\cite{bhaskar1975}. L'auditologie se veut une science du
design~\cite{hevner2004design}~: elle construit des artefacts (instruments,
métriques, grilles) et les évalue dans leur environnement. Elle adopte enfin
le falsificationnisme de Lakatos~\cite{lakatos1970}~: un noyau dur
(axiomes et piliers) et une ceinture protectrice d'hypothèses et de
prédictions testables, qui peuvent être réfutées.

% ============================================================================
\section{L'Instrument et ses Métriques}
\label{sec:c16-instrument}
% ============================================================================

\subsection{Quatre conditions pour \NEK{}}

L'article A1 formalise l'instrument trichotomique~\cite{minka2026nek}. Pour
chaque question $q$, l'espace des réponses est
$\Omega_q = \{\OUI{}, \NON{}, \NEK{}\}$, et A1 précise quatre conditions
qui distinguent un \NEK{} légitime d'un mauvais usage.

\begin{tableaucours}{L{3.4cm}L{4.8cm}L{4.0cm}}{%
    \tableauHeader{Condition (A1)} &
    \tableauHeader{Sens} &
    \tableauHeader{Dans le cours}
}
C1. Spécificité de l'objet & \NEK{} porte sur l'objet précis de la
question & Questions atomiques (Chapitre~\ref{chap:conclure}) \\
C2. Asymétrie d'accès & L'information existe, mais le répondant n'y a pas
accès & Le visa de l'Agence, les configurations des GAB \\
C3. Non-report & Si l'on peut vérifier dans le délai, on répond & Les \NEK{}
levés avant l'audit de l'ANTIC (Chapitre~\ref{chap:droit}) \\
C4. Neutralité & Aucun barème ne rend \NEK{} équivalent à \OUI{} ou à
\NON{} & Pas de «~demi-point~» \\
\end{tableaucours}
\vspace{-0.8em}
\begin{center}{\small\itshape Tableau~16.3. Les conditions de la réponse
\NEK{} (d'après A1)}\end{center}

\subsection{Décomposer l'entropie}

L'entropie trichotomique mêle deux sources de dispersion~: l'ignorance, et
le désaccord entre ce que l'on sait. La propriété de groupement de
l'entropie~\cite{shannon1948} les sépare exactement. En notant
$H_b(x) = -x\log_2 x - (1-x)\log_2(1-x)$ l'entropie binaire~:
\[
    H \;=\; H_b(U) \;+\; (1 - U)\, H_b\!\left(\frac{\mu}{1 - U}\right).
\]
Le premier terme est la part de l'ignorance~; le second, pondéré par la
part de ce qui est connu, est la part du désaccord entre \OUI{} et \NON{}.

\begin{boxlecture}{Ignorance ou désaccord~? Cinq mesures du cours}
\begin{tableaucours}{L{3.6cm}cccccc}{%
    \tableauHeader{Ensemble} &
    \tableauHeader{O/N/K} & \tableauHeader{$U$} &
    \tableauHeader{$H_b(U)$} & \tableauHeader{Désacc.} &
    \tableauHeader{$H$} & \tableauHeader{Ignor.}
}
Agences, prologue & 2/4/4 & 0,40 & 0,97 & 0,55 & 1,52 & 64~\% \\
Registre 2025 (Chap.~\ref{chap:risque}) & 6/2/4 & 0,33 & 0,92 & 0,54 & 1,46 &
63~\% \\
Déclaration d'applicabilité (Chap.~\ref{chap:iso27000}) & 3/3/3 & 0,33 & 0,92 &
0,67 & 1,58 & 58~\% \\
Mission de mars 2028 (Chap.~\ref{chap:audit}) & 3/3/2 & 0,25 & 0,81 & 0,75 &
1,56 & 52~\% \\
Infrastructure LiZbiZ 2028 (Chap.~\ref{chap:audit-technique}) & 14/4/2 & 0,10 &
0,47 & 0,69 & 1,16 & 41~\% \\
\end{tableaucours}

\medskip
\textbf{Constatation.} Les cinq ensembles ont une entropie proche de
1,5~bit, sauf le dernier~; mais la part de l'ignorance décroît de 64~\% à
41~\% à mesure que le cours avance dans le temps et que les preuves
s'accumulent.

\medskip
\textbf{Analyse} (confiance élevée sur le calcul, faible sur la tendance~:
les ensembles mesurés ne sont pas identiques). Une même entropie peut
appeler deux actions opposées~: chercher la preuve quand l'ignorance
domine, chercher la cause du désaccord quand elle ne domine plus. La
décomposition rend quantitative la règle «~$H$ élevé, $U$ faible~» du
Chapitre~\ref{chap:conclure}.
\end{boxlecture}

\subsection{Comparer deux cycles}

L'article A1 propose de comparer deux cycles par la \termecle{divergence de
Kullback-Leibler}\index{divergence de Kullback-Leibler}~\cite{kullback1951},
calculée dans le sens du temps~:
\[
    D_{\mathrm{KL}}\!\left(p^{(t+1)} \,\middle\|\, p^{(t)}\right) =
    \sum_{m} p^{(t+1)}(m) \log_2 \frac{p^{(t+1)}(m)}{p^{(t)}(m)}.
\]
Elle vaut zéro si la distribution n'a pas changé. Elle exige que toute
modalité observée au cycle $t+1$ l'ait été au cycle $t$.

\begin{boxscenario}{Opération LIZBIZ. L'infrastructure entre 2026 et 2028}
\textbf{Calcul.} En 2026~: $(0{,}40~;~0{,}45~;~0{,}15)$~; en 2028~:
$(0{,}70~;~0{,}20~;~0{,}10)$.
\[
D_{\mathrm{KL}} = 0{,}7\log_2\tfrac{0{,}7}{0{,}4} + 0{,}2\log_2\tfrac{0{,}2}{0{,}45}
  + 0{,}1\log_2\tfrac{0{,}1}{0{,}15} \approx 0{,}565 - 0{,}234 - 0{,}059
  \approx 0{,}27~\text{bit}.
\]

\textbf{Analyse} (confiance moyenne). Le changement est net, porté surtout
par le passage de \NON{} à \OUI{}. Mais A0 rappelle que deux cycles ne
distinguent pas encore une amélioration réelle d'un artefact de mesure~:
les exports de 2028 ne couvrent pas exactement les mêmes équipements qu'en
2026. C'est l'objet des travaux sur la trajectoire (Section~\ref{sec:c16-statut}).
\end{boxscenario}

% ============================================================================
\section{L'Usage Stratégique de \NEK{}}
\label{sec:c16-sui}
% ============================================================================

Un \NEK{} est \emph{authentique} s'il traduit une vraie limite de
connaissance~; il est \emph{stratégique} si le répondant sait, mais évite
les conséquences d'une réponse nette. Pour une réponse isolée, les deux
sont indiscernables. A1 propose de les distinguer sur le \emph{profil} d'un
répondant~: l'ignorance authentique ne dépend pas, en moyenne, de l'enjeu
de la question~; l'ignorance stratégique se concentre sur les questions à
fort enjeu~\cite{minka2026nek}.

Pour chaque question $q$, un score d'enjeu $s_q \in [0,1]$ est fixé
\emph{avant} la collecte par des experts indépendants. Pour chaque
répondant $r$, A1 estime un modèle logistique
\[
    \log \frac{\Pr[\NEK{}_{r,q} = 1]}{1 - \Pr[\NEK{}_{r,q} = 1]}
    \;=\; \alpha_r + \beta_r\, s_q + \gamma_r^{\top} \mathbf{z}_q,
\]
où $\mathbf{z}_q$ contient des variables de contrôle, et définit l'indice
d'usage stratégique $\mathrm{SUI}_r = \hat\beta_r$.

\begin{boxscenario}{Un atelier de cotation des risques}
Lors d'un atelier de cotation, un directeur répond \NEK{} à toutes les
questions qui portent sur son propre périmètre et dont l'enjeu est élevé
(acceptation d'un risque intolérable, défaut d'un prestataire qu'il a
choisi), et répond nettement partout ailleurs. Le profil est celui que
l'indice cherche à détecter~; il ne suffit pas à le calculer~: il faut des
scores d'enjeu fixés à l'avance par d'autres, et assez de questions par
répondant.
\end{boxscenario}

\begin{boxalert}
\textbf{Trois réserves d'A1, à ne jamais oublier.}
\begin{enumerate}
    \item Un indice positif signale une corrélation entre \NEK{} et enjeu,
          qui peut avoir d'autres causes (questions plus difficiles, accès
          plus restreint).
    \item Le score d'enjeu doit être fixé hors des répondants~; sinon,
          l'indice n'est pas identifié.
    \item L'indice décrit une population, \emph{pas} une faute
          individuelle~: il n'établit pas l'intention et ne fonde aucune
          sanction (Chapitre~\ref{chap:droit}, «~constater n'est pas
          qualifier~»).
\end{enumerate}
\end{boxalert}

% ============================================================================
\section{Établi, Déposé, en Préparation}
\label{sec:c16-statut}
% ============================================================================

\subsection{Les travaux en préparation}

Quatre articles prolongent le programme. Ils ne sont pas publiés~; le
cours les présente d'après leur objet, sans en citer de résultat comme
acquis.
\begin{itemize}
    \item \textbf{La trajectoire (A2)}~: un score d'amélioration entre
          cycles, un descripteur topologique des trajectoires et une chaîne
          de Markov sur les niveaux de maturité, pour distinguer une
          amélioration réelle d'un artefact de mesure.
    \item \textbf{Le signalement stratégique (A3)}~: l'interaction entre
          auditeur et audité comme un jeu de signalement, où \NEK{} est un
          signal~; extension à l'audit de sécurité face à un attaquant.
    \item \textbf{La décomposition (A4)}~: la justification formelle du
          gabarit canonique à huit niveaux, sous une borne cognitive et des
          conditions d'exhaustivité et de non-redondance.
    \item \textbf{L'instrument (A5)}~: un outil logiciel qui génère les
          grilles à partir du gabarit et calcule les mesures en continu.
\end{itemize}

\begin{tableaucours}{L{4.0cm}L{3.4cm}L{4.8cm}}{%
    \tableauHeader{Élément} &
    \tableauHeader{Statut} &
    \tableauHeader{Ce qui reste à faire}
}
Entropie, divergence, groupement & Établi (théorie de l'information) &
· \\
Programme de l'auditologie (A0) & Prépublication HAL, 2026 & Discussion par
la communauté~; tests des prédictions \\
Instrument, métriques, indice d'usage stratégique (A1) & Prépublication HAL,
2026 & Expérience comparant instruments binaire et trichotomique \\
A2, A3, A4, A5 & En préparation & Dépôt, puis validation \\
Applications de ce cours au risque cyber (base de preuve des registres,
couplage avec FAIR) & Questions ouvertes & Formulation et test (sujets
ci-dessous) \\
Cas SANAGA, LiZbiZ, ANPP & Fictifs, pédagogiques & Ils illustrent~; ils ne
prouvent rien \\
\end{tableaucours}
\vspace{-0.8em}
\begin{center}{\small\itshape Tableau~16.4. Le statut de ce que le cours a
présenté}\end{center}

\subsection{Questions ouvertes pour le risque cyber}

Ce cours a appliqué l'instrument à des objets que les articles ne traitent
pas directement. Ces applications ouvrent des questions, que le cours ne
tranche pas.
\begin{itemize}
    \item \textbf{La base de preuve des registres}~: coter \OUI{}~/ \NON{}~/
          \NEK{} la base de chaque cote d'un registre des risques
          (Chapitre~\ref{chap:risque}) améliore-t-il les décisions de
          traitement~? Les registres dont $U$ est élevé sont-ils suivis de
          plus d'incidents imprévus~?
    \item \textbf{Le couplage avec la quantification}~: comment combiner la
          distribution des pertes de FAIR (incertitude aléatoire, sous
          hypothèses) et la cote des hypothèses (incertitude épistémique),
          au-delà de l'analyse de sensibilité du Chapitre~\ref{chap:fair}~?
    \item \textbf{L'écosystème}~: la concentration des \NEK{} sur la
          maturité et la confiance des parties prenantes
          (Chapitre~\ref{chap:ebios}) est-elle générale~? Que devient le
          niveau de menace quand on la prend en compte~?
    \item \textbf{Les tests automatiques}~: un contrôle continu silencieux
          (Chapitre~\ref{chap:grc-pratique}) doit-il voir sa cote se
          dégrader en \NEK{} avec le temps, et à quel rythme~?
\end{itemize}

% ============================================================================
\section{Formuler un Sujet de Recherche}
\label{sec:c16-sujet}
% ============================================================================

\begin{boxprocedure}
\textbf{Un sujet de mémoire falsifiable.} Un sujet recevable précise~:
\begin{itemize}
    \item la \textbf{question}, en une phrase, et l'axiome ou le pilier
          auquel elle se rattache~;
    \item une \textbf{hypothèse} énoncée de sorte que des données puissent
          la réfuter («~les registres dont $U \geq 0{,}30$ sont suivis de
          plus d'incidents non prévus que les autres~»)~;
    \item les \textbf{données}~: lesquelles, chez qui, avec quelle
          autorisation et quelle protection des personnes et des
          organisations (Chapitre~\ref{chap:droit})~;
    \item la \textbf{méthode} de mesure et de comparaison~;
    \item ce qui \textbf{réfuterait} l'hypothèse, écrit avant la collecte~;
    \item la \textbf{limite} principale, et la manière dont on la contrôle.
\end{itemize}
\end{boxprocedure}

\begin{boxinfo}
\textbf{Pourquoi le statut compte.} L'auditologie s'applique à elle-même ce
qu'elle demande aux auditeurs. Un cours qui présenterait comme acquis des
travaux en préparation ferait exactement ce qu'il reproche au registre de
2025~: afficher une cote sans base. Un étudiant qui reprend ces travaux
dans un mémoire cite les prépublications comme telles, et ne leur attribue
que ce qu'elles contiennent.
\end{boxinfo}

\separateur

\begin{boxresume}
\textbf{Ce qu'il faut retenir de ce chapitre~:}
\begin{itemize}
    \item \textbf{L'auditologie} étudie la réduction de l'incertitude sur une
          mission $\langle \mathcal{S}, \mathcal{A}, \mathcal{I},
          \mathcal{G}, \mathcal{T} \rangle$, sous six axiomes.
    \item \textbf{L'instrument} \OUI{}~/ \NON{}~/ \NEK{} suppose quatre
          conditions~; ses mesures se lisent ensemble.
    \item \textbf{L'entropie se décompose} exactement en ignorance et en
          désaccord~; les deux appellent des actions différentes.
    \item \textbf{La divergence} compare deux cycles, sans distinguer encore
          progrès réel et artefact de mesure.
    \item \textbf{L'usage stratégique de \NEK{}} se lit sur un profil, avec
          trois réserves~; il ne fonde aucune sanction.
    \item \textbf{Tout n'a pas le même statut}~: établi, déposé, en
          préparation, question ouverte.
\end{itemize}
\end{boxresume}

\begin{boxinfo}
\textbf{Points de contrôle.} Vous avez compris ce chapitre si vous savez~:
\begin{itemize}
    \item énoncer les six axiomes et dire où le cours les a appliqués~;
    \item décomposer une entropie et dire quelle action elle appelle~;
    \item calculer une divergence entre deux cycles et en dire la limite~;
    \item formuler une hypothèse de recherche réfutable.
\end{itemize}
\end{boxinfo}

% ============================================================================
\section*{Exercices}
% ============================================================================

\begin{boxexo}[title={\ding{45}\quad Exercice~16.1. Quel axiome~? \niveaubase}]
Associez chaque situation à l'axiome de A0 qu'elle illustre le mieux.
\begin{enumerate}[(a)]
    \item Le prestataire des GAB refuse l'accès aux configurations.
    \item Une grille de 25 questions sans décomposition produit des oublis.
    \item La DSI corrige la journalisation pendant la mission~; le constat
          est maintenu pour la période.
    \item Un directeur répond \NEK{} sur les questions qui engagent sa
          responsabilité.
    \item Chaque cote est accompagnée de la preuve qui permet de la défendre
          devant le comité ou le superviseur.
\end{enumerate}
\end{boxexo}

\begin{boxexo}[title={\ding{45}\quad Exercice~16.2. Décomposer \niveaumoyen}]
\begin{enumerate}[(a)]
    \item Décomposez l'entropie de la base commune de l'ANPP (4/3/3,
          Chapitre~\ref{chap:referentiels}) et du registre des obligations de
          la BANQUE SANAGA (7/2/3, Chapitre~\ref{chap:grc-pratique}).
    \item Pour chacun, dites si l'action prioritaire est de chercher la
          preuve ou de chercher la cause d'un désaccord.
    \item Montrez, à partir de la formule, qu'une ignorance unanime donne
          $H = 0$.
\end{enumerate}
\end{boxexo}

\begin{boxexo}[title={\ding{45}\quad Exercice~16.3. Deux cycles \niveaumoyen}]
\begin{enumerate}[(a)]
    \item Calculez la divergence entre les agences de la BANQUE SANAGA au
          prologue (2/4/4) et lors de la mission de mars 2028 (3/3/2).
    \item Comparez-la à celle de l'infrastructure de LiZbiZ (0,27~bit).
    \item Pourquoi ces deux divergences ne sont-elles pas directement
          comparables~? Citez au moins deux raisons.
\end{enumerate}
\end{boxexo}

\begin{boxexo}[title={\ding{45}\quad Exercice~16.4. Un sujet de mémoire \niveauexpert}]
Choisissez l'une des questions ouvertes de la Section~\ref{sec:c16-statut}
et rédigez un projet de sujet de mémoire d'une page, selon la procédure de
la Section~\ref{sec:c16-sujet}~: question, hypothèse réfutable, données et
autorisations, méthode, condition de réfutation, limite principale. Dites
quels travaux vous citez, et avec quel statut.
\end{boxexo}

\bigskip

\begin{boxinfo}
\textbf{Pour aller plus loin.}
\begin{itemize}
    \item L'article \emph{Vers l'auditologie}~\cite{minka2026auditologie}.
    \item L'article sur l'incertitude épistémique~\cite{minka2026nek}.
    \item Shannon~\cite{shannon1948} et Kullback et
          Leibler~\cite{kullback1951}, pour les mesures.
    \item Lakatos~\cite{lakatos1970} et Hevner et
          collègues~\cite{hevner2004design}, pour la méthode.
\end{itemize}
\end{boxinfo}

\begin{boxinfo}
\textbf{Pour conclure le cours.} Le prologue s'ouvrait sur un certificat
accroché au mur et des guichets fermés. Tout ce qui a suivi revient à une
seule exigence~: savoir ce que l'on sait, savoir ce que l'on ignore, et le
dire à ceux qui décident. Les normes, les méthodes de risque, les audits et
le droit en sont les instruments~; la mesure de l'inconnu en est la
discipline.
\end{boxinfo}


% --------------------------------------------------------------------------
% ANNEXES
% --------------------------------------------------------------------------
\appendix

    % ============================================================================
% Fichier  : annexes/Ann_A_glossaire.tex
% Titre    : Glossaire
% Version  : 1.0 -- novembre 2026
% ============================================================================
% Définitions reprises des encadrés « Définition » des chapitres (termes
% marqués \termecle), abrégées ; le chapitre de référence est indiqué.
% ============================================================================

\chapter{Glossaire}
\label{ann:glossaire}

Ce glossaire reprend, sous une forme abrégée, les définitions données dans
les chapitres. Le numéro renvoie au chapitre où le terme est défini et
illustré~; c'est là qu'il faut lire la définition complète.

\begin{description}
    \item[Acceptation du risque] Décision argumentée, prise par l'instance
          qui en a l'autorité, de ne pas traiter davantage un risque
          résiduel, pour une durée et à des conditions données
          (Ch.~\ref{chap:iso27005}).
    \item[Appétence au risque] Niveau de risque qu'une organisation accepte
          pour atteindre ses objectifs~; fixée par l'organe de direction
          (Ch.~\ref{chap:grc}, \ref{chap:risque}).
    \item[Attributs de mesure] Étiquettes associées à chaque mesure de
          l'ISO/IEC~27002:2022 (type, propriétés, concepts, capacités,
          domaines), qui permettent de trier les mesures
          (Ch.~\ref{chap:iso27000}).
    \item[Audit] Processus méthodique, indépendant et documenté, qui obtient
          des preuves et les évalue au regard de critères
          (Ch.~\ref{chap:audit}).
    \item[Audit continu] Analyse répétée de toute une population de
          données ou de transactions pour y détecter les anomalies
          (Ch.~\ref{chap:emergents}).
    \item[Auditologie] Science appliquée de la réduction de l'incertitude
          sur les missions auditologiques (A0~; Ch.~\ref{chap:auditologie}).
    \item[Base commune de contrôles] Liste de contrôles propres à
          l'organisation, chacun relié aux référentiels qui l'exigent~;
          testée une fois, projetée sur plusieurs référentiels
          (Ch.~\ref{chap:referentiels}).
    \item[Base de preuve] Ce sur quoi repose une cote~: mesures internes,
          données externes, avis motivé, ou rien~; elle se cote elle-même
          \OUI{}, \NON{} ou \NEK{} (Ch.~\ref{chap:risque}).
    \item[Bien support] Composant sur lequel repose une valeur métier~:
          matériel, logiciel, réseau, personne, prestataire
          (Ch.~\ref{chap:risque}).
    \item[Certificat] Attestation, par un organisme tiers accrédité, qu'un
          système de management satisfait une norme, sur un périmètre et
          pour une durée donnés (Ch.~\ref{chap:prologue}).
    \item[Conclusion d'audit] Résultat de l'audit, tiré de l'ensemble des
          constats, au regard des objectifs de la mission
          (Ch.~\ref{chap:audit}).
    \item[Conformité] Satisfaction prouvée d'une exigence (norme, loi,
          contrat, politique) (Ch.~\ref{chap:prologue}, \ref{chap:grc}).
    \item[Constat d'audit] Résultat de l'évaluation des preuves au regard
          des critères~: conforme, non conforme, ou \NEK{}
          (Ch.~\ref{chap:audit}).
    \item[Contradictoire] Présentation des constats à l'audité, qui peut
          contester un fait, apporter une preuve ou expliquer une cause
          (Ch.~\ref{chap:audit}).
    \item[Contrôle continu de la conformité] Test automatique et répété
          d'un contrôle par une extraction comparée à une règle
          (Ch.~\ref{chap:grc-pratique}).
    \item[Couple SR/OV] Association d'une source de risque et d'un objectif
          visé, dans EBIOS Risk Manager (Ch.~\ref{chap:ebios}).
    \item[Criticité] Produit vraisemblance $\times$ impact, de 1 à 16~; elle
          classe, elle ne décide pas (Ch.~\ref{chap:risque}).
    \item[Critère d'audit] Exigence de référence~: norme, politique, loi,
          contrat, plan (Ch.~\ref{chap:audit}).
    \item[Critères d'acceptation] Règles, fixées à l'avance, qui disent quel
          résiduel peut être accepté, par qui, à quelles conditions
          (Ch.~\ref{chap:iso27005}).
    \item[Déclaration d'applicabilité] Document exigé par l'ISO/IEC~27001 qui
          recense les mesures nécessaires, justifie leur inclusion et
          l'exclusion des autres, et indique leur mise en œuvre
          (Ch.~\ref{chap:iso27000}).
    \item[Divergence de Kullback-Leibler] Mesure du changement de la
          distribution des réponses entre deux cycles
          (Ch.~\ref{chap:auditologie}).
    \item[Échantillon raisonné] Échantillon choisi selon les risques~; il
          trouve des défauts, il ne permet pas d'extrapoler un taux
          (Ch.~\ref{chap:audit}).
    \item[Entropie] Dispersion des réponses \OUI{}, \NON{}, \NEK{}, en bits,
          de 0 à $\log_2 3$~; elle se décompose en ignorance et désaccord
          (Ch.~\ref{chap:conclure}, \ref{chap:auditologie}).
    \item[Estimation calibrée] Minimum, valeur la plus probable et maximum,
          choisis pour encadrer la vraie valeur avec une confiance d'environ
          90~\% (Ch.~\ref{chap:fair}).
    \item[Événement redouté] Atteinte à une valeur métier, exprimée du point
          de vue de l'organisation (Ch.~\ref{chap:risque}).
    \item[Export figé] Copie d'une configuration ou d'une liste, extraite à
          une date et scellée par une empreinte (Ch.~\ref{chap:audit-technique}).
    \item[FAIR] Modèle de quantification du risque d'information~: fréquence
          des événements de perte et ampleur des pertes (Ch.~\ref{chap:fair}).
    \item[Fiche d'assertion] Description d'un contrôle par son objectif, son
          critère, sa question, son risque et sa conséquence
          (Ch.~\ref{chap:conclure}).
    \item[Gabarit canonique] Décomposition à huit niveaux~: domaine,
          sous-domaine, processus, activité, contrôle, test, preuve, trace
          (A0~; Ch.~\ref{chap:conclure}).
    \item[Gestion des risques] Activités coordonnées pour identifier,
          analyser, évaluer, traiter et suivre les risques
          (Ch.~\ref{chap:grc}, \ref{chap:iso27005}).
    \item[Gouvernance de la sécurité] Responsabilités et pratiques par
          lesquelles l'organe de direction fixe les objectifs, le niveau de
          risque acceptable, les rôles, et surveille les résultats
          (Ch.~\ref{chap:grc}).
    \item[GRC] Gouvernance, risque et conformité pilotés comme un seul
          dispositif, avec un référentiel commun (Ch.~\ref{chap:grc},
          \ref{chap:grc-pratique}).
    \item[Incertitude aléatoire] Part de l'incertitude qui tient au hasard
          des événements~; elle se décrit, elle ne se réduit pas
          (Ch.~\ref{chap:risque}).
    \item[Incertitude épistémique] Part de l'incertitude qui tient à ce que
          l'on ne sait pas~; elle se réduit par des preuves
          (Ch.~\ref{chap:risque}).
    \item[Indicateur de risque] Signal qui annonce un changement
          d'exposition, lu avec des seuils qui déclenchent une revue
          (Ch.~\ref{chap:grc-pratique}).
    \item[Matrice RACI] Tableau des rôles~: réalise, répond, consulté,
          informé~; un seul A par activité (Ch.~\ref{chap:grc}).
    \item[Matrice risques-contrôles] Tableau qui croise contrôles et
          risques, avec la dernière cote de chaque contrôle
          (Ch.~\ref{chap:grc-pratique}).
    \item[Menace] Cause potentielle d'un incident (Ch.~\ref{chap:risque}).
    \item[Mission auditologique] Quintuplet $\langle \mathcal{S},
          \mathcal{A}, \mathcal{I}, \mathcal{G}, \mathcal{T} \rangle$ (A0~;
          Ch.~\ref{chap:auditologie}).
    \item[Modèle des trois lignes] Répartition des rôles entre organe de
          gouvernance, première ligne, deuxième ligne et audit interne (IIA,
          2020~; Ch.~\ref{chap:grc}).
    \item[\NEK{}] Réponse «~connaissance insuffisante~»~: ni \OUI{}, ni
          \NON{}~; quatre conditions la rendent légitime (A1~;
          Ch.~\ref{chap:prologue}, \ref{chap:auditologie}).
    \item[Non-conformité majeure] Non-conformité qui affecte la capacité du
          système de management à atteindre ses résultats
          (Ch.~\ref{chap:smsi}).
    \item[Non-conformité mineure] Non-conformité qui n'affecte pas cette
          capacité (Ch.~\ref{chap:smsi}).
    \item[Objectif visé] Ce que la source de risque cherche à obtenir
          (Ch.~\ref{chap:risque}, \ref{chap:ebios}).
    \item[Organe de direction] Personne ou groupe qui répond en dernier
          ressort de l'organisation~: le conseil d'administration
          (Ch.~\ref{chap:grc}).
    \item[Périmètre de complaisance] Périmètre choisi pour faciliter la
          certification plutôt que pour couvrir les risques
          (Ch.~\ref{chap:prologue}).
    \item[Piste de progrès] Constat qui n'est pas une non-conformité, mais
          une occasion d'améliorer (Ch.~\ref{chap:smsi}).
    \item[Preuve d'audit] Information vérifiable et pertinente au regard des
          critères~; sa fiabilité se classe en cinq rangs
          (Ch.~\ref{chap:audit}).
    \item[Propriétaire de risque] Personne qui a la responsabilité et
          l'autorité de gérer un risque (Ch.~\ref{chap:grc}).
    \item[Rapport d'assurance] Rapport d'un auditeur indépendant sur les
          contrôles d'un prestataire (ISAE~3402, SOC~2)~; seul le type~2 dit
          s'ils ont fonctionné (Ch.~\ref{chap:referentiels}).
    \item[Registre des obligations] Liste des exigences qui s'imposent à
          l'organisation, avec leur source, leurs contrôles, leur preuve et
          leur statut (Ch.~\ref{chap:grc-pratique}).
    \item[Risque] Effet de l'incertitude sur les objectifs (ISO~31000~;
          Ch.~\ref{chap:risque}).
    \item[Risque cyber] Risque dont la source est une action sur les
          systèmes ou les données, et dont les conséquences touchent les
          activités (Ch.~\ref{chap:risque}).
    \item[Risque inhérent] Niveau du risque avant les mesures
          (Ch.~\ref{chap:risque}).
    \item[Risque résiduel] Niveau du risque compte tenu des seules mesures
          dont l'efficacité est établie (Ch.~\ref{chap:risque}).
    \item[Scénario opérationnel] Suite d'actions élémentaires (connaître,
          rentrer, trouver, exploiter) qui réalise un scénario stratégique
          (Ch.~\ref{chap:ebios}).
    \item[Scénario stratégique] Chemin d'attaque de haut niveau, d'un couple
          SR/OV à un événement redouté, éventuellement par l'écosystème
          (Ch.~\ref{chap:ebios}).
    \item[Sécurité effective] Capacité réelle de l'organisation à résister
          aux attaques et à s'en relever (Ch.~\ref{chap:prologue}).
    \item[Simulation de Monte-Carlo] Tirage d'un grand nombre d'années selon
          les estimations, pour obtenir la distribution des pertes
          (Ch.~\ref{chap:fair}).
    \item[SMSI] Système de management de la sécurité de l'information~:
          politiques, processus, rôles et contrôles sur un périmètre, en
          amélioration continue (Ch.~\ref{chap:prologue}, \ref{chap:smsi}).
    \item[Source de risque] Élément, personne ou groupe susceptible
          d'engendrer le risque (Ch.~\ref{chap:risque}).
    \item[Table de correspondance] Tableau qui relie les exigences de deux
          référentiels, en qualifiant chaque lien (Ch.~\ref{chap:referentiels}).
    \item[Taux de conformité ($\mu$)] Part des réponses \OUI{}
          (Ch.~\ref{chap:prologue}, \ref{chap:conclure}).
    \item[Taux d'incertitude ($U$)] Part des réponses \NEK{}
          (Ch.~\ref{chap:prologue}, \ref{chap:conclure}).
    \item[Valeur métier] Activité, processus ou information qui a de la
          valeur pour l'organisation (Ch.~\ref{chap:risque}).
    \item[Vulnérabilité] Faiblesse d'un bien support qu'une menace peut
          exploiter (Ch.~\ref{chap:risque}).
\end{description}

    % ============================================================================
% Fichier  : annexes/Ann_B_outils.tex
% Titre    : Référentiel des Outils
% Version  : 1.0 -- novembre 2026
% ============================================================================
% Outils cités à titre d'exemples courants, sans recommandation commerciale ;
% versions et licences à vérifier avant usage. Tout outil de test technique
% ne s'emploie qu'avec une autorisation écrite (Chapitre 12, loi n° 2010/012,
% art. 69).
% ============================================================================

\chapter{Référentiel des Outils}
\label{ann:outils}

Les outils ci-dessous sont des exemples courants, cités pour situer les
tâches du cours. Aucun n'est imposé~; aucun ne remplace la méthode. Les
versions et les licences sont à vérifier avant usage.

\begin{boxalert}
\textbf{Règle.} Tout outil qui interroge un réseau ou un système (scan,
test d'accès, test d'intrusion) ne s'emploie qu'avec une autorisation écrite
du responsable du système, et de l'hébergeur s'il y a lieu
(Chapitre~\ref{chap:audit-technique}, loi n°~2010/012, article~69).
\end{boxalert}

\section*{Gestion et quantification des risques}

\begin{tableaucours}{L{3.6cm}L{5.0cm}c}{%
    \tableauHeader{Tâche} &
    \tableauHeader{Outils} &
    \tableauHeader{Chap.}
}
Registre, matrice, critères & Tableur partagé, à défaut d'un outil GRC &
\ref{chap:risque}, \ref{chap:grc-pratique} \\
Étude EBIOS Risk Manager & Modèle de tableur mis à disposition par
l'ANSSI~; outils labellisés par l'ANSSI & \ref{chap:ebios} \\
Simulation de Monte-Carlo & Python (bibliothèque numpy) ou tableur~; le
générateur \texttt{fair\_sanaga\_C05.py} de NDAP SUKULU & \ref{chap:fair} \\
Mesures $\mu$, $U$, $H$, $D_{\mathrm{KL}}$ & Tableur, Python &
\ref{chap:conclure}, \ref{chap:auditologie} \\
\end{tableaucours}

\section*{Audit et audit technique}

\begin{tableaucours}{L{3.6cm}L{5.0cm}c}{%
    \tableauHeader{Tâche} &
    \tableauHeader{Outils} &
    \tableauHeader{Chap.}
}
Sceller un export & Calcul d'empreinte SHA-256 (\texttt{sha256sum},
\texttt{Get-FileHash}) & \ref{chap:audit-technique} \\
Comparer à une référence & Guides CIS Benchmarks et outil d'évaluation
associé~; scripts de comparaison & \ref{chap:audit-technique} \\
Inventaire des services & Scanner de ports (par exemple Nmap) &
\ref{chap:audit-technique} \\
Vulnérabilités & Scanner authentifié (par exemple Greenbone OpenVAS)~;
calculateur CVSS du FIRST & \ref{chap:audit-technique} \\
Comptes et journaux & Extractions de l'annuaire et du collecteur de
journaux~; tableur pour les rapprochements & \ref{chap:audit},
\ref{chap:audit-technique} \\
Audit continu & Scripts de rapprochement planifiés, détenus par l'audit &
\ref{chap:emergents} \\
\end{tableaucours}

\section*{NDAP SUKULU}

Les fiches d'auto-évaluation et les ateliers du cours sont accessibles sur
NDAP SUKULU, sous le code \codeUE{}. Ils s'appuient sur des pièces et des
exports figés fictifs, générés pour le cours~: aucun test n'y porte sur un
système réel.

    % ============================================================================
% Fichier  : annexes/Ann_C_templates.tex
% Titre    : Modèles de Documents Opérationnels
% Version  : 1.0 -- novembre 2026
% ============================================================================
% Rubriques reprises des encadrés « Fiche de risque » et « Papier de
% travail » des chapitres ; chaque modèle renvoie au chapitre qui l'explique.
% ============================================================================

\chapter{Modèles de Documents Opérationnels}
\label{ann:templates}

Cette annexe rassemble les rubriques des documents que le cours a
construits. Chaque modèle renvoie au chapitre où il est expliqué et
illustré.

\section*{Vue d'ensemble}

\begin{tableaucours}{cL{4.2cm}L{5.4cm}c}{%
    \tableauHeader{N°} &
    \tableauHeader{Document} &
    \tableauHeader{Usage} &
    \tableauHeader{Chap.}
}
M1 & Fiche de risque & Décrire, coter et décider un risque &
\ref{chap:risque} \\
M2 & Décision d'acceptation & Accepter un résiduel selon les critères &
\ref{chap:iso27005} \\
M3 & Ligne de plan de traitement & Planifier une action et sa preuve &
\ref{chap:iso27005} \\
M4 & Ligne de déclaration d'applicabilité & Justifier une mesure &
\ref{chap:iso27000} \\
M5 & Lettre de mission & Fixer objectifs, périmètre, critères, accès &
\ref{chap:audit} \\
M6 & Note d'autorisation de test & Encadrer un test technique &
\ref{chap:audit-technique} \\
M7 & Papier de travail d'un test & Rendre un test refaisable &
\ref{chap:audit} \\
M8 & Fiche d'assertion & Définir un contrôle et sa cote &
\ref{chap:conclure} \\
M9 & Observation en sept éléments & Rédiger une non-conformité &
\ref{chap:conclure} \\
M10 & Fiche d'indicateur & Définir un indicateur et sa fiabilité &
\ref{chap:smsi} \\
M11 & Page GRC de la gouvernance & Une page pour décider &
\ref{chap:grc-pratique} \\
\end{tableaucours}

\section*{M1. Fiche de risque}

\begin{tableaucours}{L{3.6cm}L{8.6cm}}{%
    \tableauHeader{Rubrique} &
    \tableauHeader{Contenu attendu}
}
Référence et intitulé & Code unique (R-XX-nn), scénario en une ligne \\
Valeur métier, biens supports & Ce que l'on protège, ce par quoi on est
attaqué \\
Source de risque, objectif visé & Qui, pour quoi faire \\
Événement redouté & L'atteinte, du point de vue de l'organisation \\
Vraisemblance et impact & Cotes de 1 à 4, avec la base de preuve de chacune
(\OUI{}, \NON{}, \NEK{}) \\
Mesures existantes & Par étape du scénario, avec la preuve de leur test \\
Résiduel & Criticité selon la convention de réduction \\
Propriétaire, décision, revue & Qui répond, ce qui est décidé, quand on
revoit \\
\end{tableaucours}

\section*{M2 et M3. Décision d'acceptation, plan de traitement}

\begin{tableaucours}{L{3.6cm}L{8.6cm}}{%
    \tableauHeader{Rubrique} &
    \tableauHeader{Contenu attendu}
}
Décision d'acceptation & Risque, propriétaire, résiduel et base de preuve,
mesures en place, décision («~accepté~» ou «~accepté en l'état des
connaissances~»), conditions, durée, date de revue, signatures \\
Ligne de plan & Action, risque et étape visés, responsable nommé, échéance,
coût, preuve attendue, effet attendu sur la cote \\
\end{tableaucours}

\section*{M4. Ligne de déclaration d'applicabilité}

Mesure (numéro et intitulé de l'ISO/IEC~27002:2022)~; incluse ou exclue~;
justification (risque du registre ou obligation)~; mise en œuvre (oui,
partielle, à vérifier) avec la référence de la preuve~; pour une exclusion,
les faits vérifiés et l'approbation datée.

\section*{M5 et M6. Lettre de mission, note d'autorisation}

\begin{tableaucours}{L{3.6cm}L{8.6cm}}{%
    \tableauHeader{Document} &
    \tableauHeader{Rubriques}
}
Lettre de mission & Commanditaire~; objectif~; périmètre et période~;
critères~; équipe et indépendance~; accès aux systèmes, documents,
personnes et prestataires~; destinataires du rapport \\
Note d'autorisation & Signataire habilité et accord de l'hébergeur~;
périmètre inclus et exclu~; méthodes permises et interdites~; fenêtres~;
contact et procédure d'arrêt~; confidentialité et destruction des
résultats \\
\end{tableaucours}

\section*{M7. Papier de travail d'un test}

Référence du test et critère~; population, échantillon et méthode de
sélection~; technique et date~; preuves obtenues, rang de fiabilité et lieu
de conservation (avec empreinte pour un export)~; résultat et cote, avec,
pour un \NEK{}, sa nature (preuve insuffisante ou limitation de
l'étendue)~; auteur et réviseur.

\section*{M8 et M9. Fiche d'assertion, observation}

\begin{tableaucours}{L{3.6cm}L{8.6cm}}{%
    \tableauHeader{Document} &
    \tableauHeader{Rubriques}
}
Fiche d'assertion & Position dans le gabarit~; objectif~; un critère et une
question par condition (atomicité)~; test, preuve, trace~; risque~;
conséquence \\
Observation & 1.~Norme~; 2.~Constatation chiffrée, preuve et cote~;
3.~Risques~; 4.~Contradictoire daté~; 5.~Analyse et confiance~;
6.~Conclusion~; 7.~Recommandations qui traitent la cause \\
\end{tableaucours}

\section*{M10 et M11. Indicateur, page GRC}

\begin{tableaucours}{L{3.6cm}L{8.6cm}}{%
    \tableauHeader{Document} &
    \tableauHeader{Rubriques}
}
Fiche d'indicateur & Objectif servi~; famille (mise en œuvre, efficacité,
risque)~; définition et formule~; source et collecte~; fréquence~; cible et
seuils~; responsable~; cote de fiabilité de la source \\
Page GRC & 1.~Risques au-delà de l'appétence~; 2.~Obligations~;
3.~Contrôles clés~; 4.~Indicateurs de risque~; 5.~Ce que nous ne savons
pas~; 6.~Décisions demandées \\
\end{tableaucours}

    % ============================================================================
% Fichier  : annexes/Ann_D_textes.tex
% Titre    : Textes et Normes de Référence
% Version  : 1.0 -- novembre 2026
% ============================================================================
% Extraits identiques à ceux des chapitres. Loi n° 2010/012 : texte français
% officiel lu (Ann. D d'ACSI). COBAC R-2016/04 : art. 1 et 2 relevés sur une
% base en ligne, à collationner. Normes ISO : paraphrasées (texte sous
% licence) ; éditions vérifiées sur iso.org le 06/11/2026 pour 27000, 27001
% (et Amd 1:2024) et 27002.
% ============================================================================

\chapter{Textes et Normes de Référence}
\label{ann:textes}

Cette annexe rassemble les textes cités dans le cours, avec leur degré de
vérification. Seuls les extraits lus dans leur version officielle sont
reproduits entre guillemets (Chapitre~\ref{chap:droit}).

\section*{Degré de vérification}

\begin{tableaucours}{L{4.4cm}L{4.0cm}L{3.8cm}}{%
    \tableauHeader{Texte} &
    \tableauHeader{Degré de vérification} &
    \tableauHeader{Chapitres}
}
Loi n°~2010/012 (cybersécurité) & Texte officiel lu~; extraits cités &
\ref{chap:grc}, \ref{chap:iso27005}, \ref{chap:audit}, \ref{chap:droit},
\ref{chap:audit-technique} \\
Règlement COBAC R-2016/04 & Art.~1 et~2 relevés, à collationner &
\ref{chap:grc}, \ref{chap:droit} \\
Convention de Malabo (2014) & Présentée~; statut camerounais à vérifier &
\ref{chap:droit} \\
Loi n°~2024/017 (données personnelles) & À vérifier au Journal officiel &
\ref{chap:grc}, \ref{chap:droit}, \ref{chap:emergents} \\
Droit OHADA des sociétés & Présenté, à vérifier & \ref{chap:droit} \\
\end{tableaucours}

\section*{Loi n° 2010/012 relative à la cybersécurité et à la cybercriminalité}

\begin{boxjuris}
\textbf{Article~4, point~9}~: \textit{« Audit de sécurité~: examen
méthodique des composantes et des acteurs de la sécurité, de la politique,
des mesures, des solutions, des procédures et des moyens mis en oeuvre par
une organisation, pour sécuriser son environnement, effectuer des contrôles
de conformité, des contrôles d'évaluation de l'adéquation des moyens
(organisationnels, techniques, humains, financiers) investis au regard des
risques encourus, d'optimisation, de rationalité et de performance. »}

\medskip
\textbf{Article~7, alinéa~2} (extrait)~: l'Agence a notamment pour mission
\textit{« de s'assurer de la régularité, de l'effectivité des audits de
sécurité des systèmes d'information suivant les normes en la matière, des
organismes publics et des autorités de certification »}.

\medskip
\textbf{Article~13, alinéa~1}~: \textit{« Sont soumis à un audit de sécurité
obligatoire, les réseaux de communications électroniques et les systèmes
d'information, des opérateurs, des autorités de certification et des
fournisseurs de services de communications électroniques. »}

\medskip
\textbf{Article~14}~: \textit{« Le personnel de l'Agence et les experts
commis en vue d'accomplir des opérations d'audit sont astreints au secret
professionnel. »}

\medskip
\textbf{Article~26, alinéa~1}~: \textit{« Les exploitants des systèmes
d'information prennent toutes les mesures techniques et administratives afin
de garantir la sécurité des services offerts. A cet effet, ils se dotent de
systèmes normalisés leur permettant d'identifier, d'évaluer, de traiter et
de gérer continûment les risques liés à la sécurité des systèmes
d'information dans le cadre des services offerts directement ou
indirectement. »}

\medskip
\textbf{Article~29, alinéa~1}~: \textit{« Les exploitants des systèmes
d'information ont l'obligation de conserver les données de connexion et de
trafic de leurs systèmes d'information pendant une période de dix (10)
ans. »}

\medskip
\textbf{Article~32, alinéas~1 à~3}~: \textit{« (1) Les réseaux de
communications électroniques et les systèmes d'information sont soumis à un
régime d'audit de sécurité obligatoire et périodique de leurs systèmes de
sécurité par l'Agence. (2) L'audit de sécurité et les mesures d'impact de
gravité sont effectués chaque année ou lorsque les circonstances
l'exigent. (3) Les rapports d'audit sont confidentiels et adressés au
Ministre chargé des Télécommunications. »}

\medskip
\textbf{Article~61, alinéa~1}~: \textit{« Sont punis d'un emprisonnement de
trois (03) mois à trois (03) ans et d'une amende de 20.000 (vingt mille) à
100.000 (cent mille) F CFA, les personnels de l'Agence et les experts des
personnes morales chargés des audits qui révèlent sans autorisation, des
informations confidentielles dont ils ont eu connaissance à l'occasion
d'un audit de sécurité. »}

\medskip
\textbf{Article~69}~: \textit{« Est puni d'un emprisonnement de cinq (05) à
dix (10) ans et d'une amende de 10.000.000 (dix millions) à 100.000.000
(cent millions) F CFA ou de l'une de ces deux peines seulement, celui qui
accède sans droit, et en violation des mesures de sécurité, à l'ensemble ou
à une partie d'un réseau de communications électroniques, d'un système
d'information ou d'un équipement terminal, afin d'obtenir des informations
ou des données, en relation avec un système d'information connecté à un
autre système d'information. »}

\medskip
\textbf{Résumés} (texte officiel lu, non cité)~: article~26, alinéas~2
et~3 (mécanismes de sécurité visés par l'Agence)~; article~29, alinéa~2
(surveillance des accès aux données)~; article~30, alinéa~1 (évaluation et
révision des systèmes de sécurité)~; article~61, alinéa~3 (obstacle aux
audits)~; article~70 (attaque par saturation).
\end{boxjuris}

\section*{Règlement COBAC R-2016/04}

\begin{boxjuris}
Texte relevé sur une base réglementaire en ligne, à collationner avec la
version du Secrétariat général de la COBAC~:

\medskip
\textbf{Article~1}~: \textit{« Le présent règlement fixe les exigences
minimales relatives au système de contrôle interne que les établissements
assujettis doivent mettre en œuvre. »}

\medskip
\textbf{Article~2}, risque opérationnel~: \textit{« risque de pertes
résultant de l'inadéquation ou de la défaillance des processus internes,
des personnes et des systèmes, ou d'évènements externes. »}
\end{boxjuris}

\section*{Normes et cadres}

\begin{tableaucours}{L{4.8cm}L{3.8cm}c}{%
    \tableauHeader{Norme ou cadre} &
    \tableauHeader{Statut dans le cours} &
    \tableauHeader{Chap.}
}
ISO~31000:2018 & Paraphrasée & \ref{chap:risque}, \ref{chap:iso27005} \\
ISO/IEC~27000:2018 & Paraphrasée~; édition vérifiée & \ref{chap:grc} \\
ISO/IEC~27001:2022 et amendement~1:2024 & Paraphrasée~; éditions vérifiées
& \ref{chap:iso27000}, \ref{chap:smsi} \\
ISO/IEC~27002:2022 & Paraphrasée~; édition vérifiée~; répartition des
mesures à confirmer & \ref{chap:iso27000} \\
ISO/IEC~27004, 27005, 27007, 27014, 27017, 27018 & Paraphrasées &
\ref{chap:grc} à \ref{chap:emergents} \\
ISO~19011:2018, ISO/IEC~17021-1:2015 & Paraphrasées & \ref{chap:audit},
\ref{chap:smsi} \\
ISO/IEC~42001:2023, NIST AI RMF & Présentés & \ref{chap:emergents} \\
EBIOS Risk Manager (ANSSI) & Ateliers vérifiés sur le site de l'ANSSI~;
quatre détails à confirmer dans le guide & \ref{chap:ebios} \\
Open FAIR (O-RT 3.0.1, O-RA 2.0.1) & Versions vérifiées & \ref{chap:fair} \\
NIST CSF~2.0, SP~800-30, 800-53, 800-53A & Présentés & \ref{chap:grc},
\ref{chap:referentiels} \\
CIS Controls v8, CIS Benchmarks, PCI DSS v4.0, CVSS v4.0 & Présentés &
\ref{chap:referentiels}, \ref{chap:audit-technique} \\
\end{tableaucours}

    % ============================================================================
% Fichier  : annexes/Ann_E_bibliographie.tex
% Titre    : Bibliographie Commentée
% Version  : 1.0 -- novembre 2026
% ============================================================================
% Les textes de l'auteur sont des prépublications ; les travaux en
% préparation (A2 à A5) ne sont pas cités (cadre de cohérence du corpus).
% ============================================================================

\chapter{Bibliographie Commentée}
\label{ann:bibliographie}

Pour chaque partie, quelques lectures, avec ce qu'elles apportent. Les
références complètes figurent dans la bibliographie du livre. Les textes de
l'auteur sont des prépublications~; les travaux en préparation ne sont pas
cités.

\section*{Pour commencer}

\begin{itemize}
    \item \textbf{Le programme de l'auditologie}~\cite{minka2026auditologie}~:
          objet, axiomes, piliers et méthode~; à lire après le prologue, et à
          relire avec le Chapitre~\ref{chap:auditologie}.
    \item \textbf{L'instrument trichotomique}~\cite{minka2026nek}~: la
          réponse \NEK{}, ses conditions, les mesures et l'indice d'usage
          stratégique.
\end{itemize}

\section*{Partie I · Fondements}

\begin{itemize}
    \item L'IIA, le modèle des trois lignes~\cite{iia2020troislignes}~: la
          répartition des rôles, en une dizaine de pages.
    \item L'ISO/IEC~27014~\cite{iso27014_2020} et la fonction \emph{Gouverner}
          du NIST CSF~2.0~\cite{nist2024csf}~: la gouvernance de la sécurité.
    \item L'ISO~31000~\cite{iso31000_2018}~: le vocabulaire du risque.
    \item Cox, sur les matrices de risque~\cite{cox2008matrices}~: ce que
          les couleurs cachent.
\end{itemize}

\section*{Partie II · Méthodes}

\begin{itemize}
    \item L'ISO/IEC~27005:2022~\cite{iso27005_2022}~: le processus appliqué à
          la sécurité de l'information.
    \item Le guide EBIOS Risk Manager de l'ANSSI~\cite{anssi2018ebios} et ses
          fiches méthodes~: la méthode et ses ateliers.
    \item Freund et Jones~\cite{freund2015fair}~: l'exposé de référence de
          FAIR~; Hubbard et Seiersen~\cite{hubbard2016measure}~: le calibrage
          des experts.
    \item L'ISO~19011~\cite{iso19011_2018} et le NIST SP~800-53A~\cite{nist2022sp80053a}~:
          les principes de l'audit et les méthodes d'évaluation.
\end{itemize}

\section*{Partie III · Cadre normatif et juridique}

\begin{itemize}
    \item L'ISO/IEC~27001:2022~\cite{iso27001_2022} et l'ISO/IEC~27002:2022~\cite{iso27002_2022}~:
          à lire ensemble, l'une pour les exigences, l'autre pour les
          mesures.
    \item Les CIS Controls~\cite{cis2021controls}~: une hygiène de base
          priorisée.
    \item La loi n°~2010/012~\cite{loi2010012}, dans son texte officiel, et le
          règlement COBAC R-2016/04~\cite{cobac2016r04}.
    \item Les normes de l'IIA~\cite{iia2024gias} et le code de
          l'IESBA~\cite{iesba2023code}~: la déontologie.
\end{itemize}

\section*{Partie IV · Pratique}

\begin{itemize}
    \item L'ISO/IEC~27007~\cite{iso27007_2020} et l'ISO/IEC~27004~\cite{iso27004_2016}~:
          auditer et mesurer un SMSI.
    \item Les CIS Benchmarks~\cite{cis2026benchmarks} et la spécification
          CVSS~4.0~\cite{first2023cvss4}~: l'audit technique.
\end{itemize}

\section*{Parties V et VI · Défis avancés et recherche}

\begin{itemize}
    \item L'ISO/IEC~42001~\cite{iso42001_2023} et le NIST AI RMF~\cite{nist2023airmf}~:
          les cadres de l'IA.
    \item L'ISO/IEC~27017~\cite{iso27017_2015}~: la sécurité du nuage.
    \item Shannon~\cite{shannon1948}, Kullback et Leibler~\cite{kullback1951},
          Lakatos~\cite{lakatos1970}~: les fondements des mesures et de la
          méthode.
\end{itemize}

    % ============================================================================
% Fichier  : annexes/Ann_F_contacts.tex
% Titre    : Contacts et Ressources
% Version  : 1.0 -- novembre 2026
% ============================================================================
% Adresses de sites institutionnels données à titre indicatif, à vérifier
% avant diffusion (reprises d'ACSI, Ann. F, complétées). Contact du Cyber
% Range : celui fixé par l'auteur pour le portail (ACSI).
% ============================================================================

\chapter{Contacts et Ressources}
\label{ann:contacts}

Les adresses ci-dessous sont celles des sites institutionnels à la date de
rédaction~; elles peuvent changer et sont à vérifier avant usage.

\section*{Autorités et institutions camerounaises et régionales}

\begin{tableaucours}{L{4.4cm}L{4.6cm}L{3.2cm}}{%
    \tableauHeader{Institution} &
    \tableauHeader{Rôle dans le cours} &
    \tableauHeader{Site}
}
ANTIC & Régulation de la sécurité électronique, audits de sécurité
(loi n°~2010/012) & \texttt{www.antic.cm} \\
MINPOSTEL & Destinataire des rapports d'audit de sécurité (art.~32) &
\texttt{www.minpostel.gov.cm} \\
COBAC (Secrétariat général) & Contrôle interne des établissements de
crédit de la CEMAC & \texttt{www.sgcobac.org} \\
Chambre des comptes de la Cour suprême & Contrôle juridictionnel des
comptes publics & · \\
Union africaine & Convention de Malabo & \texttt{au.int} \\
\end{tableaucours}

\section*{Organisations normatives et professionnelles}

\begin{tableaucours}{L{4.4cm}L{4.6cm}L{3.2cm}}{%
    \tableauHeader{Organisation} &
    \tableauHeader{Ressources} &
    \tableauHeader{Site}
}
ISO & Normes ISO et ISO/IEC & \texttt{www.iso.org} \\
ANSSI (France) & EBIOS Risk Manager, fiches méthodes &
\texttt{cyber.gouv.fr} \\
NIST & CSF~2.0, SP~800-30, 800-53, AI RMF & \texttt{www.nist.gov} \\
The Open Group & Normes Open FAIR & \texttt{www.opengroup.org} \\
Center for Internet Security & CIS Controls, CIS Benchmarks &
\texttt{www.cisecurity.org} \\
FIRST & Spécification et calculateur CVSS & \texttt{www.first.org} \\
IIA & Normes mondiales d'audit interne, modèle des trois lignes &
\texttt{www.theiia.org} \\
ISACA & COBIT, ITAF, certifications & \texttt{www.isaca.org} \\
\end{tableaucours}

\section*{Recherche}

Les articles A0 (\texttt{hal-05701007}) et A1 (\texttt{hal-05701270}) sont
déposés sur l'archive ouverte HAL (\texttt{hal.science}). Les travaux en
préparation y seront déposés à leur tour.

\section*{NDAP SUKULU}

Les fiches d'auto-évaluation et les ateliers du cours sont disponibles sur
NDAP SUKULU, sous le code \codeUE{}. NDAP SUKULU est une propriété de
MaletYon, mise à la disposition des étudiants et des praticiens souhaitant
se perfectionner~: \texttt{maletyon@proton.me}.


% --------------------------------------------------------------------------
% ÉLÉMENTS FINAUX
% --------------------------------------------------------------------------
\backmatter

    \nocite{*}
    \printbibliography[heading=bibintoc, title={Références bibliographiques}]

    \printindex

% ==============================================================================
% 4ÈME DE COUVERTURE
% ==============================================================================
\clearpage
\thispagestyle{empty}
\pagecolor{CouleurPrimaire}

\begin{tikzpicture}[remember picture, overlay]
    \fill[CouleurAccent]
        (current page.south west) rectangle
        ([yshift=1.8cm] current page.south east);
    \fill[CouleurAccent]
        ([yshift=-1.4cm] current page.north west) rectangle
        ([yshift=-1.7cm] current page.north east);
\end{tikzpicture}

\vspace*{2.5cm}

\begin{center}
    {\color{CouleurAccent}\rule{0.6\linewidth}{1pt}}\\[1.5ex]
    {\color{white}\Large\bfseries À propos de ce cours}\\[1.5ex]
    {\color{CouleurAccent}\rule{0.6\linewidth}{1pt}}
\end{center}

\vspace{1cm}

\begin{center}
\begin{minipage}{0.75\textwidth}
\color{white!85!gray}\setstretch{1.4}
Un certificat n'est pas une preuve de sécurité, et une analyse de risque
sans preuve n'est qu'une opinion chiffrée. Ce cours enseigne à appliquer les
normes, à gérer les risques et à auditer la cybersécurité, en disant toujours
ce que l'on sait, comment on le sait, et ce que l'on ignore.

\medskip

Il progresse en six parties, précédées d'un prologue~: la gouvernance et le
risque, les méthodes (ISO/IEC~27005, EBIOS Risk Manager, FAIR, audit), le cadre
normatif et juridique, la pratique (SMSI, audit technique, GRC, études de
cas), les défis avancés et la recherche en auditologie.

\medskip

Il s'adresse aux étudiants du Master Professionnel en Sécurité des Systèmes
Numériques et aux
professionnels qui exercent comme RSSI, auditeurs en cybersécurité,
gestionnaires des risques cyber ou responsables GRC.
\end{minipage}
\end{center}

\vfill

\begin{center}
    {\color{CouleurAccent}\rule{0.7\linewidth}{1.2pt}}\\[1.2ex]
    {\color{white}\large\bfseries MINKA MI NGUIDJOÏ Thierry Emmanuel}\\[0.4ex]
    {\color{white!80!gray}\small
        CISA~·~CISM~·~CGEIT~·~CRISC~·~CDPSE~·~ISO 27001~·~ISO 22301~·~ISO 9001}\\[0.3ex]
    {\color{white!70!gray}\footnotesize\itshape
        Chercheur au laboratoire d'Ingénierie Mathématiques et des Systèmes d'Information}\\[0.3ex]
    {\color{white!70!gray}\footnotesize\itshape
        Chercheur~|~Cryptologue~|~Auditeur Senior~|~Expert Judiciaire~|~GRC \& Sécurité de l'Information}\\[1.2ex]
    {\color{CouleurAccent}\rule{0.7\linewidth}{1.2pt}}
\end{center}

\vspace{1cm}

\afterpage{\nopagecolor}
\end{document}
